\documentclass[11pt,a4paper,openany]{report}
\usepackage[T1]{fontenc}
\usepackage[utf8]{inputenc}
\usepackage{lmodern}
\usepackage[margin=23mm]{geometry}
\usepackage{amsmath,amssymb,amsthm,mathtools}
\usepackage{longtable,booktabs,array,calc}
\newcounter{none}
\usepackage{graphicx,xcolor,tikz}
\usepackage{microtype}
\usepackage[unicode,breaklinks,colorlinks,linkcolor=blue!45!black,urlcolor=blue!45!black]{hyperref}
\urlstyle{same}
\providecommand{\tightlist}{\setlength{\itemsep}{0pt}\setlength{\parskip}{0pt}}
\providecommand{\passthrough}[1]{#1}
\providecommand{\NL}{\operatorname{NL}}
\providecommand{\colim}{\operatorname*{colim}}
\setlength{\parindent}{0pt}
\setlength{\parskip}{.4em}
\setlength{\emergencystretch}{3em}
\setcounter{tocdepth}{1}
\setcounter{secnumdepth}{0}
\renewcommand{\chaptername}{Lesson}
\DeclareUnicodeCharacter{03BB}{\ensuremath{\lambda}}
\DeclareUnicodeCharacter{0393}{\ensuremath{\Gamma}}
\DeclareUnicodeCharacter{2192}{\ensuremath{\rightarrow}}
\DeclareUnicodeCharacter{2011}{-}
\DeclareUnicodeCharacter{202F}{\,}
\DeclareUnicodeCharacter{2212}{\ensuremath{-}}
\DeclareUnicodeCharacter{03B1}{\ensuremath{\alpha}}
\DeclareUnicodeCharacter{03B2}{\ensuremath{\beta}}
\DeclareUnicodeCharacter{03BC}{\ensuremath{\mu}}
\DeclareUnicodeCharacter{03A6}{\ensuremath{\Phi}}
\DeclareUnicodeCharacter{2260}{\ensuremath{\ne}}
\DeclareUnicodeCharacter{00D7}{\ensuremath{\times}}
\DeclareUnicodeCharacter{2082}{\ensuremath{_2}}
\DeclareUnicodeCharacter{2228}{\ensuremath{\vee}}
\DeclareUnicodeCharacter{2032}{\ensuremath{^{\prime}}}
\DeclareUnicodeCharacter{220E}{\ensuremath{\square}}
\title{Artin approximation for polynomial equations and its desingularization proof}
\author{GPT-6.1 Sol (OpenAI)\\Codex, Ultra setting}
\date{October 2026}
\begin{document}
\hypersetup{pageanchor=false}
\maketitle
\pagenumbering{roman}
\hypersetup{pageanchor=true}
\noindent Written and self-checked by the writing AI. No human or independent review is claimed.
Original contributions are dedicated under CC0 1.0. Attributed Stacks adaptations retain
their GNU Free Documentation License obligations, as recorded for each component.
The cumulative collection includes the complete version 1.2 licence below.
\par\medskip
\noindent Online course: \url{https://kokunoyumeto.github.io/open-math-courses/courses/AG-RG/}.
\begingroup\small\tableofcontents\endgroup
\clearpage
\pagenumbering{arabic}
\chapter*{Artin approximation for polynomial equations and its desingularization proof}\addcontentsline{toc}{chapter}{Artin approximation for polynomial equations and its desingularization proof}\label{AG-RG-S05-artin-approximation-for-polynomial-equations-and-its-desingularization-proof}

\emph{Written and self-checked by GPT-6.1 Sol (OpenAI), Codex, Ultra setting, 5 October 2026. This reconstruction draft includes the desingularization proof and its consumed foundations. It retains adapted GNU FDL 1.2 expression.}

The purpose of approximation is to retain exact equations while replacing coefficients in a completed local ring by coefficients on an étale neighbourhood. A coefficient need not itself belong to the local ring, and approximating each coefficient separately generally destroys its equations. Instead, retain all equations in one finite system, including those that impose the desired congruence.

The only mathematical source used to construct the new desingularization material is the freely accessible pinned Stacks exposition and its supporting chapters. Earlier programme proofs are used where their exact statements apply. The source is material for the proof written here; a link to that source does not discharge a mathematical obligation. The original theorem below remains the target until all of its consumed dependencies are proved in this lesson or an earlier programme lesson.

\section{1. The exact approximation statement}\label{AG-RG-S05-the-exact-approximation-statement}

A map of Noetherian rings is \textbf{regular} if it is flat and every fibre is geometrically regular. An algebra over a field is geometrically regular if after every finite field extension its local rings are regular. A \textbf{G-ring} is a Noetherian ring whose maps from localizations to their completions are regular. The relevant implication for a local G-ring is therefore immediate from its definition: \(R\to\widehat R\) is regular. The broader local criterion and permanence statements, including the property for local rings essentially of finite type over \(\mathbb Z\), are supplied among the foundation arguments below.

\textbf{Theorem 1.1 (Artin approximation in a pointed étale neighbourhood).} Let \((R,\mathfrak m)\) be a Noetherian local G-ring. Let

\[
f_1,\ldots,f_q\in R[X_1,\ldots,X_n],
\qquad \widehat a\in\widehat R^n,
\qquad f_j(\widehat a)=0.
\]

For every positive integer \(N\), there exist a finitely presented étale \(R\)-algebra \(R'\), a maximal ideal \(\mathfrak m'\subset R'\) over \(\mathfrak m\), and a solution \(b\in(R')^n\), such that

\[
\kappa(\mathfrak m')=\kappa(\mathfrak m),\qquad
f_j(b)=0,\qquad
\rho(b_i)-\widehat a_i\in\mathfrak m^N\widehat R.
\tag{1.1}
\]

Here \(\rho:R'\to\widehat R\) is the canonical lift of the specified residue-field point. This is a homomorphism; no assertion that the whole étale algebra embeds into the completion is required. The statement allows arbitrary characteristic, imperfect residue fields, nilpotent base elements, and nonregular local rings.

\textbf{Lemma 1.2 (the canonical comparison map).} A finitely presented étale \(R\)-algebra and a specified \(R\)-map \(R'\to R/\mathfrak m\) have a unique compatible map \(R'\to\widehat R\).

\textbf{Proof.} Starting with the specified map to \(R/\mathfrak m\), lift successively to \(R/\mathfrak m^r\). The kernel of \(R/\mathfrak m^{r+1}\to R/\mathfrak m^r\) is square zero for \(r\geq1\), so existence and uniqueness follow from formal étaleness. The lifts are compatible by uniqueness. Taking their inverse limit gives the desired map. A map into an inverse limit of rings is exactly a compatible collection of maps into the constituent rings, which also proves uniqueness. The equality \(\widehat R/\mathfrak m^r\widehat R=R/\mathfrak m^r\) used here is proved in \href{https://kokunoyumeto.github.io/open-math-courses/courses/AG-CA/completion.html\#2-topology-and-the-powers-of-the-ideal}{AG-CA, Completion, Proposition 2.2 and Theorem 3.3}. Formal étaleness of an étale algebra is the defining lifting property in \href{https://kokunoyumeto.github.io/open-math-courses/courses/AG-CA/formally-smooth-unramified-and-etale-ring-maps.html}{AG-CA, Formally smooth, unramified and étale ring maps, §6}. ∎

\textbf{Lemma 1.3 (a smooth residue-field point has an étale lift).} If \(B\) is a smooth finitely presented \(R\)-algebra and \(B\to\kappa=R/\mathfrak m\) is an \(R\)-map, there are a finitely presented étale \(R\)-algebra \(R'\), a point \(R'\to\kappa\), and an \(R\)-map \(B\to R'\) inducing the original point.

\textbf{Proof.} Use the local standard form proved in \href{https://kokunoyumeto.github.io/open-math-courses/courses/AG-CA/formally-smooth-unramified-and-etale-ring-maps.html\#5-from-a-split-sequence-to-local-coordinates}{AG-CA, Formally smooth, unramified and étale ring maps, Theorem 5.1 and Solution 8.5}. At the kernel of the given point, localize once to write

\[
B_g=R[X_1,\ldots,X_s]/(h_1,\ldots,h_c),
\qquad
\Delta=\det(\partial h_i/\partial X_j)_{i,j\leq c}
\quad\text{a unit}.
\]

This presentation includes any localization variable and its equation. Choose \(r_{c+1},\ldots,r_s\in R\) lifting the point's coordinates in the remaining variables. Substitute these elements in the equations, keeping the first \(c\) variables. Localize the resulting algebra at the substituted \(\Delta\) and at any denominators needed for the original presentation. Its number of equations and variables agree, and its square Jacobian determinant is a unit. The direct Newton calculation proving a standard étale presentation is étale appears in the same earlier lesson, Theorem 4.1 and §6: a square-zero error vector is corrected uniquely by the inverse Jacobian matrix. Thus the resulting \(R'\) is finitely presented étale. The original residue coordinates define its map to \(\kappa\), since the inverted elements have nonzero residues. The kernel of that map is maximal and has residue field \(\kappa\), because the structural map from \(R\) already surjects onto \(\kappa\). Substitution gives \(B_g\to R'\), hence \(B\to R'\). ∎

\textbf{Proof of Theorem 1.1.} The equality of finite-order quotients of a Noetherian local ring and its completion is proved at the preceding Completion locators. Choose \(c_i\in R\) representing \(\widehat a_i\) modulo \(\mathfrak m^N\widehat R\), and choose finitely many generators \(d_1,\ldots,d_M\) of \(\mathfrak m^N\). Write

\[
\widehat a_i=c_i+\sum_{l=1}^M d_l\widehat u_{il}.
\]

Put

\[
g_j(U)=f_j\left(c_1+\sum_l d_lU_{1l},\ldots,
c_n+\sum_l d_lU_{nl}\right),
\qquad A=R[U_{il}]/(g_j).
\tag{1.2}
\]

The tuple \(\widehat u\) gives an \(R\)-map \(A\to\widehat R\). Since \(R\to\widehat R\) is regular, the desingularization theorem proved by the sequence ending at \hyperref[AG-RG-S05-smoothing-theorem-popescu]{Theorem D.38}, and the elementary factorization criterion proved at \hyperref[AG-RG-S05-algebra-lemma-when-colimit]{the corresponding foundation lemma}, give a factorization

\[
A\longrightarrow B\longrightarrow\widehat R
\tag{1.3}
\]

with \(B\) smooth and finitely presented over \(R\). Compose with the residue map \(\widehat R\to\kappa\), then apply Lemma 1.3. The images \(u_{il}\in R'\) solve the \(g_j\) exactly. Set \(b_i=c_i+\sum_l d_lu_{il}\). Substitution in (1.2) gives \(f_j(b)=0\) exactly. Under the canonical map of Lemma 1.2,

\[
\rho(b_i)-\widehat a_i
=\sum_l d_l\bigl(\rho(u_{il})-\widehat u_{il}\bigr)
\in\mathfrak m^N\widehat R.
\]

This proves the prescribed congruence without demanding that the auxiliary \(u_{il}\) themselves be close to \(\widehat u_{il}\). The exact written foundations and the earlier programme proofs used in this argument are registered in §6. ∎

\textbf{Corollary 1.4 (henselian formulation).} If \(R\) is henselian, the solution can be chosen in \(R\). In general it gives a solution in the henselization \(R^h\), descending to the pointed étale neighbourhood just constructed.

\textbf{Proof.} The section criterion proved in \href{https://kokunoyumeto.github.io/open-math-courses/courses/AG-CA/henselian-local-rings-and-henselization.html\#1-roots-and-maps-with-a-prescribed-residue}{AG-CA, Henselian local rings and henselization, Theorem 1.2} gives \(R'\to R\) with the specified residue map. Its composition with \(R\to\widehat R\) is the map \(\rho\) by Lemma 1.2, so (1.1) is preserved. For \(R^h\), use the pointed étale colimit and its universal property proved in that lesson, Lemma 4.1 and Theorem 4.2. Compatibility with the map into the complete henselian ring follows from the same uniqueness. Here the principal monic chart consumed by that section-criterion proof is supplied by the full \href{https://kokunoyumeto.github.io/open-math-courses/courses/AG-RG/AG-RG-S04.html}{earlier structural étale proof, Lemma D2.2}, whose integral closure input is \href{https://kokunoyumeto.github.io/open-math-courses/courses/AG-RG/AG-RG-S04.html}{the affine conductor theorem, Theorem C4.3}. Its polynomial normality input is proved in \href{https://kokunoyumeto.github.io/open-math-courses/courses/AG-RG/AG-RG-S04.html}{the preliminary coefficient argument, Corollary B1.2}. To spell out the actual section step, localize the selected étale algebra to \((R[T]/(F))_h\), with \(F\) monic and \(F'\) nonzero at the given residue root \(\bar a\). The henselian simple-root property gives \(a\in R\) with \(F(a)=0\) and \(a\bmod\mathfrak m=\bar a\). Every denominator has nonzero residue at \(\bar a\), hence is a unit at \(a\), and evaluation defines the section. The resulting map \(R'\to R\) induces the prescribed residue map; formal étale uniqueness along the quotients \(R/\mathfrak m^n\) proves its composition with completion is \(\rho\). Thus this corollary uses these written earlier proofs rather than an external conductor proof tag. ∎

\section{2. Returning from a local model to the original base}\label{AG-RG-S05-returning-from-a-local-model-to-the-original-base}

\textbf{Lemma 2.1 (spreading finite presentation data).} Let \(R_0\) be Noetherian and \(\mathfrak p\subset R_0\). A pointed étale algebra, a finite tuple, and finitely many equations over \((R_0)_{\mathfrak p}\), as in Theorem 1.1, spread to \((R_0)_s\) for some \(s\notin\mathfrak p\). Their solution therefore gives a pointed étale neighbourhood over \(R_0\).

\textbf{Proof.} In the construction of Lemma 1.3, every polynomial coefficient, every fixed nonpivot coordinate, and every chosen inverse has a denominator outside \(\mathfrak p\). Invert their product in \(R_0\). The finitely many equations for the tuple and the finitely many witnesses verifying them likewise hold after multiplying by finitely many additional denominators; invert their product too. The standard étale presentation spreads with its invertible Jacobian determinant, so its spread algebra is étale by the same direct Newton calculation. Regard it as an \(R_0\)-algebra through the étale localization \(R_0\to(R_0)_s\); composing two formally étale, finitely presented maps remains étale because a square-zero lifting problem is solved uniquely in two successive steps. The chosen residue map survives and gives a point over \(\mathfrak p\) with unchanged residue field. The defining equations remain exact, and localization at the selected point gives the original local comparison. ∎

This proof uses the specific standard étale algebra produced above and therefore does not invoke an external spreading theorem for an arbitrary presentation. Pulling the resulting neighbourhood back to the original scheme preserves every defining equation and the specified fibre point.

\section{3. The finite system for an embedded torus}\label{AG-RG-S05-the-finite-system-for-an-embedded-torus}

The approximation step in the torus lesson begins after the formal deformation and effectivity argument has produced an affine torus \(\widehat T\) over the completed local model and a closed homomorphism \(\widehat T\hookrightarrow G_{\widehat R}\). The present lesson proves the approximation input; it does not claim to supply that preceding effectivity argument.

\phantomsection\label{AG-RG-S05-finite-complete-local-splitting}

\textbf{Lemma 3.A (finite splitting over a complete local ring).} Let \((A,\mathfrak m,\kappa)\) be a local ring with \[
A\simeq\varprojlim_{n\geq1}A/\mathfrak m^n.
\] No Noetherian, reducedness, normality or perfection hypothesis is imposed. Suppose a torus \(T/A\) has an affine etale splitting chart \(U\to\operatorname{Spec}A\) and a point of its closed fibre with finite separable residue field \(K/\kappa\). Then a finite free, faithfully flat, etale algebra \(E/A\) splits \(T\). We may choose \(E\) local and complete, with residue field a finite Galois extension \(L/\kappa\) containing \(K\), and with an action of \(\Gamma=\operatorname{Gal}(L/\kappa)\) for which \[
E\otimes_AE\xrightarrow{\sim}\prod_{\sigma\in\Gamma}E,
\qquad x\otimes y\longmapsto\bigl(x\sigma(y)\bigr)_\sigma.
\tag{3.2}
\] Here a torus is etale locally a product of copies of \(\mathbf G_m\). Starting instead with the fpqc multiplicative-type definition gives this same condition by \href{https://kokunoyumeto.github.io/open-math-courses/courses/AG-GS/AG-GS-05.html}{AG-GS-05, Theorem 7.7}; its finite torsion-free character-group step is proved in \href{https://kokunoyumeto.github.io/open-math-courses/courses/AG-RG/AG-RG-S04.html\#finite-torsion-free-lattices}{AG-RG-S04, Lemma P0.10}. Every torus therefore has such a chart and residue point: choose an affine finite-presentation part of an etale splitting cover meeting the closed fibre, and a closed point of that fibre. Its residue field is finite over \(\kappa\) by the finite-type field-algebra argument in AG-RG-S04, Lemma P0.9. The quotient sequence for differentials gives \(\Omega_{K/\kappa}=0\), because the chart is formally unramified. The finite-field differential calculation in AG-CA-17, Lemma 7.1 then proves separability.

Furthermore, a torus \(T_0/\kappa\), with character lattice split by \(L\), has a torus lift over \(A\) split by this same \(E\). Its character lattice and descent action reduce compatibly modulo every \(\mathfrak m^n\).

\textbf{Proof.} A finite separable extension embeds in a finite normal separable extension: adjoin all the roots of the minimal polynomials of a finite generating list. There are finitely many roots and each adjunction has finite separable degree. Every embedding permutes the root sets, so the resulting extension \(L/\kappa\) is normal and separable. Counting embeddings successively gives \(|\Gamma|=[L:\kappa]\). The primitive-element argument of AG-RG-S02, Lemma 4.2, applies also to finite fields, and gives \(L=\kappa(\alpha)\). Let \(\bar f\in\kappa[X]\) be the monic minimal polynomial of \(\alpha\), and lift its coefficients to a monic \(f\in A[X]\). Put \[
E=A[t]=A[X]/(f),\qquad d=\deg(f).
\] Successively subtracting leading multiples of the monic polynomial \(f\) reduces every polynomial to degree less than \(d\). A nonzero multiple \(qf\) has degree \(\deg(q)+d\), since its leading coefficient is the nonzero leading coefficient of \(q\). Thus the remainder is unique, and \(1,t,\ldots,t^{d-1}\) is a basis making \(E\) free of positive finite rank over \(A\); it is faithfully flat because tensoring a module with this free module gives a direct sum of \(d\) copies. Its quotient by \(\mathfrak mE\) is \(L\). If \(e\in E\) has nonzero residue, multiplication by \(e\) on this basis is invertible modulo \(\mathfrak m\). Its determinant is a unit of \(A\), and the cofactor calculation proved in Lemma 3.B gives an inverse to multiplication by \(e\). Thus \(e\) is a unit, \(E\) is local with maximal ideal \(\mathfrak mE\), and coordinatewise completeness gives \[
E\simeq\varprojlim_n E/\mathfrak m^nE.
\] Separability makes \(f'(t)\) a unit by the same argument. A root in a quotient by a square-zero ideal lifts uniquely by replacing any lift \(a\) with \(a-f'(a)^{-1}f(a)\). Polynomial expansion proves existence, and \(f(b)-f(a)=(b-a)(f'(a)+(b-a)H)\) proves uniqueness. Iteration gives the corresponding assertion for every nilpotent ideal. This proves formal etaleness of \(E/A\); its finite presentation makes it etale.

The residue point of \(U\) gives an algebra map \(\mathcal O(U)\to K\to L\). Formal etaleness of \(U/A\) lifts this map uniquely through \(E/\mathfrak m^{n+1}E\to E/\mathfrak m^nE\), whose kernel is square-zero for \(n\geq1\). Taking the compatible limit gives a map \(\mathcal O(U)\to E\), hence a section of \(U_E\to\operatorname{Spec}E\). Pulling back the splitting on \(U\) therefore splits \(T_E\).

We now construct (3.2), including its automorphisms. For every \(\sigma\in\Gamma\), the residue \(\sigma(\alpha)\in L\) is a simple root of \(\bar f\). In the complete local ring \(E\), Newton's recursion \[
a_{j+1}=a_j-f'(a_j)^{-1}f(a_j)
\] starting from any lift of this residue gives \(f(a_j)\in(\mathfrak mE)^{2^j}\) and \(a_{j+1}-a_j\in(\mathfrak mE)^{2^j}\). Completeness and separatedness give a root \(a_\sigma\), and the difference formula above gives its uniqueness with that residue. Distinct \(a_\sigma\) have unit differences. If a polynomial \(P\) vanishes at all these roots, monic division by \(X-a_1\) gives \(P=(X-a_1)Q\), since the remainder is \(P(a_1)=0\). At each other root \(a_\sigma\), the factor \(a_\sigma-a_1\) is a unit, so \(Q(a_\sigma)=0\). Repeating the division shows that \(\prod_\sigma(X-a_\sigma)\) divides \(P\). If \(\deg(P)<d\), monicity forces \(P=0\). The difference between \(f\) and this product has degree less than \(d\) and vanishes at every root, so \[
f(X)=\prod_{\sigma\in\Gamma}(X-a_\sigma).
\] The map \(t\mapsto a_\sigma\) defines an \(A\)-algebra endomorphism \(g_\sigma\) of \(E\). The two roots \(g_\sigma g_\tau(t)\) and \(a_{\sigma\tau}\) have the same residue, so uniqueness gives \(g_\sigma g_\tau=g_{\sigma\tau}\). Also \(a_1=t\); hence these maps are automorphisms and lift the Galois action on \(L\).

The isomorphism \(E\otimes_AE=E[X]/(f)\) and evaluation at the \(a_\sigma\) now give (3.2). Explicitly, the polynomials \[
\ell_\sigma(X)=\prod_{\tau\ne\sigma}
\frac{X-a_\tau}{a_\sigma-a_\tau}
\] give the inverse to evaluation: a tuple \((b_\sigma)_\sigma\) is represented by \(\sum_\sigma b_\sigma\ell_\sigma(X)\). Thus the isomorphism uses no division by \(|\Gamma|\). Its compatibility on triple products is exactly \(g_\sigma g_\tau=g_{\sigma\tau}\). This is the required finite etale Galois torsor.

For the final assertion choose the finite Galois splitting field of \(T_0\) supplied by AG-GS-05, Lemma 5.1, and let \(M\) be its free finite-rank character lattice, with the resulting \(\Gamma\)-action. On \(E[M]\) put \[
\sigma(ae^m)=g_\sigma(a)e^{\sigma m}.
\tag{3.3}
\] The group law and (3.2) make this a descent datum. Multiplication and every Hopf operation commute with it, since addition in \(M\) commutes with its action. The actual module and algebra effectivity proof in AG-RG-S04, Lemmas A1.1--A1.2, gives the invariant Hopf algebra \[
C=E[M]^\Gamma,\qquad E\otimes_AC\simeq E[M].
\tag{3.4}
\] Tensor products and Hopf identities descend by this isomorphism and faithful flatness. Finite presentation descends here by a finite-generator argument. Express the finitely many Laurent algebra generators over \(E\) as finite \(E\)-linear combinations of elements of \(C\). Those elements generate \(C\), since the cokernel of their generated subalgebra becomes zero over \(E\). For a resulting finite polynomial surjection \(A[Z]\to C\), its kernel becomes a finite ideal after tensoring with \(E\), because \(E[M]\) is finitely presented. To see that finiteness for this particular surjection, compare it with a finite presentation of \(E[M]\): express each presentation's generators in the other, and substitute these expressions in the finitely many relations and composite-generator identities. They generate the new kernel. Write a finite list of its generators as finite sums of elements of the original kernel with coefficients in \(E\). Their finitely many kernel components generate the original kernel by faithful flatness. Thus \(C\) is finitely presented. It defines a torus \(T_A\), because (3.4) splits it on the finite etale faithfully flat cover.

Put \(A_n=A/\mathfrak m^n\) and \(E_n=E\otimes_AA_n\). Tensoring (3.2)--(3.4) with \(A_n\) gives the same torsor and the same character action over \(E_n\): \[
E_n\otimes_{A_n}(C\otimes_AA_n)\simeq E_n[M].
\tag{3.5}
\] The algebra on the left descends the datum (3.3) modulo \(\mathfrak m^n\); uniqueness in Lemma A1.2 proves that it is its effective descent. This proves compatibility for every \(n\), even though \(A\to A_n\) need not be flat. It does not assume that taking invariants commutes with an arbitrary quotient. At \(n=1\), the descent datum is precisely the given one on \(L[M]\), so its descended torus is \(T_0\).

More generally, any compatible character data split by these \(E_n\) and identified with the same lattice have no additional nilpotent transition parameters. Comparing coefficients in \(\Delta(\sum c_me^m)=(\sum c_me^m)\otimes(\sum c_me^m)\) gives \(c_m^2=c_m\), \(c_mc_{m'}=0\) for \(m\ne m'\), and \(\sum c_m=1\). Each \(E_n\) is local, so exactly one coefficient is \(1\). A character is therefore one monomial; a torus automorphism is one integral lattice automorphism. A compatible automorphism reducing to a fixed one is that same one. Hence a compatible lattice descent action is the fixed action in (3.3), and (3.5) gives its actual compatible descent. This completes both assertions. \(\square\)

Choose a finite presentation of the Hopf algebra of \(\widehat T\). Lemma 3.A supplies a finite free, etale, faithfully flat splitting algebra \(E/\widehat R\); choose the resulting splitting isomorphism

\[
\mathcal O(\widehat T)\otimes_{\widehat R}E
\simeq E[Z_1^{\pm1},\ldots,Z_r^{\pm1}].
\tag{3.1}
\]

\phantomsection\label{AG-RG-S05-finite-free-perfect-trace}

\textbf{Lemma 3.B (perfect trace, over an arbitrary base).} Let \(A\) be a commutative ring and \(C\) a commutative \(A\)-algebra free of finite rank \(d\). If the determinant of the pairing \((c,z)\mapsto\operatorname{Tr}_{C/A}(cz)\) is a unit, then \(C/A\) is finitely presented and formally etale, hence etale. If \(d>0\), it is faithfully flat. These assertions require no Noetherian, reducedness or characteristic hypothesis.

\textbf{Proof.} First, the elementary matrix identities used here hold over every commutative ring. Define a determinant by the signed permutation sum and a cofactor by the corresponding smaller signed determinant. Expanding along a row or column gives the determinant. A matrix with two equal rows or columns has determinant zero: pair each permutation with the permutation interchanging those positions; their contributions have opposite signs, including in characteristic two. Consequently, if \(M\) is square, multiplying it on either side by its transposed cofactor matrix gives \((\det M)I\): a diagonal entry is the row or column expansion, and an off-diagonal entry is the expansion of a determinant with two equal rows or columns. A unit determinant therefore gives the inverse \((\det M)^{-1}\operatorname{adj}(M)\). For size zero the determinant is one. To verify the converse implication as well, expand \(\det(UV)\) by multilinearity in its columns. The expansion chooses, for each column of \(UV\), one column of \(U\) with its coefficient from \(V\). A tuple repeating a column of \(U\) contributes zero by the cancellation just proved. Every remaining tuple is a permutation of all the columns of \(U\), whose determinant is its permutation sign times \(\det(U)\). Summing its coefficients with these signs is precisely \(\det(V)\). Thus \(\det(UV)=\det(U)\det(V)\), including for size zero. If \(UV=I\), this equation gives \(\det(U)\det(V)=1\), so an invertible matrix has a unit determinant. For square matrices \(U,V\), the finite double sum also gives \(\operatorname{tr}(UV)=\operatorname{tr}(VU)\). Thus conjugating a multiplication matrix by an invertible change of basis preserves its trace. Multiplication matrices, their traces and their determinants commute with every scalar extension, entry by entry.

Suppose first that \(d>0\), and choose a basis \(u_1,\ldots,u_d\). Write \(u_i u_j=\sum_k m_{ijk}u_k\) and \(1=\sum_i a_i u_i\). There is an explicit finite algebra presentation \[
C=A[X_1,\ldots,X_d]/\bigl(X_iX_j-\sum_k m_{ijk}X_k,\;1-\sum_i a_iX_i\bigr).
\] Indeed, repeatedly replace a product of two variables by the displayed linear combination, reducing each monomial to degree at most one. Replace its constant term by that constant times \(\sum_i a_iX_i\). Every quotient element is therefore represented by \(\sum_i b_iX_i\). The map to \(C\) is surjective, and if this representative maps to zero, linear independence of the \(u_i\) makes every \(b_i\) zero. This proves the presentation without an assumption on ideals of \(A\).

The inverse pairing matrix supplies a trace-dual basis \(v_1,\ldots,v_d\), characterized by \(\operatorname{Tr}(u_jv_i)=\delta_{ij}\). In particular \(x=\sum_i u_i\operatorname{Tr}(v_ix)\) and \(x=\sum_i\operatorname{Tr}(u_ix)v_i\). To check the second equality, pair both sides with each \(u_j\) and use invertibility of the pairing; the first is the ordinary basis-coordinate equality. Put \[
e=\sum_i u_i\otimes v_i\in C\otimes_A C.
\] The \(A\)-linear map \[
\Phi:C\otimes_A C\longrightarrow\operatorname{End}_A(C),\qquad
\Phi(a\otimes b)(x)=a\operatorname{Tr}(bx)
\] is an isomorphism, with inverse \(F\mapsto\sum_i F(u_i)\otimes v_i\). The first coordinate equality proves \(\Phi\Phi^{-1}(F)=F\); the second proves \(\Phi^{-1}\Phi(a\otimes b)=a\otimes b\), and pure tensors generate. Since \(\Phi(e)\) is the identity, both \((c\otimes1)e\) and \((1\otimes c)e\) have image multiplication by \(c\). Thus \[
(c\otimes1)e=(1\otimes c)e\qquad(c\in C).
\tag{3.6}
\] Let \(s=\sum_i u_iv_i\). The diagonal coefficient of multiplication by \(c\) in position \(i\) is \(\operatorname{Tr}(v_i c u_i)\). Summing those coefficients gives \(\operatorname{Tr}(c)=\operatorname{Tr}(cs)\) for every \(c\). Invertibility of the pairing now gives \[
\sum_i u_iv_i=1.
\tag{3.7}
\] No division by \(d\), or assertion that \(d\) is a unit, occurs.

We prove the lifting property directly. Let \(J\subset D\) be a square-zero ideal of a commutative \(A\)-algebra, and let \(\bar f:C\to D/J\) be an \(A\)-algebra map. The free-module lifting calculation in \href{https://kokunoyumeto.github.io/open-math-courses/courses/AG-RG/AG-RG-S01.html\#1-the-algebra-of-localization}{AG-RG-S01, Lemma 1.0} gives an \(A\)-linear lift \(l:C\to D\). Its value \(l(1)\) is \(1+j\) with \(j\in J\), and has inverse \(1-j\). Multiplying the whole lift by that inverse makes \(l(1)=1\) while preserving \(A\)-linearity and its reduction. Give \(J\) its \(C\)-module action via \(\bar f\): multiply by any lift of \(\bar f(c)\), which is independent of that lift because \(J^2=0\).

The defect \(b(c,z)=l(cz)-l(c)l(z)\) lies in \(J\), is \(A\)-bilinear and has \(b(c,1)=b(1,c)=0\). Associativity in \(C\) and \(D\) gives, by expanding the four terms, \[
c\,b(z,t)-b(cz,t)+b(c,zt)-b(c,z)\,t=0.
\tag{3.8}
\] Define the \(A\)-linear correction \(h(c)=\sum_i u_i b(v_i,c)\) in \(J\). Apply the \(A\)-linear map \(a\otimes z\mapsto a b(z,t)\) to (3.6); it gives \(c h(t)=\sum_i u_i b(cv_i,t)\). Using (3.8) with \((v_i,c,t)\) and then (3.7), we obtain \[
c h(t)-h(ct)+h(c)t
=\sum_i u_i\bigl(b(v_ic,t)-b(v_i,ct)+b(v_i,c)t\bigr)
=\sum_i u_iv_i b(c,t)=b(c,t).
\] Consequently \(f=l+h\) is multiplicative: the defect of this new map is \(b(c,t)+h(ct)-c h(t)-h(c)t=0\), since \(h(c)h(t)=0\). Also \(h(1)=0\). Thus \(f\) is a unital \(A\)-algebra lift of \(\bar f\).

For uniqueness, the difference \(\delta\) of two such lifts takes values in \(J\) and satisfies \(\delta(ct)=c\delta(t)+\delta(c)t\); expand their products and use \(J^2=0\). Apply \(a\otimes z\mapsto\delta(a)z\) to (3.6). The difference between its two sides is \(\delta(c)\sum_i u_iv_i=\delta(c)\), by that identity and (3.7). It is zero, so every lift is unique. If \(J^N=0\), lift successively through \(D/J^{r+1}\to D/J^r\), for \(1\le r<N\). Its kernel is square-zero because \(2r\ge r+1\). Existence and uniqueness at each step prove formal etaleness for every nilpotent ideal. Together with the explicit finite presentation, this is etaleness by the programme definition.

For \(d=0\), \(C\) is the zero algebra with presentation \(A/(1)\). If a unital map \(C\to D/J\) exists and \(J\) is nilpotent, then \(D/J=0\), so \(1\in J\) and nilpotence forces \(D=0\). There is exactly one lift. Finally, for positive \(d\), tensoring an \(A\)-module with \(C\) gives \(d\) copies of that module; this is exact and detects whether it is zero. This proves faithful flatness. \(\square\)

This is the finite-free case of the discriminant criterion in the freely accessible \href{https://stacks.math.columbia.edu/tag/0BJF}{Stacks Project, Lemma 49.3.1, Tag 0BJF}. The complete lifting proof above supplies the result used here.

Return to the splitting algebra \(E/\widehat R\) from Lemma 3.A. Its trace determinant is a unit, as we now check rather than adding an unsupported inverse variable. By (3.2), after scalar extension from \(\widehat R\) to \(E\) the algebra becomes \(\prod_{\sigma\in\Gamma}E\). Multiplication in this product is diagonal in the coordinate idempotent basis. Its trace pairing matrix is the identity, hence invertible. The preceding trace and basis-change calculations imply that the original pairing matrix becomes invertible over \(E\), so its determinant \(\Delta\) becomes a unit there. Then \((\widehat R/\Delta\widehat R)\otimes_{\widehat R}E=0\). Faithful flatness from Lemma 3.A detects zero modules, giving \(\widehat R/\Delta\widehat R=0\), which means that \(\Delta\) is already a unit of \(\widehat R\).

Fix the finite free basis of \(E\) supplied by Lemma 3.A. Its multiplication table, unit and multiplication matrices define finitely many coefficient variables; associativity, commutativity and the unit laws are finite polynomial equations. Add a variable \(w\) and the equation \(w\Delta=1\) for the trace determinant. The original completed coefficient tuple satisfies this equation by the preceding argument. After approximation it still satisfies the same equation exactly. Lemma 3.B makes the resulting finite free algebra etale and, since its unchanged rank is positive, faithfully flat.

Add coefficients for the finite presentations, Hopf maps, the isomorphism (3.1), its inverse, and the homomorphism to \(G\). Laurent expressions can be written polynomially by adding inverse variables and relations \(Z_iW_i=1\). The maps need to be specified on finite generating sets. Relations between their images need finitely many witnesses in the finitely generated relation ideals. The two composites of the splitting isomorphism and its inverse are tested on finite generators; equality of these composites with the identity is recorded by their witnesses. The same method records comultiplication, counit, inverse, and the homomorphism identities on a finite list of generators. Coassociativity and every needed descent identity are likewise checked on finite generators, using the finite multiplication tables on the finite overlap algebras. Each expression and witness has finite support.

Retain the closed immersion in this same system. The pullback \(\mathcal O(G_{\widehat R})\to\mathcal O(\widehat T)\) is surjective. Choose preimages of a finite list of algebra generators of \(\mathcal O(\widehat T)\), and include the equations asserting that their images are those generators. After approximation these equations still prove surjectivity; consequently the resulting homomorphism is a closed immersion.

We have obtained one finite polynomial system over the local model, solved by the original completed coefficient tuple. Apply Theorem 1.1 at order one. The defining equations hold exactly after approximation, while all the chosen coefficients retain their original residues. The preserved splitting maps identify the source Hopf algebra with a Laurent polynomial Hopf algebra on the finite étale faithfully flat cover. Thus it is a torus, with the specified special fibre and closed embedding. If a prescribed higher jet is needed, use that order instead. Lemma 2.1 spreads the data and the neighbourhood back to an open part of the finite-type model; base change returns them to the original base.

The local model is essentially of finite type over \(\mathbb Z\). The G-ring permanence argument in §5 is therefore a consumed part of this application, including its coefficient-ring inputs. It must be proved before the approximation step is declared complete. An excellence label or an external Stacks link does not replace this work.

\textbf{Scope of the supporting algebra.} The approximation theorem and the desingularization theorem concern Noetherian rings. Every base ring of the finite-type algebra diagrams in their proof is Noetherian. In the foundational complete-intersection arguments below we state and prove the exact Noetherian forms consumed there; arguments explicitly stated over arbitrary rings retain that wider scope. Polynomial presentations with arbitrarily many variables occur only in the two-term cotangent comparison, whose elementary presentation and filtered-colimit proofs are written without Noetherian assumptions. This distinction preserves arbitrary characteristic, imperfect residue fields, and the full Noetherian local G-ring statement.

\section{4. Desingularization proof sequence}\label{AG-RG-S05-desingularization-proof-sequence}

The following sequence contains the actual lifting, presentation, local descent, separable, and inseparable arguments. It is adapted from the pinned free source. The source's historical bibliography is not reproduced. Its statements use the conormal definition of smoothness reproduced among the foundation statements below. The displayed arrow lists specify the maps of the original commutative diagrams.

\subsection{Singular ideals}\label{AG-RG-S05-singular-ideals}

\phantomsection\label{AG-RG-S05-smoothing-section-singular-ideal}

Let \(R \to A\) be a ring map. The singular ideal of \(A\) over \(R\) is the radical ideal in \(A\) cutting out the singular locus of the morphism \(\operatorname{Spec}(A) \to \operatorname{Spec}(R)\). Here is a formal definition.

\phantomsection\label{AG-RG-S05-smoothing-definition-singular-ideal}

\subsubsection{Definition D.1: singular ideal}\label{AG-RG-S05-definition-d.1-singular-ideal}

Reading source: \href{https://stacks.math.columbia.edu/tag/07C5}{Stacks, Tag 07C5}.

For any ring map \(R\to A\), let \(U_{A/R}\) be its smooth locus in \(\operatorname{Spec}A\), with smoothness at a prime understood as in \hyperref[AG-RG-S05-algebra-definition-smooth-at-prime]{F.7}. Its \emph{singular ideal} \(H_{A/R}\) is the radical ideal characterized by \[V(H_{A/R})=\operatorname{Spec}A\setminus U_{A/R}.\] No finite-presentation hypothesis on \(A/R\), or quasi-compactness hypothesis on \(U_{A/R}\), is part of this definition.

\textbf{Verification.} The definition of smoothness at a prime writes \[U_{A/R}=\bigcup_{\substack{g\in A\\ R\to A_g\ \mathrm{smooth}}}D(g).\] Indeed a witnessing \(g\) gives smoothness at every prime avoiding it, and every smooth prime has such a witness. Thus \(U_{A/R}\) is open. Write its closed complement as \(V(J)\) and set \(H_{A/R}=\sqrt J\). A prime contains \(J\) exactly when it contains \(\sqrt J\), since a power belonging to a prime forces its base element into that prime. This proves existence.

Here is the radical-ideal uniqueness argument used below. For any ideal \(J\) and \(x\notin\sqrt J\), the ideals containing \(J\) and avoiding \(\{1,x,x^2,\ldots\}\) have a maximal member \(L\), by the chain-union argument and Zorn's lemma. If \(uv\in L\) while \(u,v\notin L\), maximality provides a power of \(x\) in each of \(L+(u)\) and \(L+(v)\). Their product is again a power of \(x\) and belongs to \(L\), a contradiction. Hence \(L\) is a prime containing \(J\) and avoiding \(x\). It follows that \[\sqrt J=\bigcap_{\mathfrak q\in V(J)}\mathfrak q,\] where the empty intersection is \(A\). Two radical ideals with the same closed set are therefore equal. This also covers \(A=0\) and an empty singular locus. The smooth open need not have a finite subcover, as the fully proved example \hyperref[AG-RG-S05-smoothing-example-not-quasi-compact]{D.6} shows. ∎

\phantomsection\label{AG-RG-S05-smoothing-lemma-find-strictly-standard}

\subsubsection{Lemma D.2: find strictly standard}\label{AG-RG-S05-lemma-d.2-find-strictly-standard}

\emph{Adapted from the Stacks project, smoothing.tex, lines 141--181.}
\nopagebreak[4]

Let \(R\) be a ring. Let \(A = R[x_1, \ldots, x_n]/(f_1, \ldots, f_m)\). Let \(\mathfrak q \subset A\) be a prime ideal. Assume \(R \to A\) is smooth at \(\mathfrak q\). Then there exists an \(a \in A\), \(a \not \in \mathfrak q\), an integer \(c\), \(0 \leq c \leq \min(n, m)\), subsets \(U \subset \{1, \ldots, n\}\), \(V \subset \{1, \ldots, m\}\) of cardinality \(c\) such that

\[
a = a' \det(\partial f_j/\partial x_i)_{j \in V, i \in U}
\]

for some \(a' \in A\) and

\[
a f_\ell \in (f_j, j \in V) + (f_1, \ldots, f_m)^2
\]

for all \(\ell \in \{1, \ldots, m\}\).

\textbf{Proof.} Put \(P=R[x_1,\ldots,x_n]\), \(I=(f_1,\ldots,f_m)\) and \(M=I/I^2\). By \hyperref[AG-RG-S05-algebra-definition-smooth-at-prime]{F.7}, choose \(s\notin\mathfrak q\) such that \(A_s\) is smooth. The presentation comparison \hyperref[AG-RG-S05-algebra-lemma-NL-homotopy]{F.11}, localization \hyperref[AG-RG-S05-algebra-lemma-localize-NL]{F.65}, and definition \hyperref[AG-RG-S05-algebra-definition-smooth]{F.6} give the split exact conormal sequence \[0\longrightarrow M_s\xrightarrow{\mathrm d}A_s^n\longrightarrow\Omega_{A_s/R}\longrightarrow0.\] Its splitting follows because the last module is projective. Localize at \(\mathfrak q\) and tensor with \(K=\kappa(\mathfrak q)\). Splitting preserves the injection after this tensor operation. Choose \(c\) of the equation classes, indexed by \(V\), giving a basis of \(M\otimes_AK\) from its displayed finite spanning list. Their differential columns are linearly independent in \(K^n\), so \(c\le\min(n,m)\).

There are \(c\) rows, indexed by \(U\), on which these columns have nonzero determinant. Here is the finite field calculation. For \(c>0\), move a nonzero entry of the first column to the first row, then subtract appropriate multiples of that row from the other rows. Independence makes the remaining \((n-1)\times(c-1)\) block have independent columns. Inductively choose \(c-1\) rows with a nonzero minor in that block. On these rows together with the first row, the operations just performed do not change the original selected determinant, which is the nonzero pivot times that minor. For \(c=0\), use the empty minor \(1\). Thus \[\delta=\det(\partial f_j/\partial x_i)_{j\in V,\ i\in U}\in A\setminus\mathfrak q.\]

The module \(Q=M/\langle f_j:j\in V\rangle\) is finite, since the \(m\) equation classes generate \(M\). Its residue at \(\mathfrak q\) is zero, so the finite-module Nakayama theorem \hyperref[AG-RG-S05-algebra-lemma-NAK]{F.10} gives \(Q_{\mathfrak q}=0\). For each member of a finite generating list of \(Q\), the definition of localization gives an element outside \(\mathfrak q\) killing it. Their product \(t\notin\mathfrak q\) kills all of \(Q\); for the zero module take \(t=1\).

Set \(a=t\delta\) and \(a'=t\). Then \(a\notin\mathfrak q\) and the claimed determinant equality is exact in \(A\). Also \(aQ=0\), so for every \(\ell\) the conormal class \(a[f_\ell]\) is a combination of the selected classes. Lifting the coefficients and \(a\) to \(P\) gives \[\widetilde a f_\ell\in(f_j:j\in V)+I^2.\] This is precisely the asserted \(A\)-module relation; changing the lift of \(a\) adds an element of \(I^2\). The construction uses the original presentation throughout and cancels no possible zero divisor. It includes empty variable, equation and minor lists. ∎

We will use the notion of a \emph{strictly standard} element in \(A\) over \(R\). We also define an \emph{elementary standard} element to be one of the type we found in the lemma above. We compare the different types of elements in Lemma \hyperref[AG-RG-S05-smoothing-lemma-compare-standard]{D.15}.

\phantomsection\label{AG-RG-S05-smoothing-definition-strictly-standard}

\subsubsection{Definition D.3: strictly standard}\label{AG-RG-S05-definition-d.3-strictly-standard}

Reading source: \href{https://stacks.math.columbia.edu/tag/07C7}{Stacks, Tag 07C7}.

Let \(R\to A\) be finitely presented. A witness for either standard-element condition consists of \[P=R[x_1,\ldots,x_n],\qquad I=(f_1,\ldots,f_m),\qquad A=P/I,\qquad 0\leq c\leq\min(n,m).\] All Jacobian entries below are reduced to \(A\). For \(a\in A\), choose any lift \(\widetilde a\in P\) when writing a polynomial-ideal condition.

The element \(a\) is \emph{elementary standard} if some such witness and \(a'\in A\) satisfy both

\phantomsection\label{AG-RG-S05-equation-elementary-standard-one}

\[a=a'\det\left(\frac{\partial f_j}{\partial x_i}\right)_{i,j=1}^{c}\tag{D.3e1}\]

and

\phantomsection\label{AG-RG-S05-equation-elementary-standard-two}

\[\widetilde a f_{c+j}\in(f_1,\ldots,f_c)+I^2\quad(1\leq j\leq m-c).\tag{D.3e2}\]

The element \(a\) is \emph{strictly standard} if some witness and coefficients \(a_U\in A\), indexed by the \(c\)-element subsets of \(\{1,\ldots,n\}\), satisfy both

\phantomsection\label{AG-RG-S05-equation-strictly-standard-one}

\[a=\sum_{|U|=c}a_U\det\left(\frac{\partial f_j}{\partial x_i}\right)_{j=1,\ldots,c,\ i\in U}\tag{D.3s1}\]

and

\phantomsection\label{AG-RG-S05-equation-strictly-standard-two}

\[\widetilde a f_{c+j}\in(f_1,\ldots,f_c)+I^2\quad(1\leq j\leq m-c).\tag{D.3s2}\]

Use increasing order on each selected variable set in the determinants. Changing the lift \(\widetilde a\) adds an element of \(I\); multiplying this difference by \(f_{c+j}\in I\) puts it in \(I^2\). Thus both relation conditions are independent of the lift. Equivalently they assert \(a[f_{c+j}]\in\langle[f_1],\ldots,[f_c]\rangle\) in the \(A\)-module \(I/I^2\). The minor condition and this relation condition are separate requirements.

For \(c=0\) there is one minor, the empty determinant \(1\), and the first equation imposes no restriction on \(a\). The second equation then reads \(\widetilde a I\subset I^2\). For \(m=c\) the second equation has an empty list of tests. These conventions include empty polynomial presentations and the zero algebra. ∎

\phantomsection\label{AG-RG-S05-smoothing-lemma-parse-equation-strictly-standard-one}

\subsubsection{Lemma D.4: parse equation strictly standard one}\label{AG-RG-S05-lemma-d.4-parse-equation-strictly-standard-one}

Reading source: \href{https://stacks.math.columbia.edu/tag/07ET}{Stacks, Tag 07ET}.

In the presentation of \hyperref[AG-RG-S05-smoothing-definition-strictly-standard]{D.3}, the strict minor condition for \(a\) gives an \(A\)-linear map \(\psi:A^n\to A^c\) for which the composite \[A^c\xrightarrow{e_j\mapsto[f_j]}I/I^2\xrightarrow{\mathrm d}A^n\xrightarrow{\psi}A^c\] is \(a1_c\). Conversely the existence of this map gives the strict minor condition for \(a^c\). These assertions concern the minor condition alone; the remaining-equation condition still has to be imposed independently.

\textbf{Proof.} Identify \(A^n\) with the free differential module on \(dx_1,\ldots,dx_n\). Its conormal map is well-defined: for \(u,v\in I\), the product rule gives \(d(uv)=u\,dv+v\,du\), which vanishes after reduction to \(A\). The composite of its first two maps has the \(n\times c\) matrix \[M_{ij}=\frac{\partial f_j}{\partial x_i}\bmod I.\] For a selected variable set \(U\), let \(E_U:A^n\to A^c\) be coordinate selection and put \(M_U=E_UM\). The determinant in D.3s1 is \(\det M_U\): transposition preserves the determinant by reversing the permutation indexing in its signed expansion.

The arbitrary-ring adjugate and determinant identities are proved in \hyperref[AG-RG-S05-algebra-lemma-matrix-left-inverse]{F.70}. If \(a=\sum_U a_U\det M_U\), they give the explicit map with matrix \[B=\sum_U a_U\operatorname{adj}(M_U)E_U,\qquad BM=a1_c.\] This is the required \(\psi\). Conversely write \(B\) for any proposed \(\psi\). The same proved multilinear expansion gives \[a^c=\det(BM)=\sum_{|U|=c}\det(B^U)\det(M_U),\] where \(B^U\) selects the columns indexed by \(U\). This is the required minor expression for \(a^c\), without cancellation or an invertibility assumption. If \(c=0\), both modules are zero, the map is unique, and \(a^0=1\) has its expression using the empty minor. Zero rings obey the same formulas. ∎

\phantomsection\label{AG-RG-S05-smoothing-lemma-elkik}

\subsubsection{Lemma D.5: elkik}\label{AG-RG-S05-lemma-d.5-elkik}

Reading source: \href{https://stacks.math.columbia.edu/tag/07CA}{Stacks, Tag 07CA}.

For a finitely presented \(R\)-algebra \(A\), let \(J_e\) and \(J_s\) be the ideals generated by all its elementary standard and strictly standard elements, respectively. Then \[H_{A/R}=\sqrt{J_e}=\sqrt{J_s}.\] In particular \(A_a\) is smooth over \(R\) whenever \(a\) is strictly standard.

\textbf{Proof.} First fix a strictly standard witness from D.3 and put \(B=A_a\) and \(N=(I/I^2)\otimes_A B\). Its relation condition says that the first \(c\) equation classes generate \(N\): divide each relation by the unit \(a\). Write \(q:B^c\to N\) for this surjection and \(d:N\to B^n\) for the conormal differential. D.4 gives \(\psi:B^n\to B^c\) with \[\psi d q=a1_c.\] If \(q(v)=0\) this identity gives \(av=0\), hence \(v=0\). Consequently \(q\) is an isomorphism. The map \(\chi=q\,a^{-1}\psi:B^n\to N\) satisfies \(\chi d=1_N\). Thus \[B^n=d(N)\oplus\ker\chi.\] Its quotient by \(d(N)\) is a finite projective module, being a direct summand of \(B^n\), and \(d\) has zero kernel.

The actual localization still has a finite presentation over \(R\): if \(h\in P\) lifts \(a\), it is \(P[z]/(f_1,\ldots,f_m,hz-1)\). The quotient isomorphism follows by adjoining the specified inverse of \(h\). The full presentation comparison \hyperref[AG-RG-S05-algebra-lemma-NL-homotopy]{F.11} and localization calculation \hyperref[AG-RG-S05-algebra-lemma-localize-NL]{F.65} identify its first cotangent homology and differential module with the kernel and cokernel just calculated. Therefore the smoothness definition \hyperref[AG-RG-S05-algebra-definition-smooth]{F.6} proves that \(R\to B\) is smooth. If \(B=0\), these modules are zero and the same finite presentation and criterion apply.

Each singular prime contains every such \(a\), since a prime avoiding \(a\) has the smooth neighbourhood \(A_a\). Hence \(V(H_{A/R})\subset V(J_s)\). Setting every other minor coefficient to zero makes an elementary witness a strict witness, so \(J_e\subset J_s\). If a prime is smooth, \hyperref[AG-RG-S05-smoothing-lemma-find-strictly-standard]{D.2} supplies an elementary witness outside it; relabel its selected variables and equations to place them first, absorbing any ordering sign in the coefficient. Such a prime cannot contain \(J_e\). These statements give \[V(J_e)\subset V(H_{A/R})\subset V(J_s)\subset V(J_e).\] The radical/prime equality proved in D.1 now proves both asserted equalities of ideals. Neither this argument nor the conormal splitting assumes that the base is Noetherian or that the smooth locus is quasi-compact. ∎

\phantomsection\label{AG-RG-S05-smoothing-example-not-quasi-compact}

\subsubsection{Example D.6: not quasi compact}\label{AG-RG-S05-example-d.6-not-quasi-compact}

\emph{Adapted from the Stacks project, smoothing.tex, lines 285--292.}
\nopagebreak[4]

The set of points where a finitely presented ring map is smooth needn't be a quasi-compact open. For example, let \(k\) be a field and set \[
R=k[x,y_1,y_2,y_3,\ldots]/(xy_i\mid i\geq1),\qquad A=R/(x).
\] The distinguished opens below are in \(\operatorname{Spec}(A)\). The smooth locus of \(R\to A\) is \[
U=\bigcup_{i\geq1}D(y_i)=\operatorname{Spec}(A)\setminus\{\mathfrak m\},
\qquad \mathfrak m=(y_1,y_2,\ldots),
\] and this open is not quasi-compact.

\textbf{Proof.} The quotient identifies \(A\) with \(k[y_1,y_2,\ldots]\). The map \(R\to A\) is of finite presentation: it is the quotient by the single generator \(x\), with no algebra variables required. The base ring \(R\) is not assumed Noetherian. Evaluation at all \(y_i=0\) gives \(A/\mathfrak m=k\), so \(\mathfrak m\) is maximal and is the only prime containing every \(y_i\).

Here is the ideal calculation. Setting all \(y_i\) to zero gives a retraction \(R\to k[x]\) of the natural inclusion \(k[x]\to R\), so that inclusion is injective. For \(I=(x)\subset R\), the relations \(xy_i=0\) give \(xr=xr(x,0,0,\ldots)\) for every \(r\in R\). Therefore \[
I=xk[x],\qquad I^2=x^2k[x],\qquad I/I^2=k\,\bar x\simeq A/\mathfrak m.
\] In the last identification every \(y_i\) annihilates \(\bar x\). For the zero-variable polynomial presentation \(R\twoheadrightarrow A\), the conormal complex is \([I/I^2\to0]\). The presentation comparison in Lemma \hyperref[AG-RG-S05-algebra-lemma-NL-homotopy]{F.11} identifies it with the naive cotangent complex used in Definition \hyperref[AG-RG-S05-algebra-definition-smooth]{F.6}.

First suppose \(y_i\notin\mathfrak q\). Inverting \(y_i\) in \(R\) makes \(x=0\), hence \(A_{y_i}=R_{y_i}\). As an \(R\)-algebra this is \[
R[t]/(h),\qquad h=y_it-1.
\] The element \(h\) is a nonzerodivisor over this arbitrary base ring. Indeed, if \(h\sum_{n=0}^Nb_nt^n=0\), its constant coefficient gives \(b_0=0\), and its successive coefficients \(y_ib_{n-1}-b_n=0\) give every \(b_n=0\). Consequently \((h)/(h^2)\) is free on \(\bar h\) over \(R[t]/(h)\): multiplication by \(h\) gives a surjection onto it, and cancellation of the nonzerodivisor \(h\) shows that its kernel is exactly \((h)\). An \(R\)-derivation of \(R[t]\) is determined by its value on \(t\); the product rule gives \(d(t^n)=nt^{n-1}dt\). Thus \(\Omega_{R[t]/R}=R[t]dt\), and this presentation has conormal complex \[
A_{y_i}\,\bar h\xrightarrow{\ y_i\ }A_{y_i}\,dt.
\] Its differential is an isomorphism, since \(y_i\) is a unit. Both first homology and the differential module are zero, and the zero module is finite projective. The displayed presentation is finite, so Definition \hyperref[AG-RG-S05-algebra-definition-smooth]{F.6} proves \(R\to A_{y_i}\) smooth. Definition \hyperref[AG-RG-S05-algebra-definition-smooth-at-prime]{F.7} proves smoothness at every point of \(D(y_i)\).

We must also exclude the exceptional point \(\mathfrak m\) from every smooth neighbourhood. Let \(g\in A\setminus\mathfrak m\). Its constant coefficient \(c=g(0,0,\ldots)\) is nonzero. Regard its polynomial expression as a lift \(\widetilde g\in R\). A presentation of \(A_g\) over \(R\) is \[
P=R[t]\twoheadrightarrow A_g,\qquad
J=\ker(P\to A_g)=(x,\widetilde g t-1).
\] The class of \(x\) in \(J/J^2\) has relative differential zero, because \(x\) belongs to the base ring \(R\). This class is nonzero. To see this, put \(C=(R/(x^2))_{\widetilde g}\). The \(R\)-algebra map \(P\to C\) sending \(t\) to \(\widetilde g^{-1}\) maps \(J\) into the square-zero ideal \((x)C\), and maps \(J^2\) to zero. It therefore maps the class of \(x\) in \(J/J^2\) to \(x\in C\). The assignments \(x\mapsto\epsilon\) and \(y_i\mapsto0\) define \[
R/(x^2)\longrightarrow k[\epsilon]/(\epsilon^2).
\] They respect all the relations, and send \(\widetilde g\) to the unit \(c\). The map extends to \(C\) and sends \(x\) to the nonzero element \(\epsilon\); its nonzeroness follows directly from the degree bound for a multiple of \(\epsilon^2\). Hence \(x\ne0\) in \(C\), and the class of \(x\) is a nonzero kernel element in the conormal complex \([J/J^2\to A_gdt]\). By the same presentation comparison, \(H_1(\operatorname{NL}_{A_g/R})\ne0\). Definition \hyperref[AG-RG-S05-algebra-definition-smooth]{F.6} shows that \(R\to A_g\) is not smooth. This holds for every \(g\notin\mathfrak m\), so Definition \hyperref[AG-RG-S05-algebra-definition-smooth-at-prime]{F.7} shows that \(R\to A\) is not smooth at \(\mathfrak m\).

Every other prime omits some \(y_i\), proving that the smooth locus is exactly the displayed open \(U\). In the notation of Definition \hyperref[AG-RG-S05-smoothing-definition-singular-ideal]{D.1}, this also gives \(H_{A/R}=\mathfrak m\).

Finally the open cover \(U=\bigcup_{i\geq1}D(y_i)\) has no finite subcover. Given a finite set \(F\) of indices, choose \(j\notin F\). Evaluation \(y_j\mapsto1\) and \(y_i\mapsto0\) for \(i\ne j\) defines a surjection \(A\to k\) with prime kernel \(\mathfrak q_j\). This point lies in \(D(y_j)\subset U\) but in none of the \(D(y_i)\) with \(i\in F\). Thus no finite subfamily covers \(U\), and \(U\) is not quasi-compact. \(\square\)

\phantomsection\label{AG-RG-S05-smoothing-lemma-strictly-standard-base-change}

\subsubsection{Lemma D.7: strictly standard base change}\label{AG-RG-S05-lemma-d.7-strictly-standard-base-change}

Reading source: \href{https://stacks.math.columbia.edu/tag/07CC}{Stacks, Tag 07CC}.

Let \(R\to A\) be finitely presented and \(R\to R'\) any ring map. An elementary standard element \(a\) remains elementary standard as \(a\otimes1\) in \(A\otimes_R R'\). The same assertion holds for strictly standard elements. No flatness assumption on \(R'\) is needed.

\textbf{Proof.} Take the finite presentation and witness from D.3 and set \(P'=R'[x_1,\ldots,x_n]\). The polynomial base-change map \(P\otimes_R R'\to P'\) sends \(x^E\otimes r'\) to \(r'x^E\). Its inverse sends the coefficient of every monomial to the same tensor; the polynomial monomials are the free \(R\)-basis of \(P\), so these maps are inverse. Tensor right exactness identifies \[A\otimes_R R'=P'/I',\qquad I'=(f'_1,\ldots,f'_m),\] where \(f'_j\) is obtained by changing coefficients along \(R\to R'\).

For a monomial \(r x^E\), differentiation multiplies its coefficient by the integer \(E_i\) and reduces the \(i\)th exponent by one. Changing coefficients commutes with this operation, including when that integer becomes zero in \(R'\). Consequently every selected Jacobian entry and its determinant is the image of the old one. The elementary single-minor identity, or the strict sum-of-minors identity, maps to exactly the required new identity with coefficients \(a'\otimes1\) or \(a_U\otimes1\).

For the relation condition, choose a polynomial lift \(\widetilde a\) and write each old relation with its finitely many coefficients in \((f_1,\ldots,f_c)+I^2\). Its image expresses \(\widetilde a' f'_{c+j}\) in \((f'_1,\ldots,f'_c)+(I')^2\), because products of two elements of \(I\) map to products of two elements of \(I'\). These are exactly D.3e2 or D.3s2 for the new witness. No injectivity of a tensor map was used. Empty minors, empty relation lists and a zero base-changed algebra satisfy the same identities. ∎

\phantomsection\label{AG-RG-S05-smoothing-lemma-final-solve}

\subsubsection{Lemma D.8: final solve}\label{AG-RG-S05-lemma-d.8-final-solve}

Reading source: \href{https://stacks.math.columbia.edu/tag/07EU}{Stacks, Tag 07EU}.

Suppose \(R\to A\xrightarrow{\rho}\Lambda\), with \(A\) finitely presented over \(R\) and \(H_{A/R}\Lambda=\Lambda\). Then \(\rho\) factors through an \(R\)-smooth algebra: \[A\longrightarrow B\longrightarrow\Lambda.\]

\textbf{Proof.} If \(\Lambda=0\), choose \(B=0\). Its presentation \(R/(1)\) is finite and its cotangent homology and differential module are zero, so F.6 makes it smooth. Assume henceforth \(\Lambda\ne0\).

Let \(J_s\) be the ideal generated by strictly standard elements. D.5 gives \(H_{A/R}=\sqrt{J_s}\). We first check that extending the smaller ideal also gives the unit ideal. Choose \(h_1,\ldots,h_t\in H_{A/R}\) and \(\lambda_i\in\Lambda\) with \(\sum_i\rho(h_i)\lambda_i=1\). For each \(i\) some \(n_i\geq1\) has \(h_i^{n_i}\in J_s\). Raise that identity to \[N=1+\sum_i(n_i-1).\] Every term in its distributive expansion uses some \(h_i\) at least \(n_i\) times: otherwise its total degree would be at most \(N-1\). Each term therefore belongs to \(J_s\Lambda\). Thus \(1\in J_s\Lambda\).

Expanding the finite expressions in this ideal gives strictly standard elements \(a_1,\ldots,a_r\) and \(\mu_i\in\Lambda\) with \(\sum_i\rho(a_i)\mu_i=1\). Here \(r>0\), since the target is nonzero. Define \[B=A[T_1,\ldots,T_r]/\left(\sum_i a_iT_i-1\right).\] The map sending \(T_i\) to \(\mu_i\) gives \(B\to\Lambda\) and extends the prescribed \(\rho\). Appending these finite variables and this one relation to a finite presentation of \(A\) proves that \(B/R\) is finitely presented.

In \(B\), the elements \(a_i\) generate one. Their principal opens therefore cover its spectrum: a prime containing all of them would contain the displayed unit. For each \(i\) there is the explicit isomorphism \[B_{a_i}\cong A_{a_i}[T_j:j\ne i].\] The map from right to left includes the remaining variables; its inverse substitutes \(T_i=a_i^{-1}(1-\sum_{j\ne i}a_jT_j)\). Substitution kills the relation and both composites fix their generators. D.5 proves \(A_{a_i}/R\) smooth. A polynomial algebra is standard smooth using its empty equation list, by \hyperref[AG-RG-S05-algebra-lemma-standard-smooth]{F.85}, and smooth composition \hyperref[AG-RG-S05-algebra-lemma-compose-smooth]{F.29} therefore makes \(B_{a_i}/R\) smooth. The proved finite principal-cover criterion \hyperref[AG-RG-S05-algebra-lemma-locally-smooth]{F.67} now makes \(B/R\) smooth. Zero localizations create empty charts and are already included in those criteria. This proves the required factorization over an arbitrary base ring. ∎

\subsection{Presentations of algebras}\label{AG-RG-S05-presentations-of-algebras}

\phantomsection\label{AG-RG-S05-smoothing-section-presentations}

Some of the results in this section are due to Elkik. Note that the algebra \(C\) in the following lemma is a symmetric algebra over \(A\). Moreover, if \(R\) is Noetherian, then \(C\) is of finite presentation over \(R\).

\phantomsection\label{AG-RG-S05-smoothing-lemma-improve-presentation}

\subsubsection{Lemma D.9: improve presentation}\label{AG-RG-S05-lemma-d.9-improve-presentation}

Reading source: \href{https://stacks.math.columbia.edu/tag/07CE}{Stacks, Tag 07CE}.

For a finitely presented algebra \(A\) over an arbitrary ring \(R\), there is an \(R\)-algebra map \(A\to C\) of finite type with a retraction \(C\to A\) such that:

\begin{enumerate}
\def\labelenumi{\arabic{enumi}.}
\tightlist
\item
  Whenever \(R\to A_a\) is a local complete intersection, \(A_a\to C_a\) is smooth and \(C_a\) has a finite polynomial presentation over \(R\) with free conormal module.
\item
  Whenever \(A_a/R\) is smooth, \(\Omega_{C_a/R}\) is free.
\end{enumerate}

\textbf{Proof.} Fix \(P=R[X_1,\ldots,X_n]\twoheadrightarrow A\) with \(I=(f_1,\ldots,f_m)\) and let \[M=I/I^2,\qquad K=\ker(A^m\longrightarrow M),\qquad C=\operatorname{Sym}_A(M).\] The map \(A^m\to M\) sends \(e_j\) to \([f_j]\). The symmetric algebra and its degree-zero augmentation are constructed in \hyperref[AG-RG-S05-more-algebra-lemma-symmetric-algebra-smooth]{F.128}; their composite on \(A\) is identity. Its universal property identifies \(C\) with the quotient of \(P[Y_1,\ldots,Y_m]\) by an ideal \(J\) generated by \(I\) and the linear forms whose coefficient classes belong to \(K\). In fact maps out of this quotient are exactly maps \(A\to D\) together with linear maps \(A^m\to D\) killing \(K\), hence linear maps \(M\to D\). The resulting maps are inverse because they fix \(A\) and every \(Y_j\). This proves the kernel description and finite type even when \(K\) is not finite.

Fix \(a\in A\) and write \(B=A_a\). Suppose first \(R\to B\) is a local complete intersection. The arbitrary-ring presentation comparison \hyperref[AG-RG-S05-more-algebra-lemma-lci-presentation-independent]{F.94.A} supplies local Koszul-regular equation lists for the inverse-variable polynomial presentation of \(B\). Such a list gives a conormal basis: correct a relation modulo the square of its ideal by coefficients in the ideal; the resulting exact relation is a first Koszul cycle, hence a boundary, whose coefficients belong to the ideal. This proves independence as well as generation of its equation classes. \hyperref[AG-RG-S05-algebra-lemma-principal-localization-NL]{F.72} identifies this locally free conormal module with \(M_a\oplus B\). Its summand \(M_a\) is locally finite flat and finitely presented, since both summands are generated by projections of a finite basis. The local finite-flat proof \href{https://kokunoyumeto.github.io/open-math-courses/courses/AG-RG/AG-RG-S01.html}{S01, Lemma 1.2} therefore makes \(M_a\) free at each local ring; its finite presentation permits extending that basis to a principal neighbourhood. Thus \(M_a\) is locally finite free.

Here are the finiteness and global splitting steps. Choose finitely many principal charts \(D(b_i)\) covering \(\operatorname{Spec}B\) on which \(M_a\) is free. A finite subcover exists because the covering elements generate the unit ideal, witnessed by a finite expression for \(1\). On each chart the surjection \(B_{b_i}^m\to(M_a)_{b_i}\) splits and its kernel \((K_a)_{b_i}\) is a finite direct summand. Lift its finite generating list to fractions of elements of \(K_a\) and collect their numerators in a finite submodule \(K_0\subset K_a\). Then \((K_a/K_0)_{b_i}=0\) for every \(i\). Any module element vanishing on these charts is zero: a power of each \(b_i\) kills it, and these powers generate \(1\), since a prime containing them would contain all the \(b_i\). Hence \(K_0=K_a\), so \(M_a\) is finitely presented.

For each local splitting \(s_i:M_{ab_i}\to B_{b_i}^m\), finite presentation identifies \[\operatorname{Hom}_B(M_a,B^m)_{b_i}=\operatorname{Hom}_{B_{b_i}}(M_{ab_i},B_{b_i}^m).\] To verify this identity, present \(M_a\) by finite free modules; both Hom modules are the kernels of the same finite matrix map after exact localization. Clear denominators of \(s_i\), then multiply by a further power of \(b_i\) to kill the error on a finite generating list. This gives \(t_i:M_a\to B^m\) with \(\pi t_i=b_i^{e_i}1_{M_a}\). Choose \(r_i\) satisfying \(\sum_i r_i b_i^{e_i}=1\). The map \(\sum_i r_it_i\) splits \(\pi\). Consequently \[M_a\oplus K_a\cong B^m,\] and both summands are finite projective. This argument also handles \(B=0\) directly by its zero modules.

The identities in F.128 give \(C_a=\operatorname{Sym}_B(M_a)\), which is smooth and finitely presented over \(B\). The inverse-variable presentation makes \(B/R\) finitely presented, and finite substitution \hyperref[AG-RG-S05-algebra-lemma-compose-finite-type]{F.28} therefore makes \(C_a/R\) finitely presented as well. The conormal calculation \hyperref[AG-RG-S05-more-algebra-lemma-cotangent-complex-symmetric-algebra]{F.102} identifies the kernel conormal of \(B[Y]\to C_a\) with \(K_a\otimes_B C_a\). On principal charts where the split modules are free, this kernel is generated by a subset of polynomial coordinates. Those coordinates form a regular sequence: multiplication by each new variable is injective by polynomial coefficient comparison, even over an arbitrary ring. The regular-list cone calculation in \hyperref[AG-RG-S05-more-algebra-lemma-noetherian-finite-all-equivalent]{F.118} gives their Koszul exactness in positive degrees; that direction requires no Noetherian hypothesis. Thus \hyperref[AG-RG-S05-more-algebra-lemma-transitive-lci-at-end]{F.130} supplies the actual left injection in the conormal transitivity sequence.

If \(h\in P\) lifts \(a\), introduce an inverse variable \(X_0\) and use the equation ideals \(I'=(I,hX_0-1)\) and \(J'=(J,hX_0-1)\). The explicit localization calculation \hyperref[AG-RG-S05-algebra-lemma-principal-localization-NL]{F.72} and the conormal sequence just justified give \[0\longrightarrow C_a\oplus(M_a\otimes_B C_a)\longrightarrow J'/(J')^2\longrightarrow K_a\otimes_B C_a\longrightarrow0.\] The last module is projective because a split finite free presentation stays split after tensoring. Hence \[J'/(J')^2\cong C_a\oplus((M_a\oplus K_a)\otimes_B C_a)\cong C_a^{m+1}.\] Together with the explicit polynomial quotient for \(C_a\) this proves (1). This conormal calculation uses only finite projectivity of \(M_a\) and \(K_a\), so it also applies when those properties are supplied by the smooth branch below.

For (2), suppose \(B/R\) is smooth, without first assuming (1). \hyperref[AG-RG-S05-algebra-definition-smooth]{F.6}, the full presentation comparison \hyperref[AG-RG-S05-algebra-lemma-NL-homotopy]{F.11} and localization \hyperref[AG-RG-S05-algebra-lemma-localize-NL]{F.65} give the split exact conormal sequence for the original \(n\) variables: \[0\longrightarrow M_a\longrightarrow B^n\longrightarrow\Omega_{B/R}\longrightarrow0.\] Thus \(M_a\) is finite projective and \(M_a\oplus\Omega_{B/R}\cong B^n\). F.128 gives smoothness of \(C_a/B\); its derivation universal property also identifies \(\Omega_{C_a/B}=M_a\otimes_B C_a\), since a derivation from the symmetric algebra is exactly a linear map on its degree-one generators. The full cotangent transitivity calculation \hyperref[AG-RG-S05-algebra-lemma-exact-sequence-NL]{F.42} now gives \[0\longrightarrow\Omega_{B/R}\otimes_B C_a\longrightarrow\Omega_{C_a/R}\longrightarrow M_a\otimes_B C_a\longrightarrow0.\] The right term is projective, so the sequence splits and its middle term is \(C_a^n\). This proves (2). Zero algebras and empty variable or equation lists obey the same identities. ∎

\phantomsection\label{AG-RG-S05-smoothing-proposition-lift-smooth}

\subsubsection{Proposition D.10: lift smooth}\label{AG-RG-S05-proposition-d.10-lift-smooth}

Reading source: \href{https://stacks.math.columbia.edu/tag/07M8}{Stacks, Tag 07M8}.

Let \(R\twoheadrightarrow R_0=R/I\) be a surjection of Noetherian rings. A syntomic \(R_0\)-algebra \(A_0\) has a syntomic lift \(A\) over \(R\), with a specified isomorphism \(A/IA\cong A_0\). If \(A_0/R_0\) is smooth, the lift can be chosen smooth. This is the auxiliary Noetherian lifting assertion used in desingularization.

\textbf{Proof.} If \(A_0=0\), take \(A=0\): its presentation \(R/(1)\) is finite, it is flat, its cotangent modules are zero, and its fibres are empty. Assume \(A_0\ne0\) and syntomic. Choose a finite polynomial presentation \(A_0=R_0[X]/J_0\) and put \(P_0=J_0/J_0^2\). \hyperref[AG-RG-S05-algebra-lemma-syntomic-presentation-ideal-mod-squares]{F.88} makes \(P_0\) finite projective. By \hyperref[AG-RG-S05-more-algebra-lemma-syntomic-lci]{F.129} the map is a local complete intersection, so \hyperref[AG-RG-S05-smoothing-lemma-improve-presentation]{D.9} constructs \[C_0=\operatorname{Sym}_{A_0}(P_0)\] with a free-conormal polynomial presentation and a smooth map \(A_0\to C_0\). Its finite presentation follows from \hyperref[AG-RG-S05-more-algebra-lemma-symmetric-algebra-smooth]{F.128} and finite substitution \hyperref[AG-RG-S05-algebra-lemma-compose-finite-type]{F.28}. The smooth/syntomic and composition proofs \hyperref[AG-RG-S05-algebra-lemma-smooth-syntomic]{F.83} and \hyperref[AG-RG-S05-algebra-lemma-composition-syntomic]{F.32} make \(C_0/R_0\) syntomic.

The complete equation-basis argument \hyperref[AG-RG-S05-algebra-lemma-huber]{F.55} yields \[C_0=R_0[Y_1,\ldots,Y_v]/(\bar f_1,\ldots,\bar f_c),\] with these classes a conormal basis. Flatness identifies that conormal after residue-field base change with the fibre conormal, by right exact tensor and the injection of the equation ideal into the polynomial ring. The fibres are locally complete intersections by syntomicness. \hyperref[AG-RG-S05-algebra-lemma-lci]{F.57} therefore makes their point-dimensions \(v-c\); applying this at their finitely many minimal primes gives dimension \(v-c\) for every nonempty fibre.

Lift each \(\bar f_j\) to the full polynomial ring \(R[Y_1,\ldots,Y_v]\) and set \(C^{\mathrm{pre}}=R[Y]/(f_1,\ldots,f_c)\). Its reduction is exactly \(C_0\). \hyperref[AG-RG-S05-algebra-lemma-localize-relative-complete-intersection]{F.66(1)} gives \(g\equiv1\pmod{IC^{\mathrm{pre}}}\) such that \(C=(C^{\mathrm{pre}})_g\) is a relative global complete intersection. Thus \(C/IC=C_0\), with no localization of the special fibre, and \hyperref[AG-RG-S05-algebra-lemma-relative-global-complete-intersection]{F.76} makes \(C/R\) syntomic.

Choose a complementary finite projective module \(Q_0\) with \(P_0\oplus Q_0=A_0^N\). Apply the complete lifting proof \hyperref[AG-RG-S05-more-algebra-lemma-lift-projective-module]{F.117} to the finite projective \(C_0\)-module \(Q_0\otimes_{A_0}C_0\). It gives a finitely presented étale \(C\to C'\) inducing \(C'/IC'\cong C_0\) and a finite projective \(C'\)-module \(Q\) reducing to that specified module. Put \(D=\operatorname{Sym}_{C'}(Q)\). It is smooth over \(C'\) by F.128, and the proved composition statements make \(D/R\) syntomic: étale maps are smooth by their zero cotangent complex and finite presentation, smooth maps are syntomic, and the base rings here are Noetherian finite-type algebras. The actual symmetric-algebra identities from F.128 give \[\begin{aligned}
D/ID&\cong\operatorname{Sym}_{C_0}(Q_0\otimes_{A_0}C_0)\\
&\cong\operatorname{Sym}_{A_0}(Q_0)\otimes_{A_0}\operatorname{Sym}_{A_0}(P_0)\\
&\cong\operatorname{Sym}_{A_0}(Q_0\oplus P_0)\cong A_0[Z_1,\ldots,Z_N].
\end{aligned}\] Choose \(z_i\in D\) lifting these variables and set \(E=D/(z_1,\ldots,z_N)\). Reduction of this quotient gives the specified \(E/IE=A_0\).

We verify the claimed syntomic neighbourhood of every \(\mathfrak q\in V(IE)\) rather than assuming it. Let \(\mathfrak t\) be its preimage in \(D\) and \(\mathfrak p=\mathfrak q\cap R\). Then \(I\subset\mathfrak p\). By \hyperref[AG-RG-S05-algebra-lemma-syntomic]{F.87} choose \(u\notin\mathfrak t\) and a relative global complete-intersection presentation \[D_u=R[W_1,\ldots,W_m]/(a_1,\ldots,a_s).\] In the fibre over \(k=\kappa(\mathfrak p)\) this is the localization of \((A_0\otimes_{R_0}k)[Z_1,\ldots,Z_N]\) at the image of \(u\). Its ordered \(Z_i\) are a regular sequence at the prime in question: multiplication by a polynomial variable is injective over any coefficient ring, successive quotients remove one variable, and exact localization preserves each injection. Their ideal stays proper at this prime because it comes from \(\mathfrak q\) of the quotient.

Lift each \(z_i\) to a polynomial \(b_i\) in \(R[W]\). By \hyperref[AG-RG-S05-algebra-lemma-relative-global-complete-intersection-conormal]{F.77}, the fibre \(\bar a_j\) form a regular sequence in the corresponding polynomial local ring. The just proved regularity of \(Z_i\) in its quotient makes \((\bar a_1,\ldots,\bar a_s,\bar b_1,\ldots,\bar b_N)\) a regular sequence there. \hyperref[AG-RG-S05-algebra-lemma-lci]{F.57} gives fibre point-dimension \(m-s-N\) for \[E_{\bar u}=R[W]/(a_1,\ldots,a_s,b_1,\ldots,b_N).\] Apply \hyperref[AG-RG-S05-algebra-lemma-localize-relative-complete-intersection]{F.66(3)} to this actual \(s+N\)-equation presentation. It supplies a relative global complete-intersection neighbourhood of \(\mathfrak q\), hence a syntomic one by F.76. A further localization of \(E_{\bar u}\) is a principal localization of \(E\): write its inverted element as \(e/\bar u^r\) and invert \(\bar u e\). This proves the claim at every point of \(V(IE)\).

Let \(W\) be the union of these syntomic neighbourhoods and write its closed complement as \(V(L)\) in \(\operatorname{Spec}E\). Since \(V(L)\cap V(IE)=\varnothing\), the ideal \(L+IE\) is all of \(E\); otherwise a containing maximal ideal would give a point of that intersection. Choose \(h\in L\) with \(h\equiv1\pmod{IE}\). Then \(D(h)\subset W\), and \(E_h/IE_h=A_0\) exactly. The finite principal-cover criterion \hyperref[AG-RG-S05-algebra-lemma-local-syntomic]{F.62} makes \(A=E_h\) syntomic over \(R\): its spectrum has a finite subcover by the already constructed syntomic charts, using the finite expression for \(1\) in their covering ideal. This proves the first assertion, including \(N=0\).

Now assume \(A_0/R_0\) smooth. For every prime \(\mathfrak q\in V(IA)\) the fibre of \(A/R\) is the same as the fibre of \(A_0/R_0\) over its contracted base prime. Those fibres are smooth by the following calculation, which does not assume that the residue-field map is flat. \hyperref[AG-RG-S05-algebra-lemma-smooth-syntomic]{F.83} covers \(A_0\) by standard smooth principal charts; choose finitely many by their unit-ideal cover. Tensor right exactness changes each polynomial quotient to the quotient by the coefficient-changed equations. Monomial differentiation commutes with changing coefficients, so its chosen Jacobian determinant is the image of the old unit and is still a unit. \hyperref[AG-RG-S05-algebra-lemma-standard-smooth]{F.85} makes every fibre chart smooth. The finite expression for \(1\) in the covering ideal remains such an expression after base change, so these charts cover the fibre; \hyperref[AG-RG-S05-algebra-lemma-locally-smooth]{F.67} makes the whole fibre smooth. Since \(A/R\) is finitely presented and flat, \hyperref[AG-RG-S05-algebra-lemma-flat-fibre-smooth]{F.45} gives a smooth neighbourhood of each such \(\mathfrak q\). Write the complement of their union as \(V(L')\). Exactly as above choose \(h'\in L'\) with \(h'\equiv1\pmod{IA}\). Then \(A_{h'}/IA_{h'}=A_0\), while F.67 makes \(A_{h'}/R\) smooth. No hypothesis that \(I\) is nilpotent or lies in the Jacobson radical is needed for either shrinking step. ∎

For Noetherian \(R\), \hyperref[AG-RG-S05-algebra-lemma-syntomic]{F.87} gives relative global complete-intersection neighbourhoods of a syntomic map. The next lemma proves the required arbitrary-base extension and constructs a symmetric-algebra bundle over \(\operatorname{Spec}(A)\) with a retraction whose total space is a relative global complete intersection over \(R\).

\phantomsection\label{AG-RG-S05-smoothing-lemma-syntomic-complete-intersection}

\subsubsection{Lemma D.11: syntomic complete intersection}\label{AG-RG-S05-lemma-d.11-syntomic-complete-intersection}

Reading source: \href{https://stacks.math.columbia.edu/tag/07CG}{Stacks, Tag 07CG}.

For a syntomic map \(R\to A\) over any ring \(R\), there is a smooth \(A\)-algebra \(C\) and a retraction \(C\to A\) such that \[C=R[x_1,\ldots,x_n]/(f_1,\ldots,f_c)\] is flat over \(R\) and every nonempty fibre has dimension \(n-c\).

\textbf{Proof.} We first prove the arbitrary-base local complete-intersection assertion needed to apply \hyperref[AG-RG-S05-smoothing-lemma-improve-presentation]{D.9}. Fix a finite presentation \(A=R[X_1,\ldots,X_N]/(u_1,\ldots,u_s)\) and a prime \(\mathfrak q\) of \(A\) over \(\mathfrak p\) of \(R\). In the polynomial local ring of its \(k=\kappa(\mathfrak p)\) fibre, syntomicness and \hyperref[AG-RG-S05-algebra-lemma-lci-at-prime]{F.58} give a complete-intersection equation ideal. Its conormal is free by \hyperref[AG-RG-S05-algebra-lemma-lci]{F.57}. Select from the \(u_i\) a list whose conormal residue classes are a basis. Nakayama and the unit determinant of its coordinate matrix make its local conormal classes a basis; the final assertion of F.57 makes that selected list generate the fibre ideal regularly.

Model the local map \(R_{\mathfrak p}\to A_{\mathfrak q}\) using \hyperref[AG-RG-S05-algebra-lemma-limit-essentially-finite-presentation]{F.59}. Explicitly, for finite-type \(\mathbf Z\)-subalgebras \(R_0\subset R\) containing all equation coefficients, put \(T_0=R_0[X]/(u_i)\) and contract the chosen primes to \(\mathfrak p_0,\mathfrak q_0\). The stages are \(R_{0,\mathfrak p_0}\to(T_0)_{\mathfrak q_0}\). Adding numerators and denominators of any finite set of fractions proves cofinality in the stated local-model system. The models are Noetherian, their transitions are exactly localized base changes, and their colimits are the given local rings. Since \(A_{\mathfrak q}\) is flat over \(R_{\mathfrak p}\) and is a finite module over itself, \hyperref[AG-RG-S05-algebra-lemma-colimit-eventually-flat]{F.21} gives a stage where \((T_0)_{\mathfrak q_0}\) is flat over \(R_{0,\mathfrak p_0}\).

We verify its closed-fibre complete-intersection property. Let \(k_0=\kappa(\mathfrak p_0)\). There is an embedding \(k_0\subset k\) by contraction of primes and the fraction-field construction. The map from the polynomial fibre local ring \(V_0\) at the contracted prime to the final polynomial fibre local ring \(V\) is local and flat, being field extension followed by localization. The selected equations generate the final fibre ideal. The finite \(V_0\)-module consisting of the model ideal modulo their span becomes zero after tensoring with \(V\), so it is zero: a finite module with zero tensor along a local flat map has zero residue-field tensor, and Nakayama kills it. Here the target residue field extends the source residue field, so that residue-field tensor is faithful on vector spaces.

The same argument applies to each positive homology module of the selected finite free Koszul complex. These modules are finite because \(V_0\) is Noetherian. Flat tensor commutes with their kernels and images, so their tensors are zero, since the final list is regular. Nakayama again gives zero homology in the model. The list lies in the maximal ideal, and \hyperref[AG-RG-S05-more-algebra-lemma-noetherian-finite-all-equivalent]{F.118} makes it regular in \(V_0\). Thus the model has precisely the closed fibre required by the Noetherian-base theorem \hyperref[AG-RG-S05-algebra-lemma-syntomic]{F.87}.

Apply that theorem to obtain \(g_0\notin\mathfrak q_0\) with \(B_0=(T_0)_{g_0}\) a relative global complete intersection over \(R_0\). Its scalar extension is exactly \(A_g\), with \(g\notin\mathfrak q\). Write \(B_0=P_0/(b_1,\ldots,b_t)\) using its finite polynomial complete-intersection presentation. By \hyperref[AG-RG-S05-more-algebra-lemma-relative-global-complete-intersection-koszul]{F.126} its augmented Koszul complex is exact, and by \hyperref[AG-RG-S05-algebra-lemma-relative-global-complete-intersection]{F.76} its last quotient \(B_0\) is \(R_0\)-flat. Every complex term is also \(R_0\)-flat, being finite free over the polynomial ring.

All successive kernels in this augmented complex are \(R_0\)-flat. For \(0\to Z\to F\to Q\to0\) with \(F,Q\) flat, quotient flatness makes tensor preserve the left injection. For any inclusion of modules, the resulting map after tensoring with \(Z\) is injective by embedding its two tensor products into those for \(F\). Thus \(Z\) is flat. Induct upwards from the flat augmentation quotient. Tensoring each short exact sequence with \(R\) now remains exact. The actual base-changed equation list for \(A_g\) is therefore Koszul-regular. This is the needed arbitrary-base assertion; it does not apply F.87 directly to a non-Noetherian base. The comparison \hyperref[AG-RG-S05-more-algebra-lemma-lci-presentation-independent]{F.94.A} transports these presentations to the fixed finite polynomial presentation of \(A\) at every prime, proving that \(R\to A\) is a local complete intersection.

Let \(I=(u_i)\) and \(M=I/I^2\). D.9 with \(a=1\) constructs \(C=\operatorname{Sym}_A(M)\) with its degree-zero augmentation \(C\to A\), proves \(M\) finite projective and \(A\to C\) smooth, and supplies a polynomial presentation over \(R\) with free conormal. The finite-presentation proof \hyperref[AG-RG-S05-more-algebra-lemma-symmetric-algebra-smooth]{F.128} and finite substitution \hyperref[AG-RG-S05-algebra-lemma-compose-finite-type]{F.28} make \(C/R\) finitely presented. The complete equation-basis construction \hyperref[AG-RG-S05-algebra-lemma-huber]{F.55} gives a presentation \(C=R[x_1,\ldots,x_n]/J\) with \(J=(f_1,\ldots,f_c)\) and those classes a basis of \(J/J^2\).

On a finite principal cover of \(\operatorname{Spec}A\) trivializing \(M\), the symmetric-algebra identities in F.128 make \(C\) a polynomial algebra over each chart. Those algebras are \(A\)-flat. Tensoring an injection with \(C\) and detecting its kernel on the cover proves \(A\)-flatness of \(C\) globally. The given flatness of \(A/R\) consequently makes \(C/R\) flat.

For a base residue field \(k\), the same identities give \(C_k=\operatorname{Sym}_{A_k}(M\otimes_A A_k)\). The original syntomic hypothesis makes \(A_k\) locally a complete intersection. Trivializing the module and restricting to complete-intersection charts makes \(C_k\) a polynomial algebra over such a chart. Its regular equations remain regular after polynomial extension and localization, because both are flat; the ambient polynomial local rings are regular. F.58 therefore proves that \(C_k\) is locally a complete intersection.

Flatness of \(C/R\) makes tensoring \(0\to J\to R[x]\to C\to0\) injective on the left. Hence \(J\otimes_R k\) is the fibre equation ideal \(\bar J\). Right exactness applied to \(J^2\to J\to J/J^2\to0\) identifies its conormal: the image of \(J^2\otimes_R k\) consists exactly of products of the fibre ideal generators. Thus \[\bar J/\bar J^2\cong(J/J^2)\otimes_R k\cong C_k^c.\] At a minimal prime of a nonzero \(C_k\), F.57 identifies this rank with the height of its polynomial preimage, since the quotient local ring has dimension zero. Each component therefore has dimension \(n-c\) by \href{https://kokunoyumeto.github.io/open-math-courses/courses/AG-CA/krull-dimension-and-noether-normalization.html}{AG-CA-09, Theorem 4.3}. The field algebra has finitely many components, so its dimension is their maximum, also \(n-c\). The augmentation fixes \(A\), and is the required retraction.

If \(A=0\), take \(C=0\) with presentation \(R/(1)\). It is flat, the identity \(A\to C\) is smooth, and all fibres are empty. Empty variable or regular lists use the same constructions. No condition has been added to the arbitrary base ring. ∎

\phantomsection\label{AG-RG-S05-smoothing-lemma-smooth-standard-smooth}

\subsubsection{Lemma D.12: smooth standard smooth}\label{AG-RG-S05-lemma-d.12-smooth-standard-smooth}

Reading source: \href{https://stacks.math.columbia.edu/tag/07CH}{Stacks, Tag 07CH}.

Every smooth algebra \(A\) over an arbitrary ring \(R\) is a retract of an \(R\)-algebra \(B\) that is standard smooth over \(R\), with the map \(A\to B\) itself smooth. Thus \(B\) has a finite polynomial presentation with an invertible square Jacobian minor on its first equation variables.

\textbf{Proof.} If \(A=0\), set \(B=0\). The presentation \(R[X]/(1)\) has Jacobian entry \(0\), which is a unit in the zero ring, and the identity is its smooth retraction. Suppose \(A\ne0\).

The smooth branch of \hyperref[AG-RG-S05-smoothing-lemma-improve-presentation]{D.9} constructs \(C=\operatorname{Sym}_A(I/I^2)\), with \(A\to C\) smooth and an augmentation \(C\to A\). In this branch the conormal sequence is split by F.6, so the same finite-projective conormal calculation in D.9 applies, independently of any syntomic implication over the base. Appending the inverse variable for \(1\) gives a finite polynomial presentation of \(C\) with free conormal. \hyperref[AG-RG-S05-algebra-lemma-huber]{F.55} replaces it by \[C=R[X_1,\ldots,X_n]/(f_1,\ldots,f_c)\] whose equation classes are a conormal basis. \hyperref[AG-RG-S05-algebra-lemma-compose-smooth]{F.29} makes \(C/R\) smooth. It is nonzero, because it retracts onto \(A\).

By \hyperref[AG-RG-S05-algebra-definition-smooth]{F.6} and the full comparison \hyperref[AG-RG-S05-algebra-lemma-NL-homotopy]{F.11}, the conormal differential in this presentation is a split injection. Write its matrix with equations as rows as \(J=(\partial f_j/\partial X_i)\), of size \(c\times n\). A left inverse of the transposed conormal map gives an \(n\times c\) matrix \(T\) over \(C\) such that \(JT=1_c\). Reduction at any maximal ideal of \(C\) gives \(c\le n\) by vector-space dimension. Lift its entries to polynomials \(\widetilde T_{i\ell}\).

Set \(E=C[Y_1,\ldots,Y_c]\) and introduce redundant coordinates \(Z_i\) by the equations \[h_i=X_i-\sum_{\ell=1}^c\widetilde T_{i\ell}(X)Y_\ell-Z_i.\] Eliminating \(Z_i\) gives inverse generator maps for \[E=R[X,Y,Z]/(f_1,\ldots,f_c,h_1,\ldots,h_n).\] Let \(G\) be the determinant of its square Jacobian on the selected variables \((X_1,\ldots,X_n,Y_1,\ldots,Y_c)\), using the displayed equation order. Modulo \(Y\) in \(C\) its matrix is \[\begin{pmatrix}J&0\\1_n&-T\end{pmatrix}.\] Interchange the two row blocks, which contributes \((-1)^{nc}\), then subtract \(J\) times the first block row from the second. The resulting matrix is block upper triangular with diagonal \(1_n,JT=1_c\). Its determinant is \(1\). Thus the image \(g\) of \(G\) in \(E\) satisfies \(g\bmod(Y)=(-1)^{nc}\).

Put \(B=E_g\). Evaluation at \(Y=0\) sends \(g\) to that unit and therefore extends the augmentation to \(B\to C\). Its composite with \(C\to B\) is identity; composing with \(C\to A\) gives the required retraction of \(A\to B\). Polynomial extension is smooth by the empty-equation case of \hyperref[AG-RG-S05-algebra-lemma-standard-smooth]{F.85}, and principal localization is standard smooth via its equation \(gU-1\) with invertible derivative \(g\). Hence F.85 and F.29 make \(A\to B\) smooth.

Finally adjoining the inverse of \(G\) writes \[B=R[X,Y,U,Z]/(f_1,\ldots,f_c,h_1,\ldots,h_n,UG-1).\] Order the selected variables as \((X,Y,U)\), before all the \(Z_i\). There are \(n+c+1\) equations and selected variables. The Jacobian has the earlier square block in its upper left, zeros above the last column, and \(G\) in its last diagonal position. Its determinant is therefore \(G^2\), a unit since \(UG=1\). Renaming the variables in this order gives the standard smooth presentation asserted. Empty blocks have determinant \(1\); the argument includes \(c=0\) and \(n=0\) over any base ring. ∎

\phantomsection\label{AG-RG-S05-smoothing-lemma-colimit-standard-smooth}

\subsubsection{Lemma D.13: colimit standard smooth}\label{AG-RG-S05-lemma-d.13-colimit-standard-smooth}

Reading source: \href{https://stacks.math.columbia.edu/tag/07CI}{Stacks, Tag 07CI}.

If an \(R\)-algebra \(\Lambda\) is a filtered colimit of smooth \(R\)-algebras, then it is a filtered colimit of standard smooth \(R\)-algebras, over an arbitrary base ring \(R\).

\textbf{Proof.} First consider a map \(A\to\Lambda\) with \(A\) finitely presented over \(R\). Choose representatives for its finite generator list at one common smooth stage. Each of its finitely many defining relations is zero at a later stage, by the equality rule for filtered colimits. A common successor therefore gives a factorization \(A\to B\to\Lambda\) with \(B/R\) smooth. By \hyperref[AG-RG-S05-smoothing-lemma-smooth-standard-smooth]{D.12}, \(B\) is a retract of a standard smooth \(R\)-algebra \(C\). The composite \(A\to B\to C\) followed by \(C\to B\to\Lambda\) is the original map. Thus every such finite-presentation map factors through a standard smooth algebra.

Take as objects finite polynomial presentations of standard smooth \(R\)-algebras together with maps to \(\Lambda\), and as arrows algebra maps respecting those maps. This is a small category: the finite coefficient lists form a set, and each map to \(\Lambda\) is determined by a finite tuple of elements. The presentation of \(R\) with empty variable and equation lists is an object. Two objects have a common successor by applying the factorization just proved to the map from their tensor product to \(\Lambda\); that tensor product has the finite presentation with their disjoint variables and both relation lists. To equalize parallel arrows \(u,v\colon B\rightrightarrows C\), choose generators \(b_1,\ldots,b_s\) of \(B\) and append \(u(b_i)-v(b_i)\) to the relations of \(C\). The resulting finitely presented quotient maps to \(\Lambda\); its factorization through a standard smooth object supplies an arrow out of \(C\) equalizing \(u,v\), since equality on generators implies equality on all polynomials in them. These constructions prove that the category is filtered.

Let \(L\) be its colimit. The canonical map \(L\to\Lambda\) is surjective: factor \(R[T]\to\Lambda\) with \(T\mapsto\lambda\) to represent any chosen \(\lambda\) at one of the objects. For injectivity, if \(b\in B\) has zero image, the finitely presented algebra \(B/(b)\) maps to \(\Lambda\) and factors through an object. The resulting successor kills \(b\), so its colimit class is zero. For two representatives with equal images, first move them to a common successor and use the same argument on their difference. Hence \(L\to\Lambda\) is an isomorphism and every stage is standard smooth. Empty lists and zero rings obey the same constructions. ∎

\phantomsection\label{AG-RG-S05-smoothing-lemma-standard-smooth-include-generators}

\subsubsection{Lemma D.14: standard smooth include generators}\label{AG-RG-S05-lemma-d.14-standard-smooth-include-generators}

Reading source: \href{https://stacks.math.columbia.edu/tag/07EY}{Stacks, Tag 07EY}.

Suppose \(A/R\) is standard smooth and \(E=\{a_1,\ldots,a_n\}\subset A\) is finite. There is a standard smooth presentation \[A=R[x_1,\ldots,x_{n+m}]/(f_1,\ldots,f_c),\qquad c\ge n,\] whose first \(n\) coordinate classes are the specified \(a_i\) in any chosen order, and whose Jacobian minor on its first \(c\) variables is a unit.

\textbf{Proof.} Start with \(A=R[y_1,\ldots,y_m]/(g_1,\ldots,g_d)\) and rename its variables so that the \(d\times d\) minor on \(y_1,\ldots,y_d\) is a unit. Choose polynomials \(h_i(y)\) mapping to \(a_i\). Adjoin variables \(x_i\) with equations \(x_i-h_i(y)\). Substitution defines an isomorphism \[R[x_1,\ldots,x_n,y_1,\ldots,y_m]/(x_i-h_i,g_j)\longrightarrow A.\] Its inverse sends the old \(y_j\) classes to themselves. The two composites fix all \(y_j\), and fix \(x_i\) because its new relation identifies it with \(h_i(y)\); they are therefore identity maps.

Order the equations by the \(n\) new graph equations followed by the \(d\) old equations. With equations as rows, the minor on \((x_1,\ldots,x_n,y_1,\ldots,y_d)\) is \[\begin{pmatrix}1_n&-(\partial h_i/\partial y_j)\\0&(\partial g_i/\partial y_j)\end{pmatrix}.\] In the signed permutation formula, each of its first \(n\) columns has only its diagonal entry nonzero. The surviving terms therefore give exactly the old \(d\times d\) determinant, which is a unit. Transposing yields the determinant convention with variables as rows and equations as columns. Now set \(c=n+d\) and rename \(y_j=x_{n+j}\), placing the unselected \(y\) variables last. Since \(d\le m\), this is the required minor on the first \(c\) variables. The first \(n\) classes are the prescribed elements. Empty blocks use determinant \(1\), so the proof includes an empty \(E\) or an empty old equation list. ∎

\phantomsection\label{AG-RG-S05-smoothing-lemma-compare-standard}

\subsubsection{Lemma D.15: compare standard}\label{AG-RG-S05-lemma-d.15-compare-standard}

\emph{Adapted from the Stacks project, smoothing.tex, lines 718--850.}
\nopagebreak[4]

Let \(R \to A\) be a ring map of finite presentation. Let \(a \in A\). Consider the following conditions on \(a\):

\begin{enumerate}
\def\labelenumi{\arabic{enumi}.}
\item
  \(A_a\) is smooth over \(R\),
\item
  \(A_a\) is smooth over \(R\) and \(\Omega_{A_a/R}\) is stably free,
\item
  \(A_a\) is smooth over \(R\) and \(\Omega_{A_a/R}\) is free,
\item
  \(A_a\) is standard smooth over \(R\),
\item
  \(a\) is strictly standard in \(A\) over \(R\),
\item
  \(a\) is elementary standard in \(A\) over \(R\).
\end{enumerate}

Then we have

\begin{enumerate}
\def\labelenumi{\arabic{enumi}.}
\item
  \textbf{(a)} (4) \(\Rightarrow\) (3) \(\Rightarrow\) (2) \(\Rightarrow\) (1),
\item
  \textbf{(b)} (6) \(\Rightarrow\) (5),
\item
  \textbf{(c)} (6) \(\Rightarrow\) (4),
\item
  \textbf{(d)} (5) \(\Rightarrow\) (2),
\item
  \textbf{(e)} (2) \(\Rightarrow\) the elements \(a^e\), \(e \geq e_0\) are strictly standard in \(A\) over \(R\),
\item
  \textbf{(f)} \(A_a\) is smooth over \(R\) and \(\Omega_{A_a/R}\) has a finite basis \(\mathrm db_1,\ldots,\mathrm db_d\) with \(b_i\in A\) if and only if every sufficiently large positive power \(a^e\) is elementary standard in \(A\) over \(R\). These conditions also follow if any one positive power of \(a\) is elementary standard.
\end{enumerate}

The requirement that the differential-basis representatives in (f) belong to \(A\) is essential. Standard smoothness of \(A_a\) alone does not imply the elementary-power conclusion over arbitrary base rings; the counterexample after the proof makes this distinction explicit.

\textbf{Proof of (a)--(e).} By \hyperref[AG-RG-S05-algebra-lemma-standard-smooth]{F.85}, a standard smooth algebra is smooth and its differentials are free. A free module is stably free by adding the zero free module; forgetting this condition gives (a). For (b), take the coefficients of the other minors to be zero in D.3s1; transposing the one square minor does not change its determinant. D.3e2 and D.3s2 are the same relation condition.

For (c), use an elementary presentation \(A=P/I\), where \(P=R[x_1,\ldots,x_n]\), \(I=(f_1,\ldots,f_m)\) and the first \(c\) equations have minor \(\Delta\). Lift \(a\) to \(h\in P\). Write \[B=A_a=P[x_0]/J,\qquad J=(x_0h-1,f_1,\ldots,f_m),\] and \(K=(x_0h-1,f_1,\ldots,f_c)\). The localization conormal calculation \hyperref[AG-RG-S05-algebra-lemma-principal-localization-NL]{F.72} and D.3e2 give \(J=K+J^2\): after inverting \(a\), the classes of the remaining \(f_j\) are combinations of the first \(c\) classes. The module \(J/K\) is finite, since \(J\) has the displayed finite generating list, and satisfies \(J/K=J(J/K)\). Choose generators \(v_i\) and write \(v_i=\sum_j b_{ij}v_j\) with \(b_{ij}\in J\). Multiplication by the adjugate of \((\delta_{ij}-b_{ij})\), whose identity is proved in \hyperref[AG-RG-S05-algebra-lemma-matrix-left-inverse]{F.70}, shows that \[g=\det(\delta_{ij}-b_{ij})\in1+J,\qquad gJ\subset K.\] For the zero module choose \(g=1\). Thus \(J_g=K_g\). Since \(g\) maps to \(1\) in \(B\), adjoining an inverse for it does not change \(B\), and \[B=P[x_0,z]/(x_0h-1,f_1,\ldots,f_c,gz-1).\] For these ordered equations select \((x_0,x_1,\ldots,x_c,z)\). The first \(c+1\) rows have zero last column. Their first \(c+1\) columns form the block matrix with upper left entry \(h\), lower left zero, and lower right Jacobian minor \(\Delta\). The last diagonal entry of the whole matrix is \(g\). Its determinant is therefore \(gh\Delta\). In \(B\), \(g=1\), \(h=a\) is a unit, and D.3e1 gives \(a=a'\Delta\), so \(\Delta^{-1}=a'/a\). This proves standard smoothness, including the empty minor case \(c=0\).

For (d), D.3s2 makes the first \(c\) equation classes generate \(M=(I/I^2)_a\). D.4 and D.3s1 give a map back from \(B^n\) whose composite with the map \(B^c\to M\to B^n\) is multiplication by the unit \(a\). Thus \(B^c\to M\) is both surjective and injective, and the conormal differential is a split injection. D.5 proves that \(B=A_a\) is smooth. The localized conormal sequence, by F.65 and F.72, now gives \[M\oplus\Omega_{B/R}\cong B^n,\qquad M\cong B^c.\] Consequently \(\Omega_{B/R}\oplus B^c\cong B^n\), exactly the required stable freeness.

For (e), start with a finite presentation \(A=R[X_1,\ldots,X_n]/I\) and \(B=A_a\). If \(B=0\), the equality \(1/1=0\) means \(a^L=0\) for some \(L\). For every \(e\ge L\), take \(c=0\) in any finite presentation: the empty minor is \(1\) and all multiplied relation classes are zero. Such powers are elementary, hence strictly standard. Assume \(B\ne0\). Smoothness, the conormal description of the naive cotangent complex, and F.65/F.72 give the split exact sequence \[0\to M=(I/I^2)_a\to B^n\to\Omega_{B/R}\to0.\] If \(\Omega_{B/R}\oplus B^s\cong B^t\), then \(M\oplus B^t\cong B^{n+s}\). Add \(t\) variables \(Y_i\) mapping to zero in \(A\). The kernel of this enlarged presentation is \(J=(I,Y_1,\ldots,Y_t)\), and \[J/J^2\cong I/I^2\oplus A^t.\] Indeed, evaluation at \(Y=0\) gives the first component, and the coefficients of the terms linear in \(Y\) modulo \(I\) give the second; their kernel consists precisely of products of two elements of \(J\). The differential sends the new generators to the independent \(\mathrm dY_i\). Thus \((J/J^2)_a\) is finite free and its differential is a split injection into \(B^{n+t}\).

Choose representatives \(f_1,\ldots,f_c\in J\) of a basis after localization: a basis element is a fraction of an element of \(J/J^2\), and multiplying each basis element by its own denominator, a unit in \(B\), still gives a basis. Extend them to a finite generating list \(f_1,\ldots,f_m\) of \(J\). Finiteness follows from \hyperref[AG-RG-S05-algebra-lemma-finite-presentation-independent]{F.44}, or directly from the old finite list and the new variables. Each remaining class is zero in the localization of \[Q=(J/J^2)/\langle f_1,\ldots,f_c\rangle.\] For each of its finitely many generators some power of \(a\) kills that class. Choose a common exponent \(L\). For a polynomial lift \(h\) of \(a\) this says \[h^L f_j\in(f_1,\ldots,f_c)+J^2\qquad(j>c).\] Multiplying by further powers of \(h\) gives D.3s2 for every \(a^e\), \(e\ge L\).

Let \(F\) be the \((n+t)\times c\) matrix of the differentials of the first \(c\) equations over \(A\). The split injection after localization has a matrix left inverse \(P\) over \(B\). Choose \(u\) and a matrix \(P_0\) over \(A\) with \((P_0)_a=a^uP\). Every entry of \(P_0F-a^u1_c\) becomes zero in \(B\), so a common further power \(a^v\) kills all those finitely many entries. Hence \[P_1F=a^{u+v}1_c,\qquad P_1=a^vP_0,\] as an exact equality over \(A\). Increase \(u+v\) to a positive integer \(k\) by multiplying \(P_1\) by an additional power of \(a\) if necessary. No power of \(a\) has been cancelled. F.70 puts \(a^{ck}\) in the ideal of the \(c\times c\) minors of \(F\). Reduction at a prime of the nonzero ring \(B\) shows \(c\le n+t\). For every \(e\ge\max(L,ck,1)\), multiply the minor coefficients by \(a^{e-ck}\) to obtain D.3s1, together with the already established D.3s2. When \(c=0\), the minor ideal is \(A\) and the same conclusion uses the empty determinant \(1\). This proves (e) over arbitrary rings, including zero divisors. ∎

\textbf{Proof of (f).} Put \(B=A_a\). If \(B=0\), some power \(a^L\) is zero. The zero differential module has the empty basis, and, as in (e), every \(a^e\) with \(e\ge\max(L,1)\) is elementary using \(c=0\) and the empty minor \(1\). Assume \(B\ne0\).

First suppose \(B\) is smooth and \(\mathrm db_1,\ldots,\mathrm db_d\) is a basis with \(b_i\in A\). Choose finite algebra generators \(z_1,\ldots,z_n\) of \(A\) over \(R\) and the surjection \[P=R[X_1,\ldots,X_n,Y_1,\ldots,Y_d]\longrightarrow A,
\qquad X_i\longmapsto z_i,\quad Y_j\longmapsto b_j.\] Its kernel \(J\) is finitely generated by \hyperref[AG-RG-S05-algebra-lemma-finite-presentation-independent]{F.44}. The smoothness definition \hyperref[AG-RG-S05-algebra-definition-smooth]{F.6}, the presentation description of the naive cotangent complex \hyperref[AG-RG-S05-algebra-lemma-NL-homotopy]{F.11}, and localization \hyperref[AG-RG-S05-algebra-lemma-localize-NL]{F.65} give the exact sequence \[0\longrightarrow M=(J/J^2)_a\xrightarrow{\mathrm d}
B^n\oplus B^d\longrightarrow\Omega_{B/R}\longrightarrow0.\] Here the two summands have coordinates \(\mathrm dX_i\) and \(\mathrm dY_j\). The map from the second summand to \(\Omega_{B/R}\) is an isomorphism, because its images are the chosen basis. Consequently projection onto the first summand restricts to an isomorphism \(M\to B^n\): for each \(v\in B^n\), there is a unique \(w\in B^d\) whose differential image cancels that of \(v\), so \((v,w)\) is in the image of \(M\); a vector of \(M\) with first component zero has second component zero as well. Injectivity of \(\mathrm d\) proves uniqueness in \(M\).

Let \(\eta_1,\ldots,\eta_n\) be the resulting basis of \(M\) whose projected differentials are the standard basis of \(B^n\). Write each \(\eta_i\) as a fraction of a class represented by \(f_i\in J\), with denominator \(a^{k_i}\). Thus the classes of \(f_i\) are \(a^{k_i}\eta_i\), still a basis over \(B\). Their square Jacobian on \(X_1,\ldots,X_n\) is diagonal after passage to \(B\), with entries \(a^{k_i}\). In particular the image of \[\delta=\det(\partial f_i/\partial X_j)_{i,j=1}^n\in A\] is a unit in \(B\). Extend \(f_1,\ldots,f_n\) to a finite generating list \(f_1,\ldots,f_m\) of \(J\). The finite module \[Q=(J/J^2)/\langle f_1,\ldots,f_n\rangle\] has zero localization at \(a\). A common power \(a^L\) kills its finite generating list. For a polynomial lift \(h\in P\) of \(a\), this means \[h^L f_j\in(f_1,\ldots,f_n)+J^2\qquad(j>n).\] All higher powers satisfy the same relation condition.

Since \(\delta\) is a unit in \(B\), choose \(u\in A\) and \(r\ge0\) with \(\delta(u/a^r)=1\) in \(B\). The equality in this localization supplies \(v\ge0\) such that \[a^v(\delta u-a^r)=0\quad\hbox{in }A,
\qquad a^{r+v}=\delta(a^v u).\] Increase the exponent to a positive \(K\) by multiplying this equality by an additional power of \(a\) if necessary. This gives \(a^K=\delta u_0\) in \(A\), without cancelling a power of a possible zero divisor. Every \(e\ge\max(K,L,1)\) then satisfies both D.3e1 and D.3e2 for the presentation \(A=P/J\), with \(c=n\) and the selected variables \(X_i\). This proves the forward direction, including \(n=0\), when \(\delta\) is the empty determinant \(1\).

Conversely, suppose \(a^r\) is elementary for some \(r\ge1\). Inverting \(a^r\) is the same as inverting \(a\), since \(a^{-1}=a^{r-1}(a^r)^{-1}\). Part (c), applied to \(a^r\), proves that \(B\) is standard smooth, hence smooth. Take its elementary presentation \(A=R[x_1,\ldots,x_N]/I\) with \(I=(g_1,\ldots,g_m)\) and selected minor on \(x_1,\ldots,x_c\). In \(B\), D.3e2 makes every remaining conormal class a combination of the first \(c\). Therefore \[\Omega_{B/R}=B^N/\langle\mathrm dg_1,\ldots,\mathrm dg_c\rangle.\] The selected \(c\times c\) Jacobian block is invertible. Subtracting a unique combination of these \(c\) differential rows eliminates the first \(c\) coordinates of any vector in \(B^N\); a combination with those coordinates zero has all its coefficients zero. Hence the images of \(\mathrm dx_{c+1},\ldots,\mathrm dx_N\) form a basis of this quotient. Their representatives are elements of \(A\). The forward direction now proves that every sufficiently large power is elementary. This proves (f). ∎

\textbf{Why standard smoothness alone is insufficient.} Take \[R=\mathbf Z,\qquad A=\mathbf Z[t^2,t^3]\subset\mathbf Z[t],\qquad a=t^2.\] This algebra has the finite presentation \[A\cong\mathbf Z[X,Y]/(Y^2-X^3),\qquad X\longmapsto t^2,\quad Y\longmapsto t^3.\] Indeed division by the monic polynomial \(Y^2-X^3\) writes every class uniquely as \(p(X)+Yq(X)\). After substitution, the first summand has even exponents and the second has odd exponents at least \(3\); their disjoint monomials prove injectivity. The image is precisely the indicated subalgebra.

Inverting \(a\) gives \[B=A_a=\mathbf Z[t,t^{-1}],\qquad t=Y/X,\quad t^{-1}=Y/X^2.\] These identities follow from \(Y^2=X^3\) and give both inclusions with the Laurent ring. The presentation \(B=\mathbf Z[T,U]/(TU-1)\) is standard smooth: select \(U\) for its one equation, whose derivative is the unit \(T\). Every derivation of \(B\) is uniquely determined by its value on \(t\), with \(\mathrm d(t^{-1})=-t^{-2}\mathrm dt\). Thus \(\Omega_{B/\mathbf Z}=B\,\mathrm dt\).

The units of \(\mathbf Z[t,t^{-1}]\) are exactly \(\pm t^k\). To check this, the lowest and highest nonzero exponents of a product add, since their coefficients are nonzero products in the domain \(\mathbf Z\). For a product equal to \(1\), each factor must therefore have width zero and be a monomial; its integer coefficient must be \(1\) or \(-1\).

Every \(b\in A\) has the form \(b_0+\sum_{k\ge2}b_k t^k\), with integer coefficients and a finite sum. Consequently \[\mathrm db=\left(\sum_{k\ge2}k b_k t^{k-1}\right)\mathrm dt.\] This coefficient cannot be a unit: if it were \(\pm t^j\), necessarily \(j\ge1\), and its sole coefficient would require \((j+1)b_{j+1}=\pm1\), impossible in \(\mathbf Z\). A basis of the free rank-one module \(B\,\mathrm dt\) must contain one generator; tensoring with \(\mathbf Q(t)\) verifies that its length is one. Hence it has no basis consisting of differentials of elements of \(A\). By the converse in (f), no positive power of \(a\) is elementary standard, despite standard smoothness of \(B\).

There is no conflict with (e). In the original cusp presentation the two one-by-one minors are \(-3X^2\) and \(2Y\), and in \(A\) \[X^3=(-X)(-3X^2)+(-Y)(2Y).\] The relation condition is empty because there is only one equation. Thus \(a^3\), and every \(a^e\) with \(e\ge3\) after multiplication by \(a^{e-3}\), is strictly standard. This exact example distinguishes the sum-of-minors condition from the single-minor condition over the same arbitrary base.

\subsection{The lifting problem}\label{AG-RG-S05-the-lifting-problem}

\phantomsection\label{AG-RG-S05-smoothing-section-lifting}

The goal in this section is to prove (Proposition \hyperref[AG-RG-S05-smoothing-proposition-lift]{D.18}) that the collection of algebras which are filtered colimits of smooth algebras is closed under infinitesimal flat deformations. The proof is elementary and only uses the results on presentations of smooth algebras from Section \hyperref[AG-RG-S05-smoothing-section-presentations]{§ presentations}.

\phantomsection\label{AG-RG-S05-smoothing-lemma-lift-once}

\subsubsection{Lemma D.16: lift once}\label{AG-RG-S05-lemma-d.16-lift-once}

Reading source: \href{https://stacks.math.columbia.edu/tag/07CK}{Stacks, Tag 07CK}.

Let \(R\to\Lambda\) be a ring map and \(I\subset R\) an ideal with \(I^2=0\). Suppose that \(\Lambda/I\Lambda\) is a filtered colimit of smooth \(R/I\)-algebras. Every map \(\varphi:A\to\Lambda\) from a finitely presented \(R\)-algebra factors as \[A\longrightarrow B/J\longrightarrow\Lambda,\] where \(B/R\) is smooth and \(J\subset IB\) is a finitely generated ideal. Neither \(R\) nor \(I\) is required to be Noetherian or finite.

\textbf{Proof.} Write \(\bar R=R/I\) and \(\bar A=A/IA\). By \hyperref[AG-RG-S05-smoothing-lemma-colimit-standard-smooth]{D.13}, the induced map \(\bar A\to\Lambda/I\Lambda\) factors through a standard smooth \(\bar R\)-algebra \(C\). This factorization occurs at a finite stage: take representatives for the finite generators of \(\bar A\) and then a common successor where all its finitely many equations vanish.

There is a finite presentation \(C=\bar A[t_1,\ldots,t_r]/(\bar g_1,\ldots,\bar g_s)\) relative to this specified map. To see the finiteness directly, choose \(C=\bar R[t]/(k_1,\ldots,k_v)\) and generators \(a_1,\ldots,a_u\) of \(\bar A\). Represent their images by polynomials \(v_i(t)\). Over \(\bar A[t]\), the finite equations \(k_j(t)\) and \(a_i-v_i(t)\) give exactly \(C\): substituting for the \(a_i\) recovers its displayed \(\bar R\)-presentation. Include these equations in the list \(\bar g\).

Lift each \(\bar g_i\) to \(g_i\in A[t]\) and the images of the \(t_j\) in \(\Lambda/I\Lambda\) to \(\lambda_j\in\Lambda\). Evaluation gives a finite expression \[\varphi(g_i)(\lambda)=\sum_{j=1}^{q_i}\epsilon_{ij}\mu_{ij},\qquad \epsilon_{ij}\in I,\quad\mu_{ij}\in\Lambda.\] Introduce one variable \(z_{ij}\) for each of these terms and set \[T=A[t,z]/\bigl(g_i-\sum_j\epsilon_{ij}z_{ij}\bigr).\] This is finitely presented over \(R\), by combining the finite presentation of \(A\) with the displayed lists. The assignments \(a\mapsto\varphi(a)\), \(t_j\mapsto\lambda_j\), \(z_{ij}\mapsto\mu_{ij}\) define \(T\to\Lambda\), and reduction gives \[T/IT\cong C[z].\] Adjoining free variables to a standard smooth presentation preserves its invertible Jacobian minor: retain the equations and their pivot variables. Thus choose \[T/IT=(R/I)[X_1,\ldots,X_n]/(\bar f_1,\ldots,\bar f_c)\] with \(\bar\Delta=\det(\partial\bar f_i/\partial X_j)_{1\leq i,j\leq c}\) invertible. Lift the \(f_i\) to \(R[X]\), put \(\Delta=\det(\partial f_i/\partial X_j)_{1\leq i,j\leq c}\), and form \[B=R[X_1,\ldots,X_n,W]/(f_1,\ldots,f_c,W\Delta-1).\] On pivots \(X_1,\ldots,X_c,W\), its Jacobian has upper right block zero and lower right entry \(\Delta\); its determinant is \(\Delta^2\), a unit in \(B\). It is therefore standard smooth, hence smooth by \hyperref[AG-RG-S05-algebra-lemma-standard-smooth]{F.85}. Its reduction is precisely \(T/IT\), since \(W\) becomes the inverse of \(\bar\Delta\).

We construct a map \(\psi:B\to T\) inducing that isomorphism. Choose lifts \(y_i\in T\) of the displayed coordinate classes in \(T/IT\), and let \(L=IT\), so \(L^2=0\). The element \(\Delta(y)\) is a unit: a lift \(d\) of its modular inverse satisfies \(\Delta(y)d=1+\ell\) with \(\ell\in L\), and \((1+\ell)^{-1}=1-\ell\). Let \(H=(\partial f_i/\partial X_j(y))_{1\leq i,j\leq c}\). Its determinant is \(\Delta(y)\), so the cofactor identity proved in \hyperref[AG-RG-S05-algebra-lemma-matrix-left-inverse]{F.70} gives an inverse matrix. Every \(f_i(y)\) lies in \(L\). Choose \[ (u_1,\ldots,u_c)^{\mathsf t}=-H^{-1}(f_1(y),\ldots,f_c(y))^{\mathsf t}\in L^c\] and set \(u_j=0\) for \(j>c\). Expanding each monomial, all terms with two correction factors vanish in \(L^2\), and hence \[ f_i(y+u)=f_i(y)+\sum_{j=1}^c\frac{\partial f_i}{\partial X_j}(y)u_j=0.\] The new determinant is still a unit because its reduction is unchanged. Set \(\psi(X_i)=y_i+u_i\) and \(\psi(W)=\Delta(y+u)^{-1}\). These assignments satisfy every equation of \(B\) and lift the required reduction map.

The map \(\psi\) is onto without a finiteness assumption on \(I\). For \(t\in T\), lift its reduction to \(b\in B\) and write \(t-\psi(b)=\sum_\nu\eta_\nu t_\nu\) with \(\eta_\nu\in I\). Lift each reduction of \(t_\nu\) to \(b_\nu\in B\). Since \(I(IT)=0\), one obtains \(t=\psi(b+\sum_\nu\eta_\nu b_\nu)\). Its kernel \(J\) lies in \(IB\), because \(B/IB\to T/IT\) is injective. To prove that \(J\) is finite, compose the finite polynomial presentation \(R[X,W]\twoheadrightarrow B\) with \(\psi\). The kernel of the resulting surjection to the finitely presented algebra \(T\) is finite by \hyperref[AG-RG-S05-algebra-lemma-finite-presentation-independent]{F.44}. Its images in \(B\) generate \(J\). Thus \(T\cong B/J\), and the maps \(A\to T\to\Lambda\) give the claimed factorization. Empty equation lists use \(\Delta=1\) and require no correction coordinates. ∎

\phantomsection\label{AG-RG-S05-smoothing-lemma-lift-twice}

\subsubsection{Lemma D.17: lift twice}\label{AG-RG-S05-lemma-d.17-lift-twice}

Reading source: \href{https://stacks.math.columbia.edu/tag/07CL}{Stacks, Tag 07CL}.

Let \(R\to\Lambda\) be flat, let \(I^2=0\), and suppose that \(\Lambda/I\Lambda\) is a filtered colimit of smooth \(R/I\)-algebras. If \(B/R\) is smooth, \(\varphi:B\to\Lambda\) is an \(R\)-map, and \(J\subset IB\) is a finitely generated ideal with \(\varphi(J)=0\), there are maps \[B\xrightarrow{\alpha}B'\xrightarrow{\beta}\Lambda\] with \(B'/R\) smooth, \(\alpha(J)=0\), and the exact equality \(\beta\alpha=\varphi\).

\textbf{Proof.} We first give the square-zero lifting calculation needed below. For any smooth \(R\)-algebra \(B\), any \(R\)-algebra \(E\), any ideal \(L\subset E\) with \(L^2=0\), and any \(R\)-map \(\gamma:B\to E/L\), there is a lift \(B\to E\). Choose a finite polynomial presentation \(P=R[Z_1,\ldots,Z_n]\twoheadrightarrow B\), with kernel \(K\), and lifts \(e_i\in E\) of \(\gamma(Z_i)\). Evaluation of a relation \(f\in K\) lies in \(L\) and vanishes on \(K^2\). It gives a \(B\)-linear map \[\delta:K/K^2\longrightarrow L,\] where \(L\) is a \(B\)-module through \(\gamma\); its \(E/L\)-module structure is well defined because \(L^2=0\). Indeed evaluation of \(pf\) is evaluation of \(p\) times evaluation of \(f\), and that multiplication on \(L\) depends only on \(p\) modulo \(K\) and \(E\) modulo \(L\).

By the definition \hyperref[AG-RG-S05-algebra-definition-smooth]{F.6} and presentation comparison \hyperref[AG-RG-S05-algebra-lemma-NL-homotopy]{F.11}, the sequence \[0\longrightarrow K/K^2\xrightarrow{d}B^n\longrightarrow\Omega_{B/R}\longrightarrow0\] is exact and its last module is finite projective. This sequence splits: realize that projective module as a direct summand of a finite free module, lift the free basis through the surjection, and restrict the lift to the summand. The resulting section defines a retraction \(r:B^n\to K/K^2\) of \(d\). Put \(\ell=\delta r:B^n\to L\) and \(c_i=\ell(dZ_i)\). Evaluate polynomials instead at \(e_i-c_i\). The monomial expansion, using \(L^2=0\), gives for each \(f\in K\) \[ f(e-c)=f(e)-\sum_i\frac{\partial f}{\partial Z_i}(e)c_i
 =\delta([f])-\ell(d[f])=0.\] In the second equality the derivative coefficients act on \(L\) through their images under \(\gamma\), so no choice of their lifts changes the product. Evaluation therefore descends to \(B\to E\) and lifts \(\gamma\). This proves the needed lifting statement in full arbitrary-ring scope.

Now suppose first that \(J=(h)\). Write \(h=\sum_{i=1}^q\epsilon_i b_i\) with \(\epsilon_i\in I\) and \(b_i\in B\). Consider the row map \(v:R^q\to R\), \(v(a)=\sum_i\epsilon_i a_i\), and let \(K_v=\ker v\), \(H_v=\operatorname{im}v\). Flatness preserves the two exact sequences \(0\to K_v\to R^q\to H_v\to0\) and \(0\to H_v\to R\). Consequently \[\ker(v\otimes_R\Lambda)=\operatorname{im}(K_v\otimes_R\Lambda\longrightarrow\Lambda^q).\] The vector \((\varphi(b_i))_i\) belongs to this kernel, since \(\varphi(h)=0\). An element of a tensor product is a finite sum of pure tensors. Thus there are finitely many \(\lambda_j\in\Lambda\) and columns \((a_{ij})_i\in K_v\) with \[\varphi(b_i)=\sum_{j=1}^m a_{ij}\lambda_j,\qquad
\sum_i\epsilon_i a_{ij}=0\quad\text{for each }j.\] This derives the finite relation directly from flatness, including the case of a zero vector and an empty tensor sum.

Set \(C=B[T_1,\ldots,T_m]/(b_i-\sum_j a_{ij}T_j)\). Its finite presentation follows by appending these variables and equations to a finite presentation of \(B\). The displayed relations define a map \(C\to\Lambda\) extending \(\varphi\), with \(T_j\mapsto\lambda_j\). By \hyperref[AG-RG-S05-smoothing-lemma-lift-once]{D.16} this map factors through \(E/J'\), where \(E/R\) is smooth and \(J'\subset IE\) is finite. In particular \((J')^2=0\). Apply the lifting calculation above to \(B\to C\to E/J'\) and obtain \(\alpha:B\to E\). Let \(\beta:E\to E/J'\to\Lambda\). Since the reduction of \(\alpha\) is the original map through \(C\), one has \(\beta\alpha=\varphi\) exactly.

Choose lifts \(\xi_j\in E\) of the images of \(T_j\). The equations in \(C\) give \(\alpha(b_i)=\sum_j a_{ij}\xi_j+\theta_i\) for some \(\theta_i\in J'\). Therefore \[\alpha(h)=\sum_j\Bigl(\sum_i\epsilon_i a_{ij}\Bigr)\xi_j
              +\sum_i\epsilon_i\theta_i=0.\] The first sum vanishes by the column relations, the second because \(IJ'\subset I^2E=0\). This proves the one-generator case, with \(B'=E\).

For \(J=(h_1,\ldots,h_s)\), apply that construction successively to the image of \(h_i\) in the smooth algebra obtained at the preceding step. Each such image lies in \(I\) times that algebra, and its map to \(\Lambda\) is zero because the composite to \(\Lambda\) remains \(\varphi\). Each new map preserves the previously killed generators. After finitely many steps their composite kills every generator, hence the whole ideal \(J\), and still factors \(\varphi\) exactly. If \(J=0\), take \(B'=B\), \(\alpha=1\), and \(\beta=\varphi\). ∎

\phantomsection\label{AG-RG-S05-smoothing-proposition-lift}

\subsubsection{Proposition D.18: lift}\label{AG-RG-S05-proposition-d.18-lift}

Reading source: \href{https://stacks.math.columbia.edu/tag/07CM}{Stacks, Tag 07CM}.

Suppose that \(R\to\Lambda\) is flat, \(I\subset R\) is nilpotent, and \(\Lambda/I\Lambda\) is a filtered colimit of smooth \(R/I\)-algebras. Then \(\Lambda\) is a filtered colimit of smooth \(R\)-algebras.

\textbf{Proof.} Begin with \(I^2=0\). Given any \(R\)-map \(A\to\Lambda\) from a finitely presented algebra, \hyperref[AG-RG-S05-smoothing-lemma-lift-once]{D.16} gives maps \(A\to B/J\to\Lambda\) with \(B/R\) smooth, \(J\subset IB\) finite. Let \(\varphi:B\to\Lambda\) be the induced map; it kills \(J\). The flat map and the same square-zero ideal meet every hypothesis of \hyperref[AG-RG-S05-smoothing-lemma-lift-twice]{D.17}. Obtain \(B\to B'\to\Lambda\) with \(B'/R\) smooth and \(J\) killed in \(B'\). The first arrow descends to \(B/J\to B'\), and its composite to \(\Lambda\) is the given quotient map because D.17 gives equality of the maps from \(B\). Thus \(A\to B'\to\Lambda\) is an exact factorization of the original map. Criterion \hyperref[AG-RG-S05-algebra-lemma-when-colimit]{F.90}, whose filtered-category construction and colimit isomorphism are proved there, gives the assertion for square-zero ideals.

Induct on a positive integer \(n\) with \(I^n=0\). For \(n=1\) the ideal is zero and the assertion is the hypothesis. For \(n\geq2\), set \(J=I^{n-1}\), so \(J^2=0\) because \(2n-2\geq n\). Over \(\bar R=R/J\) put \(\bar\Lambda=\Lambda/J\Lambda\) and \(\bar I=I/J\). Then \(\bar I^{n-1}=0\) and \(\bar\Lambda/\bar I\bar\Lambda=\Lambda/I\Lambda\). The map \(\bar R\to\bar\Lambda\) is flat: for every \(\bar R\)-module \(M\) the maps on pure tensors give \[M\otimes_{\bar R}(\Lambda/J\Lambda)\cong M\otimes_R\Lambda.\] Here tensors involving \(J\Lambda\) vanish because \(JM=0\); this verifies the inverse maps as well. An exact sequence of \(\bar R\)-modules is exact over \(R\), so the original flatness makes the displayed tensor functor exact. The induction hypothesis therefore expresses \(\bar\Lambda\) as a filtered colimit of smooth \(\bar R\)-algebras. Apply the square-zero case already proved to \(J\subset R\) and the original flat map \(R\to\Lambda\). Its quotient is exactly this smooth colimit, hence it expresses \(\Lambda\) as a filtered colimit of smooth \(R\)-algebras. This completes the induction without assuming that the nilpotent ideal is finitely generated. ∎

\subsection{The lifting lemma}\label{AG-RG-S05-the-lifting-lemma}

\phantomsection\label{AG-RG-S05-smoothing-section-lifting-lemma}

The next lemma constructs a finite presentation lifting the smooth locus.

\phantomsection\label{AG-RG-S05-smoothing-lemma-lifting}

\subsubsection{Lemma D.19: lifting}\label{AG-RG-S05-lemma-d.19-lifting}

Reading source: \href{https://stacks.math.columbia.edu/tag/07CP}{Stacks, Tag 07CP}.

Let \(R\) be Noetherian, let \(\pi\in R\), let \(\Lambda\) be an \(R\)-algebra, and assume \[
\operatorname{Ann}_R(\pi)=\operatorname{Ann}_R(\pi^2),
\qquad
\operatorname{Ann}_\Lambda(\pi)=\operatorname{Ann}_\Lambda(\pi^2).
\] Given a finitely presented \(R/\pi^2R\)-algebra \(\bar C\) and a map \(\bar C\to\Lambda/\pi^2\Lambda\), there is a finitely presented \(R\)-algebra \(D\), a map \(D\to\Lambda\), and commuting maps \[
\begin{array}{ccccc}
R/\pi^2R&\longrightarrow&\bar C&\longrightarrow&\Lambda/\pi^2\Lambda\\
\downarrow&&\downarrow&&\downarrow\\
R/\pi R&\longrightarrow&D/\pi D&\longrightarrow&\Lambda/\pi\Lambda.
\end{array}
\] The map \(R\to D\) is smooth at every prime avoiding \(\pi\). It is also smooth at each prime containing \(\pi\) whose inverse image in \(\bar C\) is smooth over \(R/\pi^2R\). At these latter primes, \(\bar C/\pi\bar C\to D/\pi D\) is smooth as well.

\textbf{Proof.} Put \(P=R[x_1,\ldots,x_n]\) and choose lifts \(f_1,\ldots,f_m\in P\) of a finite relation list for \[
\bar C=(R/\pi^2R)[x_1,\ldots,x_n]/(\bar f_1,\ldots,\bar f_m).
\] Write \(\bar P=(R/\pi^2R)[x]\) and \(\bar J=(\bar f_1,\ldots,\bar f_m)\). Thus the kernel of \(P\to\bar C\) is \((f_1,\ldots,f_m)+\pi^2P\); \(\pi^2\) need not be in the chosen list.

At a smooth prime, \hyperref[AG-RG-S05-smoothing-lemma-find-strictly-standard]{D.2} supplies a principal neighbourhood, selected equation indices \(E\), and equally many variable indices \(U\). On that neighbourhood the selected Jacobian minor \(\delta\) is a unit and the selected classes generate \(\bar J/\bar J^2\). They are a basis: their composite with the conormal differential projected onto the \(U\)-coordinates is invertible, so a relation among them is zero. Shrink the polynomial neighbourhood to make \(\bar J\) itself generated by these selected equations. Here are the details of that shrinking. The finite module \[
M=(\bar J/(\bar f_j:j\in E))_a
\] satisfies \(M=\bar J_aM\). The determinant form of \hyperref[AG-RG-S05-algebra-lemma-NAK]{F.10} gives \(b'\in\bar J_a\) with \((1+b')M=0\). Write \(b'=b/a^N\), with \(b\in\bar J\). Replacing the polynomial \(a\) by \(a(a^N+b)\) inverts \(1+b'\). Its image in \(\bar C\) is the old \(a^{N+1}\), so this replacement preserves the principal open in \(\operatorname{Spec}\bar C\).

These principal opens cover the smooth locus. They have a finite subcover for an algebraic reason: \(\bar C\) is Noetherian by \href{https://kokunoyumeto.github.io/open-math-courses/courses/AG-RG/AG-RG-S01.html}{S01, Lemma 1.3}, so the ideal generated by all their defining elements is generated by finitely many of those elements. To obtain such a list, express a finite ideal generating list as finite sums of the original elements and collect those that occur. Their principal opens have the same union. Index the resulting charts by \(k=1,\ldots,r\). An empty smooth locus uses \(r=0\).

Clear the finitely many denominators in the equality of localized relation ideals on each chart and replace its \(a_k\) by a sufficiently large positive power. There are then exact identities in \(P\) \[
a_kf_\ell=\sum_{j\in E_k}h_{k\ell}^jf_j+\pi^2g_{k\ell}.
\tag{D.19a}
\] For \(\ell\in E_k\), take \(h_{k\ell}^j=a_k\delta_{\ell j}\) and \(g_{k\ell}=0\). The selected conormal classes remain a basis on \(\bar C_{a_k}\), and the selected Jacobian minor \(\delta_k\) remains invertible there.

In \(Q=P[z_1,\ldots,z_m]\), put \[
F_\ell=f_\ell-\pi z_\ell,\qquad
p_{k\ell}=a_kz_\ell-\sum_{j\in E_k}h_{k\ell}^jz_j-\pi g_{k\ell},
\] and define \(D=Q/(F_\ell,p_{k\ell})_{\ell,k}\). All lists are finite. Choose lifts \(\lambda_i\in\Lambda\) of the prescribed images of \(x_i\), and write \(f_\ell(\lambda)=\pi^2\mu_\ell\). Send \(x_i\) to \(\lambda_i\) and \(z_\ell\) to \(\pi\mu_\ell\). The \(F_\ell\) vanish. Evaluation of (D.19a) shows that \[
\pi^2\left(a_k(\lambda)\mu_\ell-
\sum_{j\in E_k}h_{k\ell}^j(\lambda)\mu_j-g_{k\ell}(\lambda)\right)=0.
\] The annihilator hypothesis in \(\Lambda\) makes its product with \(\pi\) zero; that product is exactly the image of \(p_{k\ell}\). This defines \(D\to\Lambda\). Reduction of \(F_\ell=0\) modulo \(\pi\) defines \(\bar C\to D/\pi D\), and the variable assignments verify every square in the displayed diagram.

The polynomial identity \[
\pi p_{k\ell}=-a_kF_\ell+\sum_{j\in E_k}h_{k\ell}^jF_j
\tag{D.19b}
\] follows by substituting (D.19a). After inverting \(\pi\), all \(p_{k\ell}\) follow from the \(F_\ell\), and \(F_\ell=0\) eliminates \(z_\ell=f_\ell/\pi\). Thus \(D_\pi=R_\pi[x]\), a polynomial localization smooth over \(R\) by the empty-equation and inverse-variable cases of \hyperref[AG-RG-S05-algebra-lemma-standard-smooth]{F.85}.

Fix \(k\), and retain only its selected equations and its own compatibility equations: \[
D_k=Q/(F_j:j\in E_k,\ p_{k\ell}:\ell\notin E_k).
\] There is a surjection \(D_k\to D\). Let \(\mathfrak q\subset D\) contain \(\pi\) and lie over a smooth prime of \(\bar C\); choose \(k\) with \(a_k\notin\mathfrak q\), and let \(\mathfrak q_k\) be its inverse image in \(D_k\). In \(D_k/\pi D_k\), after inverting \(a_k\), each \(z_\ell\) for \(\ell\notin E_k\) is uniquely \[
z_\ell=a_k^{-1}\sum_{j\in E_k}h_{k\ell}^jz_j.
\] Substitution and its reverse on the remaining variables give the actual isomorphism \[
(D_k/\pi D_k)_{a_k}
 \cong(\bar C/\pi\bar C)_{a_k}[z_j:j\in E_k].
\tag{D.19c}
\] Here the equality of the relation ideals modulo \(\pi^2\) also gives their equality modulo \(\pi\).

There are \(m\) equations in this presentation of \(D_k\). Choose as pivot variables \(x_i\), \(i\in U_k\), followed by \(z_\ell\), \(\ell\notin E_k\). Their Jacobian modulo \(\pi\), with selected equations first, has blocks \[
\begin{pmatrix}
(\partial f_j/\partial x_i)_{j\in E_k,i\in U_k}&0\\
*&a_k1_{m-|E_k|}
\end{pmatrix}.
\] Its determinant is \(\delta_k a_k^{m-|E_k|}\), invertible at \(\mathfrak q_k\). If \(\Delta_k\) denotes the determinant before reduction, \((D_k)_{a_k\Delta_k}\) is standard smooth: adjoining an inverse of \(a_k\Delta_k\) adds one equation and one pivot, whose block determinant is \((a_k\Delta_k)\Delta_k\). \hyperref[AG-RG-S05-algebra-lemma-standard-smooth]{F.85} proves smoothness. \hyperref[AG-RG-S05-algebra-lemma-smooth-syntomic]{F.83}, with the Noetherian base \(R\), proves flatness of this smooth chart. Localization is exact by \href{https://kokunoyumeto.github.io/open-math-courses/courses/AG-RG/AG-RG-S01.html}{S01, Lemma 1.1}, so \(T=(D_k)_{\mathfrak q_k}\) is flat over \(R\).

We now show that every additional relation of \(D_k\to D\) vanishes in \(T\). First (D.19b), with the \(k\)-relations zero and \(a_k\) invertible, gives \(F_\ell=0\) for all \(\ell\). For another chart \(k'\), set \[
b_{kk'\ell}^j=a_{k'}h_{k\ell}^j-
\sum_{j'\in E_{k'}}h_{k'\ell}^{j'}h_{kj'}^j
\quad(j\in E_k).
\] Reduce (D.19a) modulo \(\pi^2\) and express \(a_ka_{k'}[\bar f_\ell]\) in the conormal basis indexed by \(E_k\) in the two possible orders. The difference is \(\sum_{j\in E_k}b_{kk'\ell}^j[\bar f_j]=0\). Basis independence therefore makes every coefficient zero in \(\bar C_{a_k}\). Equivalently, \[
b_{kk'\ell}^j\in
((f_j:j\in E_k)+\pi^2P)_{a_k}.
\] Its image in \((D_k)_{a_k}\) belongs to \(\pi(D_k)_{a_k}\), since every selected \(f_j\) equals \(\pi z_j\). Directly expanding the definitions, modulo \(\pi\), gives \[
a_kp_{k'\ell}-a_{k'}p_{k\ell}
+\sum_{j'\in E_{k'}}h_{k'\ell}^{j'}p_{kj'}
=\sum_{j\in E_k}b_{kk'\ell}^jz_j.
\] Consequently \(p_{k'\ell}\in\pi T\). Identity (D.19b) for \(k'\) also gives \(\pi p_{k'\ell}=0\) in \(T\).

Flat tensor applied to the multiplication-kernel sequences identifies \[
\operatorname{Ann}_T(\pi^e)
=\operatorname{im}(\operatorname{Ann}_R(\pi^e)\otimes_RT\to T)
\qquad(e=1,2).
\] Indeed, factor multiplication through its image and tensor its short exact kernel sequence and that image's injection into \(R\); flatness preserves both. The assumed equality over \(R\) thus gives equality of these annihilators in \(T\). For \(p=p_{k'\ell}=\pi u\), the equality \(\pi p=0\) means \(\pi^2u=0\); it follows that \(p=\pi u=0\). This kills all additional generators.

Finally this local calculation supplies genuine neighbourhoods. The kernel of \(D_k\to D\) is generated by the finite lists of additional \(F_\ell\) and \(p_{k'\ell}\). For each generator, vanishing at \(\mathfrak q_k\) gives an element outside \(\mathfrak q_k\) killing it. Their product \(s\) kills the entire kernel. The surjection therefore becomes an isomorphism after inverting \(a_k\Delta_ks\), which still avoids \(\mathfrak q_k\). This identifies a principal neighbourhood of \(\mathfrak q\) with a localization of the standard smooth chart, proving the required smoothness of \(R\to D\) by F.85 and \hyperref[AG-RG-S05-algebra-definition-smooth-at-prime]{F.7}. Modulo \(\pi\), (D.19c) identifies the same neighbourhood with a localization of a polynomial algebra over \(\bar C/\pi\bar C\): denominators in the \(a_k\)-localized polynomial algebra can be cleared by a power of \(a_k\). F.85 proves its smoothness over that base. This gives the second claimed smoothness assertion as well. Empty equation or chart lists use determinant and denominator product \(1\); the same constructions cover zero algebras. ∎

\subsection{The desingularization lemma}\label{AG-RG-S05-the-desingularization-lemma}

\phantomsection\label{AG-RG-S05-smoothing-section-desingularization-lemma}

The next construction improves the singular ideal.

\phantomsection\label{AG-RG-S05-smoothing-lemma-desingularize}

\subsubsection{Lemma D.20: desingularize}\label{AG-RG-S05-lemma-d.20-desingularize}

Reading source: \href{https://stacks.math.columbia.edu/tag/07CR}{Stacks, Tag 07CR}.

Let \(R\) be Noetherian, let \(\pi\in R\), and let \(\Lambda\) be an \(R\)-algebra with \(\operatorname{Ann}_\Lambda(\pi)=\operatorname{Ann}_\Lambda(\pi^2)\). Suppose \(A\) is finitely presented over \(R\), a map \(A\to\Lambda\) is given, the image of \(\pi\) is strictly standard in \(A/R\), and a section \[
\rho:A/\pi^4A\longrightarrow R/\pi^4R
\] agrees with the prescribed map to \(\Lambda/\pi^4\Lambda\). Then there is a factorization \(A\to B\to\Lambda\), with \(B/R\) finitely presented, such that \[
\mathfrak aB\subset H_{B/R},\qquad
\mathfrak a=\operatorname{Ann}_R
\bigl(\operatorname{Ann}_R(\pi^2)/\operatorname{Ann}_R(\pi)\bigr).
\] No annihilator equality in \(R\) is assumed.

\textbf{Proof.} Choose the strict witness from \hyperref[AG-RG-S05-smoothing-definition-strictly-standard]{D.3}: \[
A=R[x_1,\ldots,x_n]/(f_1,\ldots,f_m),\quad
0\le c\le\min(n,m),\quad I=(f_1,\ldots,f_m).
\] It gives a sum-of-minors identity for \(\pi\) and the relations \(\pi f_\ell\in(f_1,\ldots,f_c)+I^2\) for \(\ell>c\). Subtract lifts in \(R\) of the finitely many values \(\rho(x_i)\) from the variables. This translation is a polynomial automorphism, preserves both witness identities, and makes \(\rho(x_i)=0\). Hence \(f_j(0)\in\pi^4R\), and the actual image of \(x_i\) in \(\Lambda\) has the form \(\pi^4\lambda_i\).

Write the linear terms of \(f_j\) as \(\sum_i r_{ji}x_i\), and let \(M=(r_{ji})_{j\le c,i\le n}\). Evaluate the sum-of-minors identity at \(\rho\). There are coefficients \(d_U\in R\) and \(u\in1+\pi^3R\) such that \[
u\pi=\sum_{|U|=c}d_U\det(M_U).
\] Apply \hyperref[AG-RG-S05-algebra-lemma-matrix-left-inverse]{F.70} to \(M^{\mathsf t}\) and transpose its resulting identity. This supplies an \(n\times c\) matrix \(S=(s_{ij})\) with \(MS=u\pi1_c\).

Introduce \(v_1,\ldots,v_c,w_1,\ldots,w_n\), set \(L_i=\sum_j s_{ij}v_j+\pi w_i\), and substitute \(x_i=\pi^2L_i\). For \(j\le c\), constant terms are divisible by \(\pi^4\), linear terms become \[
\pi^2\sum_i r_{ji}L_i
=\pi^3u v_j+\pi^3\sum_i r_{ji}w_i,
\] and every monomial of degree \(d\ge2\) becomes divisible by \(\pi^{2d}\). Dividing these specified powers in the coefficient expressions, rather than cancelling a ring element, constructs a polynomial \(G_j\) and an exact identity \[
f_j(\pi^2L)=\pi^3g_j,\qquad
g_j=u v_j+\sum_i r_{ji}w_i+\pi G_j.
\tag{D.20a}
\] For example, a constant term \(\pi^4a_j\) contributes \(a_j\) to \(G_j\), and a degree-\(d\) term with coefficient \(b_\alpha\) contributes \(b_\alpha\pi^{2d-4}L^\alpha\). Thus this construction is valid even when \(\pi\) is a zero divisor.

Let \(h_i=x_i-\pi^2L_i\), and define \[
B=R[x,v,w]/(h_i,\ g_j:j\le c,\ f_\ell:1\le\ell\le m).
\] Eliminating \(x_i\) using the \(h_i\) is an actual inverse substitution isomorphism. The two required maps send the old \(x_i\) to their classes in \(B\), and send \[
x_i\longmapsto\pi^4\lambda_i,\qquad
v_j\longmapsto0,\qquad w_i\longmapsto\pi\lambda_i
\] in \(\Lambda\). The \(h_i\) and \(f_\ell\) vanish there. Each \(g_j\) has image \(t\in\pi\Lambda\) by (D.20a) modulo \(\pi\), and satisfies \(\pi^3t=0\). The annihilator equality implies \(\operatorname{Ann}_\Lambda(\pi^e)=\operatorname{Ann}_\Lambda(\pi)\) for every \(e\ge1\): if \(\pi^e b=0\) for \(e\ge2\), apply the equality to \(\pi^{e-2}b\) and reduce the exponent successively. Writing \(t=\pi b\) now gives \(\pi^4b=0\), hence \(t=0\). All relations therefore vanish, and the factorization is exact.

After inverting \(\pi\), the \(h_i\) uniquely eliminate \(w_i\), and (D.20a) makes the \(g_j\) consequences of the \(f_j\). Consequently \[
B_\pi\cong A_\pi[v_1,\ldots,v_c].
\] \hyperref[AG-RG-S05-smoothing-lemma-elkik]{D.5} makes \(A_\pi\) smooth over \(R\); a polynomial extension is smooth by \hyperref[AG-RG-S05-algebra-lemma-standard-smooth]{F.85}, and \hyperref[AG-RG-S05-algebra-lemma-compose-smooth]{F.29} proves smoothness of \(B_\pi/R\).

For primes containing \(\pi\), first use the algebra \[
B'=R[v,w]/(g_1,\ldots,g_c).
\] The Jacobian on the pivot variables \(v_j\) is \(u1_c+\pi(\partial G_j/\partial v_k)\), so its determinant \(\Delta\) is \(1\) modulo \(\pi\). Thus \(B'_\Delta\) is standard smooth by the inverse-variable Jacobian calculation in F.85. It is flat over the Noetherian base \(R\) by \hyperref[AG-RG-S05-algebra-lemma-smooth-syntomic]{F.83}. Notice also the explicit reduction \(B'/\pi B'\cong(R/\pi R)[w]\), obtained by eliminating \(v_j=-\sum_i r_{ji}w_i\).

The natural surjection \(B'\to B\) has kernel generated by \[
F_\ell=f_\ell(\pi^2L)\quad(\ell>c);
\] the first \(c\) equations have already vanished by (D.20a). Let \(\mathfrak q\subset B\) contain \(\pi\) and avoid some \(r\in\mathfrak a\), and let \(\mathfrak q'\) be its inverse image in \(B'\). Then \(\Delta\notin\mathfrak q'\), and \(T=B'_{\mathfrak q'}\) is flat over \(R\), since it is a localization of that flat chart. Every \(F_\ell\) belongs to \(\pi^2T\): its substituted variables have this divisibility and its constant term belongs to \(\pi^4R\).

Substitute in the strict relation condition and use the first \(c\) equations already killed in \(B'\). It gives \[
\pi F_\ell=\sum_{a,b>c}b_{\ell ab}F_aF_b.
\] For the finite \(T\)-module \(N=(\pi F_\ell:\ell>c)\), divisibility of \(F_a\) by \(\pi^2\) gives \(N\subset\pi N\); the reverse inclusion holds for every ideal. Since \(\pi\) lies in the maximal ideal of \(T\), \hyperref[AG-RG-S05-algebra-lemma-NAK]{F.10} gives \(N=0\). Thus \(\pi F_\ell=0\).

Flatness identifies each multiplication kernel and preserves their quotient sequence, giving the canonical isomorphism \[
\bigl(\operatorname{Ann}_R(\pi^2)/\operatorname{Ann}_R(\pi)\bigr)
\otimes_RT
\cong
\operatorname{Ann}_T(\pi^2)/\operatorname{Ann}_T(\pi).
\] For the kernel identification, tensor the short exact sequence through the image of multiplication by \(\pi^e\), and its injection into \(R\), for \(e=1,2\). Tensor their annihilator inclusion and quotient sequence next. The element \(r\) kills the left-hand module and is a unit in \(T\), so the quotient on the right is zero. Write \(F_\ell=\pi b_\ell\). From \(\pi F_\ell=0\) and the resulting annihilator equality in \(T\), we obtain \(F_\ell=0\).

To obtain a smooth neighbourhood rather than only an isomorphism of local rings, choose for each of the finitely many \(F_\ell\) an element of \(B'\setminus\mathfrak q'\) killing it; these exist by the localization equality criterion. Their product \(s\) kills the whole kernel. Hence \(B'_{\Delta s}\to B_{\Delta s}\) is an isomorphism, and \(\Delta s\) avoids \(\mathfrak q\). This is a localization of a standard smooth algebra, so F.85 and \hyperref[AG-RG-S05-algebra-definition-smooth-at-prime]{F.7} prove smoothness of \(R\to B\) at \(\mathfrak q\).

Every prime avoiding \(\pi\) is smooth by the earlier description of \(B_\pi\). Every other prime avoiding \(\mathfrak aB\) has some such \(r\) and is smooth by the chart just constructed. Thus every nonsmooth prime contains \(\mathfrak aB\). The radical/prime intersection formula proved in \hyperref[AG-RG-S05-smoothing-definition-singular-ideal]{D.1}, applied to \(H_{B/R}\), gives \(\mathfrak aB\subset H_{B/R}\). Empty lists use determinant and product \(1\), so the argument includes \(c=0\), \(c=m\), and zero rings. ∎

\phantomsection\label{AG-RG-S05-smoothing-lemma-desingularize-strictly-standard}

\subsubsection{Lemma D.21: desingularize strictly standard}\label{AG-RG-S05-lemma-d.21-desingularize-strictly-standard}

Reading source: \href{https://stacks.math.columbia.edu/tag/07CT}{Stacks, Tag 07CT}.

Let \(R\) be Noetherian, let \(\pi\in R\), and let \(A\to\Lambda\), \(D\to\Lambda\) be \(R\)-algebra maps, with \(A/R\) and \(D/R\) finitely presented. Assume \[
\operatorname{Ann}_R(\pi)=\operatorname{Ann}_R(\pi^2),\qquad
\operatorname{Ann}_\Lambda(\pi)=\operatorname{Ann}_\Lambda(\pi^2),
\] that \(\pi\) is strictly standard in \(A/R\), and that a map \(A/\pi^4A\to D/\pi^4D\) agrees with both prescribed maps to \(\Lambda/\pi^4\Lambda\). There are compatible maps \(A\to B\), \(D\to B\), \(B\to\Lambda\), with \(B/R\) finitely presented, such that \[
H_{D/R}B\subset H_{B/D},
\qquad
H_{D/R}B\subset H_{B/R}.
\]

\textbf{Proof.} Let \(C=A\otimes_RD\), with map to \(\Lambda\) sending \(a\otimes d\) to the product of their prescribed images. A finite presentation \(A=R[x]/(f)\) gives \(C=D[x]/(f)\), so \(C/D\) is finitely presented. The strict witness changes coefficients to \(D\) by \hyperref[AG-RG-S05-smoothing-lemma-strictly-standard-base-change]{D.7}. The given map \(\phi:A/\pi^4A\to D/\pi^4D\) gives the section \[
C/\pi^4C\longrightarrow D/\pi^4D,\qquad
a\otimes d\longmapsto\phi(\bar a)\bar d.
\] The tensor relations make this well-defined; on the \(D/\pi^4D\) factor it is the identity. Its composition to \(\Lambda/\pi^4\Lambda\) is the original product map by the assumed compatibility.

The ring \(D\) is Noetherian by \href{https://kokunoyumeto.github.io/open-math-courses/courses/AG-RG/AG-RG-S01.html}{S01, Lemma 1.3}. Apply \hyperref[AG-RG-S05-smoothing-lemma-desingularize]{D.20} with base \(D\) and algebra \(C\). Its hypotheses require the annihilator equality in \(\Lambda\), already assumed, and require no such equality in \(D\). We obtain \(C\to B\to\Lambda\), with \(B/D\) finitely presented, and \[
\mathfrak aB\subset H_{B/D},\qquad
M=\operatorname{Ann}_D(\pi^2)/\operatorname{Ann}_D(\pi),\quad
\mathfrak a=\operatorname{Ann}_D(M).
\] The two tensor-factor maps give \(A\to B\) and \(D\to B\), with the desired exact composites to \(\Lambda\). If \(D=R[u]/(g)\) and \(B=D[v]/(h)\), choose polynomial lifts \(\widetilde h\) in \(R[u,v]\). Substitution in both directions gives \(B=R[u,v]/(g,\widetilde h)\), proving finite presentation over \(R\).

Fix a prime \(\mathfrak q\subset D\) where \(R\to D\) is smooth. A smooth principal neighbourhood is flat over the Noetherian ring \(R\) by \hyperref[AG-RG-S05-algebra-lemma-smooth-syntomic]{F.83}; exact localization \href{https://kokunoyumeto.github.io/open-math-courses/courses/AG-RG/AG-RG-S01.html}{S01, Lemma 1.1} makes \(D_{\mathfrak q}\) flat as well. Tensor the multiplication-kernel sequences for \(\pi,\pi^2\) as in D.20. The equality of the original annihilators over \(R\) gives \[
\operatorname{Ann}_{D_{\mathfrak q}}(\pi)
=\operatorname{Ann}_{D_{\mathfrak q}}(\pi^2),\qquad M_{\mathfrak q}=0.
\] The module \(M\) is finite: its two annihilator ideals are finite over the Noetherian ring \(D\), and \(M\) is their quotient. Choose generators \(m_i\). For each, vanishing at \(\mathfrak q\) gives \(s_i\notin\mathfrak q\) with \(s_im_i=0\). Their product \(s\) kills every generator; hence \(s\in\mathfrak a\setminus\mathfrak q\). For \(M=0\), use \(s=1\).

At a prime \(\mathfrak r\subset B\) above \(\mathfrak q\), this element remains outside \(\mathfrak r\). Since \(s\in\mathfrak aB\subset H_{B/D}\), \hyperref[AG-RG-S05-smoothing-definition-singular-ideal]{D.1} shows that \(D\to B\) is smooth at \(\mathfrak r\). To verify smooth composition at the level of actual neighbourhoods, choose \(t\in D\setminus\mathfrak q\) with \(D_t/R\) smooth. Inside a smooth \(D\)-neighbourhood of \(\mathfrak r\), F.83 supplies a principal neighbourhood \(B_g\) that is standard smooth over \(D\). The denominator conversion in its proof makes \(g\) an element of \(B\setminus\mathfrak r\). Base change to \(D_t\) preserves the standard smooth presentation and its unit determinant by \hyperref[AG-RG-S05-algebra-lemma-standard-smooth]{F.85(5)}, so \(B_{tg}/D_t\) is smooth. \hyperref[AG-RG-S05-algebra-lemma-compose-smooth]{F.29} now makes \(B_{tg}/R\) smooth. Thus \(R\to B\) is smooth at every such \(\mathfrak r\).

An element of \(H_{D/R}\) avoided by a prime of \(B\) is avoided by its contraction in \(D\), so that contraction is smooth. We have just proved that the prime of \(B\) is smooth for both \(B/D\) and \(B/R\). Consequently every nonsmooth prime for either map contains \(H_{D/R}B\). Intersecting these primes and using the radical/prime formula in D.1 proves both claimed ideal inclusions. Empty singular loci and zero algebras use the same convention that the intersection of no primes is the unit ideal. ∎

\phantomsection\label{AG-RG-S05-smoothing-lemma-desingularize-lifting-apply}

\subsubsection{Lemma D.22: desingularize lifting apply}\label{AG-RG-S05-lemma-d.22-desingularize-lifting-apply}

Reading source: \href{https://stacks.math.columbia.edu/tag/07F0}{Stacks, Tag 07F0}.

Let \(R\) be Noetherian, let \(\pi\in R\), and let \(A\to\Lambda\) be an \(R\)-algebra map with \(A/R\) finitely presented. Assume \[
\operatorname{Ann}_R(\pi)=\operatorname{Ann}_R(\pi^2),\qquad
\operatorname{Ann}_\Lambda(\pi)=\operatorname{Ann}_\Lambda(\pi^2),
\] and that \(\pi\) is strictly standard in \(A/R\). Given a factorization \[
A/\pi^8A\longrightarrow\bar C\longrightarrow\Lambda/\pi^8\Lambda
\] with \(\bar C\) finitely presented over \(R/\pi^8R\), there is a factorization \(A\to B\to\Lambda\), with \(B/R\) finitely presented, such that \(B_\pi/R_\pi\) is smooth and \[
H_{\bar C/(R/\pi^8R)}(\Lambda/\pi^8\Lambda)
\subset \bigl(\sqrt{H_{B/R}\Lambda}+\pi^8\Lambda\bigr)/\pi^8\Lambda.
\]

\textbf{Proof.} If multiplication by \(u\) on a module has \(\ker u=\ker u^2\), its kernels for all positive powers agree. Indeed \(u^{r+1}v=0\) puts \(u^{r-1}v\) in \(\ker u^2=\ker u\), hence \(u^rv=0\); repeating lowers the exponent to one. Thus the annihilators of \(\pi^4\) and \(\pi^8\) agree in both \(R\) and \(\Lambda\).

Use \hyperref[AG-RG-S05-smoothing-lemma-lifting]{D.19} with parameter \(\pi^4\) and algebra \(\bar C\). It constructs \(D/R\) finitely presented, a map \(D\to\Lambda\), and commuting maps \[
\bar C/\pi^4\bar C\longrightarrow D/\pi^4D
\longrightarrow\Lambda/\pi^4\Lambda.
\] It proves smoothness of \(D/R\) away from \(\pi\), and also at a prime over \(\pi\) if its contraction to \(\bar C\) is smooth over \(R/\pi^8R\). Compose the reduced map from \(A\) with this diagram. The resulting \(A/\pi^4A\to D/\pi^4D\) agrees with the prescribed maps to \(\Lambda/\pi^4\Lambda\). Therefore \hyperref[AG-RG-S05-smoothing-lemma-desingularize-strictly-standard]{D.21} constructs compatible maps \(A\to B\), \(D\to B\), \(B\to\Lambda\), with \(B/R\) finitely presented and \[
H_{D/R}B\subset H_{B/R}.
\tag{D.22a}
\] If a prime \(\mathfrak b\subset B\) avoids \(\pi\), its contraction to \(D\) is smooth and hence avoids \(H_{D/R}\). Inclusion (D.22a) makes \(\mathfrak b\) avoid \(H_{B/R}\) as well, so \(B/R\) is smooth at that prime by \hyperref[AG-RG-S05-smoothing-definition-singular-ideal]{D.1}. Its smooth principal neighbourhood can be localized further at \(\pi\) by \hyperref[AG-RG-S05-algebra-lemma-locally-smooth]{F.67}. Viewed over \(R_\pi\), it is smooth: a polynomial presentation over \(R\) in which \(\pi\) is invertible localizes to the same presentation over \(R_\pi\). Exact localization \href{https://kokunoyumeto.github.io/open-math-courses/courses/AG-RG/AG-RG-S01.html}{S01, Lemma 1.1} identifies its kernel ideal and conormal quotient with their localizations. Since \(\pi\) is already invertible on the presented algebra, its two conormal terms and their differential are unchanged. Definition \hyperref[AG-RG-S05-algebra-definition-smooth]{F.6} and comparison \hyperref[AG-RG-S05-algebra-lemma-NL-homotopy]{F.11} give the assertion. Every prime of \(B_\pi\) has such a neighbourhood, and F.67 gives global smoothness over \(R_\pi\).

Set \(J=H_{B/R}\Lambda\). For each prime \(\mathfrak l\subset\Lambda\) containing \(J\), let \(\mathfrak b,\mathfrak d\) be its contractions to \(B,D\). They contain \(H_{B/R},H_{D/R}\), respectively, so both contracted points are nonsmooth. The preceding conclusion forces \(\pi\in\mathfrak l\). The prime \(\mathfrak d/\pi^4D\) now contracts along the quotient diagram to a prime \(\mathfrak c\) of \(\bar C\). If this were smooth, D.19 would make \(\mathfrak d\) smooth, a contradiction. Consequently \(H_{\bar C/(R/\pi^8R)}\subset\mathfrak c\).

The contraction \(\mathfrak c\) is also the inverse image of \(\mathfrak l/\pi^8\Lambda\) under \(\bar C\to\Lambda/\pi^8\Lambda\). To check this equality, reduce that map modulo \(\pi^4\) and use the commuting diagram through \(D/\pi^4D\); both quotient primes are the contractions of the same \(\mathfrak l\), which contains \(\pi\). Thus the ideal generated by the images of \(H_{\bar C/(R/\pi^8R)}\) lies in every such \(\mathfrak l/\pi^8\Lambda\).

For an explicit radical argument, let \(K\subset\Lambda\) be the inverse image of that generated ideal. The preceding paragraph and \(\pi\in\mathfrak l\) give \(K\subset\mathfrak l\) for every prime containing \(J\). The prime-intersection proof in D.1 yields \(K\subset\sqrt J\), as well as \(\pi\in\sqrt J\). In particular \(\pi^8\Lambda\subset\sqrt J\); taking the image of \(K\subset\sqrt J\) in \(\Lambda/\pi^8\Lambda\) is precisely the claimed inclusion. If no prime contains \(J\), D.1 gives \(\sqrt J=\Lambda\), and the same argument holds. This proves the result for an arbitrary target \(\Lambda\), including zero quotients. ∎

\subsection{Warmup: reduction to a base field}\label{AG-RG-S05-warmup-reduction-to-a-base-field}

\phantomsection\label{AG-RG-S05-smoothing-section-reduction}

In this section we apply the lemmas in the previous sections to prove that it suffices to prove the main result when the base ring is a field, see Lemma \hyperref[AG-RG-S05-smoothing-lemma-reduce-to-field]{D.26}.

\phantomsection\label{AG-RG-S05-smoothing-situation-global}

\subsubsection{Situation D.23: global}\label{AG-RG-S05-situation-d.23-global}

\emph{Adapted from the Stacks project, smoothing.tex, lines 1915--1918.}
\nopagebreak[4]

Here \(R \to \Lambda\) is a regular ring map of Noetherian rings.

Let \(R \to \Lambda\) be as in Situation \hyperref[AG-RG-S05-smoothing-situation-global]{D.23}. We say \emph{PT holds for \(R \to \Lambda\)} if \(\Lambda\) is a filtered colimit of smooth \(R\)-algebras.

\phantomsection\label{AG-RG-S05-smoothing-lemma-product}

\subsubsection{Lemma D.24: product}\label{AG-RG-S05-lemma-d.24-product}

Reading source: \href{https://stacks.math.columbia.edu/tag/07F3}{Stacks, Tag 07F3}.

For \(i=1,2\), let \(R_i\to\Lambda_i\) be a regular map of Noetherian rings as in \hyperref[AG-RG-S05-smoothing-situation-global]{D.23}. If each \(\Lambda_i\) is a filtered colimit of smooth \(R_i\)-algebras, then \(\Lambda_1\times\Lambda_2\) is a filtered colimit of smooth \(R_1\times R_2\)-algebras.

\textbf{Proof.} Put \(R=R_1\times R_2\), with orthogonal idempotents \(e_1=(1,0)\), \(e_2=(0,1)\). For any \(R\)-module \(M\), the maps \[
M\longrightarrow e_1M\times e_2M,\quad m\longmapsto(e_1m,e_2m),
\qquad (m_1,m_2)\longmapsto m_1+m_2
\] are inverse, since \(e_1+e_2=1\) and \(e_1e_2=0\). Tensor products split by the same rule: \[
M\otimes_R(\Lambda_1\times\Lambda_2)
\cong(e_1M\otimes_{R_1}\Lambda_1)\times
      (e_2M\otimes_{R_2}\Lambda_2).
\] The maps send a pure tensor to its two component tensors; summing the inclusions gives the inverse, and cross tensors vanish because the idempotent can be moved to either factor. This verifies the formula by \href{https://kokunoyumeto.github.io/open-math-courses/courses/AG-RG/AG-RG-S01.html}{S01, Lemma 1.0}. Exact sequences split into their two component sequences, so component flatness proves flatness of the product map.

Every ideal in a finite product splits in this fashion into two ideals. Finite generator lists for those component ideals give a finite generator list for the original ideal. Hence \(R\) and \(\Lambda_1\times\Lambda_2\) are Noetherian. A prime contains exactly one of \(e_1,e_2\): it contains at least one since their product is zero, and cannot contain both since their sum is one. It therefore comes from exactly one component, with that component's residue field. Tensoring with this residue field kills the other component and identifies the fibre with the corresponding fibre of \(R_i\to\Lambda_i\). Geometric regularity of these fibres proves regularity of the product map, so it is again in Situation D.23.

We prove smoothness of each product algebra used below. Let \(C_i/R_i\) be smooth and let \(C=C_1\times C_2\). Then \(C_{e_i}=C_i\), with inverse maps given by projection and inclusion of the selected component; the other component vanishes when \(e_i\) is made invertible. Also \(R_{e_i}=R_i\). Indeed inverting an idempotent forces it to equal one, so this localization is the quotient \(R/(1-e_i)\). The inverse-variable presentation \(R[T]/(e_iT-1)\), whose derivative \(e_i\) is a unit, proves \(R_i/R\) smooth by \hyperref[AG-RG-S05-algebra-lemma-standard-smooth]{F.85}. Thus \hyperref[AG-RG-S05-algebra-lemma-compose-smooth]{F.29} makes \(C_i/R\) smooth. The two elements \(e_i\) generate one in \(C\), so \hyperref[AG-RG-S05-algebra-lemma-cover-upstairs]{F.35} gives finite presentation of \(C/R\) from that of the \(C_{e_i}/R\). At every prime one \(e_i\) is avoided; the corresponding localization is smooth over \(R\). The full gluing result \hyperref[AG-RG-S05-algebra-lemma-locally-smooth]{F.67} therefore proves \(C/R\) smooth.

Choose the assumed presentations \[
\Lambda_i=\operatorname*{colim}_{a\in I_i}C_{i,a}.
\] The category \(I_1\times I_2\) is filtered: it is nonempty, successors of two objects are chosen in each coordinate, and parallel arrows are equalized in each coordinate. The preceding argument makes every \(C_{1,a}\times C_{2,b}\) smooth over \(R\). Component maps induce \[
\operatorname*{colim}_{(a,b)\in I_1\times I_2}
(C_{1,a}\times C_{2,b})\longrightarrow\Lambda_1\times\Lambda_2.
\] Any pair on the right has representatives at a pair of stages, proving surjectivity. If a represented pair maps to zero, each component becomes zero at a later stage by the filtered equality criterion. The pair of these stages kills the represented pair, proving injectivity. The same criterion applied to differences proves equality detection, and the maps preserve sums, products and units componentwise. The displayed map is consequently an \(R\)-algebra isomorphism. Iterating gives the same conclusion for any finite nonempty product; zero components obey these identical idempotent calculations. ∎

\phantomsection\label{AG-RG-S05-smoothing-lemma-delocalize-base}

\subsubsection{Lemma D.25: delocalize base}\label{AG-RG-S05-lemma-d.25-delocalize-base}

Reading source: \href{https://stacks.math.columbia.edu/tag/07F4}{Stacks, Tag 07F4}.

Let \(R\to A\to\Lambda\) be ring maps with \(A/R\) finitely presented, and let \(S\subset R\) be multiplicative. Suppose \(S^{-1}A\to B'\to S^{-1}\Lambda\) is a factorization with \(B'/S^{-1}R\) smooth. Then the original map factors as \(A\to B\to\Lambda\), with \(B/R\) finitely presented and the image of some \(s\in S\) elementary standard in \(B/R\) in the sense of \hyperref[AG-RG-S05-smoothing-definition-strictly-standard]{D.3}.

\textbf{Proof.} If \(0\in S\), take \(B=A\) and \(s=0\). In its finite presentation choose \(c=0\): the empty minor is one with coefficient zero, and multiplying every extra equation by zero gives zero. This verifies both elementary standard conditions. Hence assume \(0\notin S\).

The retraction \hyperref[AG-RG-S05-smoothing-lemma-smooth-standard-smooth]{D.12} replaces \(B'\) by a standard smooth algebra: compose the original map into its retract inclusion, and compose its retraction with the original target map. These composites preserve the original factorization. Write \(A=R[x_1,\ldots,x_n]/(g_1,\ldots,g_t)\), with coordinate images \(\lambda_i\in\Lambda\). The graph-equation construction \hyperref[AG-RG-S05-smoothing-lemma-standard-smooth-include-generators]{D.14}, applied to this ordered list of images, gives \[
B'=S^{-1}R[x_1,\ldots,x_{n+m}]/(f_1,\ldots,f_c),
\qquad n\le c\le n+m,
\] with the first \(n\) variables representing the given images and with the minor on its first \(c\) variables invertible. Repeated images are permitted: each graph equation has its own variable and the same identity Jacobian block.

For \(i>n\), write the image of the old coordinate \(x_i\) in \(S^{-1}\Lambda\) as \(\lambda_i/d_i\), with \(d_i\in S\). Introduce the new coordinate \(y_i=d_ix_i\), leaving \(y_i=x_i\) for \(i\le n\). Express each old equation in the new coordinates by substituting \(x_i=d_i^{-1}y_i\). The reverse substitution is \(y_i=d_ix_i\), so these changes give inverse polynomial-algebra maps over \(S^{-1}R\). The new coordinates map to \(\lambda_i/1\). Differentiating the substituted equations multiplies the column for an added selected variable by \(d_i^{-1}\); the determinant remains a unit. Clear the finite coefficient denominators in each equation by multiplying it by \(u_j\in S\). This does not change the localized relation ideal and multiplies the determinant by the unit \(\prod_j u_j\). Rename the new coordinates \(x_i\) and the scaled equations \(f_j\). Now \[
P=R[x_1,\ldots,x_{n+m}],\qquad f_j\in P,\qquad
\Delta=\det(\partial f_j/\partial x_i)_{1\le j,i\le c}
\] has invertible image in \(B'\).

A polynomial representative of its inverse gives \(b_0,\ldots,b_c\in S^{-1}P\) with \(1=b_0\Delta+\sum_j b_jf_j\). Choose \(d\in S\) and \(c_j\in P\) representing all \(db_j\). The polynomial \(d-c_0\Delta-\sum_jc_jf_j\) is zero after localization; some \(e\in S\) therefore annihilates it in \(P\). With \(s_0=ed\) and \(a_j=ec_j\), we have the exact identity \[
s_0=a_0\Delta+\sum_{j=1}^c a_jf_j\quad\text{in }P.
\tag{D.25a}
\] The annihilating multiplier is required even when \(P\) has \(S\)-torsion; no element of \(S\) has been cancelled in \(P\).

Each original relation \(g_k\) is zero in \(S^{-1}(P/(f_1,\ldots,f_c))\), so choose \(s_k\in S\) and \(v_{kj}\in P\) with \(s_kg_k=\sum_jv_{kj}f_j\). Each \(f_j(\lambda)\) is zero in \(S^{-1}\Lambda\), so choose \(t_j\in S\) with \(t_jf_j(\lambda)=0\) in \(\Lambda\). Set \[
h_j=t_jf_j,\quad
B=P/(h_1,\ldots,h_c,g_1,\ldots,g_t),\quad
s=s_0\prod_{k=1}^t s_k\prod_{j=1}^c t_j.
\] These finite relations give finite presentation. The original coordinates define \(A\to B\); evaluation of all coordinates at \(\lambda_i\) defines \(B\to\Lambda\), since it kills both relation lists. Their composite is the given map.

Here is the full witness for elementary standardness. For each \(k\), multiply its relation identity by \(s_0\prod_{\ell\ne k}s_\ell\prod_jt_j\). Replacing \(t_jf_j\) by \(h_j\) gives \[
sg_k=\sum_{j=1}^c
s_0\Bigl(\prod_{\ell\ne k}s_\ell\Bigr)
\Bigl(\prod_{q\ne j}t_q\Bigr)v_{kj}h_j
\quad\text{in }P.
\] Thus \(sg_k\in(h_1,\ldots,h_c)\), which implies D.3's extra-relation condition even without its allowed \(I^2\) term. Multiplication of (D.25a) by \(\prod_jt_j\) gives \[
s_0\prod_jt_j
=a_0\det(\partial h_j/\partial x_i)_{1\le j,i\le c}
+\sum_{j=1}^c a_j\Bigl(\prod_{q\ne j}t_q\Bigr)h_j.
\] The determinant equality follows by factoring \(t_j\), a base scalar, from equation row \(j\). In \(B\), the second sum is zero. Multiplying by \(\prod_ks_k\) expresses the image of \(s\) as this selected minor times the image of \(a_0\prod_ks_k\), exactly D.3's other condition. Every product with an empty index set is one. This proves the assertion also for empty equation lists and zero target algebras. ∎

\phantomsection\label{AG-RG-S05-smoothing-lemma-reduce-to-field}

\subsubsection{Lemma D.26: reduce to field}\label{AG-RG-S05-lemma-d.26-reduce-to-field}

Reading source: \href{https://stacks.math.columbia.edu/tag/07F5}{Stacks, Tag 07F5}.

If PT holds for every regular map \(K\to\Lambda\) of Noetherian rings with \(K\) a field, then PT holds for every regular map \(R\to\Lambda\) of Noetherian rings, as in \hyperref[AG-RG-S05-smoothing-situation-global]{D.23}.

\textbf{Proof.} Suppose a map fails PT. Each quotient \(R/I\to\Lambda/I\Lambda\) is regular by \hyperref[AG-RG-S05-more-algebra-lemma-regular-base-change]{F.123}; its rings are Noetherian because ideals in a quotient lift to ideals in the original ring. The set of failing quotient ideals is nonempty and has a maximal member by the ascending-chain condition. Replace the map by that quotient. Hence every further nonzero ideal gives PT. The new \(R\) is nonzero: a unital map from the zero ring has zero target, which is already a smooth algebra over itself and a constant filtered colimit.

The nilradical \(N\) is finitely generated because \(R\) is Noetherian. If its generators satisfy \(a_i^{n_i}=0\), every monomial of degree \(1+\sum_i(n_i-1)\) contains some \(a_i^{n_i}\), so that power of \(N\) is zero. If \(N\ne0\), maximality gives PT after quotienting by \(N\), and flatness of the original regular map allows \hyperref[AG-RG-S05-smoothing-proposition-lift]{D.18} to lift PT back to \(R\). This is a contradiction. Thus \(R\) is reduced.

We supply the finite-prime decomposition needed for its total fraction ring. Every radical ideal of a Noetherian ring is a finite intersection of prime ideals. Otherwise take a maximal counterexample \(I\) by the ascending-chain condition. It is proper and not prime, so there are \(a,b\notin I\) with \(ab\in I\). Both \(\sqrt{I+(a)}\) and \(\sqrt{I+(b)}\) are strictly larger and hence finite intersections of primes (the unit ideal is the empty intersection). Their intersection is \(I\): if \(x\) lies in both, then \(x^m\in I+(a)\) and \(x^n\in I+(b)\) for some positive \(m,n\), whence \(x^{m+n}\in I+(ab)=I\); radicalness gives \(x\in I\). Combining the two finite lists contradicts the choice of \(I\).

Apply this to the radical zero ideal of the nonzero reduced ring, and remove repeated primes and any prime containing another in the list. We obtain \[
0=\bigcap_{i=1}^r\mathfrak p_i
\] with a nonempty, pairwise incomparable list. These are exactly the minimal primes. Indeed the product of the listed ideals is contained in their intersection, so any prime over zero contains at least one listed prime: otherwise choose an element from each \(\mathfrak p_i\) outside it, whose product contradicts primality. If a prime lies strictly inside \(\mathfrak p_i\), the same observation puts a listed \(\mathfrak p_j\) inside it, contradicting incomparability. This proves finiteness, minimality and the zero intersection directly.

For each \(i\), choose \[
y_i\in\bigcap_{j\ne i}\mathfrak p_j\setminus\mathfrak p_i
\] by multiplying one element of \(\mathfrak p_j\setminus\mathfrak p_i\) for each \(j\ne i\); for \(r=1\) take \(y_1=1\). An element outside all \(\mathfrak p_i\) is a nonzerodivisor: test an annihilated product in the domains \(R/\mathfrak p_i\), then use their zero intersection. Conversely \(a\in\mathfrak p_i\) satisfies \(ay_i=0\), while \(y_i\ne0\), and is a zerodivisor. Hence the multiplicative set of nonzerodivisors is \(S=R\setminus\bigcup_i\mathfrak p_i\).

Let \(d=\sum_i y_i\in S\), \(Q=S^{-1}R\), and \(e_i=y_i/d\). Fractions define an injection \[
Q\longrightarrow\prod_i\operatorname{Frac}(R/\mathfrak p_i):
\] the denominators are nonzero in every component, and a fraction in its kernel has numerator in their zero intersection. In component \(i\), \(e_i\) equals one; in all other components it equals zero. Therefore \(e_i^2=e_i\), \(e_ie_j=0\) for \(i\ne j\), and \(\sum_i e_i=1\). Projection by these idempotents and summation give inverse maps \(Q\cong\prod_i e_iQ\).

Each \(K_i=e_iQ\) is a field. A nonzero \(e_i a/s\) has \(a\notin\mathfrak p_i\). The element \[
u=e_i a+(1-e_i)
 =\frac{y_i a+\sum_{j\ne i}y_j}{d}
\] is a unit of \(Q\): its numerator avoids \(\mathfrak p_i\) because \(y_i a\) does, and avoids every other \(\mathfrak p_j\) because only \(y_j\) survives there. Its numerator is thus in \(S\), giving the fraction inverse. The inverse of \(e_i a/s\) inside \(e_iQ\) is \(e_i s u^{-1}\). This proves the claimed field decomposition with actual inverse ring maps.

The localization \(Q\to\Lambda_Q=S^{-1}\Lambda\) is a regular map of Noetherian rings. To verify Noetherianness of either localization, extend the contraction of an ideal: each fraction in the ideal has its numerator in that contraction, and finite generators of the contraction generate the localized ideal. For a \(Q\)-module \(M\), the tensor maps of \href{https://kokunoyumeto.github.io/open-math-courses/courses/AG-RG/AG-RG-S01.html}{S01, Lemma 1.0} give \[
M\otimes_Q\Lambda_Q\cong M\otimes_R\Lambda.
\] Here elements of \(S\) are already invertible on \(M\), so the maps on pure tensors and their inverse respect the fraction relations. Exactness of the right side gives flatness. For a prime \(\mathfrak q\subset Q\), its contraction \(\mathfrak p\subset R\) avoids \(S\), and \(\mathfrak q=S^{-1}\mathfrak p\): a fraction belongs to \(\mathfrak q\) exactly when its numerator does, because its denominator is a unit. Conversely a prime disjoint from \(S\) extends to a prime, since \(S^{-1}(R/\mathfrak p)\) is a domain. The quotient \(Q/\mathfrak q\) is this localized domain and its fraction field is \(\operatorname{Frac}(R/\mathfrak p)\), so its residue field equals \(\kappa(\mathfrak p)\). Tensor associativity from S01, Lemma 1.0 identifies its fibre with \(\Lambda\otimes_R\kappa(\mathfrak p)\), geometrically regular by hypothesis. The idempotents split \(\Lambda_Q=\prod_i e_i\Lambda_Q\). Each component \(K_i\to e_i\Lambda_Q\) is flat by the same component tensor calculation in \hyperref[AG-RG-S05-smoothing-lemma-product]{D.24}, and has the corresponding regular fibre. Its rings are Noetherian, since it is a quotient of the respective localization. The assumed field case and the finite-product conclusion of D.24 give PT for \(Q\to\Lambda_Q\).

Take any finitely presented \(R\)-algebra \(A\to\Lambda\). Localizing its finite generator and relation lists gives a finitely presented \(Q\)-algebra. Criterion \hyperref[AG-RG-S05-algebra-lemma-when-colimit]{F.90} factors it through a smooth \(Q\)-algebra. By \hyperref[AG-RG-S05-smoothing-lemma-delocalize-base]{D.25} there is \(A\to A_1\to\Lambda\), with \(A_1/R\) finitely presented and \(\pi\in S\) elementary standard in \(A_1/R\). It is strictly standard by \hyperref[AG-RG-S05-smoothing-definition-strictly-standard]{D.3}: put its elementary coefficient at the selected variable set and put all other minor coefficients equal to zero.

Since \(R\ne0\) and multiplication by \(\pi\) is injective, \(\pi^8\ne0\). Maximality of the failing quotient therefore gives PT for \(R/\pi^8R\to\Lambda/\pi^8\Lambda\), including when this quotient is zero. Apply F.90 to obtain \[
A_1/\pi^8A_1\longrightarrow\bar C\longrightarrow\Lambda/\pi^8\Lambda
\] with \(\bar C/(R/\pi^8R)\) smooth. All powers of \(\pi\) act injectively on \(R\), and flatness makes them injective on \(\Lambda\) by tensoring \(0\to R\xrightarrow{\pi^n}R\), using S01, Lemma 1.0. Thus both annihilator equalities hold for parameters \(\pi\) and \(\pi^4\).

Apply \hyperref[AG-RG-S05-smoothing-lemma-lifting]{D.19} with parameter \(\pi^4\). It supplies \(D/R\) finitely presented, \(D\to\Lambda\), and the compatible map \[
A_1/\pi^4A_1\longrightarrow\bar C/\pi^4\bar C
\longrightarrow D/\pi^4D\longrightarrow\Lambda/\pi^4\Lambda.
\] Every prime of \(D\) avoiding \(\pi\) is smooth by that lemma. Every prime containing \(\pi\) contracts to a smooth point of \(\bar C\), since the whole map \(\bar C/(R/\pi^8R)\) is smooth, and is therefore smooth as well. The actual neighbourhood criterion \hyperref[AG-RG-S05-algebra-lemma-locally-smooth]{F.67} makes \(D/R\) globally smooth.

Finally \hyperref[AG-RG-S05-smoothing-lemma-desingularize-strictly-standard]{D.21}, applied to \(A_1,D,\pi\) and the displayed quotient map, constructs \(A_1\to B\to\Lambda\), with \(B/R\) finitely presented and \(H_{D/R}B\subset H_{B/R}\). Definition \hyperref[AG-RG-S05-smoothing-definition-singular-ideal]{D.1} gives \(H_{D/R}=D\), so \(H_{B/R}=B\). Thus every prime of \(B\) is smooth, and F.67 makes \(B/R\) smooth. Composing with \(A\to A_1\) gives the required smooth factorization of every finitely presented input. F.90 now gives PT for the supposedly failing map, a contradiction. Consequently no failing quotient exists and PT holds for the original map. ∎

\subsection{Local tricks}\label{AG-RG-S05-local-tricks}

\phantomsection\label{AG-RG-S05-smoothing-section-local}

\phantomsection\label{AG-RG-S05-smoothing-situation-local}

\subsubsection{Situation D.27: local}\label{AG-RG-S05-situation-d.27-local}

Reading source: \href{https://stacks.math.columbia.edu/tag/07F7}{Stacks, Tag 07F7}.

Fix a Noetherian ring \(R\), an \(R\)-algebra \(A\) of finite presentation, an \(R\)-algebra map \(A\to\Lambda\), and a prime \(\mathfrak q\) of \(\Lambda\). The target need not be Noetherian. For any intermediate \(R\)-algebra \(T\) with a specified map to \(\Lambda\), write \[
\mathfrak h_T=\sqrt{H_{T/R}\Lambda},
\] where \(H_{T/R}\) is the singular ideal of \hyperref[AG-RG-S05-smoothing-definition-singular-ideal]{D.1} and its extension is generated by its images in \(\Lambda\).

A \emph{resolution at \(\mathfrak q\)} is a factorization \(A\to B\to\Lambda\) through an \(R\)-algebra \(B\) of finite presentation for which \[
\mathfrak h_A\subset\mathfrak h_B,
\qquad \mathfrak h_B\not\subset\mathfrak q.
\] We say that the specified local problem is \emph{resolvable} when such a factorization exists. Both ideals in this condition belong to the same target ring \(\Lambda\); finite presentation of \(B\) is over the fixed base \(R\). The second condition means that some element of \(\mathfrak h_B\) lies outside \(\mathfrak q\), while the first requires preservation of the entire original singular radical.

\phantomsection\label{AG-RG-S05-smoothing-lemma-lift-solution}

\subsubsection{Lemma D.28: lift solution}\label{AG-RG-S05-lemma-d.28-lift-solution}

Reading source: \href{https://stacks.math.columbia.edu/tag/07F8}{Stacks, Tag 07F8}.

Work in \hyperref[AG-RG-S05-smoothing-situation-local]{D.27}. Let \(r\ge1\), and choose \(\pi_1,\ldots,\pi_r\in R\) whose images lie in \(\mathfrak q\). For each \(i\), suppose that \(\pi_i\) is strictly standard in \(A/R\), and that multiplication by \(\pi_i\) and by \(\pi_i^2\) has the same kernel on each of \[
R/(\pi_1^8,\ldots,\pi_{i-1}^8)R,
\qquad
\Lambda/(\pi_1^8,\ldots,\pi_{i-1}^8)\Lambda.
\] If the problem obtained by quotienting \(R,A,\Lambda\) by all \(\pi_i^8\), at the image of \(\mathfrak q\), is resolvable, then the original problem is resolvable.

\textbf{Proof.} First establish a base-change inclusion used in the single-element step. If \(R\to R'\) is any map and \(A'=A\otimes_R R'\), then \[
H_{A/R}A'\subset H_{A'/R'}.
\tag{D.28a}
\] Indeed, a prime of \(A'\) whose contraction is smooth has a standard smooth principal neighbourhood after taking a smooth neighbourhood of that contraction and applying \hyperref[AG-RG-S05-algebra-lemma-smooth-syntomic]{F.83}. The coefficient base-change construction \hyperref[AG-RG-S05-algebra-lemma-standard-smooth]{F.85(5)} makes its inverse image standard smooth over \(R'\). Thus every nonsmooth prime of \(A'\) contracts to a nonsmooth prime of \(A\), and contains the images of \(H_{A/R}\). The prime-intersection verification in D.1 gives (D.28a), including an empty nonsmooth locus. No flatness was required.

Consider \(r=1\), and abbreviate \(\pi=\pi_1\). Put \(\bar R=R/\pi^8R\), \(\bar A=A/\pi^8A\), \(\bar\Lambda=\Lambda/\pi^8\Lambda\), and \(\bar{\mathfrak q}=\mathfrak q/\pi^8\Lambda\). The assumed resolution provides \(\bar A\to\bar C\to\bar\Lambda\), with \(\bar C/\bar R\) finitely presented, such that \[
\sqrt{H_{\bar A/\bar R}\bar\Lambda}
\subset\sqrt{H_{\bar C/\bar R}\bar\Lambda}
\not\subset\bar{\mathfrak q}.
\tag{D.28b}
\] The two kernel equalities and strict standardness are exactly the hypotheses of \hyperref[AG-RG-S05-smoothing-lemma-desingularize-lifting-apply]{D.22}. Let \(A\to B\to\Lambda\) be its factorization, and set \(\mathfrak h_B=\sqrt{H_{B/R}\Lambda}\). The proof of D.22 shows that every prime containing \(H_{B/R}\Lambda\) contains \(\pi\); D.1 consequently gives \(\pi\in\mathfrak h_B\). Its other conclusion is \[
H_{\bar C/\bar R}\bar\Lambda
\subset \mathfrak h_B/\pi^8\Lambda.
\tag{D.28c}
\] The ideal on the right is radical. To check this directly, lift a class whose power belongs to it to \(x\in\Lambda\). Then \(x^a\in\mathfrak h_B+\pi^8\Lambda=\mathfrak h_B\), so \(x\in\mathfrak h_B\). Taking radicals in (D.28c) is therefore legitimate.

Write \(\rho:\Lambda\to\bar\Lambda\) for the quotient. Inclusion (D.28a) gives \[
\rho(\mathfrak h_A)
\subset\sqrt{H_{\bar A/\bar R}\bar\Lambda}
\subset\mathfrak h_B/\pi^8\Lambda:
\] an element of \(\mathfrak h_A\) has a power in \(H_{A/R}\Lambda\), whose image belongs to \(H_{\bar A/\bar R}\bar\Lambda\). Since \(\ker\rho=\pi^8\Lambda\subset\mathfrak h_B\), this proves \(\mathfrak h_A\subset\mathfrak h_B\). If \(\mathfrak h_B\subset\mathfrak q\), then (D.28c), after taking radicals, would put \(\sqrt{H_{\bar C/\bar R}\bar\Lambda}\) inside \(\bar{\mathfrak q}\), contradicting (D.28b). Thus \(\mathfrak h_B\not\subset\mathfrak q\). D.22 supplies the required finite presentation of \(B/R\), so this is a resolution with both radical conditions verified.

For \(r>1\), quotient first by \(\pi_1^8\). The base \(R'=R/\pi_1^8R\) is Noetherian: any ascending chain of its ideals lifts to one in \(R\). The finite polynomial presentation of \(A\) changes coefficients to give the finite presentation of \(A'=A/\pi_1^8A\) over \(R'\). The image \(\mathfrak q'\) is prime because \(\pi_1^8\Lambda\subset\mathfrak q\). For the remaining elements \(\pi_2,\ldots,\pi_r\), each successive quotient is precisely the original quotient by \(\pi_1^8,\ldots,\pi_{i-1}^8\); hence its two kernel hypotheses are unchanged. Strict standardness survives by \hyperref[AG-RG-S05-smoothing-lemma-strictly-standard-base-change]{D.7}. The induction hypothesis gives a resolution of this first quotient problem. Apply the single-element step just proved to \(\pi_1\) to lift it to the original problem. This completes the induction, with no Noetherian assumption on \(\Lambda\). ∎

\phantomsection\label{AG-RG-S05-smoothing-lemma-delocalize-weak}

\subsubsection{Lemma D.29: delocalize weak}\label{AG-RG-S05-lemma-d.29-delocalize-weak}

\emph{Adapted from the Stacks project, smoothing.tex, lines 2147--2227.}
\nopagebreak[4]

Let \(R \to A \to \Lambda \supset \mathfrak q\) be as in Situation \hyperref[AG-RG-S05-smoothing-situation-local]{D.27}. Let \(\mathfrak p = R \cap \mathfrak q\). Assume that \(\mathfrak q\) is minimal over \(\mathfrak h_A\) and that \(R_\mathfrak p \to A_\mathfrak p \to \Lambda_\mathfrak q
\supset \mathfrak q\Lambda_\mathfrak q\) can be resolved. Then there exists a factorization \(A \to C \to \Lambda\) with \(C\) of finite presentation such that \(H_{C/R} \Lambda \not \subset \mathfrak q\).

\textbf{Proof.} In the assumed resolution of the localized problem, the singular radical of the intermediate algebra is outside the maximal ideal of \(\Lambda_{\mathfrak q}\). Its extended singular ideal is therefore the unit ideal: if that ideal were contained in the maximal prime, its radical would be contained there too. \hyperref[AG-RG-S05-smoothing-lemma-final-solve]{D.8} replaces the intermediate algebra by a smooth one, and \hyperref[AG-RG-S05-smoothing-lemma-smooth-standard-smooth]{D.12} gives a standard smooth factorization, both compatible with the prescribed map from \(A\) to \(\Lambda_{\mathfrak q}\).

Write \(A=R[x_1,\ldots,x_n]/(g_1,\ldots,g_t)\), with \(x_i\mapsto\lambda_i\in\Lambda\). The identity-block construction \hyperref[AG-RG-S05-smoothing-lemma-standard-smooth-include-generators]{D.14} includes this prescribed generator list in a presentation of the standard smooth intermediate algebra: \[C=R_{\mathfrak p}[x_1,\ldots,x_n,u_1,\ldots,u_m]/(f_1,\ldots,f_c),\qquad c\ge n.\] Its selected \(c\)-variable Jacobian determinant \(\delta\) is a unit. The construction adjoins a separate equation for each member of the list, so repeated generator images cause no difficulty. Multiply each \(f_j\) by a denominator in \(R\setminus\mathfrak p\) to make its coefficients lie in \(R\). These factors are units over \(R_{\mathfrak p}\), and multiply \(\delta\) by their product, so it remains a unit. Put \(P=R[x,u]\).

Choose a polynomial representative of \(\delta^{-1}\) and of its identity with \(\delta\) modulo the equations over \(R_{\mathfrak p}\). Clear the finitely many coefficient denominators in that identity; equality of fractions supplies one further element outside \(\mathfrak p\) annihilating its error. Their product gives \(s_0\notin\mathfrak p\) for which \(\delta\) is a unit in \(P_{s_0}/(f_j)\). For each \(g_i\), its class vanishes in \((P/(f_j))_{\mathfrak p}\), so a localization witness \(s_i\notin\mathfrak p\) gives \(s_ig_i\in(f_j)\) in \(P\). With \(s=s_0\prod_i s_i\), the algebra \[C_s=P_s/(f_j)\] is standard smooth over \(R_s\), all \(g_i\) vanish there, and it is smooth over \(R\) by principal localization and smooth composition. No assertion that \(\delta\) fails to be a unit before this localization is needed.

Choose a common denominator \(h\in\Lambda\setminus\mathfrak q\) and numerators \(\nu_\ell\in\Lambda\) representing the target images of the finitely many \(u_\ell\) as \(\nu_\ell/h\). For an empty list take \(h=1\). Homogenize only these variables: for sufficiently large integers \(d_j\ge0\), put \[F_j(x,U,v)=v^{d_j}f_j(x,U/v)\in R[x,U,v].\] Every \(F_j(\lambda,\nu,h)\) is zero in \(\Lambda_{\mathfrak q}\). The product of their finite localization annihilator witnesses gives \(\lambda_0\notin\mathfrak q\) killing every one of them in \(\Lambda\).

Set \[B=R[w,x_1,\ldots,x_n,U_1,\ldots,U_m,v]/(g_i(x),\ wF_j(x,U,v)).\] Map it to \(\Lambda\) by \(w\mapsto\lambda_0\), \(x_i\mapsto\lambda_i\), \(U_\ell\mapsto\nu_\ell\), \(v\mapsto h\). Every equation vanishes exactly in \(\Lambda\), and the \(g_i\) relations give the desired factorization through \(A\). Let \(b=wvs\in B\). Its image is outside \(\mathfrak q\), because each of its three factors is outside that prime.

In \(B_b\) the factors \(w,v,s\) are units. Substitute \(u_\ell=U_\ell/v\), with inverse \(U_\ell=vu_\ell\). Each equation \(wF_j=0\) becomes \(f_j=0\), since \(wv^{d_j}\) is a unit. The equations \(g_i=0\) are already consequences of those equations in \(P_s\). These substitutions give inverse \(R\)-algebra maps and the exact isomorphism \[B_b\cong C_s[w,w^{-1},v,v^{-1}].\] This is smooth over \(R\) by polynomial extension, principal localization and F.29. Hence \(b\in H_{B/R}\), by D.1 and F.67. Since its image is outside \(\mathfrak q\), \(H_{B/R}\Lambda\not\subset\mathfrak q\). The displayed finite generator and equation lists make \(B\) finitely presented, proving the lemma without adding any Noetherian hypothesis on the target. ∎

\phantomsection\label{AG-RG-S05-smoothing-lemma-delocalize-height-zero}

\subsubsection{Lemma D.30: delocalize height zero}\label{AG-RG-S05-lemma-d.30-delocalize-height-zero}

\emph{Adapted from the Stacks project, smoothing.tex, lines 2229--2274.}
\nopagebreak[4]

Let \(R \to A \to \Lambda \supset \mathfrak q\) be as in Situation \hyperref[AG-RG-S05-smoothing-situation-local]{D.27}. Let \(\mathfrak p = R \cap \mathfrak q\). Assume

\begin{enumerate}
\def\labelenumi{\arabic{enumi}.}
\item
  \(\mathfrak q\) is minimal over \(\mathfrak h_A\),
\item
  \(R_\mathfrak p \to A_\mathfrak p \to \Lambda_\mathfrak q
  \supset \mathfrak q\Lambda_\mathfrak q\) can be resolved, and
\item
  \(\dim(\Lambda_\mathfrak q) = 0\).
\end{enumerate}

Then \(R \to A \to \Lambda \supset \mathfrak q\) can be resolved.

\textbf{Proof.} The Noetherian hypothesis is on \(R\); no Noetherian hypothesis on \(\Lambda\) is needed. By \href{https://kokunoyumeto.github.io/open-math-courses/courses/AG-CA/noetherian-and-artinian-rings.html}{AG-CA-03, Theorem 2.1}, the finite type \(R\)-algebra \(A\) is Noetherian. Choose \(H_{A/R}=(a_1,\ldots,a_r)\).

A zero-dimensional local ring has only its maximal prime: a strictly smaller prime would give a chain of length one. Thus the nilradical of \(\Lambda_{\mathfrak q}\) is \(\mathfrak q\Lambda_{\mathfrak q}\), by the prime-intersection argument in \hyperref[AG-RG-S05-smoothing-lemma-elkik]{D.5}. Since \(\mathfrak q\) contains \(\mathfrak h_A\), every \(a_i\) has nilpotent image in this localization. Choose one positive exponent \(N\) killing their finite list. The definition of localization gives \(\lambda_i\notin\mathfrak q\) with \(\lambda_i a_i^N=0\) in \(\Lambda\). Their product \(\lambda\notin\mathfrak q\) satisfies \[\lambda a_i^N=0\qquad(1\le i\le r).\] For an empty list choose \(N=1\) and \(\lambda=1\). This uses only these finitely many nilpotence witnesses; it does not assert that the local ring is Artinian or that its entire maximal ideal is nilpotent.

Apply \hyperref[AG-RG-S05-smoothing-lemma-delocalize-weak]{D.29} to the assumed localized resolution. It gives \(A\to C\to\Lambda\), with \(C\) finitely presented over \(R\) and \(\mathfrak h_C\not\subset\mathfrak q\). There is a finite \(A\)-presentation \(C=A[x_1,\ldots,x_n]/(f_1,\ldots,f_m)\): adjoining the generators of a finite \(R\)-presentation makes it a finite type \(A\)-algebra, and the polynomial ring over the Noetherian \(A\) has finite ideal kernels by the same Hilbert theorem. Choose \(c\in H_{C/R}\) whose image is outside \(\mathfrak q\). Such an element exists, since otherwise the radical of \(H_{C/R}\Lambda\) would be contained in this prime.

Define \[B=A[x_1,\ldots,x_n,y_1,\ldots,y_r,z,t_{ij}]/
\left(f_j-\sum_{i=1}^r y_i t_{ij},\ zy_i\right),\] with \(1\le j\le m\). This is finitely presented over \(R\). It maps to \(\Lambda\) by the given map on \(C\), by \(x_i\) going to the corresponding generator image, by \(y_i\mapsto a_i^N\), \(z\mapsto\lambda\) and \(t_{ij}\mapsto0\). The displayed annihilator equality makes every equation \(zy_i\) vanish, and the other equations reduce to \(f_j=0\).

Inverting \(z\) forces every \(y_i=0\), giving the exact elimination \[B_z\cong C[z,z^{-1},t_{ij}].\] Consequently \(B_{zc}\) is smooth over \(R\). Here and below a principal open lying in the smooth locus is globally smooth by \hyperref[AG-RG-S05-algebra-lemma-locally-smooth]{F.67}. For each \(\ell\), inverting \(y_\ell\) instead forces \(z=0\) and solves every equation uniquely for \(t_{\ell j}\): \[t_{\ell j}=y_\ell^{-1}\left(f_j-\sum_{i\ne\ell}y_i t_{ij}\right).\] Hence \[B_{y_\ell}\cong A[x_j,y_i,y_\ell^{-1},t_{ij}\ (i\ne\ell)],\] a polynomial extension and principal localization over \(A\). Since \(A_{a_\ell}\) is smooth over \(R\), \(B_{a_\ell y_\ell}\) is smooth over \(R\) as well. Thus \(zc\) and all \(a_\ell y_\ell\) belong to \(H_{B/R}\).

The image of \(zc\) in \(\Lambda\) is \(\lambda c\notin\mathfrak q\), so \(\mathfrak h_B\not\subset\mathfrak q\). The image of \(a_\ell y_\ell\) is \(a_\ell^{N+1}\), so each image \(a_\ell\) belongs to the radical \(\mathfrak h_B\). Taking radicals gives \(\mathfrak h_A\subset\mathfrak h_B\). The finite presentation and the two required radical conditions prove that \(A\to B\to\Lambda\) is the desired resolution. ∎

\subsection{Separable residue fields}\label{AG-RG-S05-separable-residue-fields}

\phantomsection\label{AG-RG-S05-smoothing-section-separable}

In this section we explain how to solve a local problem in the case of a separable residue field extension.

\phantomsection\label{AG-RG-S05-smoothing-lemma-ogoma}

\subsubsection{Lemma D.31: ogoma}\label{AG-RG-S05-lemma-d.31-ogoma}

Reading source: \href{https://stacks.math.columbia.edu/tag/07FC}{Stacks, Tag 07FC}.

Let \(A\) be Noetherian, \(M\) a finite \(A\)-module, \(S\subset A\) multiplicative, and \(\pi\in A\). Suppose that multiplication by \(\pi\) and \(\pi^2\) has equal kernels on \(S^{-1}M\). There is one \(s\in S\) such that, for every positive integer \(n\), \[
\ker(s^n\pi:M\to M)=\ker((s^n\pi)^2:M\to M).
\]

\textbf{Proof.} Let \[
N=\{m\in M:\pi^2m/1=0\text{ in }S^{-1}M\}.
\] It is a submodule of the finite module \(M\), so \href{https://kokunoyumeto.github.io/open-math-courses/courses/AG-RG/AG-RG-S01.html}{S01, Lemma 1.3} gives a finite generating list \(u_1,\ldots,u_t\). By the assumed equality of localized kernels, \(\pi u_j/1=0\) for every \(j\). The fraction-zero calculation of \href{https://kokunoyumeto.github.io/open-math-courses/courses/AG-RG/AG-RG-S01.html}{S01, Lemma 1.1} gives \(s_j\in S\) with \(s_j\pi u_j=0\) in \(M\). Put \(s=\prod_j s_j\), using \(s=1\) for an empty list. Multiplication by \(s\pi\) kills every generator and hence all of \(N\).

Fix \(n>0\) and suppose \((s^n\pi)^2m=0\). Since \(s\) is a unit after localization, this implies \(\pi^2m/1=0\), so \(m\in N\). Therefore \(s\pi m=0\), and multiplication by \(s^{n-1}\) gives \(s^n\pi m=0\). The reverse inclusion follows by applying \(s^n\pi\) a second time to an element of its kernel. The same \(s\) works for all \(n\). The argument also covers \(0\in S\) and \(M=0\), without cancelling \(s\) inside \(M\). ∎

\phantomsection\label{AG-RG-S05-smoothing-lemma-find-sequence}

\subsubsection{Lemma D.32: find sequence}\label{AG-RG-S05-lemma-d.32-find-sequence}

Reading source: \href{https://stacks.math.columbia.edu/tag/07FD}{Stacks, Tag 07FD}.

Let \(\Lambda\) be Noetherian, let \(I\subset\mathfrak q\) with \(\mathfrak q\) prime, and fix positive integers \(n,e\). Suppose that \(L=\Lambda_{\mathfrak q}\) is regular local of dimension \(d\) and that \(\mathfrak q^nL\subset IL\). There are \(\pi_1,\ldots,\pi_d\in\Lambda\) satisfying \[
(\pi_1,\ldots,\pi_d)L=\mathfrak qL,
\qquad \pi_i^n\in I,
\] and, for every \(1\le i\le d\), \[
\operatorname{Ann}_{\Lambda/(\pi_1^e,\ldots,\pi_{i-1}^e)\Lambda}(\pi_i)
=\operatorname{Ann}_{\Lambda/(\pi_1^e,\ldots,\pi_{i-1}^e)\Lambda}(\pi_i^2).
\] The prescribed exponents \(n,e\) are retained.

\textbf{Proof.} Set \(S=\Lambda\setminus\mathfrak q\). Regularity says that \(\mathfrak qL/(\mathfrak qL)^2\) has dimension \(d\). Lift a basis to \(\mathfrak qL\); finite-module Nakayama, \href{https://kokunoyumeto.github.io/open-math-courses/courses/AG-RG/AG-RG-S01.html}{S01, Lemma 1.2}, makes these \(d\) lifts generate \(\mathfrak qL\). Each lift is a fraction with numerator in \(\mathfrak q\) and denominator in \(S\). Multiplying each by its denominator replaces it by \(a_i\in\mathfrak q\) and preserves the minimal generating property, since these denominators are units in \(L\).

The containment \(\mathfrak q^nL\subset IL\) gives \(a_i^n/1=0\) in \(S^{-1}(\Lambda/I)\). Exact localization and the fraction-zero criterion \href{https://kokunoyumeto.github.io/open-math-courses/courses/AG-RG/AG-RG-S01.html}{S01, Lemma 1.1} provide \(v_i\in S\) with \(v_i a_i^n\in I\). Replace \(a_i\) by \(v_i a_i\). Its \(n\)th power is \(v_i^{n-1}(v_i a_i^n)\in I\), and the list still minimally generates the localized maximal ideal. No change of \(n\) was made. Later multiplication of an entry by an element of \(S\) also preserves these two properties.

Now choose the entries in order. Suppose \(\pi_1,\ldots,\pi_{i-1}\) have been fixed, with their required annihilator equalities, and keep the remaining entries as their current \(a_j\). The whole list is a minimal generating list of \(\mathfrak qL\), so it is regular by the full parameter proof \href{https://kokunoyumeto.github.io/open-math-courses/courses/AG-CA/regular-sequences-depth-and-cohen-macaulay-modules.html}{AG-CA-12, Theorem 6.1}. The power theorem \href{https://kokunoyumeto.github.io/open-math-courses/courses/AG-CA/regular-sequences-depth-and-cohen-macaulay-modules.html}{AG-CA-12, Proposition 1.3}, applied with exponents \(e\) on the first \(i-1\) entries and exponent one on the rest, therefore makes \(a_i\) injective on \[
L/(\pi_1^e,\ldots,\pi_{i-1}^e)L.
\] Its square is injective too, so their kernels are both zero.

Take the finite \(\Lambda\)-module \(M_i=\Lambda/(\pi_1^e,\ldots,\pi_{i-1}^e)\Lambda\). S01, Lemma 1.1 identifies its localization with the displayed quotient of \(L\). Apply \hyperref[AG-RG-S05-smoothing-lemma-ogoma]{D.31} to \(M_i,S,a_i\). It supplies \(w_i\in S\) for which multiplication by \(w_i a_i\) and its square has the same kernel on \(M_i\); choose the positive power one in that lemma and set \(\pi_i=w_i a_i\). This is exactly the required \(i\)th annihilator equality.

The earlier equalities involve only their own entry and the entries preceding it, none of which has changed. The \(n\)th-power condition remains true because \(I\) is an ideal. In \(L\), the unit \(w_i\) preserves the generating list and its successive regularity: the preceding power ideals stay the same, and multiplication by a unit preserves injectivity on every quotient. Thus the next step has the same hypotheses. After \(d\) steps the list satisfies all assertions. If \(d=0\), Nakayama gives \(\mathfrak qL=0\); the empty list already satisfies them, with both given exponents unchanged. ∎

\phantomsection\label{AG-RG-S05-smoothing-lemma-resolve-special}

\subsubsection{Lemma D.33: resolve special}\label{AG-RG-S05-lemma-d.33-resolve-special}

Reading source: \href{https://stacks.math.columbia.edu/tag/07FE}{Stacks, Tag 07FE}.

In \hyperref[AG-RG-S05-smoothing-situation-local]{D.27}, let the base be a field \(k\), let \(\Lambda\) be Noetherian, and let \(\mathfrak q\) be minimal over \(\mathfrak h_A\). Assume that \(L=\Lambda_{\mathfrak q}\) is regular local and that its residue field \(K=\kappa(\mathfrak q)\) is separable over \(k\). Then the problem has a resolution.

\textbf{Proof.} Write \(\mathfrak m=\mathfrak qL\) and \(d=\dim L\). If \(d=0\), regularity says that \(\mathfrak m/\mathfrak m^2=0\). The finite-module Nakayama lemma \hyperref[AG-RG-S05-algebra-lemma-NAK]{F.10} makes \(\mathfrak m=0\), so \(L=K\). By \hyperref[AG-RG-S05-algebra-lemma-colimit-syntomic]{F.23}, this field is a filtered colimit of smooth \(k\)-algebras. Finite presentation and \hyperref[AG-RG-S05-algebra-lemma-when-colimit]{F.90} factor \(A\to L\) through one such algebra. The resulting localized resolution returns to \(\Lambda\) by \hyperref[AG-RG-S05-smoothing-lemma-delocalize-height-zero]{D.30}. Henceforth \(d>0\).

Set \(I=H_{A/k}\Lambda\). Minimality gives \(\sqrt{IL}=\mathfrak m\). Choose finite generators of \(\mathfrak m\), and, for each generator, a positive power belonging to \(IL\). A monomial whose degree exceeds the sum of these powers minus the number of generators contains one of those powers. Consequently there is \(n>0\) with \(\mathfrak m^n\subset IL\). The Hilbert theorem \href{https://kokunoyumeto.github.io/open-math-courses/courses/AG-CA/noetherian-and-artinian-rings.html}{AG-CA-03, Theorem 2.1} makes \(A\) Noetherian; choose \(H_{A/k}=(a_1,\ldots,a_r)\).

Introduce the polynomial base \(P=k[x_1,\ldots,x_d]\) and the finitely presented algebra \[
B=A[x_1,\ldots,x_d,z_{ij}]/(x_i^n-\sum_j a_jz_{ij}).
\] On the chart where \(a_j\) is invertible, each equation eliminates \(z_{ij}\). This chart is a polynomial algebra over \(A_{a_j}[x_1,\ldots,x_d]\), and is smooth over \(P\). On inverting \(x_i\), the equation gives \(1=\sum_j a_jx_i^{-n}z_{ij}\); thus these \(a_j\)-charts cover \(B_{x_i}\). Smooth gluing \hyperref[AG-RG-S05-algebra-lemma-locally-smooth]{F.67} proves that \(B_{x_i}/P\) is smooth.

Apply \hyperref[AG-RG-S05-smoothing-lemma-improve-presentation]{D.9} over \(P\). It supplies \(B\to C\) and a retraction \(C\to B\), with each \(C_{x_i}/P\) smooth and its differential module finite free. Its symmetric-algebra construction is smooth also on every \(a_j\)-chart. The algebra \(C\) is finitely presented, since it is finite type over the Noetherian ring \(P\). By \hyperref[AG-RG-S05-smoothing-lemma-compare-standard]{D.15(e)}, choose a common \(c>0\) such that every \(x_i^c\) is strictly standard in \(C/P\), and put \(e=8c\).

Use \hyperref[AG-RG-S05-smoothing-lemma-find-sequence]{D.32} with the fixed \(n,e,I\). Obtain elements \(\pi_i\in\Lambda\) generating \(\mathfrak m\) in \(L\), with \[
\pi_i^n=\sum_j a_j\lambda_{ij}
\quad\hbox{in }\Lambda,
\] and equal kernels for \(\pi_i\) and \(\pi_i^2\) on each quotient by the earlier \(\pi_j^e\). Sending \(x_i\) to \(\pi_i\) and \(z_{ij}\) to \(\lambda_{ij}\) defines \(B\to\Lambda\); the retraction defines \(C\to\Lambda\).

It is enough to resolve this new problem over \(P\). In fact the smooth \(a_j\)-charts put every \(a_j\) in \(H_{C/P}\). A \(P\)-resolution through \(T\) therefore preserves \(\mathfrak h_A\). A smooth \(P\)-neighbourhood is smooth over \(k\), so \(H_{T/P}\subset H_{T/k}\). Finite presentation over \(P\) is finite presentation over \(k\), by adjoining the finitely many \(x_i\) to a presentation. These observations give both conditions of an original \(k\)-resolution.

For the new problem apply \hyperref[AG-RG-S05-smoothing-lemma-lift-solution]{D.28} to \(x_1^c,\ldots,x_d^c\). On the source quotients multiplication by a remaining polynomial variable is injective, as equality of polynomial coefficients shows. On a target quotient, equality \(\ker u=\ker u^2\) implies \(\ker u^s=\ker u\) for all \(s\ge1\): if \(u^{s+1}v=0\), applying the equality to \(u^{s-1}v\) gives \(u^sv=0\), and induction finishes the argument. Thus D.32 supplies the equal kernels for \(\pi_i^c\) and \(\pi_i^{2c}\). Since \((x_i^c)^8=x_i^e\), D.28 reduces the task to the quotient by these \(e\)-th powers.

Localize that quotient at its target prime. The base and target are \[
S=P_{(x_1,\ldots,x_d)}/(x_1^e,\ldots,x_d^e),
\qquad L_0=L/(\pi_1^e,\ldots,\pi_d^e).
\] The contracted prime is \((x_1,\ldots,x_d)\), because reduction of a polynomial in the \(\pi_i\) is its constant in \(k\subset K\). The polynomial local base is regular by iterating \href{https://kokunoyumeto.github.io/open-math-courses/courses/AG-CA/regular-local-rings.html}{AG-CA-14, Proposition 3.3} from the field. The regular local ring \(L\) has dimension \(d\), and its maximal ideal is generated by the \(d\) elements \(\pi_i\). The minimal-parameter theorem \href{https://kokunoyumeto.github.io/open-math-courses/courses/AG-CA/regular-sequences-depth-and-cohen-macaulay-modules.html}{AG-CA-12, Theorem 6.1} makes this list regular. The criterion \hyperref[AG-RG-S05-algebra-lemma-flat-over-regular]{F.46} gives flatness of \(P_{(x)}\to L\); quotient base change gives flatness of \(S\to L_0\).

The maximal ideals of \(S\) and \(L_0\) are nilpotent: every monomial of degree \(d(e-1)+1\) contains an \(e\)-th power. The closed fibre of \(S\to L_0\) is \(K/k\). F.23 expresses it as a smooth colimit, and the nilpotent lifting theorem \hyperref[AG-RG-S05-smoothing-proposition-lift]{D.18} does the same for \(L_0/S\). By F.90, the finitely presented algebra \(C\otimes_P S\) factors through a smooth \(S\)-algebra.

To return to the unlocalized quotient, let its target prime be \(\overline{\mathfrak q}\). Its target localization has dimension zero by the preceding nilpotence calculation. If the quotient input's singular radical escapes \(\overline{\mathfrak q}\), its identity factorization already resolves the problem. Otherwise \(\overline{\mathfrak q}\) is minimal over that radical, and D.30 delocalizes the smooth factorization. D.28 lifts the quotient resolution to \(P\to C\to\Lambda\). The change-of-base argument above then gives the required resolution over \(k\). ∎

\subsection{Inseparable residue fields}\label{AG-RG-S05-inseparable-residue-fields}

\phantomsection\label{AG-RG-S05-smoothing-section-inseparable}

In this section we explain how to solve a local problem in the case of an inseparable residue field extension.

\phantomsection\label{AG-RG-S05-smoothing-lemma-helper}

\subsubsection{Lemma D.34: helper}\label{AG-RG-S05-lemma-d.34-helper}

Reading source: \href{https://stacks.math.columbia.edu/tag/07FG}{Stacks, Tag 07FG}.

Let \(k\) be a field of characteristic \(p>0\), and let \((\Lambda,\mathfrak m,K)\) be an Artinian local \(k\)-algebra. If \(H_1(L_{K/k})\) is finite-dimensional over \(K\), then \(\Lambda\) is a directed union of Artinian local \(k\)-subalgebras \(A\), essentially of finite type over \(k\), such that \(A\to\Lambda\) is faithfully flat and \(\mathfrak m_A\Lambda=\mathfrak m\).

\textbf{Proof.} Keep the specified map \(i:k\to\Lambda\) throughout. Its reduction identifies \(k\) with a subfield of \(K\). By \href{https://kokunoyumeto.github.io/open-math-courses/courses/AG-CA/noetherian-and-artinian-rings.html}{AG-CA-03, Theorem 4.2}, choose generators \(u_1,\ldots,u_d\) of \(\mathfrak m\) and a least \(N\ge1\) with \(\mathfrak m^N=0\). We use the following internal prime-field lifting calculation. For a field \(E\) of characteristic \(p\), \hyperref[AG-RG-S05-more-algebra-lemma-p-basis]{F.119} proves the p-basis presentation \[
E=E^p[T_b:b\in\mathcal B]/(T_b^p-b^p:b\in\mathcal B).
\] In a square-zero test \(f:E\to C/J\), define a map on \(E^p\) by sending \(a^p\) to the pth power of any lift of \(f(a)\). The pth root in \(E\) is unique, and the value is independent of the lift because \(j^p=0\) for \(j\in J\). The Frobenius identities prove addition and multiplication for this map. Choose any lifts of the images of the basis elements \(b\). Their pth powers are exactly the prescribed values of \(b^p\), so the displayed presentation gives an \(\mathbb F_p\)-map \(E\to C\) lifting \(f\). This proves formal smoothness for every such field, including an infinite p-basis.

It lifts maps across a nilpotent ideal by its successive powers: the kernel from the quotient by \(J^{a+1}\) to the quotient by \(J^a\) is square zero for \(a\ge1\). In particular a residue section \(s:K\to\Lambda\) exists. Define \(\psi_s:K[X_1,\ldots,X_d]\to\Lambda\) by this section and \(X_j\mapsto u_j\).

We first prove that the section can be chosen so that, for a field \(k\subset F_0\subset K\) finitely generated over \(k\), the subring \(\psi_s(F_0[X])\) contains \(i(k)\). This is automatic for \(N=1\), with \(s\) the identity and \(F_0=k\). Suppose \(N>1\) and put \(I=\mathfrak m^{N-1}\). Then \(I^2=0\) and \(\mathfrak m I=0\), so scalar multiplication on \(I\) factors through \(K\). Applying the induction assertion to \(\Lambda/I\) gives a section \(s_0\) and a finitely generated field \(F\supset k\) containing the prescribed image in \(s_0(F)[\bar u]\). Lift \(s_0\) to \(s\). The ring \[
T=F[X_1,\ldots,X_d]/(X_1,\ldots,X_d)^N
\] maps onto \(s_0(F)[\bar u]\). The kernel is contained in the ideal of positive-degree terms, because a nonzero constant from \(F\) has nonzero residue. It is therefore nilpotent. Formal smoothness of \(k/\mathbb F_p\) lifts its specified map to a map \(v:k\to T\). The difference \(e=i-\psi_s v:k\to I\) is an \(\mathbb F_p\)-derivation: expand a product, use agreement modulo \(I\), and discard the product of the two errors in \(I^2\).

Choose the absolute differential basis \((dz_j)_{j\in J}\) supplied by \hyperref[AG-RG-S05-more-algebra-lemma-p-basis]{F.119}. Formal smoothness and \hyperref[AG-RG-S05-algebra-lemma-characterize-formally-smooth-field-extension]{F.17} give \(H_1(L_{K/\mathbb F_p})=0\). Consequently \hyperref[AG-RG-S05-more-algebra-lemma-transitivity-gamma]{F.131} identifies the kernel in \[
\Omega_{k/\mathbb F_p}\otimes_kK\longrightarrow\Omega_{K/\mathbb F_p}
\] with \(H_1(L_{K/k})\). A finite basis of this kernel uses only finitely many coordinates \(dz_j\); let \(J_0\) contain their supports. The images of the other basis vectors are linearly independent, since a relation supported outside \(J_0\) would also have to be supported inside it. Define a \(K\)-linear map on their span by \(dz_j\mapsto e(z_j)\). Extend a basis of this span to a basis of \(\Omega_{K/\mathbb F_p}\), assign zero on the extra vectors, and use the universal differential map to obtain a derivation \(\delta:K\to I\).

The formula \(s_1=s+\delta\) gives another ring section: the derivation identity accounts for the two cross terms, and \(I^2=0\) discards the last term. In evaluating \(v(a)\), the constant coefficient is \(a\in k\subset F\). Changes to positive-degree coefficients are annihilated by their monomials, using \(\mathfrak m I=0\). Thus \(\psi_{s_1}v(a)-\psi_s v(a)=\delta(a)\). The remaining error has the form \[
e(a)-\delta(a)=\sum_{j\in J_0}a_j w_j,
\qquad da=\sum_j a_j dz_j,\quad a_j\in k,
\quad w_j=e(z_j)-\delta(z_j)\in I.
\] Each \(w_j\) is a linear combination of the finitely many degree-\((N-1)\) monomials \(u^E\), with coefficients in \(K\). Indeed these monomials generate \(I\), and a coefficient may be replaced by its residue because \(\mathfrak m I=0\). Adjoin all these coefficients to \(F\) to obtain \(F_0\). When multiplying a \(w_j\), the scalar \(i(a_j)\) can equally be represented by \(s_1(a_j)\), since their difference belongs to \(\mathfrak m\). The displayed error and \(\psi_{s_1}v(a)\) therefore both lie in \(\psi_{s_1}(F_0[X])\). This proves the assertion for the actual structure map. Fix this section henceforth and call it \(s\).

The map \(\psi_s\) is surjective. Subtract a residue representative, write the remainder using the \(u_j\), and do the same for its coefficients. After \(N\) steps the remainder lies in \(\mathfrak m^N=0\). This produces a polynomial expression for each element. By the Hilbert basis theorem \href{https://kokunoyumeto.github.io/open-math-courses/courses/AG-CA/noetherian-and-artinian-rings.html}{AG-CA-03, Theorem 2.1}, its kernel has finitely many generators \(g_1,\ldots,g_b\). Enlarge \(F_0\) by their finitely many coefficients to a field \(F_*\). For each finitely generated intermediate field \(F_*\subset F\subset K\), set \[
A_F=F[X_1,\ldots,X_d]/(g_1,\ldots,g_b).
\] The presentation gives \(K\otimes_F A_F=\Lambda\). Choose a basis of \(K/F\) containing \(1\). As an \(A_F\)-module this tensor product is a direct sum of copies of \(A_F\); its \(1\)-coordinate makes the natural map injective, exact, and faithful on every module. Thus \(A_F\subset\Lambda\) and this inclusion is faithfully flat. It contains \(i(k)\) by the preceding assertion.

The ideal generated by the \(X_j\) has \(N\)th power zero, as may be checked in the injective image. Its quotient is \(F\), since each \(g_j\) has zero constant coefficient. The monomials of degree below \(N\) span \(A_F\) over \(F\). Finite dimension makes it Artinian; the nilpotent ideal is its unique maximal ideal, since an element with nonzero constant term is invertible by the finite geometric series. That ideal extends to \(\mathfrak m\). The residue field \(F/k\) is finitely generated, hence is the fraction field of the finite-type domain generated by a field-generator list. \hyperref[AG-RG-S05-algebra-lemma-essentially-of-finite-type-into-artinian-local]{F.40} now makes \(A_F\) essentially of finite type for its actual, possibly different, \(k\)-structure map.

Two finite field-generator lists have a common successor, their joint generated field. Every polynomial expression in \(\Lambda\) uses finitely many coefficients of \(K\); adjoining them to \(F_*\) places that element in some \(A_F\). These actual subrings therefore form a directed union equal to \(\Lambda\), with all the required properties. ∎

\phantomsection\label{AG-RG-S05-smoothing-artinian-denominator-enlargement}

\subsubsection{Lemma D.34A: Artinian denominator enlargement}\label{AG-RG-S05-lemma-d.34a-artinian-denominator-enlargement}

Reading sources: the unit enlargement in \href{https://stacks.math.columbia.edu/tag/07FI}{Stacks, Tag 07FI} and the finite étale construction in \href{https://stacks.math.columbia.edu/tag/04GK}{Tag 04GK}. The auxiliary statement and the calculations below are written here.

Let \(k\) have characteristic \(p>0\), and let \(A\to B\) be a flat local map of Noetherian Artinian local \(k\)-algebras with \(\mathfrak m_AB=\mathfrak m_B\) and \(\mathfrak m_B^N=0\). Denote their residue fields by \(F\subset K\). Given finitely many units \(b_j\in B\), there is a local factorization \(A\to A'\to B\) such that \(A\to A'\) is essentially smooth, \(A'\to B\) is faithfully flat, \(\mathfrak m_{A'}B=\mathfrak m_B\), and a \(p\)-power of each \(b_j\) belongs to the subring \(A'\subset B\). The residue extension is a finite tower of simple transcendental and finite separable extensions. At a transcendental step the unit itself is adjoined.

\textbf{Proof.} The local faithful-flatness theorem \href{https://kokunoyumeto.github.io/open-math-courses/courses/AG-CA/faithful-flatness-and-the-local-criterion-for-flatness.html}{AG-CA-08, Theorem 2.1} makes the given map faithfully flat, hence injective. For one unit \(b\), write \(\beta\in K^\times\) for its residue.

If \(\beta\) is transcendental over \(F\), take \(A'=A[T]_{\mathfrak m_AA[T]}\) and send \(T\) to \(b\). Every denominator has a nonzero polynomial residue at \(\beta\), hence maps to a unit of \(B\). This defines a local, essentially smooth map with maximal ideal \(\mathfrak m_AA'\) and residue \(F(\beta)\). The ring is Noetherian by the Hilbert basis theorem and localization. Its maximal ideal is nilpotent. Each successive ideal-power quotient is a finite-dimensional vector space over its residue field, since each power is a finite module. These quotients give finite length, so \(A'\) is Artinian.

In the remaining algebraic case, take \(A'=A\) if \(\beta\in F\). Otherwise let \(f\) be its monic irreducible polynomial. If \(f'=0\), every positive exponent occurring in \(f\) is divisible by \(p\), so \(f(T)=f_1(T^p)\). The polynomial \(f_1\) is irreducible, since a factorization would give one of \(f\) by substitution. Repeat this operation while the derivative vanishes. The degrees decrease, so it stops at an irreducible polynomial with nonzero derivative. Its roots include \(\beta^{q_0}\) for a \(p\)-power \(q_0\). Bézout with the derivative proves this root separable. Thus \(F_1=F(\beta^{q_0})\) is finite separable over \(F\).

Lift its monic minimal polynomial to \(h\in A[T]\), and put \(A'=A[T]/(h)\). Monic division makes \(A'\) finite free over \(A\). Its reduction is the field \(F_1\); the ideal \(\mathfrak m_AA'\) is nilpotent and is its unique maximal ideal. It is Artinian by the same finite-length argument. Bézout for \(\bar h,\bar h'\) makes \(h'\) invertible in the reduction, and a finite geometric series lifts its inverse across the nilpotent ideal. The one-equation Jacobian proof \href{https://kokunoyumeto.github.io/open-math-courses/courses/AG-CA/formally-smooth-unramified-and-etale-ring-maps.html}{AG-CA-17, Proposition 6.1} makes \(A'/A\) étale.

Here is the actual map to \(B\). Choose \(r_0\in B\) reducing to \(\beta^{q_0}\). Its derivative value is a unit, and \(h(r_0)\in\mathfrak m_B\). Define \[
r_{a+1}=r_a-h'(r_a)^{-1}h(r_a).
\] Polynomial expansion shows that \(h(r_{a+1})\) belongs to the square of the ideal containing \(h(r_a)\): the constant and linear terms cancel, and all remaining terms contain at least two copies of the correction. Consequently \(h(r_a)\in\mathfrak m_B^{2^a}\). The residues stay fixed and the derivatives stay units. Once \(2^a\ge N\), this gives an exact root and the required homomorphism \(A'\to B\). Two roots with that residue are equal: their difference times the divided-difference polynomial is zero, and that polynomial has the nonzero derivative as residue and is a unit. This is the simple-root construction without an external lifting step.

In all three cases \(A'\) is \(A\)-flat and \(A'/\mathfrak m_AA'=F'\), its residue field. Tensor a free \(A\)-resolution of \(F\) with \(A'\); flatness makes it a free \(A'\)-resolution of \(F'\). Tensoring further with \(B\) gives \[
\operatorname{Tor}^{A'}_1(F',B)=\operatorname{Tor}^{A}_1(F,B)=0.
\] Apply the Noetherian local criterion \href{https://kokunoyumeto.github.io/open-math-courses/courses/AG-CA/faithful-flatness-and-the-local-criterion-for-flatness.html}{AG-CA-08, Theorem 4.2} with local source \(A'\), local target \(B\), and the finite \(B\)-module \(B\). It proves \(A'\to B\) flat. Locality makes it faithfully flat and injective, so \(A'\) is the asserted subring. Its maximal ideal still generates \(\mathfrak m_B\).

In the two algebraic cases choose a unit \(x\in A'\) with residue \(\beta^{-q_0}\), taking \(q_0=1\) if \(\beta\in F\). Then \(xb^{q_0}-1\in\mathfrak m_B\). For a \(p\)-power \(q_1\ge N\), the characteristic-\(p\) binomial identity yields \[
(xb^{q_0})^{q_1}=1,
\qquad b^{q_0q_1}=x^{-q_1}\in A'.
\] In the transcendental case \(b\) itself already belongs to \(A'\). Repeat for the finite unit list. Each earlier power remains in every later subring, and the maps remain essentially smooth by composition.

For use in D.35, retain the natural differential information. A simple transcendental residue step has \[
\Omega_{F(\beta)/k}
=(\Omega_{F/k}\otimes_F F(\beta))\oplus F(\beta)\,d\beta,
\] by polynomial differentials and the inverse rule in localization. In a finite separable step, differentiating the monic equation solves its new differential uniquely, because its derivative is invertible; it leaves precisely the old differential space after scalar extension. Both residue extensions are formally smooth by \hyperref[AG-RG-S05-algebra-lemma-formally-smooth-extensions-easy]{F.49}, so their relative first homology vanishes by \hyperref[AG-RG-S05-algebra-lemma-characterize-formally-smooth-field-extension]{F.17}. The exact field sequence \hyperref[AG-RG-S05-more-algebra-lemma-transitivity-gamma]{F.131} therefore gives the natural isomorphism \[
H_1(L_{F'/k})=H_1(L_{F/k})\otimes_F F'.
\] Composition gives these formulas for the whole finite tower. In particular its first-homology image in \(H_1(L_{K/k})\) is unchanged. ∎

\phantomsection\label{AG-RG-S05-smoothing-lemma-solution-modulo}

\subsubsection{Lemma D.35: solution modulo}\label{AG-RG-S05-lemma-d.35-solution-modulo}

Reading source: \href{https://stacks.math.columbia.edu/tag/07FH}{Stacks, Tag 07FH}. The global representatives and their denominators are constructed explicitly below.

Let \(k\) have characteristic \(p>0\), let \(\Lambda\) be a Noetherian geometrically regular \(k\)-algebra, and choose a prime \(\mathfrak q\), an integer \(n\ge1\), and a finite subset \(E\subset\Lambda_{\mathfrak q}/\mathfrak q^n\Lambda_{\mathfrak q}\). There is a map \(\varphi:P=k[y_1,\ldots,y_m]\to\Lambda\), for some \(m\ge0\), such that, with \(\mathfrak p=\varphi^{-1}(\mathfrak q)\):

\begin{enumerate}
\def\labelenumi{\arabic{enumi}.}
\tightlist
\item
  \(P_{\mathfrak p}\to\Lambda_{\mathfrak q}\) is flat and \(\mathfrak p\Lambda_{\mathfrak q}=\mathfrak q\Lambda_{\mathfrak q}\).
\item
  Its quotient map has a local Artinian factorization \[
  P_{\mathfrak p}/\mathfrak p^nP_{\mathfrak p}
  \longrightarrow D\longrightarrow
  \Lambda_{\mathfrak q}/\mathfrak q^n\Lambda_{\mathfrak q},
  \] with the first arrow essentially smooth and the second faithfully flat.
\item
  The image of \(D\) contains \(E\).
\end{enumerate}

\textbf{Proof.} Write \(L=\Lambda_{\mathfrak q}\), \(\mathfrak m=\mathfrak qL\), \(K=L/\mathfrak m\), and \(B=L/\mathfrak m^n\). The injection in \hyperref[AG-RG-S05-more-algebra-proposition-characterization-geometrically-regular]{F.133} embeds \(H_1(L_{K/k})\) in the finite-dimensional space \(\mathfrak m/\mathfrak m^2\). By \hyperref[AG-RG-S05-smoothing-lemma-helper]{D.34}, choose a local Artinian subring \(A\subset B\) containing \(E\), essentially of finite type over \(k\), faithfully flat to \(B\), with \(\mathfrak m_AB=\mathfrak m_B\). Let \(F\subset K\) be its residue field. This field is finitely generated over \(k\): present \(A\) as a localization of a finite-type algebra, take its residue quotient, and express the resulting field as fractions of the image of the finite algebra-generator list.

The differentials of such a field-generator list span \(\Omega_{F/k}\), since the product and inverse rules differentiate every rational expression. Select a basis \(d\theta_1,\ldots,d\theta_t\) from this spanning list and lift its elements to \(A\). Each lift has a representative \(a_j/b_j\in L\), with \(a_j\in\Lambda\) and \(b_j\notin\mathfrak q\).

By the left injection in \hyperref[AG-RG-S05-more-algebra-lemma-transitivity-gamma]{F.131}, \(H_1(L_{F/k})\otimes_FK\) embeds in \(H_1(L_{K/k})\). Denote its image in \(\mathfrak m/\mathfrak m^2\) by \(V\). If \(n\ge2\), the elements of \(\mathfrak m_A\) span this cotangent space over \(K\): its identification with the cotangent space of \(B\) and the equality of extended maximal ideals prove that assertion. Choose elements whose images are a basis of a complement to \(V\). Lift them to fractions \(c_i/d_i\), with \(c_i\in\mathfrak q\) and \(d_i\notin\mathfrak q\). For \(n=1\), select instead representatives \(c_i\in\mathfrak q\) of a complement in \(\mathfrak m/\mathfrak m^2\) and put \(d_i=1\). Their images in \(B=K\) are zero, so are already in \(A\).

Apply \hyperref[AG-RG-S05-smoothing-artinian-denominator-enlargement]{D.34A} to the finitely many units represented by \(b_j,d_i\) in \(B\). Obtain \(A\subset A'\subset B\), with \(A\to A'\) essentially smooth, \(A'\to B\) faithfully flat, and unchanged extended maximal ideal. For appropriate \(p\)-powers \(q_j,r_i\), all four kinds of global elements \[
a_jb_j^{q_j-1},\quad b_j^{q_j},\quad
c_id_i^{r_i-1},\quad d_i^{r_i}
\] have images in \(A'\). The first and third are the fractional lifts multiplied by the indicated powers; the second and fourth are supplied by the enlargement. Record in addition every original denominator which was adjoined at a transcendental residue step. It itself has its image in \(A'\).

Let \(F'\) be the residue field of \(A'\). The residues of the first two kinds of elements, together with these recorded transcendental denominators, have differentials spanning \(\Omega_{F'/k}\). Indeed their quotient is \(\theta_j\), so the inverse rule expresses each old \(d\theta_j\) in their span. D.34A adds just one new differential at each transcendental step and none at a finite separable step. Select a differential basis from this finite list, and call its actual global representatives \(\lambda_1,\ldots,\lambda_s\in\Lambda\). This argument does not set \(d(b_j^{q_j})\) to zero in \(F'\): the root \(b_j\) need not belong to that field. Every selected representative has its class in \(A'\).

Define \(P'=k[y_1,\ldots,y_s]\to\Lambda\) by these elements, let \(\mathfrak p'\) be the inverse image of \(\mathfrak q\), and set \(K_0=\kappa(\mathfrak p')\subset F'\). Since all the basis differentials come from \(K_0\), \(\Omega_{F'/K_0}=0\). The extension is finitely generated, because \(F'/k\) is finitely generated. The equality proved in \hyperref[AG-RG-S05-more-algebra-lemma-cartier-equality]{F.98} gives \[
-\dim_{F'}H_1(L_{F'/K_0})
=\operatorname{trdeg}_{K_0}F'.
\] Its two nonnegative dimensions must both be zero. Finitely many algebraic field generators give a finite extension by multiplying their degrees; the zero first homology gives separability by \hyperref[AG-RG-S05-algebra-proposition-characterize-separable-field-extensions]{F.91}. Thus \(F'/K_0\) is finite separable. The field sequence F.131, whose relative first homology and differential module here vanish, identifies \[
H_1(L_{K_0/k})\otimes_{K_0}K
=H_1(L_{F'/k})\otimes_{F'}K
=H_1(L_{F/k})\otimes_FK.
\] The last identity is the natural one retained in D.34A, so its image in \(\mathfrak m/\mathfrak m^2\) is still \(V\).

Apply \hyperref[AG-RG-S05-more-algebra-lemma-geometrically-regular-over-field]{F.109} using the selected differential basis. It gives a flat map \(P'_{\mathfrak p'}\to L\) and a regular quotient \(L/\mathfrak p'L\). Its written conormal square identifies the image of the source cotangent space with the preceding first-homology image \(V\). Append the global elements \(\lambda_{s+i}=c_id_i^{r_i-1}\) to form \(P=k[y_1,\ldots,y_m]\to\Lambda\). Their cotangent classes are the originally chosen complementary classes multiplied by the nonzero scalars \(\bar d_i^{r_i}\), and therefore still form a complementary basis. Their residues are zero. Hence \(P_{\mathfrak p}\) is the localization of \(P'_{\mathfrak p'}[y_{s+1},\ldots,y_m]\) at the prime generated by \(\mathfrak p'\) and these variables.

Polynomial local regularity is proved in \href{https://kokunoyumeto.github.io/open-math-courses/courses/AG-CA/regular-local-rings.html}{AG-CA-14, Proposition 3.3}. The old regular parameters and the appended variables are a regular parameter list for this source: its maximal ideal is generated by them and its cotangent space is the old cotangent space plus one independent class for each new variable. Their images are a regular sequence in \(L\). The old list remains regular by the flat map from \(P'_{\mathfrak p'}\), applying flat base change to each successive multiplication injection. The new list is a minimal parameter list on the regular local quotient \(L/\mathfrak p'L\), hence is regular by \href{https://kokunoyumeto.github.io/open-math-courses/courses/AG-CA/regular-sequences-depth-and-cohen-macaulay-modules.html}{AG-CA-12, Theorem 6.1}. Concatenating these lists proves the assertion. \hyperref[AG-RG-S05-algebra-lemma-flat-over-regular]{F.46} now makes \(P_{\mathfrak p}\to L\) flat. Their cotangent classes span \(\mathfrak m/\mathfrak m^2\); Nakayama \hyperref[AG-RG-S05-algebra-lemma-NAK]{F.10} gives \(\mathfrak pL=\mathfrak m\). This proves (1).

All polynomial generators have their classes in \(A'\), so the map to \(B\) factors through \(A'\). An element outside \(\mathfrak p\) has nonzero residue and maps to a unit of this local ring. Moreover \(\mathfrak p^n\) maps to zero in \(B\), hence in \(A'\subset B\). We obtain \[
R=P_{\mathfrak p}/\mathfrak p^nP_{\mathfrak p}
\longrightarrow A'\longrightarrow B.
\] The ring \(R\) is Artinian: it is Noetherian local with nilpotent maximal ideal, and its finite ideal-power quotients have finite dimension over its residue field. Its composite map to \(B\) is the flat base change of (1). Since \(A'\to B\) is faithfully flat, \hyperref[AG-RG-S05-algebra-lemma-flatness-descends-more-general]{F.48} makes \(R\to A'\) flat. The ideals \(\mathfrak pA'\) and \(\mathfrak m_{A'}\) become equal in \(B\); faithful flatness detects equality by tensoring their quotient, so they are equal. The residue extension is the finite separable extension \(F'/K_0\) just proved. \hyperref[AG-RG-S05-algebra-lemma-essentially-of-finite-type-into-artinian-local]{F.40} makes \(A'\) a finite \(R\)-module. Since \(R\) is Noetherian, it is a finitely presented algebra, by taking finite module generators and a finite kernel in their polynomial presentation. Its only fibre is the finite separable residue field. The proved étale criterion \hyperref[AG-RG-S05-algebra-lemma-characterize-etale]{F.16} therefore makes \(R\to A'\) étale, in particular essentially smooth. Take \(D=A'\). It retains \(E\subset A\), proving (2) and (3). ∎

\phantomsection\label{AG-RG-S05-smoothing-lemma-enlarge-solution-modulo}

\subsubsection{Lemma D.36: enlarge solution modulo}\label{AG-RG-S05-lemma-d.36-enlarge-solution-modulo}

Reading source: \href{https://stacks.math.columbia.edu/tag/07FI}{Stacks, Tag 07FI}.

Use the notation and a local Artinian factorization of D.35: \[
R=k[y_1,\ldots,y_m]_{\mathfrak p}/\mathfrak p^n
\longrightarrow D\longrightarrow
B=\Lambda_{\mathfrak q}/\mathfrak q^n\Lambda_{\mathfrak q},
\] where \(R\to D\) is essentially smooth and \(D\to B\) is flat. For every \(\lambda\in\Lambda\setminus\mathfrak q\), there are a positive \(p\)-power \(q\) and a local Artinian factorization \(R\to D\to D'\to B\) with \(D\to D'\) essentially smooth, \(D'\to B\) faithfully flat, and \(\lambda^q\) in the image of \(D'\).

\textbf{Proof.} The polynomial-source maximal ideal extends to that of \(B\), by D.35(1). We verify that the maximal ideal of \(D\) does also, even for a factorization other than the particular étale one constructed there. The closed fibre \(D/\mathfrak m_RD\) is Artinian local and essentially smooth over the field \(R/\mathfrak m_R\). The regularity theorem \href{https://kokunoyumeto.github.io/open-math-courses/courses/AG-CA/smooth-algebras-over-a-field-and-the-jacobian-criterion.html}{AG-CA-18, Theorem 2.1}, followed by localization, makes it regular. An Artinian local ring has dimension zero. For a regular local ring of dimension zero its cotangent dimension is zero; its finite maximal ideal equals its square, and Nakayama makes that ideal zero. Thus this fibre is a field and \(\mathfrak m_D=\mathfrak m_RD\). Extending to \(B\) gives \(\mathfrak m_DB=\mathfrak m_B\).

All rings required by D.34A are Noetherian. The source \(R\) is Noetherian, and the essentially smooth algebra \(D\) is a localization of a finitely presented \(R\)-algebra, so the Hilbert basis theorem and localization make it Noetherian. The ring \(B\) is the displayed Artinian quotient of the Noetherian local ring \(\Lambda_{\mathfrak q}\), and \(\mathfrak m_B^n=0\). The flat local map \(D\to B\) is faithfully flat by the local criterion. The image of \(\lambda\) is a unit because its residue is nonzero. Apply the one-unit case of \hyperref[AG-RG-S05-smoothing-artinian-denominator-enlargement]{D.34A}. It supplies \(D\subset D'\subset B\), an essentially smooth first map, a faithfully flat last map, and the exact power containment \(\lambda^q\in D'\). Composition with \(R\to D\) gives the required factorization. ∎

\phantomsection\label{AG-RG-S05-smoothing-lemma-resolve-general}

\subsubsection{Lemma D.37: resolve general}\label{AG-RG-S05-lemma-d.37-resolve-general}

Reading source: \href{https://stacks.math.columbia.edu/tag/07FJ}{Stacks, Tag 07FJ}.

Let \(k\) have characteristic \(p>0\). In \hyperref[AG-RG-S05-smoothing-situation-local]{D.27}, suppose that \(\Lambda\) is Noetherian and geometrically regular over \(k\), and that \(\mathfrak q\) is minimal over \(\mathfrak h_A\). Then the problem has a resolution.

\textbf{Proof.} Put \(L=\Lambda_{\mathfrak q}\), \(\mathfrak m=\mathfrak qL\), \(d=\dim L\), and \(H=H_{A/k}=(a_1,\ldots,a_r)\). These generators exist by the Hilbert theorem \href{https://kokunoyumeto.github.io/open-math-courses/courses/AG-CA/noetherian-and-artinian-rings.html}{AG-CA-03, Theorem 2.1}. The ring \(L\) is regular by geometric regularity. As in D.33, minimality and a finite maximal-ideal generating list give \(N>0\) with \(\mathfrak m^N\subset HL\).

Over \(P_0=k[x_1,\ldots,x_d]\) form \[
B=A[x_1,\ldots,x_d,z_{ij}]/(x_i^{2N}-\sum_j a_jz_{ij}).
\] Eliminating \(z_{ij}\) on an \(a_j\)-chart proves smoothness there; the equation makes those charts cover \(B_{x_i}\). By \hyperref[AG-RG-S05-algebra-lemma-locally-smooth]{F.67}, \(B_{x_i}/P_0\) is smooth. Apply \hyperref[AG-RG-S05-smoothing-lemma-improve-presentation]{D.9} to obtain \(B\to C\) with retraction \(C\to B\). The smooth charts of \(C\) have finite free differentials. Since \(C/P_0\) is finitely presented by the Hilbert theorem, \hyperref[AG-RG-S05-smoothing-lemma-compare-standard]{D.15(e)} gives one \(c>0\) for which all \(x_i^c\) are strictly standard. If \(d=0\), take \(c=1\). Fix, before making any Artinian approximation, \[
e=8c,\qquad n=\max\{N+dc,\ d(e-1)+1\}.
\tag{D.37a}
\]

Apply \hyperref[AG-RG-S05-smoothing-lemma-solution-modulo]{D.35} to the images of a finite \(k\)-algebra generating list of \(A\) in \(\overline L=L/\mathfrak m^n\). It gives a polynomial algebra \(P=k[y_1,\ldots,y_m]\) mapping to \(\Lambda\), its contracted prime \(\mathfrak p\), and local maps \[
S_0=P_{\mathfrak p}/\mathfrak p^nP_{\mathfrak p}
\longrightarrow D\longrightarrow\overline L.
\] The first map is essentially smooth, the second is flat, and the generator images belong to \(D\). Also \(P_{\mathfrak p}\to L\) is flat and \(\mathfrak pL=\mathfrak m\). Its closed fibre is a field, so the dimension formula \href{https://kokunoyumeto.github.io/open-math-courses/courses/AG-CA/dimension-theory-of-noetherian-local-rings.html}{AG-CA-11, Theorem 5.1} gives \(\dim P_{\mathfrak p}=d\).

Choose \(\pi_1,\ldots,\pi_d\in\mathfrak p\) representing a regular parameter list of \(P_{\mathfrak p}\). Polynomial local rings are regular by iteration of \href{https://kokunoyumeto.github.io/open-math-courses/courses/AG-CA/regular-local-rings.html}{AG-CA-14, Proposition 3.3}, starting with the field. The minimal-parameter theorem \href{https://kokunoyumeto.github.io/open-math-courses/courses/AG-CA/regular-sequences-depth-and-cohen-macaulay-modules.html}{AG-CA-12, Theorem 6.1} makes their minimal generating list regular. Their images generate \(\mathfrak m\) and form a regular parameter list of \(L\), by the same theorem and \(\dim L=d\). Let \[
R=P[t_1,\ldots,t_d],\qquad \gamma_i=\pi_it_i.
\] Throughout, localization of this polynomial ring at the coefficients means \[
R_S=(P\setminus\mathfrak p)^{-1}R=P_{\mathfrak p}[t_1,\ldots,t_d].
\] It is distinct from localization at a prime of \(R\).

First arrange the source kernel conditions. By \hyperref[AG-RG-S05-algebra-lemma-regular-sequence-in-polynomial-ring]{F.75}, the \(\gamma_i\) form a regular sequence in \(R_S\). Powers of earlier entries preserve regularity by \href{https://kokunoyumeto.github.io/open-math-courses/courses/AG-CA/regular-sequences-depth-and-cohen-macaulay-modules.html}{AG-CA-12, Proposition 1.3}. Inductively, on the finite \(R\)-module \(R/(\gamma_1^e,\ldots,\gamma_{i-1}^e)\), the localized kernels for \(\gamma_i\) and \(\gamma_i^2\) are zero. \hyperref[AG-RG-S05-smoothing-lemma-ogoma]{D.31} supplies \(s_i\in P\setminus\mathfrak p\) such that replacing \(\pi_i\) by \(s_i\pi_i\) makes these kernels equal globally. Earlier equalities are unchanged, because their entries and quotient ideals have not changed. Such scaling is by local units, so both local parameter lists remain regular. Continue through all \(i\), using the modified \(\pi_i\) from now on.

Next construct units \(\delta_i\in\Lambda\setminus\mathfrak q\), coefficients \(\lambda_{ij}\in\Lambda\), and an essentially smooth Artinian enlargement \(D\to D'\) with flat local map \(D'\to\overline L\), such that \[
(\delta_i\pi_i)^{2N}=\sum_j a_j\lambda_{ij}
\quad\hbox{in }\Lambda,
\tag{D.37b}
\] and both \(\delta_i\) and \(\lambda_{ij}\) have images in \(D'\). We additionally require equal kernels for \(\delta_i\pi_i\) and its square modulo the earlier \((\delta_h\pi_h)^e\). A flat local map into the nonzero local ring \(\overline L\) is faithfully flat, and hence injective, by \href{https://kokunoyumeto.github.io/open-math-courses/courses/AG-CA/faithful-flatness-and-the-local-criterion-for-flatness.html}{AG-CA-08, Theorem 2.2}. Thus every ring in this construction may be regarded as a subring of \(\overline L\). In particular all relations among the chosen \(A\)-generator images already hold in \(D\), so there is a map \(A\to D\).

Suppose the construction has reached \(i-1\). The element \(\pi_i^N\) lies in \(H\overline L\); faithful flatness contracts this ideal to \(HD'\). Write \(\pi_i^N=\sum_j a_jd_j\) in \(D'\), and lift the \(d_j\) to \(u_j\in L\). The error lies in \(\mathfrak m^n\subset H\mathfrak m^{n-N}\), since \(n\ge N\). Thus choose \(v_j\in\mathfrak m^{n-N}\) with \[
\pi_i^N=\sum_j a_j(u_j+v_j)\quad\hbox{in }L.
\] Set \(w_j=\pi_i^N(u_j+v_j)\). Then \(\pi_i^{2N}=\sum_j a_jw_j\), and the residue of \(w_j\) in \(\overline L\) is \(\pi_i^Nd_j\), because \(\pi_i^Nv_j\in\mathfrak m^n\).

Here is the global denominator calculation. Choose \(s_0\notin\mathfrak q\) so that every \(s_0^{2N}w_j\) is represented by \(\widetilde\mu_j\in\Lambda\). The error \[
r_i=(s_0\pi_i)^{2N}-\sum_j a_j\widetilde\mu_j
\] vanishes in \(L\). Hence some \(v\notin\mathfrak q\) annihilates \(r_i\), by the fraction-zero criterion of \href{https://kokunoyumeto.github.io/open-math-courses/courses/AG-RG/AG-RG-S01.html}{S01, Lemma 1.1}. Set \(s=s_0v\) and \(\mu_j=v^{2N}\widetilde\mu_j\). Now \((s\pi_i)^{2N}=\sum_j a_j\mu_j\) in \(\Lambda\), and \(\mu_j/s^{2N}=w_j\) in \(L\).

Use \hyperref[AG-RG-S05-smoothing-lemma-enlarge-solution-modulo]{D.36} to put \(s^h\) in a further \(D'\), for some \(h>0\). Replace \(s\) by \(s^h\) and \(\mu_j\) by \(s_{\rm old}^{2N(h-1)}\mu_j\). The exact equation survives, and the new coefficient residue is \(s^{2N}\pi_i^Nd_j\in D'\). In \(L/((\delta_1\pi_1)^e,\ldots,(\delta_{i-1}\pi_{i-1})^e)\), multiplication by \(s\pi_i\) is injective, because the preceding entries are powers of a regular parameter list scaled by units. Apply D.31 on the corresponding finite \(\Lambda\)-module to obtain \(s'\notin\mathfrak q\) for which every \((s')^q s\pi_i\), \(q>0\), has equal global kernels with its square. D.36 puts one such \((s')^q\) in an enlargement of \(D'\). Take \[
\delta_i=(s')^qs,\qquad \lambda_{ij}=(s')^{2Nq}\mu_j.
\] This proves (D.37b), the residue assertions, and the required target kernel equality. Later enlargements retain all earlier assertions. The finite induction constructs the full list; when \(d=0\) this construction is empty and \(D'=D\).

Map \(R\to\Lambda\) by the given \(P\)-map and \(t_i\mapsto\delta_i\). Give \(R\) its \(P_0\)-structure by \(x_i\mapsto\gamma_i\). Equation (D.37b) defines \(B\to\Lambda\) by \(x_i\mapsto\delta_i\pi_i\) and \(z_{ij}\mapsto\lambda_{ij}\). The retraction defines \(C\to\Lambda\), and these maps induce \[
C_R=C\otimes_{P_0}R\longrightarrow\Lambda.
\] Every \(a_j\)-chart of \(C_R/R\) is smooth, so the singular radical of this new input contains \(\mathfrak h_A\). A resolution over \(R\) returns to one over \(k\): finite presentation composes with the polynomial base, and smooth composition gives \(H_{T/R}\subset H_{T/k}\).

Let \(I=(\gamma_1^e,\ldots,\gamma_d^e)\subset R\). \hyperref[AG-RG-S05-smoothing-lemma-strictly-standard-base-change]{D.7} transports the strictly standard elements \(x_i^c\) to \(\gamma_i^c\). The arranged source and target kernel equalities hold also for their \(c\)-th powers, by the kernel induction in D.33. \hyperref[AG-RG-S05-smoothing-lemma-lift-solution]{D.28} therefore lifts a resolution modulo \(I\) to an \(R\)-resolution. For \(d=0\), \(I=0\) and this lifting step is unnecessary.

We construct the quotient resolution explicitly. Put \(J=(\pi_1^e,\ldots,\pi_d^e)\subset P\), and let \(\mathfrak r\) be the inverse image of \(\mathfrak q\) in \(R\). Each \(t_i\notin\mathfrak r\), since \(\delta_i\notin\mathfrak q\). Consequently \[
IR_{\mathfrak r}=JR_{\mathfrak r},
\qquad IL=JL.
\tag{D.37c}
\] The parameters generate \(\mathfrak pP_{\mathfrak p}\) and \(\mathfrak m\). The bound \(n\ge d(e-1)+1\) in (D.37a) gives, by expansion into monomials, \[
\mathfrak p^nP_{\mathfrak p}\subset JP_{\mathfrak p},
\qquad\mathfrak m^n\subset JL.
\tag{D.37d}
\] For \(d=0\), both maximal ideals are zero, so these containments hold as well.

Set \(S_1=S_0[t_1,\ldots,t_d]\). Bars below denote elements in \(D'\subset\overline L\). Each \(\overline\delta_i\) is a unit of \(D'\): its image is a unit and the map is local. Define \[
B\longrightarrow D'[t_1,\ldots,t_d],
\qquad
x_i\longmapsto\overline\pi_it_i,\quad
z_{ij}\longmapsto
\overline\lambda_{ij}\overline\delta_i^{-2N}t_i^{2N}.
\tag{D.37e}
\] Injectivity of \(D'\to\overline L\) transports (D.37b) to \(D'\). The image of the \(i\)-th defining equation in (D.37e) is therefore \[
t_i^{2N}\left(\overline\pi_i^{2N}
-\overline\delta_i^{-2N}\sum_j\overline a_j\overline\lambda_{ij}\right)=0.
\] Its restriction to \(P_0\) is exactly the prescribed \(x_i\mapsto\pi_it_i\) in \(S_1\), so the map extends to \(B\otimes_{P_0}S_1\). Evaluation \(t_i\mapsto\overline\delta_i\) recovers the prescribed images \(\delta_i\pi_i,\lambda_{ij}\) in \(\overline L\).

The map \(S_0\to D'\) is essentially smooth by construction and composition. Express it as a localization of a smooth finitely presented \(S_0\)-algebra. Localization is the filtered colimit of its principal localizations at finite products of the inverted elements. These stages are smooth, and adjoining the \(t_i\) gives smooth \(S_1\)-stages with colimit \(D'[t_i]\). Finite presentation and \hyperref[AG-RG-S05-algebra-lemma-when-colimit]{F.90} factor (D.37e) through one smooth \(S_1\)-stage \(T\), followed by the specified map to \(\overline L\). Composing with \(C\to B\) gives the same smooth factorization for \(C\otimes_{P_0}S_1\).

By (D.37d) there are quotient maps \(S_1\to P_{\mathfrak p}[t_i]/J\) and \(\overline L\to L/JL\). Base changing \(T\) gives a smooth factorization over \(P_{\mathfrak p}[t_i]/J\). Localize further at \(\mathfrak r\). Every denominator then has a unit image in the local target, so its map extends. By (D.37c) this is a smooth factorization of the localized problem modulo \(I\).

The target \(L/JL\) has nilpotent maximal ideal: degree \(d(e-1)+1\) monomials vanish. Thus its dimension is zero. If the singular radical of the unlocalized quotient input escapes the image of \(\mathfrak q\), the identity factorization resolves that quotient. Otherwise this prime is minimal over the radical, and \hyperref[AG-RG-S05-smoothing-lemma-delocalize-height-zero]{D.30} delocalizes the smooth factorization just obtained. D.28 then lifts it before quotienting by \(I\). Finally the polynomial change of base returns the resolution to \(k\). All constructions above include \(d=0\), empty parameter lists, and the specified target maps. ∎

\subsection{The main theorem}\label{AG-RG-S05-the-main-theorem}

\phantomsection\label{AG-RG-S05-smoothing-section-main}

In this section we wrap up the discussion.

\phantomsection\label{AG-RG-S05-smoothing-theorem-popescu}

\subsubsection{Theorem D.38: popescu}\label{AG-RG-S05-theorem-d.38-popescu}

Reading source: \href{https://stacks.math.columbia.edu/tag/07GC}{Stacks, Tag 07GC}.

If \(R\to\Lambda\) is a regular map between Noetherian rings, then \(\Lambda\), as an \(R\)-algebra, is a filtered colimit of smooth \(R\)-algebras.

\textbf{Proof.} The reduction \hyperref[AG-RG-S05-smoothing-lemma-reduce-to-field]{D.26} reduces the assertion to a field \(k\) and a Noetherian geometrically regular \(k\)-algebra \(\Lambda\). By \hyperref[AG-RG-S05-smoothing-lemma-final-solve]{D.8}, it is enough that every finite type factorization \(k\to A\to\Lambda\) can be followed by \(A\to B\to\Lambda\) with \(B\) finite type and \(\mathfrak h_B=\Lambda\). All such algebras are finitely presented by the Hilbert theorem.

Assume some factorization has no such successor. Among the singular radicals of all these unsuccessful factorizations choose a maximal one, say \(\mathfrak h_A\). A nonempty collection of ideals in a Noetherian ring has a maximal member: otherwise successive strict enlargements would contradict the ascending chain condition. The ideal \(\mathfrak h_A\) is proper, since if it were \(\Lambda\), \(A\) itself would be a successor.

There is a prime \(\mathfrak q\) minimal over \(\mathfrak h_A\). Here are the existence details. Zorn's lemma first gives an ideal maximal among the proper ideals containing \(\mathfrak h_A\): a chain union is still proper, since it cannot contain \(1\). Such an ideal is prime, because if \(a\) is outside it, the ideal and \(a\) generate \(1\), and multiplication of that relation by \(b\) shows that \(ab\) in the ideal forces \(b\) into it. Now order all primes containing \(\mathfrak h_A\) by reverse inclusion. A chain intersection is a proper ideal containing \(\mathfrak h_A\), and is prime: if \(ab\) belongs to the intersection while \(a,b\) do not, two members excluding them have a smaller member excluding both, contradicting its primality. Zorn's lemma therefore supplies the desired minimal prime.

Geometric regularity makes \(\Lambda_{\mathfrak q}\) regular local. In positive characteristic \hyperref[AG-RG-S05-smoothing-lemma-resolve-general]{D.37} resolves this problem. In characteristic zero every field extension is separable in the sense of \hyperref[AG-RG-S05-algebra-definition-separable-field-extension]{F.5}: each finitely generated subextension has a finite transcendence basis followed by a finite algebraic extension, whose irreducible polynomials have nonzero derivatives and hence distinct roots. Thus \hyperref[AG-RG-S05-smoothing-lemma-resolve-special]{D.33} resolves it in that case.

In either case obtain \(A\to C\to\Lambda\) with \[
\mathfrak h_A\subset\mathfrak h_C,\qquad
\mathfrak h_C\not\subset\mathfrak q.
\] Because \(\mathfrak h_A\subset\mathfrak q\), the first inclusion is strict. Maximality of the unsuccessful ideal \(\mathfrak h_A\) means that \(C\) has a successor with unit singular radical. Composing its factorization with \(A\to C\) supplies the forbidden successor of \(A\), a contradiction. Hence every finite type input has the required successor. D.8 supplies the smooth filtered presentation, and D.26 returns the assertion to the original regular map. ∎

\section{5. Foundation statements and their free-source proofs}\label{AG-RG-S05-foundation-statements-and-their-free-source-proofs}

All resolution comparisons, balance of Tor and connecting exact sequences used in the following flatness arguments are proved for arbitrary modules in the earlier \href{https://kokunoyumeto.github.io/open-math-courses/courses/AG-RG/AG-RG-S01.html\#tor-calculus}{AG-RG-S01, Lemma 1.A}. This binds their homological inputs to a written programme proof before the approximation argument.

These proofs are included as mathematical text, rather than cited as substitutes for proofs. Their transitive dependencies are checked against the earlier programme locators and the open-obligation register in §6. Inclusion of a proof does not automatically close its unproved inputs.

\phantomsection\label{AG-RG-S05-cotangent-context}

\textbf{Cotangent-complex context.} For a polynomial presentation \(P=R[X_s]\twoheadrightarrow S\) with kernel \(I\), write \(\operatorname{NL}(P\to S)\) for the two-term complex \([I/I^2\to\bigoplus_sS\,dX_s]\) in degrees \(-1,0\). The differential takes the class of \(f\) to \(df\). The canonical presentation is \(R[X_s\mid s\in S]\to S\), \(X_s\mapsto s\). For a commuting square of ring maps \(R\to R'\), \(S\to S'\), a presentation map is a map of these polynomial algebras lifting \(S\to S'\). The full comparison and homotopy calculations below show independence of the polynomial presentation. Localization is likewise the displayed two-term localization proved at the corresponding foundation lemma, not an import of the whole cotangent chapter.

\phantomsection\label{AG-RG-S05-inseparable-exponent-context}

\textbf{The inseparable exponent calculation.} Let \(\alpha\) be algebraic over a field \(F\) of characteristic \(p\). Its irreducible polynomial is \(f(T)=g(T^{p^s})\) for a maximal \(s\geq0\). Such \(s\) exists because \(f\) has finite positive degree. Maximality gives \(g'\ne0\). Any factorization of \(g\) would give a factorization of \(f\), so \(g\) is irreducible, and its relatively prime derivative makes it separable. The element \(\alpha^{p^s}\) is a root of \(g\), hence is separable over \(F\). This proves the field-theoretic exponent choice used in the enlargement argument.

\phantomsection\label{AG-RG-S05-dimension-context}

\textbf{Dimension notation.} Krull dimension is the supremum of the lengths of finite chains of prime ideals; the height of a prime is the dimension of its localization. A system of parameters of a Noetherian local ring of dimension \(d\) is a list of \(d\) elements generating an ideal whose radical is the maximal ideal. Their existence is AG-CA-11, Theorem 2.1. Minimal primes, irreducible components and the dimension of polynomial local rings are supplied respectively by AG-CA-01, Theorems 4.1--4.3, AG-CA-09, Theorem 4.3, and AG-CA-11, Theorem 6.1. This paragraph fixes notation; it does not assert an additional dimension theorem. The separate dimension-in-a-neighbourhood and codimension arguments consumed below remain separately registered until matched or proved.

\phantomsection\label{AG-RG-S05-algebra-definition-lci}

\subsubsection{Definition F.1: lci}\label{AG-RG-S05-definition-f.1-lci}

\emph{Adapted from the Stacks project, algebra.tex, lines 37874--37881.}
\nopagebreak[4]

A ring map \(R \to S\) is called \emph{syntomic}, or we say \(S\) is a \emph{flat local complete intersection over \(R\)} if it is flat, of finite presentation, and if all of its fibre rings \(S \otimes_R \kappa(\mathfrak p)\) are local complete intersections, see Definition \hyperref[AG-RG-S05-algebra-definition-lci-field]{F.2}.

\phantomsection\label{AG-RG-S05-algebra-definition-lci-field}

\subsubsection{Definition F.2: lci field}\label{AG-RG-S05-definition-f.2-lci-field}

\emph{Adapted from the Stacks project, algebra.tex, lines 37249--37264.}
\nopagebreak[4]

Let \(k\) be a field. Let \(S\) be a finite type \(k\)-algebra.

\begin{enumerate}
\def\labelenumi{\arabic{enumi}.}
\item
  We say that \(S\) is a \emph{global complete intersection over \(k\)} if there exists a presentation \(S = k[x_1, \ldots, x_n]/(f_1, \ldots, f_c)\) such that \(\dim(S) = n - c\).
\item
  We say that \(S\) is a \emph{local complete intersection over \(k\)} if there exists a covering \(\operatorname{Spec}(S) = \bigcup D(g_i)\) such that each of the rings \(S_{g_i}\) is a global complete intersection over \(k\).
\end{enumerate}

We will also use the convention that the zero ring is a global complete intersection over \(k\).

\phantomsection\label{AG-RG-S05-algebra-definition-lci-local-ring}

\subsubsection{Definition F.3: lci local ring}\label{AG-RG-S05-definition-f.3-lci-local-ring}

\emph{Adapted from the Stacks project, algebra.tex, lines 37446--37458.}
\nopagebreak[4]

Let \(k\) be a field. Let \(S\) be a local \(k\)-algebra essentially of finite type over \(k\). We say \(S\) is a \emph{complete intersection (over \(k\))} if there exists a local \(k\)-algebra \(R\) and elements \(f_1, \ldots, f_c \in \mathfrak m_R\) such that

\begin{enumerate}
\def\labelenumi{\arabic{enumi}.}
\item
  \(R\) is essentially of finite type over \(k\),
\item
  \(R\) is a regular local ring,
\item
  \(f_1, \ldots, f_c\) form a regular sequence in \(R\), and
\item
  \(S \cong R/(f_1, \ldots, f_c)\) as \(k\)-algebras.
\end{enumerate}

\phantomsection\label{AG-RG-S05-algebra-definition-relative-global-complete-intersection}

\subsubsection{Definition F.4: relative global complete intersection}\label{AG-RG-S05-definition-f.4-relative-global-complete-intersection}

\emph{Adapted from the Stacks project, algebra.tex, lines 37943--37951.}
\nopagebreak[4]

Let \(R \to S\) be a ring map. We say that \(R \to S\) is a \emph{relative global complete intersection} if there exists a presentation \(S = R[x_1, \ldots, x_n]/(f_1, \ldots, f_c)\) and every nonempty fibre of \(\operatorname{Spec}(S) \to \operatorname{Spec}(R)\) has dimension \(n - c\). We will say ``let \(S = R[x_1, \ldots, x_n]/(f_1, \ldots, f_c)\) be a relative global complete intersection'\,' to indicate this situation.

\phantomsection\label{AG-RG-S05-algebra-definition-separable-field-extension}

\subsubsection{Definition F.5: separable field extension}\label{AG-RG-S05-definition-f.5-separable-field-extension}

\emph{Adapted from the Stacks project, algebra.tex, lines 10135--10146.}
\nopagebreak[4]

Let \(K/k\) be a field extension.

\begin{enumerate}
\def\labelenumi{\arabic{enumi}.}
\item
  We say \(K\) is \emph{separably generated over \(k\)} if there exists a transcendence basis \(\{x_i; i \in I\}\) of \(K/k\) such that the extension \(K/k(x_i; i \in I)\) is a separable algebraic extension.
\item
  We say \(K\) is \emph{separable over \(k\)} if for every subextension \(k \subset K' \subset K\) with \(K'\) finitely generated over \(k\), the extension \(K'/k\) is separably generated.
\end{enumerate}

\phantomsection\label{AG-RG-S05-algebra-definition-smooth}

\subsubsection{Definition F.6: smooth}\label{AG-RG-S05-definition-f.6-smooth}

\emph{Adapted from the Stacks project, algebra.tex, lines 38526--38533.}
\nopagebreak[4]

A ring map \(R \to S\) is \emph{smooth} if it is of finite presentation and the naive cotangent complex \(\NL_{S/R}\) is quasi-isomorphic to a finite projective \(S\)-module placed in degree \(0\): this means that \(H_1(\NL_{S/R}) = 0\) and that \(\Omega_{S/R}\) is a finite projective \(S\)-module.

\phantomsection\label{AG-RG-S05-algebra-definition-smooth-at-prime}

\subsubsection{Definition F.7: smooth at prime}\label{AG-RG-S05-definition-f.7-smooth-at-prime}

\emph{Adapted from the Stacks project, algebra.tex, lines 38953--38960.}
\nopagebreak[4]

Let \(R \to S\) be a ring map. Let \(\mathfrak q\) be a prime of \(S\). We say \(R \to S\) is \emph{smooth at \(\mathfrak q\)} if there exists a \(g \in S\), \(g \not \in \mathfrak q\) such that \(R \to S_g\) is smooth.

\phantomsection\label{AG-RG-S05-algebra-definition-unramified}

\subsubsection{Definition F.8: unramified}\label{AG-RG-S05-definition-f.8-unramified}

\emph{Adapted from the Stacks project, algebra.tex, lines 42897--42912.}
\nopagebreak[4]

Let \(R \to S\) be a ring map.

\begin{enumerate}
\def\labelenumi{\arabic{enumi}.}
\item
  We say \(R \to S\) is \emph{unramified} if \(R \to S\) is of finite type and \(\Omega_{S/R} = 0\).
\item
  We say \(R \to S\) is \emph{G-unramified} if \(R \to S\) is of finite presentation and \(\Omega_{S/R} = 0\).
\item
  Given a prime \(\mathfrak q\) of \(S\) we say that \(S\) is \emph{unramified at \(\mathfrak q\)} if there exists a \(g \in S\), \(g \not \in \mathfrak q\) such that \(R \to S_g\) is unramified.
\item
  Given a prime \(\mathfrak q\) of \(S\) we say that \(S\) is \emph{G-unramified at \(\mathfrak q\)} if there exists a \(g \in S\), \(g \not \in \mathfrak q\) such that \(R \to S_g\) is G-unramified.
\end{enumerate}

\phantomsection\label{AG-RG-S05-algebra-lemma-CM-over-regular-flat}

\subsubsection{Lemma F.9: CM over regular flat}\label{AG-RG-S05-lemma-f.9-cm-over-regular-flat}

\emph{Adapted from the Stacks project, algebra.tex, lines 34036--34074.}
\nopagebreak[4]

Let \(R \to S\) be a local homomorphism of Noetherian local rings. Assume

\begin{enumerate}
\def\labelenumi{\arabic{enumi}.}
\item
  \(R\) is regular,
\item
  \(S\) is Cohen-Macaulay,
\item
  \(\dim(S) = \dim(R) + \dim(S/\mathfrak m_R S)\).
\end{enumerate}

Then \(R \to S\) is flat.

\textbf{Proof.} Induct on \(d=\dim R\). If \(d=0\), regularity and Nakayama give \(\mathfrak m_R=0\), so \(R\) is a field and every \(R\)-module is flat. Suppose \(d>0\). Hypothesis (3) gives \(\dim S>\dim(S/\mathfrak m_RS)\). The CM associated-component theorem AG-CA-12, Theorem 4.1, identifies the associated primes of \(S\) with its finitely many minimal primes \(\mathfrak q_i\), and gives \(\dim(S/\mathfrak q_i)=\dim S\). Hence \(\mathfrak m_RS\) cannot be contained in any \(\mathfrak q_i\): such a containment would make \(S/\mathfrak q_i\) a quotient of \(S/\mathfrak m_RS\), contradicting these dimensions.

Thus the contractions \(\mathfrak p_i=R\cap\mathfrak q_i\) are proper subprimes of \(\mathfrak m_R\). By \hyperref[AG-RG-S05-algebra-lemma-silly]{F.80}, choose \(x\in\mathfrak m_R\setminus\mathfrak m_R^2\) outside all \(\mathfrak p_i\). It avoids every associated prime of \(S\), so multiplication by \(x\) on \(S\) is injective. The regular-quotient and CM dimension calculations in AG-CA-12, Sections 2 and 4, make \(S/xS\) CM of dimension \(\dim S-1\). AG-CA-14, Theorem 1.1 and Corollary 1.2, make \(x\) a nonzerodivisor on \(R\) and \(R/xR\) regular of dimension \(d-1\).

The induced quotient map is local and its closed fibre is still \(S/\mathfrak m_RS\), since \(x\in\mathfrak m_R\). Subtracting one from both total dimensions in (3) gives \[\dim(S/xS)=\dim(R/xR)+\dim(S/\mathfrak m_RS).\] Every induction hypothesis therefore holds, so \(S/xS\) is flat over \(R/xR\).

The exact free resolution \(0\to R\xrightarrow{x}R\to R/xR\to0\) identifies \(\operatorname{Tor}_1^R(R/xR,S)\) with the kernel of multiplication by \(x\) on \(S\), which is zero. Apply AG-CA-08, Corollary 4.4, to the proper ideal \((x)\subset R\) and the module \(S\), finite over the stipulated Noetherian local overring \(S\). Its quotient is flat and its first Tor vanishes, so \(S\) is flat over \(R\). ∎

\phantomsection\label{AG-RG-S05-algebra-lemma-NAK}

\subsubsection{Lemma F.10: NAK}\label{AG-RG-S05-lemma-f.10-nak}

\emph{Adapted from the Stacks project, algebra.tex, lines 3732--3813.}
\nopagebreak[4]

Let \(R\) be a ring with Jacobson radical \(\text{rad}(R)\). Let \(M\) be an \(R\)-module. Let \(I \subset R\) be an ideal.

\begin{enumerate}
\def\labelenumi{\arabic{enumi}.}
\item
  If \(IM = M\) and \(M\) is finite, then there exists an \(f \in 1 + I\) such that \(fM = 0\).
\item
  If \(IM = M\), \(M\) is finite, and \(I \subset \text{rad}(R)\), then \(M = 0\).
\item
  If \(N, N' \subset M\), \(M = N + IN'\), and \(N'\) is finite, then there exists an \(f \in 1 + I\) such that \(fM \subset N\) and \(M_f = N_f\).
\item
  If \(N, N' \subset M\), \(M = N + IN'\), \(N'\) is finite, and \(I \subset \text{rad}(R)\), then \(M = N\).
\item
  If \(N \to M\) is a module map, \(N/IN \to M/IM\) is surjective, and \(M\) is finite, then there exists an \(f \in 1 + I\) such that \(N_f \to M_f\) is surjective.
\item
  If \(N \to M\) is a module map, \(N/IN \to M/IM\) is surjective, \(M\) is finite, and \(I \subset \text{rad}(R)\), then \(N \to M\) is surjective.
\item
  If \(x_1, \ldots, x_n \in M\) generate \(M/IM\) and \(M\) is finite, then there exists an \(f \in 1 + I\) such that \(x_1, \ldots, x_n\) generate \(M_f\) over \(R_f\).
\item
  If \(x_1, \ldots, x_n \in M\) generate \(M/IM\), \(M\) is finite, and \(I \subset \text{rad}(R)\), then \(M\) is generated by \(x_1, \ldots, x_n\).
\item
  If \(IM = M\), \(I\) is nilpotent, then \(M = 0\).
\item
  If \(N, N' \subset M\), \(M = N + IN'\), and \(I\) is nilpotent then \(M = N\).
\item
  If \(N \to M\) is a module map, \(I\) is nilpotent, and \(N/IN \to M/IM\) is surjective, then \(N \to M\) is surjective.
\item
  If \(\{x_\alpha\}_{\alpha \in A}\) is a set of elements of \(M\) which generate \(M/IM\) and \(I\) is nilpotent, then \(M\) is generated by the \(x_\alpha\).
\end{enumerate}

\textbf{Proof.} Proof of (\hyperref[AG-RG-S05-algebra-lemma-NAK]{F.10 / nakayama}). Choose generators \(y_1, \ldots, y_m\) of \(M\) over \(R\). For each \(i\) we can write \(y_i = \sum z_{ij} y_j\) with \(z_{ij} \in I\) (since \(M = IM\)). In other words \(\sum_j (\delta_{ij} - z_{ij})y_j = 0\). Let \(f\) be the determinant of the \(m \times m\) matrix \(A = (\delta_{ij} - z_{ij})\). Note that \(f \in 1 + I\) (since the matrix \(A\) is entrywise congruent to the \(m \times m\) identity matrix modulo \(I\)). By Lemma \hyperref[AG-RG-S05-algebra-lemma-matrix-left-inverse]{F.70} (1), there exists an \(m \times m\) matrix \(B\) such that \(BA = f 1_{m \times m}\). Writing out we see that \(\sum_{i} b_{hi} a_{ij} = f \delta_{hj}\) for all \(h\) and \(j\); hence, \(\sum_{i, j} b_{hi} a_{ij} y_j
= \sum_{j} f \delta_{hj} y_j = f y_h\) for every \(h\). In other words, \(0 = f y_h\) for every \(h\) (since each \(i\) satisfies \(\sum_j a_{ij} y_j = 0\)). This implies that \(f\) annihilates \(M\).

By Lemma \hyperref[AG-RG-S05-algebra-lemma-contained-in-radical]{F.34} an element of \(1 + \text{rad}(R)\) is an invertible element of \(R\). Hence we see that (\hyperref[AG-RG-S05-algebra-lemma-NAK]{F.10 / nakayama}) implies (2). We obtain (3) by applying (1) to \(M/N\) which is finite as \(N'\) is finite. We obtain (4) by applying (2) to \(M/N\) which is finite as \(N'\) is finite. We obtain (5) by applying (3) to \(M\) and the submodules \(\operatorname{im}(N \to M)\) and \(M\). We obtain (6) by applying (4) to \(M\) and the submodules \(\operatorname{im}(N \to M)\) and \(M\). We obtain (7) by applying (5) to the map \(R^{\oplus n} \to M\), \((a_1, \ldots, a_n) \mapsto a_1x_1 + \ldots + a_nx_n\). We obtain (8) by applying (6) to the map \(R^{\oplus n} \to M\), \((a_1, \ldots, a_n) \mapsto a_1x_1 + \ldots + a_nx_n\).

Part (9) holds because if \(M = IM\) then \(M = I^nM\) for all \(n \geq 0\) and \(I\) being nilpotent means \(I^n = 0\) for some \(n \gg 0\). Parts (10), (11), and (12) follow from (9) by the arguments used above. ∎

\phantomsection\label{AG-RG-S05-algebra-lemma-NL-homotopy}

\subsubsection{Lemma F.11: NL homotopy}\label{AG-RG-S05-lemma-f.11-nl-homotopy}

\emph{Adapted from the Stacks project, algebra.tex, lines 36700--36815.}
\nopagebreak[4]

Let \(R\to R'\) and \(S\to S'\) be a commutative square of ring maps with structure maps \(R\to S\) and \(R'\to S'\). Let \(\alpha : P \to S\) and \(\alpha' : P' \to S'\) be presentations.

\begin{enumerate}
\def\labelenumi{\arabic{enumi}.}
\item
  There exists a morphism of presentations from \(\alpha\) to \(\alpha'\).
\item
  Any two morphisms of presentations induce homotopic morphisms of complexes \(\operatorname{NL}(\alpha) \to \operatorname{NL}(\alpha')\).
\item
  The construction is compatible with compositions of morphisms of presentations (see proof for exact statement).
\item
  If \(R \to R'\) and \(S \to S'\) are isomorphisms, then for any map \(\varphi\) of presentations from \(\alpha\) to \(\alpha'\) the induced map \(\operatorname{NL}(\alpha) \to \operatorname{NL}(\alpha')\) is a homotopy equivalence and a quasi-isomorphism.
\end{enumerate}

In particular, comparing \(\alpha\) to the canonical presentation described in the cotangent-complex context above we conclude there is a quasi-isomorphism \(\operatorname{NL}(\alpha) \to \NL_{S/R}\) well defined up to homotopy and compatible with all functorialities (up to homotopy).

\textbf{Proof.} Since \(P\) is a polynomial algebra over \(R\) we can write \(P = R[x_a, a \in A]\) for some set \(A\). As \(\alpha'\) is surjective, we can choose for every \(a \in A\) an element \(f_a \in P'\) such that \(\alpha'(f_a) = \phi(\alpha(x_a))\). Let \(\varphi : P = R[x_a, a \in A] \to P'\) be the unique \(R\)-algebra map such that \(\varphi(x_a) = f_a\). This gives the morphism in (1).

Let \(\varphi\) and \(\varphi'\) be morphisms of presentations from \(\alpha\) to \(\alpha'\). Let \(I = \ker(\alpha)\) and \(I' = \ker(\alpha')\). We have to construct the diagonal map \(h\) in the diagram

\[
\begin{gathered}
I/I^2 {\text{d}} \xrightarrow{-} \Omega_{P/R} \otimes_P S \\
I/I^2 {\text{d}} \xrightarrow{\varphi'_1} I'/(I')^2 {\text{d}} \\
I/I^2 {\text{d}} \xrightarrow{\varphi_1} I'/(I')^2 {\text{d}} \\
\Omega_{P/R} \otimes_P S \xrightarrow{\varphi'_0} \Omega_{P'/R'} \otimes_{P'} S' \\
\Omega_{P/R} \otimes_P S \xrightarrow{\varphi_0} \Omega_{P'/R'} \otimes_{P'} S' \\
\Omega_{P/R} \otimes_P S \xrightarrow{h} I'/(I')^2 {\text{d}} \\
I'/(I')^2 {\text{d}} \xrightarrow{-} \Omega_{P'/R'} \otimes_{P'} S'
\end{gathered}
\]

where the vertical maps are induced by \(\varphi\) and \(\varphi'\), and \(h\) must satisfy

\[
\varphi_1 - \varphi'_1 = h \circ \text{d}
\quad\text{and}\quad
\varphi_0 - \varphi'_0 = \text{d} \circ h
\]

Consider the map \(\varphi - \varphi' : P \to P'\). Since both \(\varphi\) and \(\varphi'\) are compatible with \(\alpha\) and \(\alpha'\) we obtain \(\varphi - \varphi' : P \to I'\). This implies that \(\varphi, \varphi' : P \to P'\) induce the same \(P\)-module structure on \(I'/(I')^2\), since \(\varphi(p)i' - \varphi'(p)i' = (\varphi - \varphi')(p)i' \in (I')^2\). Also \(\varphi - \varphi'\) is \(R\)-linear and

\[
(\varphi - \varphi')(fg) =
\varphi(f)(\varphi - \varphi')(g) + (\varphi - \varphi')(f)\varphi'(g)
\]

Hence the induced map \(D : P \to I'/(I')^2\) is an \(R\)-derivation. Thus we obtain a canonical map \(h : \Omega_{P/R} \otimes_P S \to I'/(I')^2\) such that \(D = h \circ \text{d}\). For \(f\in I\), the definition of the universal derivation gives \(h(df)=D(f)=\varphi(f)-\varphi'(f)\bmod (I')^2\), so the degree-one difference is \(h\circ d\). On a basis element \(dx_a\) of \(\Omega_{P/R}\otimes_PS\), one has \(d(h(dx_a))=d(\varphi(x_a)-\varphi'(x_a))\), exactly the difference of the degree-zero maps. These two identities verify the homotopy on generators, hence on the whole complexes.

Suppose that we have a commutative diagram

\[
\begin{gathered}
S \xrightarrow{\phi} S' \\
S' \xrightarrow{\phi'} S'' \\
R \xrightarrow{} R' \\
R \xrightarrow{} S \\
R' \xrightarrow{} S' \\
R' \xrightarrow{} R'' \\
R'' \xrightarrow{} S''
\end{gathered}
\]

and that

\begin{enumerate}
\def\labelenumi{\arabic{enumi}.}
\item
  \(\alpha : P \to S\),
\item
  \(\alpha' : P' \to S'\), and
\item
  \(\alpha'' : P'' \to S''\)
\end{enumerate}

are presentations. Suppose that

\begin{enumerate}
\def\labelenumi{\arabic{enumi}.}
\item
  \(\varphi : P \to P'\) is a morphism of presentations from \(\alpha\) to \(\alpha'\) and
\item
  \(\varphi' : P' \to P''\) is a morphism of presentations from \(\alpha'\) to \(\alpha''\).
\end{enumerate}

Then it is immediate that \(\varphi' \circ \varphi : P \to P''\) is a morphism of presentations from \(\alpha\) to \(\alpha''\) and that the induced map \(\operatorname{NL}(\alpha) \to \operatorname{NL}(\alpha'')\) of naive cotangent complexes is the composition of the maps \(\operatorname{NL}(\alpha) \to \operatorname{NL}(\alpha')\) and \(\operatorname{NL}(\alpha') \to \operatorname{NL}(\alpha'')\) induced by \(\varphi\) and \(\varphi'\).

In the simple case of complexes with 2 terms a quasi-isomorphism is just a map that induces an isomorphism on both the cokernel and the kernel of the maps between the terms. Note that homotopic maps of 2 term complexes (as explained above) define the same maps on kernel and cokernel. Hence if \(\varphi\) is a map from a presentation \(\alpha\) of \(S\) over \(R\) to itself, then the induced map \(\operatorname{NL}(\alpha) \to \operatorname{NL}(\alpha)\) is a quasi-isomorphism being homotopic to the identity by part (2). To prove (4) in full generality, consider a morphism \(\varphi'\) from \(\alpha'\) to \(\alpha\) which exists by (1). The compositions \(\operatorname{NL}(\alpha) \to \operatorname{NL}(\alpha') \to \operatorname{NL}(\alpha)\) and \(\operatorname{NL}(\alpha') \to \operatorname{NL}(\alpha) \to \operatorname{NL}(\alpha')\) are homotopic to the identity maps by (2) and (3), hence these maps are homotopy equivalences by definition. It follows formally that both maps \(\operatorname{NL}(\alpha) \to \operatorname{NL}(\alpha')\) and \(\operatorname{NL}(\alpha') \to \operatorname{NL}(\alpha)\) are quasi-isomorphisms. Explicitly, if \(uv\) and \(vu\) are homotopic to the respective identities, their maps on kernel and cokernel are identities, since a homotopy contributes \(hd\) in degree one and \(dh\) in degree zero. The maps induced by \(u\) and \(v\) on both groups are therefore inverse isomorphisms. ∎

\phantomsection\label{AG-RG-S05-algebra-lemma-Noetherian-field-extension}

\subsubsection{Lemma F.12: Noetherian field extension}\label{AG-RG-S05-lemma-f.12-noetherian-field-extension}

\emph{Adapted from the Stacks project, algebra.tex, lines 6331--6347.}
\nopagebreak[4]

If \(A\) is a Noetherian \(k\)-algebra and \(K/k\) is a finitely generated field extension, then \(A\otimes_kK\) is Noetherian.

\textbf{Proof.} Choose finitely many field generators and let \(B\subset K\) be the \(k\)-algebra they generate. Its fraction field is \(K\), so \(K=(B\setminus\{0\})^{-1}B\). A finite polynomial presentation of \(B\) identifies \(A\otimes_kB\) with a quotient of a polynomial algebra in finitely many variables over \(A\). Hilbert basis AG-CA-03, Theorem 2.1, makes it Noetherian. Tensor and the fraction universal property identify \(A\otimes_kK\) with its localization at the images of the nonzero elements of \(B\). The same Noetherian permanence proof in that earlier theorem makes this localization Noetherian. In the special case of a finite field extension, one can alternatively use its finite field basis to see a finite free \(A\)-module, whose ideals are finite submodules. This proves the full finitely generated form consumed in the regularity tests. ∎

\phantomsection\label{AG-RG-S05-algebra-lemma-another-variant-local-criterion-flatness}

\subsubsection{Lemma F.13: another variant local criterion flatness}\label{AG-RG-S05-lemma-f.13-another-variant-local-criterion-flatness}

\emph{Adapted from the Stacks project, algebra.tex, lines 24942--24990.}
\nopagebreak[4]

For a local square of Noetherian rings \(R\to S\), \(R'\to S'\) with \(S'\) a localization of \(S\otimes_RR'\), a finite \(S\)-module \(M\), and a proper ideal \(I\subset R\), assume \(M/IM\) is flat over \(R/I\) and the map \(\operatorname{Tor}_1^R(M,R/I)\to\operatorname{Tor}_1^{R'}(M\otimes_SS',R'/IR')\) is zero. Then \(M\otimes_SS'\) is \(R'\)-flat.

\textbf{Proof.} Write a free presentation \(0\to K\to F\to M\to0\) over \(R\), allowing an infinite free \(F\). Put \(T=\operatorname{Tor}_1^R(M,R/I)\) and let \(\bar K\) be the image of \(K/IK\) in \(F/IF\). The two sequences \(0\to T\to K/IK\to\bar K\to0\) and \(0\to\bar K\to F/IF\to M/IM\to0\) are exact. The assumed flatness of the last quotient makes \(\bar K\) flat by AG-CA-07, Theorem 6.1. Tensor these sequences with any \(R/I\)-module \(N\); their quotient flatness preserves their left injections and identifies \[\operatorname{Tor}_1^R(M,N)=T\otimes_{R/I}N.\] This identity follows from the same free-presentation kernel formula for Tor, so does not assume that \(K\) is flat.

After base change to \(R'\), let \(K'\) be the image of \(K\otimes_RR'\) in \(F\otimes_RR'\). The map onto \(K'\) is surjective. Tensor it with \(R'/IR'\); every cycle in \(\ker(K'\otimes R'/IR'\to F\otimes R'/IR')\) lifts to an element of \(K\otimes_RR'/IR'\), and that lift is automatically a cycle because the two free-module targets are identical. Therefore \(\operatorname{Tor}_1^R(M,R'/IR')\to\operatorname{Tor}_1^{R'}(M\otimes_RR',R'/IR')\) is surjective. By the preceding identity its source is \(T\otimes_RR'\). Localizing at the stipulated prime in the algebra commutes with this kernel calculation by exact localization. Thus \(T\otimes_RR'\) surjects onto the Tor group of \(M'=M\otimes_SS'\). The assumed zero map annihilates this surjection, so that Tor group is zero.

Base change of flatness makes \(M'/IR'M'\) flat over \(R'/IR'\). AG-CA-08, Corollary 4.4, is the proved proper-ideal local criterion under these Noetherian local and finite-over-\(S'\) hypotheses; it now makes \(M'\) flat over \(R'\). This proves the transition criterion rather than invoking an unproved change-of-rings spectral sequence. ∎

\phantomsection\label{AG-RG-S05-algebra-lemma-base-change-relative-global-complete-intersection}

\subsubsection{Lemma F.14: base change relative global complete intersection}\label{AG-RG-S05-lemma-f.14-base-change-relative-global-complete-intersection}

\emph{Adapted from the Stacks project, algebra.tex, lines 38076--38106.}
\nopagebreak[4]

Relative global complete-intersection presentations are preserved by base change, principal localization, and replacing the base by a localization already inverted in the algebra.

\textbf{Proof.} First record the required field-dimension argument. For a finite-type algebra over a field and any extension field, dimension is unchanged. Quotient by its finitely many minimal primes; their intersection is nilpotent by AG-CA-03, Proposition 2.2. This remains nilpotent after field extension. For each resulting domain, Noether normalization AG-CA-09, Corollary 3.2, gives an injective finite polynomial subalgebra. Field tensor is exact, so the injection and finiteness persist. Dimension is preserved by an injective integral map, AG-CA-05, Theorem 5.1. Thus each component and their maximum retain dimension. This proof applies to an arbitrary, possibly infinite, field extension.

The fibre after any base change is a field extension of the fibre over the contracted base prime. Consequently its dimension stays \(n-c\). For localization, each original fibre is equidimensional of dimension \(n-c\): the minimal-prime height calculation in \hyperref[AG-RG-S05-algebra-lemma-relative-global-complete-intersection-conormal]{F.77}, applied just to the polynomial algebra over the fibre field, proves this without any base Noetherianity. A nonempty principal open in an irreducible finite-type field scheme has unchanged dimension, since its coordinate domain has the same fraction field and AG-CA-09, Theorem 4.2, identifies dimension with transcendence degree. Therefore every nonempty localized fibre still has dimension \(n-c\). The presentation \(S_g=R[X,Z]/(f_1,\ldots,f_c,hZ-1)\) has one additional variable and equation, hence the same expected dimension. Finally if a base element is already inverted in \(S\), tensoring with that base localization gives back \(S\), proving the third assertion. ∎

\phantomsection\label{AG-RG-S05-algebra-lemma-change-base-NL}

\subsubsection{Lemma F.15: change base NL}\label{AG-RG-S05-lemma-f.15-change-base-nl}

\emph{Adapted from the Stacks project, algebra.tex, lines 36963--36981.}
\nopagebreak[4]

For a flat base change \(R\to R'\), the two-term conormal complex of a polynomial presentation \(P\twoheadrightarrow S\) base changes to the corresponding complex for \(P\otimes_RR'\twoheadrightarrow S\otimes_RR'\).

\textbf{Proof.} Let \(I\) be the presentation kernel. Flatness identifies the new kernel with \(I\otimes_RR'\), since it preserves the exact sequence \(0\to I\to P\to S\to0\). Its square is the image of \(I^2\otimes_RR'\), by writing products of generators. The cokernel description therefore identifies the conormal module with \((I/I^2)\otimes_RR'\). The differential module of the polynomial ring is free on the same \(dX_s\) before and after base change. Both differentials send the class of a polynomial \(f\) to \(df\), so these are an isomorphism of complexes. Comparison of polynomial presentations and its explicit homotopies, proved above, give the canonical homotopy equivalence. ∎

\phantomsection\label{AG-RG-S05-algebra-lemma-characterize-etale}

\subsubsection{Lemma F.16: characterize etale}\label{AG-RG-S05-lemma-f.16-characterize-etale}

\emph{Adapted from the Stacks project, algebra.tex, lines 40942--40975.}
\nopagebreak[4]

Let \(R \to S\) be a ring map. Let \(\mathfrak q\) be a prime of \(S\) lying over a prime \(\mathfrak p\) of \(R\). If

\begin{enumerate}
\def\labelenumi{\arabic{enumi}.}
\item
  \(R \to S\) is of finite presentation,
\item
  \(R_{\mathfrak p} \to S_{\mathfrak q}\) is flat
\item
  \(\mathfrak p S_{\mathfrak q}\) is the maximal ideal of the local ring \(S_{\mathfrak q}\), and
\item
  the field extension \(\kappa(\mathfrak q)/\kappa(\mathfrak p)\) is finite separable,
\end{enumerate}

then \(R \to S\) is étale at \(\mathfrak q\).

\textbf{Proof.} Put \(k=\kappa(\mathfrak p)\) and \(L=\kappa(\mathfrak q)\). The local fibre ring is \(S_{\mathfrak q}/\mathfrak pS_{\mathfrak q}=L\) by (3). Since \(L/k\) is finite, \hyperref[AG-RG-S05-algebra-lemma-isolated-point-fibre]{F.56} supplies \(g\notin\mathfrak q\) for which \(\mathfrak q\) is the only point of the fibre of \(S_g\). Replace \(S\) by \(S_g\); finite presentation is preserved by \hyperref[AG-RG-S05-algebra-lemma-cover-upstairs]{F.35}. The fibre \(B=S\otimes_R k\) has a unique prime, so is local: every element outside its unique maximal ideal is a unit. Its localization there is \(S_{\mathfrak q}/\mathfrak pS_{\mathfrak q}\), hence \(B=L\).

Here is the finite separable fibre calculation, including its finite presentation. Choose a finite generating tower \(k=E_0\subset E_1\subset\cdots\subset E_t=L\), with \(E_i=E_{i-1}(\alpha_i)\). A finite \(k\)-basis provides such a generating list. Let \(p_i(T)\) be the monic minimal polynomial of \(\alpha_i\) over \(E_{i-1}\). Polynomial division gives \(E_i=E_{i-1}[T]/(p_i)\); the quotient is a field by polynomial Bezout. These full field calculations are proved in \href{https://kokunoyumeto.github.io/open-math-courses/courses/AG-RG/AG-RG-S02.html}{AG-RG-S02, Lemma 4.A}. Represent the coefficients of \(p_i\) by polynomials in the preceding generators, obtaining \(P_i(T_1,\ldots,T_i)\). Successively taking these monic quotients gives the finite presentation \[
L=k[T_1,\ldots,T_t]/(P_1,\ldots,P_t).
\] Put \(I=(P_1,\ldots,P_t)\subset k[T_1,\ldots,T_t]\). The minimal polynomial of \(\alpha_i\) over \(E_{i-1}\) divides its separable minimal polynomial over \(k\). A repeated root of the former would be a repeated root of the latter, so \(p_i'(\alpha_i)\ne0\), as also proved in Lemma 4.A. The Jacobian of this presentation over \(L\) is triangular with these nonzero diagonal entries, hence invertible. The map \(L^t\to I/I^2\) given by the equation classes is surjective, and its composite with the conormal differential to \(L^t\) is this invertible matrix. Its kernel is therefore zero, and the conormal differential is an isomorphism. Thus \(H_1(\NL_{L/k})=0\) and \(\Omega_{L/k}=0\), proving that \(L/k\) is smooth by \hyperref[AG-RG-S05-algebra-definition-smooth]{Definition F.6}. The empty tower gives the same assertion for \(L=k\).

The stated local flatness and finite presentation now satisfy \hyperref[AG-RG-S05-algebra-lemma-flat-fibre-smooth]{F.45}, so another principal neighbourhood of \(\mathfrak q\) is smooth over \(R\). By the arbitrary-base standard-chart assertion of \hyperref[AG-RG-S05-algebra-lemma-smooth-syntomic]{F.83}, shrink once more to a standard smooth presentation \[
S=R[X_1,\ldots,X_n]/(f_1,\ldots,f_c).
\] Every element inverted in these shrinkings has nonzero image in the field \(L\), so the fibre remains \(L\). By \hyperref[AG-RG-S05-algebra-lemma-standard-smooth]{F.85(2)}, \(\Omega_{S/R}\) is free on \(dX_{c+1},\ldots,dX_n\). The complete base-change calculation in AG-CA-16, Theorem 2.1, identifies its tensor with \(L\) with \(\Omega_{L/k}=0\). Hence \(L^{n-c}=0\); a coordinate vector in a nonzero field module shows that this forces \(n=c\).

For clarity, the square Jacobian gives étaleness directly. Given an \(R\)-map \(S\to D/J\) with \(J^2=0\), lift its variable images to a tuple \(a\in D^n\). Its error vector \(f(a)\) lies in \(J^n\). The Jacobian matrix \(A(a)\) is invertible: its determinant is a unit modulo \(J\), and a lift of a modular inverse gives \(1+j\), whose inverse is \(1-j\); the cofactor identity of \hyperref[AG-RG-S05-algebra-lemma-matrix-left-inverse]{F.70} then gives the matrix inverse. Put \(b=a-A(a)^{-1}f(a)\). Polynomial expansion gives \(f(b)=f(a)+A(a)(b-a)=0\), since terms with two corrections vanish in \(J^2\). Any other lifting root is \(b+\epsilon\) with \(\epsilon\in J^n\); the same expansion gives \(A(b)\epsilon=0\), so \(\epsilon=0\). Thus every square-zero lifting test has a unique solution. Finite presentation makes this algebra étale by the programme definition, so the original map is étale at \(\mathfrak q\). Empty variable lists use determinant one and the unique empty tuple. No Noetherian or perfect-field assumption was used. ∎

\phantomsection\label{AG-RG-S05-algebra-lemma-characterize-formally-smooth-field-extension}

\subsubsection{Lemma F.17: characterize formally smooth field extension}\label{AG-RG-S05-lemma-f.17-characterize-formally-smooth-field-extension}

\emph{Adapted from the Stacks project, algebra.tex, lines 45959--45969.}
\nopagebreak[4]

For a field extension \(K/k\), formal smoothness is equivalent to \(H_1(\operatorname{NL}_{K/k})=0\).

\textbf{Proof.} Use any polynomial presentation, including an infinite one. Its first homology is the kernel of the conormal differential by definition. If it vanishes, the conormal sequence is short exact and its quotient \(\Omega_{K/k}\) is a vector space, hence free and projective; choosing and lifting a basis splits it. AG-CA-17, Theorem 3.1, with arbitrary variable sets then proves formal smoothness. Conversely that theorem gives split exactness and hence zero kernel. Polynomial comparison already proved identifies this homology with the presentation-independent \(H_1(L_{K/k})\) used in the statement. ∎

\phantomsection\label{AG-RG-S05-algebra-lemma-characterize-smooth-kbar}

\subsubsection{Lemma F.18: characterize smooth kbar}\label{AG-RG-S05-lemma-f.18-characterize-smooth-kbar}

\emph{Adapted from the Stacks project, algebra.tex, lines 39962--40064.}
\nopagebreak[4]

Let \(k\) be an algebraically closed field. Let \(S\) be a finite type \(k\)-algebra. Let \(\mathfrak m \subset S\) be a maximal ideal. The following are equivalent:

\begin{enumerate}
\def\labelenumi{\arabic{enumi}.}
\item
  The ring \(S_{\mathfrak m}\) is a regular local ring.
\item
  We have \(\dim_{\kappa(\mathfrak m)} \Omega_{S/k} \otimes_S \kappa(\mathfrak m)
  \leq \dim(S_{\mathfrak m})\).
\item
  We have \(\dim_{\kappa(\mathfrak m)} \Omega_{S/k} \otimes_S \kappa(\mathfrak m)
  = \dim(S_{\mathfrak m})\).
\item
  There exists a \(g \in S\), \(g \not \in \mathfrak m\) such that \(S_g\) is smooth over \(k\). In other words \(S/k\) is smooth at \(\mathfrak m\).
\end{enumerate}

\textbf{Proof.} By AG-CA-18, Proposition 3.1, the residue field at this closed point is \(k\) and \[\Omega_{S/k}\otimes_S k\cong\mathfrak m/\mathfrak m^2.\] AG-CA-14, Section 1, gives \(\dim S_{\mathfrak m}\le\dim_k\mathfrak m/\mathfrak m^2\), with equality precisely for a regular local ring. Consequently (1), (2) and (3) are equivalent. The same complete earlier Proposition 3.1 proves that smoothness at this rational closed point implies regularity, so (4) implies (1).

We write the converse construction explicitly. Put \(d=\dim S_{\mathfrak m}\) and choose a presentation \(S=k[x_1,\ldots,x_n]/I\) after subtracting the residue values of the generators. Its polynomial preimage is \(\mathfrak n=(x_1,\ldots,x_n)\). The conormal sequence at the rational point is \[I/\mathfrak nI\longrightarrow\mathfrak n/\mathfrak n^2\longrightarrow\mathfrak m/\mathfrak m^2\longrightarrow0.\] Its right space has dimension \(d\) by regularity. Choose \(c=n-d\) polynomials \(f_1,\ldots,f_c\in I\) whose linear classes form a basis of the kernel. Let \(J=(f_1,\ldots,f_c)\) and \(T=k[x]/J\), with maximal ideal \(\mathfrak m'=\mathfrak n/J\). The same sequence for \(J\) shows that the cotangent space of \(T_{\mathfrak m'}\) has dimension \(d\). Its surjection onto \(S_{\mathfrak m}\) implies \(\dim T_{\mathfrak m'}\ge d\), while the embedding-dimension bound gives the reverse inequality. Thus \(T_{\mathfrak m'}\) is regular of dimension \(d\).

This local ring is a domain by AG-CA-14, Theorem 1.1. Its surjection to \(S_{\mathfrak m}\) has zero kernel. Indeed, if a nonzero proper kernel existed, every prime containing it would be nonzero. Any chain of length \(d\) in the quotient would lift to such a chain in the domain, and prepending \((0)\) would make a chain of length \(d+1\), contradicting its dimension \(d\). A local Noetherian ring of finite dimension has a chain attaining that integer dimension, so the quotient could not have dimension \(d\).

The Jacobian rows of the \(f_j\) at \(x=0\) are their independent linear coefficient vectors. Some \(c\times c\) minor \(\Delta\) is therefore nonzero there; for \(c=0\) take \(\Delta=1\). After reordering variables, \(T_\Delta\) is standard smooth: adjoin an inverse variable \(u\) and the equation \(u\Delta-1\). The selected Jacobian block on the original \(c\) pivot variables and \(u\) has determinant \(\Delta^2\), which is a unit. This is the principal-localization calculation of \hyperref[AG-RG-S05-algebra-lemma-standard-smooth]{F.85}, so \(T_\Delta\) is smooth over \(k\).

It remains to identify this neighbourhood with one in \(S\). The kernel of \(T\to S\) is finite because \(T\) is Noetherian, and its localization at \(\mathfrak m'\) is zero by the preceding local isomorphism. For each of its finitely many generators choose an annihilating denominator outside \(\mathfrak m'\). Their product \(a\notin\mathfrak m'\) kills the entire kernel. Surjectivity and exact localization give \(T_a\cong S_{\bar a}\). Both \(a\) and \(\Delta\) avoid the point, so \[T_{a\Delta}\cong S_{\overline{a\Delta}}.\] Principal localization preserves the standard smooth presentation by F.85(4). Therefore the right side is smooth over \(k\), proving (4) with \(g=\overline{a\Delta}\notin\mathfrak m\). This also covers the empty equation list. All assertions retain the stipulated algebraically closed base field. ∎

\phantomsection\label{AG-RG-S05-algebra-lemma-charpoly}

\subsubsection{Lemma F.19: charpoly}\label{AG-RG-S05-lemma-f.19-charpoly}

Reading source: \href{https://stacks.math.columbia.edu/tag/00DX}{Stacks, Tag 00DX}.

For a square matrix \(M\) of size \(n\) over any commutative ring \(R\), let \(P(T)=\det(T1_n-M)\). Then \(P(M)=0\) as an endomorphism of \(R^n\).

\textbf{Proof.} Define determinants by the signed permutation formula. Grouping its terms according to the entry in a fixed row gives the cofactor expansion. If that row is replaced by another row, the same expansion is zero: permutations pair with the transposition of the two equal rows, and the paired terms have opposite signs. The column calculation is identical. Thus, for every square matrix \(C\), the transpose of its signed cofactor matrix, denoted \(\operatorname{adj}(C)\), satisfies \[
C\operatorname{adj}(C)=\operatorname{adj}(C)C=(\det C)1_n.
\] This argument uses addition and multiplication only, so it applies over \(R[T]\) as well.

Assume \(n>0\), write \[
\operatorname{adj}(T1_n-M)=\sum_{j=0}^{n-1}C_jT^j,
\qquad P(T)=\sum_{j=0}^n p_jT^j,\quad p_n=1.
\] The degree bound follows because each cofactor is a determinant of size \(n-1\). Comparing coefficients in the first adjugate identity gives \[
-MC_0=p_0 1_n,\qquad
C_{j-1}-MC_j=p_j1_n\ (1\leq j<n),\qquad C_{n-1}=1_n.
\] Multiply the equation indexed by \(j\) on the left by \(M^j\), the constant equation by \(1_n\), and the last equation by \(M^n\). Adding them cancels every \(M^{j+1}C_j\) against its next occurrence. The result is \(\sum_{j=0}^n p_jM^j=0\). No substitution into a polynomial with noncommuting matrix coefficients, and no division, is involved. If \(n=0\), the empty determinant gives \(P=1\); its value is the identity of the zero module, which is the zero endomorphism. ∎

\phantomsection\label{AG-RG-S05-algebra-lemma-codimension}

\subsubsection{Lemma F.20: codimension}\label{AG-RG-S05-lemma-f.20-codimension}

\emph{Adapted from the Stacks project, algebra.tex, lines 29578--29596.}
\nopagebreak[4]

For a surjection of finite-type \(k\)-algebras \(S'\to S\) and corresponding primes, \(\dim_{x'}\operatorname{Spec}S'-\dim_x\operatorname{Spec}S=\operatorname{height}(\mathfrak p')-\operatorname{height}(\mathfrak p)\).

\textbf{Proof.} AG-CA-09, Theorem 6.1, gives point dimension as local dimension plus the transcendence degree of the residue field. The two residue fields here are equal, since the map is a surjection and the primes correspond. The local dimensions are the heights of those primes, by their chain definition. Subtracting the two identities cancels exactly the equal transcendence degrees, giving the asserted formula. ∎

\phantomsection\label{AG-RG-S05-algebra-lemma-colimit-eventually-flat}

\subsubsection{Lemma F.21: colimit eventually flat}\label{AG-RG-S05-lemma-f.21-colimit-eventually-flat}

\emph{Adapted from the Stacks project, algebra.tex, lines 34105--34145.}
\nopagebreak[4]

Let \(R_\lambda\to S_\lambda\) be a compatible directed system of local maps of Noetherian local rings, with local transition maps, and put \(R=\operatorname*{colim}R_\lambda\), \(S=\operatorname*{colim}S_\lambda\). Assume that for every \(\mu\ge\lambda\), \(S_\mu\) is the stipulated localization of \(S_\lambda\otimes_{R_\lambda}R_\mu\). Let \(M_\lambda\) be finite \(S_\lambda\)-modules with the compatible identifications \(M_\mu=M_\lambda\otimes_{S_\lambda}S_\mu\), and put \(M=\operatorname*{colim}M_\lambda\). The constructions \hyperref[AG-RG-S05-algebra-lemma-limit-essentially-finite-presentation]{F.59} and \hyperref[AG-RG-S05-algebra-lemma-limit-module-essentially-finite-presentation]{F.60} supply exactly such models. If \(M\) is flat over \(R\), then some model \(M_\mu\) is flat over \(R_\mu\).

\textbf{Proof.} Fix one stage \(\lambda\) and set \(T=\operatorname{Tor}_1^{R_\lambda}(M_\lambda,R_\lambda/\mathfrak m_\lambda)\). The free-ideal kernel formula identifies it with \(\ker(\mathfrak m_\lambda\otimes_{R_\lambda}M_\lambda\to M_\lambda)\). It is finite over \(S_\lambda\): the ideal and module are finite, and kernels of maps of finite modules over that Noetherian algebra are finite. Choose finitely many generators. In the limit their images in \((\mathfrak m_\lambda R)\otimes_RM\) vanish, because multiplication into \(M\) is injective by its \(R\)-flatness. Tensor and filtered colimits commute by their finite-expression universal properties. Consequently each such vanishing holds at a later stage; a common successor \(\mu\) kills all generators in \((\mathfrak m_\lambda R_\mu)\otimes_{R_\mu}M_\mu\).

Thus the transition map from \(T\) to \(\operatorname{Tor}_1^{R_\mu}(M_\mu,R_\mu/\mathfrak m_\lambda R_\mu)\) is zero. The module \(M_\lambda/\mathfrak m_\lambda M_\lambda\) is flat over the field \(R_\lambda/\mathfrak m_\lambda\). The preceding transition criterion applies to the local square \(R_\lambda\to S_\lambda\), \(R_\mu\to S_\mu\), whose localized tensor description was proved in the model construction, with ideal \(I=\mathfrak m_\lambda\). It makes \(M_\mu\) flat over \(R_\mu\), as required. ∎

\phantomsection\label{AG-RG-S05-algebra-lemma-colimit-formally-etale}

\subsubsection{Lemma F.22: colimit formally etale}\label{AG-RG-S05-lemma-f.22-colimit-formally-etale}

\emph{Adapted from the Stacks project, algebra.tex, lines 42732--42747.}
\nopagebreak[4]

A filtered colimit of formally etale \(R\)-algebras is formally etale over \(R\).

\textbf{Proof.} Given a map from the colimit to \(A/J\) with \(J^2=0\), restrict to each stage. Formal etaleness gives a unique lift from that stage to \(A\). Uniqueness makes these maps compatible with every transition arrow, so their colimit supplies a lift. Two such lifts coincide on every stage and hence on the colimit. No finite-presentation assumption is used. ∎

\phantomsection\label{AG-RG-S05-algebra-lemma-colimit-syntomic}

\subsubsection{Lemma F.23: colimit syntomic}\label{AG-RG-S05-lemma-f.23-colimit-syntomic}

\emph{Adapted from the Stacks project, algebra.tex, lines 46090--46123.}
\nopagebreak[4]

Every field extension \(K/k\) is a filtered union of global complete-intersection \(k\)-subalgebras. If it is separable, it is a filtered union of smooth subalgebras.

\textbf{Proof.} Given a finite subset \(E\subset K\), form its finitely generated subfield \(L/k\). Choose a transcendence basis \(T_1,\ldots,T_d\) and a finite algebraic generator tower \(y_1,\ldots,y_r\). At the \(i\)th step the monic minimal polynomial has coefficients in the preceding field, which is spanned over \(k(T)\) by the standard monomials in the preceding generators. Represent those coefficients by polynomials in the preceding \(y_j\) with rational-function coefficients in \(T\). Clearing a common denominator \(h(T)\) puts all of these monic polynomials \(P_i\) in \(k[T,1/h,Y]\). Their successive quotients are free over the preceding coefficient ring on \(1,Y_i,\ldots,Y_i^{\deg P_i-1}\), by monic division. In particular each monic polynomial is a nonzerodivisor in the relevant polynomial ring: its highest nonzero coefficient in a putative product cannot disappear. The list is therefore regular.

The resulting algebra \(C\) is finite free over \(k[T,1/h]\). After passing to \(k(T)\) it is exactly the field \(L\), by the minimal-polynomial tower. Since the coefficient ring is a domain and \(C\) is free, \(C\to C\otimes k(T)=L\) is injective. Each member of \(E\) is a linear combination of the finite standard monomial basis with coefficients in \(k(T)\); enlarging \(h\) clears those finitely many denominators and puts \(E\) in \(C\). This finite injective extension has dimension \(d\) by AG-CA-05, Theorem 5.1, and AG-CA-09, Theorem 4.2. Adjoining the localization variable with \(hZ-1\) gives a polynomial presentation with \(d+r+1\) variables and \(r+1\) equations and dimension \(d\), so it is a global complete intersection. The collection of these subalgebras is directed: to dominate two, apply the construction to their finite generator lists. Every field element is included, proving that the union is \(K\).

In the separable case, Definition \hyperref[AG-RG-S05-algebra-definition-separable-field-extension]{F.5} makes the finitely generated subfield containing \(E\) separably generated. The forward construction in \hyperref[AG-RG-S05-algebra-lemma-localization-smooth-separable]{F.64} gives a smooth finite-type domain inside this subfield with the same fraction field. Express the finitely many elements of \(E\) as fractions and invert the product of their nonzero denominators. This principal localization remains smooth and contains \(E\). Applying the same construction to two finite generator lists proves directedness, and every field element is included. This proves the second assertion using the definition directly. ∎

\phantomsection\label{AG-RG-S05-algebra-lemma-colimits-NL}

\subsubsection{Lemma F.24: colimits NL}\label{AG-RG-S05-lemma-f.24-colimits-nl}

\emph{Adapted from the Stacks project, algebra.tex, lines 36983--36999.}
\nopagebreak[4]

The canonical two-term conormal complex commutes with filtered colimits of ring maps.

\textbf{Proof.} Use the canonical polynomial presentation with one variable for each element of the target ring. Every polynomial is a finite sum in finitely many variables and coefficients, so it occurs at one stage of the filtered system. Its being in the kernel is a finite equality in the target, hence holds at a later stage. An element in the square of the kernel is a finite sum of products of such kernel elements, and that equality also holds at a common later stage. These facts prove that the polynomial rings, their kernels, their squared kernels, and their conormal quotients have the indicated filtered colimits. Polynomial differentials are the direct sums on the same variables, and their derivatives are finite expressions, so the differentials commute with these identifications. Here is the exactness calculation for the module colimits being used. For a compatible system of short exact sequences, an element from the left whose image is zero in the middle colimit has zero image at a later stage; injectivity at that stage makes its left representative zero as well. A middle representative mapping to zero in the right colimit maps to zero at a later stage, where exactness lifts it from the left module. Every right representative lifts from the middle module at its own stage by surjectivity. Addition, scalar multiplication and comparison of representatives are checked at a common stage. Thus filtered module colimits preserve the entire short exact sequence even when transition maps are not injective. Applied to the kernels and quotients just described, this gives the asserted equality of two-term complexes. ∎

\phantomsection\label{AG-RG-S05-algebra-lemma-complete-local-Noetherian-domain-finite-over-regular}

\subsubsection{Lemma F.25: complete local Noetherian domain finite over regular}\label{AG-RG-S05-lemma-f.25-complete-local-noetherian-domain-finite-over-regular}

\emph{Adapted from the Stacks project, algebra.tex, lines 46654--46712.}
\nopagebreak[4]

A complete Noetherian local domain \(A\) is finite over a regular complete local subring \(A_0\) with the same residue field. One may take \(A_0\) to be a power-series ring over a field or a Cohen ring.

\textbf{Proof.} The coefficient-ring existence proof is AG-CA-19S, Theorem 5.1. Its characteristic-\(p\) lifting argument uses the full \(p\)-basis construction proved in Lemma \hyperref[AG-RG-S05-more-algebra-lemma-p-basis]{F.119}. In particular, that proof divides \(T^p-a^p=(T-a)^p\) inside the extension field and uses the coefficient \(-da\) of a putative proper factor, with \(1\le d<p\), to prove the required degree-\(p\) assertion. This supplies the field step for arbitrary residue fields. Since \(A\) is a domain its coefficient image is a field in equal characteristic, and is an embedded Cohen DVR in mixed characteristic. Let \(D=\dim A\). In the field case choose \(D\) parameters \(x_i\), by AG-CA-11, Theorem 2.1, and set \(A_0=k[[X_1,\ldots,X_D]]\). In the mixed case \(p\) is a nonzerodivisor; AG-CA-11, Theorem 3.2, gives \(\dim A/pA=D-1\). Choose parameters \(x_1,\ldots,x_{D-1}\) of that quotient and set \(A_0=C[[X_1,\ldots,X_{D-1}]]\). Here \(C\) is its coefficient Cohen ring. In either case the image \(I\) of the maximal ideal of \(A_0\) is primary to the maximal ideal of \(A\). A power of the latter lies in \(I\) by AG-CA-03, Proposition 2.2, making the two filtrations cofinal. Thus \(A\) is \(I\)-adically complete and separated. Evaluation of a power series at the \(x_i\) is well defined, by the same convergence and multiplication proof as AG-CA-19S, Theorem 6.1.

Choose lifts \(a_1,\ldots,a_s\) of a basis of the finite-dimensional residue quotient \(A/I\) over \(k\). Every \(a\in A\) can be expressed modulo \(I\) as a coefficient combination of them. Express the error in \(I\) as a sum of its chosen generators times elements of \(A\), and repeat the same residue decomposition on those elements. After \(n\) steps this constructs coefficients \(c_{jn}\in A_0\), compatible modulo its \(n\)-th maximal-ideal power, such that \(a-\sum_jc_{jn}a_j\in I^n\). Completeness of \(A_0\) gives limits \(c_j\) and separatedness gives \(a=\sum_jc_ja_j\). Hence \(A\) is finite as an \(A_0\)-module. The quotient \(A/I\) is finite-dimensional because its maximal ideal is nilpotent and its finitely many successive maximal-ideal layers are finite-dimensional, by AG-CA-03, Theorem 4.2.

The power-series ring \(A_0\) is the maximal-ideal completion of the regular polynomial local ring over its coefficient field or DVR, as is checked directly on the finite quotients. Regularity of the polynomial local ring is AG-CA-14, Proposition 3.3, with the DVR base regular of dimension one; its dimension is the number of displayed parameters, including \(p\) in mixed characteristic, by AG-CA-11, Theorem 5.1. Completion preserves dimension and regularity by AG-CA-19, Theorem 4.1. Thus \(A_0\) is a regular local domain of dimension \(D\). If the evaluation kernel contained \(0\ne f\), the quotient \(A_0/(f)\) would have dimension \(D-1\) by AG-CA-11, Theorem 3.2. The injective finite integral map \(A_0/\ker\to A\) preserves dimension by AG-CA-05, Theorem 5.1, but \(\dim(A_0/\ker)\le D-1\) would contradict \(\dim A=D\). The evaluation is therefore injective. Its residue map is the coefficient-field identification, completing the proof. For \(D=0\) the domain is a field and the assertion takes \(A_0=A\). ∎

\phantomsection\label{AG-RG-S05-algebra-lemma-completion-at-quasi-finite-prime}

\subsubsection{Lemma F.26: completion at quasi finite prime}\label{AG-RG-S05-lemma-f.26-completion-at-quasi-finite-prime}

\emph{Adapted from the Stacks project, algebra.tex, lines 32412--32464.}
\nopagebreak[4]

Let \(R\) be Noetherian, \(R\to S\) finite type, and \(q\) a quasi-finite prime over \(p\). Then \(\widehat{R_p}\otimes_RS\cong\widehat{S_q}\times D\), with its first projection the natural completed local map.

\textbf{Proof.} The earlier full affine conductor theorem \href{https://kokunoyumeto.github.io/open-math-courses/courses/AG-RG/AG-RG-S04.html}{Zariski Main, Theorem C4.3} gives an integral subalgebra \(C\subset S\) and \(g\in C\setminus q\) with \(C_g=S_g\). Choose finitely many numerators from \(C\) that, together with \(g^{-1}\), generate \(S_g\) over \(R\). Let \(S'\subset C\) be generated by them and \(g\). Its generators are integral, hence it is a finite \(R\)-module algebra, and \(S'_g=S_g\). Put \(q'=q\cap S'\). Localizing this equality again gives \(S'_{q'}=S_q\), so their maximal-adic completions agree.

The finite-extension completion argument of \hyperref[AG-RG-S05-algebra-lemma-completion-finite-extension]{F.27} gives a product \[\widehat{R_p}\otimes_RS'=A\times D',\qquad A=\widehat{S'_{q'}}=\widehat{S_q}.\] Let \(e\) be the central idempotent of this selected factor. Its image splits \(T=\widehat{R_p}\otimes_RS\) as \(eT\times(1-e)T\). The element \(g\) maps to a unit in \(A\), because it is outside \(q'\), so it is a unit in \(eT\) through the map \(A\to eT\). Consequently \(eT=(eT)_g\). But localizing the tensor map at \(g\) is an isomorphism, since \(S'_g=S_g\). Its selected \(e\)-factor therefore identifies \((eT)_g\) with \(A_g=A\). This proves \(eT=A\) and the desired product decomposition, including its adic factor.

To check its first projection is the natural map, it agrees on \(\widehat{R_p}\) and on \(S'\) by the finite-extension completion construction. For any \(s\in S\), equality after localizing at \(g\) supplies an exponent \(a\) with \(g^as\in S'\): write its fraction using \(S'_g\), then multiply by one further power to kill the localization error in \(S\). Since \(g\) is a unit in \(\widehat{S_q}\), agreement on \(g^as\) and on \(g\) forces agreement on \(s\). Pure tensors generate \(T\), so its projection is precisely the natural completed-local map. This is an explicit adic comparison, not an identification of finite scheme completion with ring completion. ∎

\phantomsection\label{AG-RG-S05-algebra-lemma-completion-finite-extension}

\subsubsection{Lemma F.27: completion finite extension}\label{AG-RG-S05-lemma-f.27-completion-finite-extension}

\emph{Adapted from the Stacks project, algebra.tex, lines 24144--24184.}
\nopagebreak[4]

For a finite map \(R\to S\) of Noetherian rings and a prime \(\mathfrak p\) of \(R\), with primes \(\mathfrak q_i\) above it, one has \(\widehat{R_{\mathfrak p}}\otimes_RS=\prod_i\widehat{S_{\mathfrak q_i}}\).

\textbf{Proof.} Localize the base at \(\mathfrak p\). All maximal ideals of the finite algebra lie above this maximal ideal, by AG-CA-05, Theorem 3.3, and there are finitely many by Theorem 3.4. For each \(n\), the finite algebra \(S/\mathfrak p^nS\) is a finite module over the Artinian ring \(R/\mathfrak p^n\), hence Artinian by AG-CA-03, Theorem 3.3 and its finite maximal-ideal-layer proof. The product theorem 4.2 of that lesson gives \(S/\mathfrak p^nS=\prod_iS_{\mathfrak q_i}/\mathfrak p^nS_{\mathfrak q_i}\). On each factor \(\mathfrak pS_{\mathfrak q_i}\) has radical \(\mathfrak q_iS_{\mathfrak q_i}\). A power of the latter ideal lies in the former, by AG-CA-03, Proposition 2.2; their adic filtrations are cofinal, so their inverse limits are the same. Taking the finite product of the inverse limits now gives the asserted product of local completions. Finally the completion-tensor equality for the finite \(R_{\mathfrak p}\)-module \(S_{\mathfrak p}\) is AG-CA-19, Theorem 3.1. ∎

\phantomsection\label{AG-RG-S05-algebra-lemma-compose-finite-type}

\subsubsection{Lemma F.28: compose finite type}\label{AG-RG-S05-lemma-f.28-compose-finite-type}

\emph{Adapted from the Stacks project, algebra.tex, lines 547--571.}
\nopagebreak[4]

Finite-type and finitely presented ring maps compose. If \(R\to S'\to S\) and \(S\) is finite type over \(R\), it is finite type over \(S'\). If additionally \(S/R\) is finitely presented and \(S'/R\) finite type, then \(S/S'\) is finitely presented.

\textbf{Proof.} Generators of \(S'/R\), followed by generators of \(S/S'\), generate the composite. For presentations \(S'=R[Y]/(a)\) and \(S=S'[X]/(b)\), lift the finitely many coefficients of \(b\) to \(R[Y]\) to obtain the finite composite presentation \(R[Y,X]/(a,\widetilde b)\). The same generators of \(S/R\) generate it over \(S'\). Finally write \(S=R[X]/(f_1,\ldots,f_m)\) and \(S'=R[Y]/I\), with finitely many \(Y_i\) mapping to polynomials \(h_i(X)\) in \(S\). Then \(S=S'[X]/(f_j,Y_i-h_i(X))\). To check this equality, substitute \(Y_i=h_i(X)\); every relation in \(I\) vanishes in \(S\) and so belongs to the ideal of the \(f_j\). This proves both inverse maps directly without assuming \(I\) finitely generated. ∎

\phantomsection\label{AG-RG-S05-algebra-lemma-compose-smooth}

\subsubsection{Lemma F.29: compose smooth}\label{AG-RG-S05-lemma-f.29-compose-smooth}

\emph{Adapted from the Stacks project, algebra.tex, lines 39039--39053.}
\nopagebreak[4]

A composition of smooth ring maps is smooth.

\textbf{Proof.} Let \(R\to S\to T\) be smooth maps. Lemma \hyperref[AG-RG-S05-algebra-lemma-compose-finite-type]{F.28} gives finite presentation of \(T/R\). The Jacobi--Zariski sequence \hyperref[AG-RG-S05-algebra-lemma-exact-sequence-NL]{F.42} places \(H_1(\operatorname{NL}_{T/R})\) between \(H_1(\operatorname{NL}_{S/R}\otimes_ST)\) and \(H_1(\operatorname{NL}_{T/S})\). The latter is zero by smoothness. For the former, the last assertion of F.42 identifies it with \(H_1(\operatorname{NL}_{S/R})\otimes_ST=0\): its two Tor hypotheses hold because \(\Omega_{S/R}\) is projective. Hence \(H_1(\operatorname{NL}_{T/R})=0\). The same sequence gives

\[0\longrightarrow\Omega_{S/R}\otimes_ST\longrightarrow\Omega_{T/R}\longrightarrow\Omega_{T/S}\longrightarrow0.\]

The last module is projective, so this sequence splits. The first is finite projective because a finite direct-summand presentation for \(\Omega_{S/R}\) remains such a presentation after tensoring with \(T\); the last is finite projective by smoothness of \(T/S\). Their direct sum is therefore finite projective. Definition \hyperref[AG-RG-S05-algebra-definition-smooth]{F.6} proves smoothness of \(T/R\). ∎

\phantomsection\label{AG-RG-S05-algebra-lemma-compose-standard-smooth}

\subsubsection{Lemma F.30: compose standard smooth}\label{AG-RG-S05-lemma-f.30-compose-standard-smooth}

\emph{Adapted from the Stacks project, algebra.tex, lines 38815--38878.}
\nopagebreak[4]

A composite of standard smooth ring maps is standard smooth.

\textbf{Proof.} Write \(S=R[X_1,\ldots,X_n]/(f_1,\ldots,f_c)\) and \(T=S[Y_1,\ldots,Y_m]/(g_1,\ldots,g_d)\), with their selected Jacobian minors invertible. Lift the coefficients of the \(g_i\) to \(R[X]\). The resulting \(R\)-presentation of \(T\) has equations \(f_i,g_j\) and pivot variables the \(c\) selected \(X\) variables followed by the \(d\) selected \(Y\) variables. Its Jacobian matrix is block lower triangular, since \(\partial f_i/\partial Y_j=0\). Its determinant is the product of the two given unit determinants. This is the required standard smooth presentation. Any principal localization is encoded by one more variable \(z\) and equation \(hz-1\), whose additional diagonal Jacobian entry \(h\) is a unit. Thus the assertion covers the localized presentations as well. ∎

\phantomsection\label{AG-RG-S05-algebra-lemma-composition-essentially-of-finite-type}

\subsubsection{Lemma F.31: composition essentially of finite type}\label{AG-RG-S05-lemma-f.31-composition-essentially-of-finite-type}

\emph{Adapted from the Stacks project, algebra.tex, lines 13129--13138.}
\nopagebreak[4]

Essentially finite type, and essentially finite presentation, are preserved by composition.

\textbf{Proof.} Write the first algebra as \(U^{-1}R[Y]/(a)\), and the second before its final localization as a polynomial algebra over it modulo finitely many chosen relations in the finite-presentation case. Only finitely many denominators from \(U\) occur in those relations; absorb them in one element \(u\in U\). This finite algebra then descends to \((R[Y]/(a))_u\), and its further localization gives the original composite. The composite is therefore a localization of a finite-type \(R\)-algebra, or of a finitely presented one in the second case. The remaining elements of \(U\) and the second localization set are simply included in the final localization set. The finite-type argument only needs the finite generator list and uses the same construction with the possibly infinite relation ideal. ∎

\phantomsection\label{AG-RG-S05-algebra-lemma-composition-syntomic}

\subsubsection{Lemma F.32: composition syntomic}\label{AG-RG-S05-lemma-f.32-composition-syntomic}

\emph{Adapted from the Stacks project, algebra.tex, lines 38404--38453.}
\nopagebreak[4]

Let \(R \to S\), \(S \to S'\) be ring maps.

\begin{enumerate}
\def\labelenumi{\arabic{enumi}.}
\item
  If \(R\) is Noetherian and \(R \to S\) and \(S \to S'\) are syntomic, then \(R \to S'\) is syntomic.
\item
  If \(R \to S\) and \(S \to S'\) are relative global complete intersections, then \(R \to S'\) is a relative global complete intersection.
\end{enumerate}

\textbf{Proof.} Proof of (2). Say \(R \to S\) and \(S \to S'\) are relative global complete intersections and we have presentations \(S =  R[x_1, \ldots, x_n]/(f_1, \ldots, f_c)\) and \(S' = S[y_1, \ldots, y_m]/(h_1, \ldots, h_d)\) as in Definition \hyperref[AG-RG-S05-algebra-definition-relative-global-complete-intersection]{F.4}. Then

\[
S' \cong
R[x_1, \ldots, x_n, y_1, \ldots, y_m]/(f_1, \ldots, f_c, h'_1, \ldots, h'_d)
\]

for some lifts \(h_j' \in R[x_1, \ldots, x_n, y_1, \ldots, y_m]\) of the \(h_j\). Hence it suffices to bound the dimensions of the fibre rings. Thus we may assume \(R = k\) is a field. In this case we see that we have a ring, namely \(S\), which is of finite type over \(k\) and equidimensional of dimension \(n - c\), and a finite type ring map \(S \to S'\) all of whose nonempty fibre rings are equidimensional of dimension \(m - d\). Then, by Lemma \href{https://kokunoyumeto.github.io/open-math-courses/courses/AG-CA/dimension-theory-of-noetherian-local-rings.html}{AG-CA-11, Theorem 5.1} for example applied to localizations at maximal ideals of \(S'\), we see that \(\dim(S') \leq n - c + m - d\) as desired.

We will reduce part (1) to part (2). Assume \(R\) Noetherian and \(R \to S\) and \(S \to S'\) are syntomic. Their finite presentation makes \(S\) and \(S'\) Noetherian by the earlier Hilbert basis theorem, so both applications of \hyperref[AG-RG-S05-algebra-lemma-syntomic]{F.87} have their stated base hypothesis. Let \(\mathfrak q' \subset S'\) be a prime ideal lying over \(\mathfrak q \subset S\). By Lemma \hyperref[AG-RG-S05-algebra-lemma-syntomic]{F.87} there exists a \(g' \in S'\), \(g' \not \in \mathfrak q'\) such that \(S \to S'_{g'}\) is a relative global complete intersection. Similarly, we find \(g \in S\), \(g \not \in \mathfrak q\) such that \(R \to S_g\) is a relative global complete intersection. By Lemma \hyperref[AG-RG-S05-algebra-lemma-base-change-relative-global-complete-intersection]{F.14} the ring map \(S_g \to S'_{gg'}\) is a relative global complete intersection. By part (2) we see that \(R \to S'_{gg'}\) is a relative global complete intersection and \(gg' \not \in \mathfrak q'\). Since \(\mathfrak q'\) was arbitrary combining Lemmas \hyperref[AG-RG-S05-algebra-lemma-syntomic]{F.87} and \hyperref[AG-RG-S05-algebra-lemma-local-syntomic]{F.62} we see that \(R \to S'\) is syntomic (this also uses that the spectrum of \(S'\) is quasi-compact, see Lemma \href{https://kokunoyumeto.github.io/open-math-courses/courses/AG-CA/spectra-of-rings.html}{AG-CA-01, Proposition 2.3}). ∎

\phantomsection\label{AG-RG-S05-algebra-lemma-conormal-module-localize}

\subsubsection{Lemma F.33: conormal module localize}\label{AG-RG-S05-lemma-f.33-conormal-module-localize}

\emph{Adapted from the Stacks project, algebra.tex, lines 37203--37225.}
\nopagebreak[4]

For finite polynomial presentations \(R[X_1,\ldots,X_n]\twoheadrightarrow S\) and \(R[Y_1,\ldots,Y_m]\twoheadrightarrow S_g\) with kernels \(I,J\), there is an isomorphism \((I/I^2)_g\oplus S_g^m\cong J/J^2\oplus S_g^n\).

\textbf{Proof.} Put \(C=S_g\), \(P=R[X_1,\ldots,X_n]\), and \(Q=R[Y_1,\ldots,Y_m]\). Choose \(f\in P\) representing \(g\). Then \(P_f\twoheadrightarrow C\) has kernel \(I_f\), by exact localization. Choose \(u_j\in P_f\) lifting the images of the \(Y_j\) in \(C\), and choose \(v_i\in Q\) lifting the images of the \(X_i\). The latter lifts exist because \(Q\twoheadrightarrow C\) is surjective.

Map \(W=P[Y_1,\ldots,Y_m]\) to \(C\) using both specified presentations. Its kernel is \(L=(J,X_i-v_i(Y))\): quotienting by this ideal eliminates the \(X_i\) and gives exactly \(Q/J=C\). Translation makes \(W=Q[s_1,\ldots,s_n]\), with \(s_i=X_i-v_i(Y)\), and \(L=(J,s)\). Expanding a polynomial by its \(s\) degree proves \[
L/L^2\cong J/J^2\oplus C^n.
\] Its constant part is read modulo \(J^2\), its linear coefficients modulo \(J\), and all terms of degree at least two belong to \(L^2\). These rules define inverse maps, and justify the formula over every base ring.

The image of \(f\) in \(C\) is the unit \(g\), so localization leaves this conormal module unchanged. Ideal products and quotients localize by their fraction formulas, giving \(L_f/L_f^2=L/L^2\). On the other hand \(W_f=P_f[Y]\), and its kernel for the same map to \(C\) is \((I_f,Y_j-u_j)\). Eliminating the \(Y_j\) gives \(P_f/I_f=C\), which proves this kernel equality. Translating these variables and applying the identical constant/linear calculation gives \[
L_f/L_f^2\cong (I_f/I_f^2)\oplus C^m
\cong (I/I^2)_g\oplus C^m.
\] Combine the two displayed isomorphisms. No free summand is cancelled and no flatness of \(C\) is assumed. The formulas also hold for empty variable lists and \(C=0\). ∎

\phantomsection\label{AG-RG-S05-algebra-lemma-contained-in-radical}

\subsubsection{Lemma F.34: contained in radical}\label{AG-RG-S05-lemma-f.34-contained-in-radical}

Reading source: \href{https://stacks.math.columbia.edu/tag/0AME}{Stacks, Tag 0AME}.

For an ideal \(I\) of a commutative ring \(R\), the inclusion \(I\subset\operatorname{Jac}(R)\) is equivalent to invertibility of every \(1+i\), \(i\in I\). Under either condition, every element invertible modulo \(I\) is already invertible in \(R\). Here the Jacobson radical is the intersection of all maximal ideals.

\textbf{Proof.} We first justify the maximal-ideal test for a unit. A nonunit \(a\) generates a proper ideal. The partially ordered set of proper ideals containing \((a)\) has an upper bound for every chain, its union: that union cannot contain \(1\), since then a member would contain \(1\). Zorn's lemma supplies a maximal proper ideal containing \(a\). Conversely no unit belongs to a proper ideal. Therefore an element is a unit exactly when it belongs to no maximal ideal.

Suppose \(I\) is contained in every maximal ideal. If a maximal ideal contained \(1+i\), it would contain its difference with \(i\), namely \(1\). This is impossible, so the unit test makes every \(1+i\) invertible. Conversely, if a maximal ideal \(\mathfrak m\) fails to contain \(I\), maximality gives \(I+\mathfrak m=R\). Write \(1=i+b\) with \(i\in I\) and \(b\in\mathfrak m\). The element \(b=1-i\) has the stipulated form, yet is not a unit. This contradiction proves the inclusion.

If \(a\bmod I\) is a unit, lift one inverse to \(b\in R\). Then \(ab=1+i\) for some \(i\in I\). The element \(b(ab)^{-1}\) is an actual inverse of \(a\), proving the last assertion. In the zero ring the empty intersection is the whole ring and its unique element is its identity and a unit; the conclusions hold there too. ∎

\phantomsection\label{AG-RG-S05-algebra-lemma-cover-upstairs}

\subsubsection{Lemma F.35: cover upstairs}\label{AG-RG-S05-lemma-f.35-cover-upstairs}

\emph{Adapted from the Stacks project, algebra.tex, lines 4346--4405.}
\nopagebreak[4]

Let \(g_1,\ldots,g_r\in S\) generate the unit ideal. Finite type and finite presentation of \(S/R\) can be checked on the finitely many \(S_{g_i}\).

\textbf{Proof.} First record the module detection calculation on a finite principal cover. Let \(t_1,\ldots,t_r\) generate one in a ring \(T\), say \(\sum c_it_i=1\), and suppose a \(T\)-module \(M\) has \(M_{t_i}=0\) for each \(i\). For any \(m\in M\), choose positive integers \(e_i\) with \(t_i^{e_i}m=0\), using the equality criterion for localization. Expand \[
1=\left(\sum c_it_i\right)^N,\qquad N=1+\sum_i(e_i-1).
\] Every monomial contains some \(t_i^{e_i}\), so multiplying by \(m\) gives \(m=0\). Thus \(M=0\), with no finiteness hypothesis on \(M\). An empty cover forces \(T=0\), when this assertion also holds.

Choose \(h_i\in S\) with \(\sum h_ig_i=1\). For finite type, let \(S_0\subset S\) be the \(R\)-subalgebra generated by all \(h_i,g_i\) and the numerators of finite \(R\)-algebra generator lists for the \(S_{g_i}\). The maps \((S_0)_{g_i}\to S_{g_i}\) are injective by the localization equality criterion and surjective by these chosen generators. The \(g_i\) generate one in \(S_0\) as well. The cokernel \(S/S_0\), viewed as an \(S_0\)-module, vanishes on this principal cover, so the preceding calculation gives \(S=S_0\).

For finite presentation first use finite type to write \(S=P/J\), \(P=R[X_1,\ldots,X_n]\). Lift \(g_i,h_i\) to \(u_i,v_i\in P\). The surjection \(P[Z]\to S_{g_i}\) has kernel \(JP[Z]+(u_iZ-1)\), since its quotient is \(S[Z]/(g_iZ-1)=S_{g_i}\) by the fraction universal property. This kernel is finitely generated because \(S_{g_i}\) is finitely presented. Here is the precise presentation-independence calculation: given \(C=R[Y_1,\ldots,Y_m]/(a_1,\ldots,a_s)\) and a finite polynomial surjection \(R[W_1,\ldots,W_t]\to C\), choose polynomials \(U_j(W)\) representing the images of \(Y_j\) and \(V_l(Y)\) representing the images of \(W_l\). The ideal \[
\big(a_j(U(W)),\ W_l-V_l(U(W))\big)_{j,l}
\] is its kernel. Indeed substitution \(Y_j\mapsto U_j(W)\) defines a map from \(C\) to the displayed quotient, inverse to its map to \(C\): the first list makes substitution well-defined, the second makes the composite on the \(W_l\) the identity, and the chosen images make the other composite on the \(Y_j\) the identity.

Expand each finite kernel generator for \(P[Z]\to S_{g_i}\) as a finite combination of elements of \(J\) and \(u_iZ-1\). Collect the finitely many elements \(a_{ij}\in J\) that occur. These elements together with \(u_iZ-1\) generate exactly that kernel: they belong to it and generate each of its chosen generators. Set \(J_0=(\sum v_iu_i-1,a_{ij})\subset J\). The classes of the \(u_i\) generate one in \(P/J_0\). The equality of the preceding kernels shows that every map \((P/J_0)_{u_i}\to(P/J)_{u_i}\) is an isomorphism. Its global kernel \(J/J_0\) therefore vanishes by the module detection calculation, so \(J=J_0\) is finitely generated.

Conversely, if \(S\) is of finite type, then \(S_{g_i}\) is generated by its finite algebra generators and \(g_i^{-1}\). If \(S=P/(a_1,\ldots,a_s)\) is finitely presented and \(u_i\in P\) lifts \(g_i\), then \(S_{g_i}=P[Z]/(a_1,\ldots,a_s,u_iZ-1)\) is finitely presented. This proves both directions, including zero rings and empty generator lists. ∎

\phantomsection\label{AG-RG-S05-algebra-lemma-criterion-flatness-fibre}

\subsubsection{Lemma F.36: criterion flatness fibre}\label{AG-RG-S05-lemma-f.36-criterion-flatness-fibre}

\emph{Adapted from the Stacks project, algebra.tex, lines 34374--34496.}
\nopagebreak[4]

Let \(R\), \(S\), \(S'\) be local rings and let \(R \to S \to S'\) be local ring homomorphisms. Let \(M\) be an \(S'\)-module. Let \(\mathfrak m \subset R\) be the maximal ideal. Assume

\begin{enumerate}
\def\labelenumi{\arabic{enumi}.}
\item
  The ring maps \(R \to S\) and \(R \to S'\) are essentially of finite presentation.
\item
  The module \(M\) is of finite presentation over \(S'\).
\item
  The module \(M\) is not zero.
\item
  The module \(M/\mathfrak mM\) is a flat \(S/\mathfrak mS\)-module.
\item
  The module \(M\) is a flat \(R\)-module.
\end{enumerate}

Then \(S\) is flat over \(R\) and \(M\) is a flat \(S\)-module.

\textbf{Proof.} Use the local Noetherian models of \hyperref[AG-RG-S05-algebra-lemma-limit-essentially-finite-presentation]{F.59}: \(R=\operatorname*{colim}R_\lambda\), with local transition maps, maximal ideals \(\mathfrak p_\lambda\), and each \(R_\lambda\) essentially of finite type over \(\mathbf Z\). The finite presentations of \(S\) and \(S'\) and the finitely many coefficients giving \(S\to S'\) can be put into a common model. Explicitly, present \(S\) by a finite variable and relation list over \(R\) before its localization. Present \(S'\) by another such finite list. For every generator of \(S\) choose in \(S'\) a polynomial representative with its finitely many localization denominators. Put those denominators in the localization set, and add the finite equations identifying those representatives with the chosen \(S\)-generators. Eliminating the latter recovers the presentation of \(S'\) over \(R\). Thus we obtain compatible local models \[S_\lambda=(R_\lambda[x]/(f_\lambda))_{\mathfrak q_\lambda},\qquad S'_\lambda=(R_\lambda[x,y]/(f_\lambda,g_\lambda))_{\mathfrak q'_\lambda}\] with both relation lists included, and with the primes contracted from the actual local rings. The finite equality witnesses required for the maps hold after one common successor, by the finite-expression colimit calculation in F.59. All these rings are Noetherian.

Choose a finite presentation \((S')^s\xrightarrow{A}(S')^t\to M\to0\). The model construction \hyperref[AG-RG-S05-algebra-lemma-limit-module-essentially-finite-presentation]{F.60} puts the finitely many matrix entries into one stage. Set \(M_\lambda=\operatorname{coker}((S'_\lambda)^s\xrightarrow{A_\lambda}(S'_\lambda)^t)\). For \(\mu\ge\lambda\), the maps \[S_\lambda\otimes_{R_\lambda}R_\mu\to S_\mu,\qquad S'_\lambda\otimes_{S_\lambda}S_\mu\to S'_\mu\] are localizations, and \(M_\lambda\otimes_{S'_\lambda}S'_\mu\cong M_\mu\). The same descriptions hold with the final rings. In particular \(M_\lambda=0\) would imply \(M=0\), so every such model is nonzero. Each \(M_\lambda\) is finite over \(S'_\lambda\), as required for the Noetherian local criterion.

Apply \hyperref[AG-RG-S05-algebra-lemma-colimit-eventually-flat]{F.21} to the \(R\)-flat module \(M\). It gives a stage where \(M_\lambda\) is flat over \(R_\lambda\). That property persists at every later stage: the displayed module and ring identities express the later module as a base change followed by localization, both of which preserve flatness by their exact tensor functors.

The local transitions give \(\mathfrak m=\operatorname*{colim}\mathfrak p_\lambda\), and right exact tensor products give the quotient colimits \[S/\mathfrak mS=\operatorname*{colim}S_\lambda/\mathfrak p_\lambda S_\lambda,\qquad S'/\mathfrak mS'=\operatorname*{colim}S'_\lambda/\mathfrak p_\lambda S'_\lambda,\qquad M/\mathfrak mM=\operatorname*{colim}M_\lambda/\mathfrak p_\lambda M_\lambda.\] Reducing the same polynomial presentations and matrices gives, for \(\mu\ge\lambda\), the localization \[\frac{S'_\lambda}{\mathfrak p_\lambda S'_\lambda}\otimes_{S_\lambda/\mathfrak p_\lambda S_\lambda}\frac{S_\mu}{\mathfrak p_\mu S_\mu}\longrightarrow\frac{S'_\mu}{\mathfrak p_\mu S'_\mu}\] and the corresponding exact module base-change isomorphism. These are precisely the quotient local finite-presentation models in F.60. Apply F.21 to the assumed flat module \(M/\mathfrak mM\) over \(S/\mathfrak mS\). After taking a common successor, both model flatness assertions hold: \(M_\lambda\) is \(R_\lambda\)-flat and \(M_\lambda/\mathfrak p_\lambda M_\lambda\) is flat over \(S_\lambda/\mathfrak p_\lambda S_\lambda\).

Now \hyperref[AG-RG-S05-algebra-lemma-criterion-flatness-fibre-Noetherian]{F.37} applies to the local Noetherian maps \(R_\lambda\to S_\lambda\to S'_\lambda\) and the nonzero finite module \(M_\lambda\). It gives \(R_\lambda\)-flatness of \(S_\lambda\) and \(S_\lambda\)-flatness of \(M_\lambda\). The final ring \(S\) is a localization of \(S_\lambda\otimes_{R_\lambda}R\), and the final module \(M\) is the corresponding localization of \(M_\lambda\otimes_{S_\lambda}S\). Tensor base change and localization give the two asserted flatness conclusions. No Noetherian hypothesis has been imposed on the original rings. ∎

\phantomsection\label{AG-RG-S05-algebra-lemma-criterion-flatness-fibre-Noetherian}

\subsubsection{Lemma F.37: criterion flatness fibre Noetherian}\label{AG-RG-S05-lemma-f.37-criterion-flatness-fibre-noetherian}

\emph{Adapted from the Stacks project, algebra.tex, lines 24999--25045.}
\nopagebreak[4]

Let \(R\to S\to T\) be local maps of Noetherian local rings, and \(M\ne0\) finite over \(T\). If \(M\) is \(R\)-flat and \(M/\mathfrak m_RM\) is flat over \(S/\mathfrak m_RS\), then \(M\) is faithfully flat over \(S\) and \(S\) is flat over \(R\).

\textbf{Proof.} Put \(I=\mathfrak m_RS\). The map \(\mathfrak m_R\otimes_RM\to I\otimes_SM\) is surjective, because the images of \(\mathfrak m_R\) generate \(I\). Its composite to \(M\) is injective by \(R\)-flatness, so both maps are injective. Thus \(\operatorname{Tor}_1^S(S/I,M)=0\). The ideal form of the local flatness criterion, AG-CA-08, Corollary 4.4, applied to the proper ideal \(I\) gives \(S\)-flatness. Nakayama gives \(M/\mathfrak m_TM\ne0\), hence \(M/\mathfrak m_SM\ne0\); the faithful-flatness criterion for modules, AG-CA-08, Theorem 1.1, gives faithfulness over \(S\).

Tensor the exact sequence \(0\to\operatorname{Tor}_1^R(R/\mathfrak m_R,S)\to\mathfrak m_R\otimes_RS\to I\to0\) with this flat \(S\)-module. Its first term maps to the kernel of the already injective map \(\mathfrak m_R\otimes_RM\to I\otimes_SM\), hence is zero. Faithfulness kills the original Tor module. AG-CA-08, Theorem 4.2, then gives \(R\)-flatness of \(S\). All modules to which the local criterion is applied are finite over the stipulated Noetherian local overring. ∎

\phantomsection\label{AG-RG-S05-algebra-lemma-descent-regular}

\subsubsection{Lemma F.38: descent regular}\label{AG-RG-S05-lemma-f.38-descent-regular}

\emph{Adapted from the Stacks project, algebra.tex, lines 48353--48371.}
\nopagebreak[4]

Regularity descends along faithfully flat maps of Noetherian rings.

\textbf{Proof.} Fix a source prime and choose a target prime above it, which exists by faithful flatness. The localized map \(A\to B\) is flat local, hence faithfully flat by AG-CA-08, Theorem 2.1. The target is regular local. Take a free resolution of the source residue field \(k\) over \(A\). After tensoring with \(B\) it remains a free resolution of \(k\otimes_AB\), by flatness. Consequently \[\operatorname{Tor}_i^A(k,k)\otimes_AB=\operatorname{Tor}_i^B(k\otimes_AB,k\otimes_AB).\] AG-CA-14, Theorem 2.2, gives finite global dimension of \(B\), so these groups vanish for all sufficiently large \(i\). Faithfulness gives vanishing of the original source Tor groups. The minimal-free-resolution calculation AG-CA-13, Theorem 2.2, makes the source residue field have finite projective dimension; AG-CA-14, Theorem 2.2, then makes \(A\) regular. Doing this at every prime proves the assertion. ∎

\phantomsection\label{AG-RG-S05-algebra-lemma-dimension-fibres-bounded-open-upstairs}

\subsubsection{Lemma F.39: dimension fibres bounded open upstairs}\label{AG-RG-S05-lemma-f.39-dimension-fibres-bounded-open-upstairs}

\emph{Adapted from the Stacks project, algebra.tex, lines 32674--32690.}
\nopagebreak[4]

For a finite-type map of Noetherian rings \(R\to S\), if the fibre dimension at a point \(q\) is \(n\), some principal neighbourhood of \(q\) has fibre point-dimension at most \(n\) everywhere. This is the exact form used here.

\textbf{Proof.} Let \(p=q\cap R\). Choose a principal neighbourhood whose \(\kappa(p)\)-fibre has dimension exactly \(n\), using the point-dimension definition and the finite-type field component formula. Apply Noether normalization to that fibre. The coordinate changes in AG-CA-09, Lemma 2.1 and Corollary 3.2, use polynomials with integer coefficients in the algebra generators: at each step replace a generator by it minus a sufficiently high power of the last generator, so one selected relation has a constant nonzero leading coefficient; dividing that coefficient makes the last generator integral. Inducting yields \(n\) such parameters and finiteness over their polynomial algebra. Thus these parameters lift to elements \(t_i\in S\), and \(R[T_1,\ldots,T_n]\to S\), \(T_i\mapsto t_i\), has a finite fibre over \(p\). At the chosen point it is quasi-finite: its further fibre over the contracted polynomial prime is a finite algebra over that prime's residue field.

The earlier affine conductor proof \href{https://kokunoyumeto.github.io/open-math-courses/courses/AG-RG/AG-RG-S04.html}{Zariski Main, Theorem C5.1} proves the open quasi-finite locus from its full Theorem C4.3. Shrink once more around \(q\) so this polynomial map is quasi-finite everywhere. In a fibre over any base prime \(p'\), its algebra \(U\) is quasi-finite over \(\kappa(p')[T_1,\ldots,T_n]\). At a minimal prime \(\eta\) of \(U\), the residue field is finite over that of its contracted polynomial prime, by the isolated-fibre criterion proved above. Its transcendence degree over \(\kappa(p')\) is therefore at most \(n\). AG-CA-09, Theorem 4.2, gives \(\dim(U/\eta)\) equal to that transcendence degree. Taking the maximum over the finitely many minimal primes gives \(\dim U\le n\). Every fibre point-dimension is at most that fibre dimension. This proves the claimed neighbourhood bound. Its earlier conductor input includes the polynomial-normality proof in \href{https://kokunoyumeto.github.io/open-math-courses/courses/AG-RG/AG-RG-S04.html}{Zariski Main over an arbitrary base, Corollary B1.2}, so no old external proof tag is being substituted. ∎

\phantomsection\label{AG-RG-S05-algebra-lemma-essentially-of-finite-type-into-artinian-local}

\subsubsection{Lemma F.40: essentially of finite type into artinian local}\label{AG-RG-S05-lemma-f.40-essentially-of-finite-type-into-artinian-local}

Reading source: \href{https://stacks.math.columbia.edu/tag/07DT}{Stacks, Tag 07DT}.

Let \(R\to S\) be any ring map, with \(S\) Artinian local, maximal ideal \(\mathfrak m\), and residue field \(K=S/\mathfrak m\). Each of the following properties of \(S/R\) is equivalent to the same property of \(K/R\): finite, of finite type, and essentially of finite type.

\textbf{Proof.} The Artinian length theorem \href{https://kokunoyumeto.github.io/open-math-courses/courses/AG-CA/noetherian-and-artinian-rings.html}{AG-CA-03, Theorem 4.2} gives a finite composition series of the \(S\)-module \(S\). Every simple factor is \(K\): a nonzero element of a simple module generates it, and its annihilator is the unique maximal ideal. If \(K\) is finite over \(R\), each factor is therefore a finite \(R\)-module. Finite generation passes through an extension directly: lift finitely many generators of the quotient and adjoin finitely many generators of the submodule. Induction on the composition series makes \(S\) finite over \(R\). The converse follows by taking the images of module generators in \(K\). This proves the finite assertion for an arbitrary source ring.

For finite type, a quotient keeps the images of any finite list of algebra generators. Conversely, lift finite \(R\)-algebra generators of \(K\) to \(s_1,\ldots,s_a\in S\). The polynomial map \(P=R[X_1,\ldots,X_a]\to S\), \(X_i\mapsto s_i\), is surjective on residues. Thus \(K\) is cyclic as a \(P\)-module. The finite assertion just proved makes \(S\) a finite \(P\)-module. Choose its module generators \(u_1,\ldots,u_b\). Every element is a sum of polynomials in the \(s_i\) times the \(u_j\), so \(s_i,u_j\) together generate \(S\) as an \(R\)-algebra.

Finally, a quotient of a localization of a finite-type algebra is the localization of its quotient, by the fraction formula. Hence essential finite type passes from \(S\) to \(K\). For the converse write \(K=T^{-1}C\), with \(C\) a finite-type \(R\)-algebra, and let \(C_0\) be the image of \(C\) in \(K\). The domain \(C_0\) has fraction field \(K\): every element of the displayed localization is a fraction of its elements, and the field \(K\) contains the inverse of every nonzero element of \(C_0\).

Lift a finite list of \(R\)-algebra generators of \(C_0\) to \(S\), and let \(A\subset S\) be the \(R\)-subalgebra they generate. Put \(\mathfrak p=A\cap\mathfrak m\). Then \(A/\mathfrak p=C_0\), so \(\kappa(\mathfrak p)=K\). Every element of \(A\setminus\mathfrak p\) maps to a unit in \(S\); thus the inclusion extends to \(A_{\mathfrak p}\to S\). Applying the finite assertion to this map makes \(S\) a finite \(A_{\mathfrak p}\)-module. Choose generators \(v_1,\ldots,v_c\) and set \(B=A[v_1,\ldots,v_c]\subset S\). This is a finite-type \(R\)-algebra. Localizing \(B\) at \(A\setminus\mathfrak p\) maps injectively into \(S\), because those denominators are units there, and surjectively because the \(v_j\) generate over \(A_{\mathfrak p}\). Therefore \(S\) is that localization of \(B\). This proves the last equivalence without a finiteness assumption on \(R\). ∎

\phantomsection\label{AG-RG-S05-algebra-lemma-etale-standard-smooth}

\subsubsection{Lemma F.41: etale standard smooth}\label{AG-RG-S05-lemma-f.41-etale-standard-smooth}

\emph{Adapted from the Stacks project, algebra.tex, lines 40737--40760.}
\nopagebreak[4]

Any étale ring map is standard smooth. More precisely, if \(R \to S\) is étale, then there exists a presentation \(S = R[x_1, \ldots, x_n]/(f_1, \ldots, f_n)\) such that the image of \(\det(\partial f_j/\partial x_i)\) is invertible in \(S\).

\textbf{Proof.} If \(S=0\), use \(S=R[X]/(1)\). Its one-by-one Jacobian is zero, a unit in the zero ring, giving the required presentation. Suppose \(S\ne0\).

Choose a finite polynomial presentation \(S=P/I\), \(P=R[x_1,\ldots,x_n]\). The complete formal/conormal comparison in AG-CA-17, Theorem 3.1 and Section 4, makes its conormal differential \[I/I^2\longrightarrow S^{\oplus n},\qquad\bar f\longmapsto\sum_i\frac{\partial f}{\partial x_i}dx_i\] an isomorphism: étaleness means smoothness and zero differentials, and the split conormal sequence therefore has zero cokernel and zero kernel. In particular \(I/I^2\) is finite free.

Apply \hyperref[AG-RG-S05-algebra-lemma-huber]{F.55}. It supplies a possibly different finite presentation \(S=R[y_1,\ldots,y_m]/J\), with \(J=(g_1,\ldots,g_c)\) and these equation classes a conormal basis. Apply the same étale conormal criterion to this new presentation as well. It identifies \(S^c\cong J/J^2\) with \(S^m\). Reducing this isomorphism at any maximal ideal of the nonzero ring \(S\) gives an isomorphism of finite-dimensional vector spaces, hence \(c=m\).

In these bases the conormal isomorphism is the square Jacobian matrix \(H\) of the actual new equations. A matrix inverse exists over \(S\), so multiplicativity of determinants gives \(\det(H)\det(H^{-1})=1\). Thus its determinant is a unit. Relabeling the \(m\) variables and equations gives exactly the asserted standard smooth presentation, with equal numbers of variables and equations. ∎

\phantomsection\label{AG-RG-S05-algebra-lemma-exact-sequence-NL}

\subsubsection{Lemma F.42: exact sequence NL}\label{AG-RG-S05-lemma-f.42-exact-sequence-nl}

\emph{Adapted from the Stacks project, algebra.tex, lines 36838--36913.}
\nopagebreak[4]

Let \(A \to B \to C\) be ring maps. Choose a presentation \(\alpha : A[x_s, s \in S] \to B\) with kernel \(I\). Choose a presentation \(\beta : B[y_t, t \in T] \to C\) with kernel \(J\). Let \(\gamma : A[x_s, y_t] \to C\) be the induced presentation of \(C\) with kernel \(K\). Then we get a canonical commutative diagram

\[
\begin{gathered}
0 \xrightarrow{} \Omega_{A[x_s]/A} \otimes C \\
\Omega_{A[x_s]/A} \otimes C \xrightarrow{} \Omega_{A[x_s, y_t]/A} \otimes C \\
\Omega_{A[x_s, y_t]/A} \otimes C \xrightarrow{} \Omega_{B[y_t]/B} \otimes C \\
\Omega_{B[y_t]/B} \otimes C \xrightarrow{} 0 \\
I/I^2 \otimes C \xrightarrow{} K/K^2 \\
I/I^2 \otimes C \xrightarrow{} \Omega_{A[x_s]/A} \otimes C \\
K/K^2 \xrightarrow{} J/J^2 \\
K/K^2 \xrightarrow{} \Omega_{A[x_s, y_t]/A} \otimes C \\
J/J^2 \xrightarrow{} 0 \\
J/J^2 \xrightarrow{} \Omega_{B[y_t]/B} \otimes C
\end{gathered},
\]

with exact rows. We get the following exact sequence of homology groups

\[
H_1(\NL_{B/A} \otimes_B C) \to
H_1(L_{C/A}) \to
H_1(L_{C/B}) \to
C \otimes_B \Omega_{B/A} \to
\Omega_{C/A} \to
\Omega_{C/B} \to 0
\]

of \(C\)-modules extending the sequence of Lemma \href{https://kokunoyumeto.github.io/open-math-courses/courses/AG-CA/kahler-differentials.html}{AG-CA-16, Theorem 3.1}. If \(\text{Tor}_1^B(\Omega_{B/A}, C) = 0\) and \(\text{Tor}_2^B(\Omega_{B/A}, C) = 0\), then \(H_1(\NL_{B/A} \otimes_B C) = H_1(L_{B/A}) \otimes_B C\).

\textbf{Proof.} On differentials, the first map sends \(dx_s\) to \(dx_s\), the second sends every \(dx_s\) to zero and \(dy_t\) to \(dy_t\). On conormal modules the first map sends \((f\bmod I^2)\otimes c\) to \(cf\bmod K^2\), using any lift of \(c\) to \(A[x_s,y_t]\); the ambiguity changes the answer by \(K^2\). The second sends \(g\bmod K^2\) to its image in \(J/J^2\). Derivatives of these representatives verify every commutative square. The connecting homology map is obtained by lifting a cycle from \(J/J^2\) to \(K/K^2\), taking its derivative (which lies in the image of the left differential module), and reducing modulo the image of \(I/I^2\otimes C\). Changing the lift changes this by a boundary. The same lifting calculation proves exactness at each homology term. The exactness of the top row follows as the \(\text{d}x_s\), \(\text{d}y_t\) form a basis for the middle module. The map \(\gamma\) factors

\[
A[x_s, y_t] \to B[y_t] \to C
\]

with surjective first arrow and second arrow equal to \(\beta\). Thus we see that \(K \to J\) is surjective. Moreover, the kernel of the first displayed arrow is \(IA[x_s, y_t]\). Hence \(I/I^2 \otimes C\) surjects onto the kernel of \(K/K^2 \to J/J^2\). Finally, we can use Lemma \hyperref[AG-RG-S05-algebra-lemma-NL-homotopy]{F.11} to identify the terms as homology groups of the naive cotangent complexes.

The final assertion is a statement in homological algebra. Recall that \(\NL_{B/A} = (N^{-1} \to N^0)\) is a two term complex of \(B\)-modules with \(N^0\) free and cohomology modules \(H^0 = \Omega_{B/A}\) and \(H^{-1} = H_1(L_{B/A})\). Write \(M \subset N^0\) for the image of the differential. If \(\text{Tor}_1^B(H^0, C) = 0\), then we have an exact sequence

\[
0 \to M \otimes_B C \to N^0 \otimes_B C \to H^0 \otimes_B C \to 0
\]

Since \(N^0\) is free, we also see that \(\text{Tor}_2^B(H^0, C) =
\text{Tor}_1^B(M, C)\). Hence if \(\text{Tor}_2^B(H^0, C) = 0\) then we also have an exact sequence

\[
0 \to H^{-1} \otimes_B C \to N^{-1} \otimes_B C \to M \otimes_B C \to 0
\]

Putting everything together we see that if \(\text{Tor}_1^B(H^0, C) = 0\) and \(\text{Tor}_2^B(H^0, C) = 0\), then \(H^{-1} \otimes_B C\) is the kernel of \(N^{-1} \otimes_B C \to N^0 \otimes_B C\) as desired. ∎

\phantomsection\label{AG-RG-S05-algebra-lemma-extension}

\subsubsection{Lemma F.43: extension}\label{AG-RG-S05-lemma-f.43-extension}

Reading source: \href{https://stacks.math.columbia.edu/tag/0519}{Stacks, Tag 0519}.

Let \(R\) be any ring and let \(0\to M_1\to M_2\to M_3\to0\) be exact. If the two end modules are finite, so is the middle module; a quotient of a finite module is finite. If both end modules are finitely presented, so is the middle module. A quotient of a finitely presented module by a finite submodule is finitely presented. Finally, if \(M_2\) is finite and \(M_3\) is finitely presented, then \(M_1\) is finite.

\textbf{Proof.} Identify \(M_1\) with its image in \(M_2\). To generate the middle module, take generators of \(M_1\) and lifts of generators of \(M_3\). Subtracting a combination of the latter from an arbitrary element leaves an element of \(M_1\). Images of any generating list generate a quotient. These observations prove the two assertions about finite modules.

We first establish the presentation calculation needed in the other assertions. Suppose \(q : G\twoheadrightarrow N\) is a finite free presentation with finitely generated kernel \(L\), and \(p : F\twoheadrightarrow N\) is any surjection from a finite free module. Choose maps \(u : F\to G\) and \(v : G\to F\) by lifting their finite bases, so that \(qu=p\) and \(pv=q\). For every \(z\in\ker p\), \[
z=v(uz)+(1-vu)z,
\qquad uz\in L.
\] Both \(v(L)\) and \((1-vu)F\) lie in \(\ker p\), since \(pv=q\) and \(pvu=qu=p\). Thus that kernel is generated by the images under \(v\) of a finite generating list of \(L\), together with the finitely many columns of \(1-vu\). In particular the kernel of every finite free surjection onto a finitely presented module is finite. The choices of lifts here use only the defining property of a free basis.

If \(M_2\) is finite, choose a finite free surjection \(F\to M_2\). When \(M_3\) is finitely presented, the kernel of the composite \(F\to M_3\) is finite by this calculation. Its image in \(M_2\) is exactly \(M_1\): an element mapping to zero in \(M_3\) has any chosen lift in that kernel. Hence \(M_1\) is finite.

Next write a finitely presented \(M_2\) as \(F/K\), with \(F\) finite free and \(K\) finite. If \(M_1\) is generated by \(a_1,\ldots,a_h\), choose lifts \(\widetilde a_i\in F\). The kernel of \(F\to M_3\) is \[
K+R\widetilde a_1+\cdots+R\widetilde a_h.
\] Indeed an element in that kernel maps to a combination of the \(a_i\) in \(M_2\); subtracting the corresponding lifted combination puts it in \(K\). This is a finite kernel and gives the asserted finite presentation of the quotient.

Finally choose finite free presentations \(p_i : F_i\twoheadrightarrow M_i\), with finite kernels \(K_i\), for \(i=1,3\). Lift the basis of \(F_3\) to obtain \(\ell : F_3\to M_2\), and use \[
F_1\oplus F_3\longrightarrow M_2,
\qquad (a,b)\longmapsto p_1(a)+\ell(b).
\] It is onto by the first generator argument. For each generator \(b_j\) of \(K_3\), the element \(\ell(b_j)\) lies in \(M_1\); choose \(a_j\in F_1\) with \(p_1(a_j)=-\ell(b_j)\). Its kernel is generated by \((K_1,0)\) and the finitely many pairs \((a_j,b_j)\). To see this, the second coordinate of any kernel element lies in \(K_3\). Subtract a combination of these pairs to make that coordinate zero; the first coordinate then lies in \(K_1\). We have constructed a finite generating list of relations, so \(M_2\) is finitely presented. Empty generating lists and zero modules satisfy the same formulas. No Noetherian assumption is used. ∎

\phantomsection\label{AG-RG-S05-algebra-lemma-finite-presentation-independent}

\subsubsection{Lemma F.44: finite presentation independent}\label{AG-RG-S05-lemma-f.44-finite-presentation-independent}

Reading source: \href{https://stacks.math.columbia.edu/tag/00R2}{Stacks, Tag 00R2}.

If an \(R\)-algebra \(S\) is finitely presented, the kernel of every surjection \(\alpha : R[X_1,\ldots,X_n]\twoheadrightarrow S\) is a finitely generated ideal.

\textbf{Proof.} Fix one presentation \[
S=R[Y_1,\ldots,Y_m]/(f_1,\ldots,f_h).
\] Choose \(g_i(Y)\) representing \(\alpha(X_i)\), and, using the surjectivity of \(\alpha\), choose \(h_j(X)\) representing the class of \(Y_j\). In the polynomial ring \(P=R[X_1,\ldots,X_n]\), form the finite ideal \[
J=\bigl(f_1(h(X)),\ldots,f_h(h(X)),\;
X_i-g_i(h(X))\ (1\leq i\leq n)\bigr).
\] Each generator maps to zero under \(\alpha\), by the choices of representatives, so \(J\subset\ker\alpha\).

For the reverse inclusion, take \(F(X)\in\ker\alpha\). The polynomial \(F(g(Y))\) represents zero in the fixed presentation; hence there are polynomials \(c_l(Y)\) with \[
F(g(Y))=\sum_{l=1}^h c_l(Y)f_l(Y).
\] Substitute \(Y=h(X)\). This puts \(F(g(h(X)))\) in \(J\). Also \[
F(X)-F(g(h(X)))\in
\bigl(X_i-g_i(h(X)):1\leq i\leq n\bigr).
\] Here is the polynomial identity behind this last inclusion, valid over an arbitrary ring. Replace the variables one at a time and telescope the differences. At a single variable, every monomial difference is divisible by its coordinate difference, because \(U^a-V^a=(U-V)\sum_{b=0}^{a-1}U^{a-1-b}V^b\) for \(a>0\); the difference is zero for \(a=0\). Adding the monomials proves the inclusion without division by any coefficient. Thus \(F\in J\), proving \(\ker\alpha=J\). The argument includes either empty variable list, an empty relation list and the zero algebra. ∎

\phantomsection\label{AG-RG-S05-algebra-lemma-flat-fibre-smooth}

\subsubsection{Lemma F.45: flat fibre smooth}\label{AG-RG-S05-lemma-f.45-flat-fibre-smooth}

\emph{Adapted from the Stacks project, algebra.tex, lines 39115--39149.}
\nopagebreak[4]

Let \(R \to S\) be a ring map. Let \(\mathfrak q \subset S\) be a prime lying over the prime \(\mathfrak p\) of \(R\). Assume

\begin{enumerate}
\def\labelenumi{\arabic{enumi}.}
\item
  there exists a \(g \in S\), \(g \not\in \mathfrak q\) such that \(R \to S_g\) is of finite presentation,
\item
  the local ring homomorphism \(R_{\mathfrak p} \to S_{\mathfrak q}\) is flat,
\item
  the fibre \(S \otimes_R \kappa(\mathfrak p)\) is smooth over \(\kappa(\mathfrak p)\) at the prime corresponding to \(\mathfrak q\).
\end{enumerate}

Then \(R \to S\) is smooth at \(\mathfrak q\).

\textbf{Proof.} We prove the stated arbitrary-ring assertion. First replace \(S\) by the given \(S_s\), with \(s\notin\mathfrak q\), and keep the corresponding prime denoted by \(\mathfrak q\). Choose a finite presentation \[
P=R[X_1,\ldots,X_n]\longrightarrow S=P/I,\qquad
I=(a_1,\ldots,a_m).
\] Let \(\mathfrak r\subset P\) be the inverse image of \(\mathfrak q\), put \(A=R_{\mathfrak p}\), \(Q=P_{\mathfrak r}\), and \(B=S_{\mathfrak q}\). The exact sequence \(0\to I_{\mathfrak r}\to Q\to B\to0\) is a sequence of \(A\)-modules, and its quotient \(B\) is \(A\)-flat by hypothesis. We will use only finite generation of \(I\), not a chain condition on any of these rings.

Here is the tensor-injection calculation needed for that sequence. More generally, suppose \(0\to K\to M\to F\to0\) is exact over a ring \(A\) and \(F\) is flat; \(M\) need not be flat or free. For an arbitrary \(A\)-module \(N\), choose a free presentation \(F_1\xrightarrow{d}F_0\to N\to0\), and put \(T=\ker d\) and \(L=\operatorname{im}d\). The free presentations, direct-sum exactness and right-exact tensor calculations are written in \href{https://kokunoyumeto.github.io/open-math-courses/courses/AG-RG/AG-RG-S01.html\#1-the-algebra-of-localization}{AG-RG-S01, Lemma 1.0}. If \(z\in K\otimes_A N\) maps to zero in \(M\otimes_A N\), lift it to \(x\in K\otimes_A F_0\). Right exactness gives \(y\in M\otimes_A F_1\) whose image in \(M\otimes_A F_0\) is the image of \(x\). The image of \(y\) in \(F\otimes_A F_1\) is killed by \(1\otimes d\). Tensoring \(0\to T\to F_1\to L\to0\) with the flat \(F\), and using the injection \(L\to F_0\), shows that this kernel is the image of \(F\otimes_A T\). Lift such a preimage to \(M\otimes_A T\), using the surjection \(M\to F\) and right exactness, and subtract its image from \(y\). The resulting \(y'\) still maps to the image of \(x\), and its image in \(F\otimes_A F_1\) is zero. Since \(F_1\) is free, coordinatewise exactness lifts \(y'\) from \(K\otimes_A F_1\). Its image in \(K\otimes_A F_0\) is \(x\): this follows from the injection \(K\otimes_A F_0\to M\otimes_A F_0\), again coordinatewise because \(F_0\) is free. Thus \(x\) comes from \(K\otimes_A F_1\), so \(z=0\). We have proved the tensor injection with no Tor assertion.

Set \(k=\kappa(\mathfrak p)=A/\mathfrak pA\). Apply that calculation to \(N=k\). Right exactness and the quotient-tensor formula identify \[
J:=\ker(Q/\mathfrak pQ\to B/\mathfrak pB)
\simeq I_{\mathfrak r}/\mathfrak pI_{\mathfrak r}.
\] Write \(\overline Q=Q/\mathfrak pQ\) and \(\overline B=B/\mathfrak pB\). Then \[
\overline Q=k[X_1,\ldots,X_n]_{\overline{\mathfrak r}},
\qquad
\overline B=\overline Q/J.
\] Indeed localization commutes with quotient and coefficient change; every polynomial over \(k\) lifts coefficientwise to \(A\), and being outside \(\overline{\mathfrak r}\) is exactly being outside the prime before reduction. The ring \(\overline Q\) is local, \(J\) is contained in its maximal ideal, and \(\overline B\) is the local ring of the fibre at the specified point. In particular \(J\) is finite, generated by the images of the \(a_j\).

The smooth-fibre hypothesis makes \(\overline B\) formally smooth over \(k\). Here the equivalence of the definitions is used at its proved scope: Definition F.6 and the presentation comparison in Lemma F.11 give, on a smooth fibre neighbourhood, an injective conormal map with projective differential cokernel. Projectivity splits that quotient, and the full conormal criterion \href{https://kokunoyumeto.github.io/open-math-courses/courses/AG-CA/formally-smooth-unramified-and-etale-ring-maps.html}{AG-CA-17, Theorem 3.1} gives formal smoothness. Its Theorem 1.2 proves that arbitrary further localizations preserve formal smoothness: a lift on the unlocalized algebra sends the inverted elements to units, since units lift across a square-zero ideal, and hence extends to the localization. Thus this step does not assume that a local ring is itself finitely presented.

We obtain the needed conormal splitting in this same localized polynomial presentation, without changing presentations. Formal smoothness lifts the identity of \(\overline B\) to a \(k\)-algebra section \[
\sigma:\overline B\longrightarrow\overline Q/J^2
\] of its square-zero quotient. The map \[
D(q)=q\bmod J^2-\sigma(q\bmod J)
\quad:\quad\overline Q\longrightarrow J/J^2
\] is a \(k\)-derivation: writing \(q=\sigma(\overline q)+D(q)\), and discarding products of two elements of \(J/J^2\), gives \(D(qq')=\overline qD(q')+\overline q'D(q)\). For \(j\in J\), it has \(D(j)=j\bmod J^2\). The universal differential construction therefore supplies a left inverse to \[
J/J^2\longrightarrow
\Omega_{\overline Q/k}\otimes_{\overline Q}\overline B
=\overline B^{\,n},\qquad
[j]\longmapsto\left(\frac{\partial j}{\partial X_\nu}\right)_\nu.
\] The free differential basis and its localization here are precisely \href{https://kokunoyumeto.github.io/open-math-courses/courses/AG-CA/kahler-differentials.html}{AG-CA-16, Theorems 5.1 and 2.2}; localization extends a derivation by \(D(q/t)=t^{-1}D(q)-qt^{-2}D(t)\). Thus \(J/J^2\) is a finite direct summand of \(\overline B^{\,n}\).

We spell out the finite local-module argument, including the required Nakayama uses. Over any local ring \((D,\mathfrak n)\), a finite module \(M\) with \(M=HM\) for \(H\subset\mathfrak n\) is zero: choose generators and write \(m_i=\sum_j h_{ij}m_j\) with \(h_{ij}\in H\). The cofactor identity of Lemma 3.B applied to \(I-(h_{ij})\) shows that its determinant annihilates every generator. That determinant is congruent to one modulo \(\mathfrak n\), hence is a unit. An element outside the unique maximal ideal is a unit, because a nonunit generates a proper ideal, which lies in a maximal ideal by the proper-ideal Zorn argument. Therefore \(M=0\). It follows also that lifts of a basis of \(M/\mathfrak nM\) generate a finite \(M\), by applying this assertion to their quotient.

Choose from the finite list \(a_1,\ldots,a_m\) elements \(f_1,\ldots,f_c\) whose conormal classes give a basis of \((J/J^2)\otimes_{\overline B}\kappa(\mathfrak q)\); a finite spanning list over a field contains a basis by successively adjoining independent members. Nakayama over \(\overline B\) makes these classes generate \(J/J^2\). They are in fact a free basis. To see this, the surjection \(\overline B^c\to J/J^2\) splits: lift the projection \(\overline B^n\to J/J^2\) through it by choosing lifts of the \(n\) coordinate images, and restrict that lift to the direct summand \(J/J^2\). Its kernel is then a finite direct summand of \(\overline B^c\), generated by the projections of its coordinate vectors. Modulo the maximal ideal the surjection is an isomorphism by our basis choice, so this kernel has zero residue quotient. The just-proved Nakayama argument kills it.

Let \(H\) be the \(c\) by \(n\) equation-by-variable matrix \((\partial f_i/\partial X_\nu)\), with coefficients reduced to \(\overline B\). The split conormal injection has matrix \(H^{\mathsf T}\) and a left inverse \(V\), so \(VH^{\mathsf T}=I_c\). We write the minor argument here. Expanding \(\det(VH^{\mathsf T})\) by multilinearity in its columns chooses \(c\) columns of \(V\) with their coefficients from \(H^{\mathsf T}\). Repeated choices have determinant zero by the signed-permutation cancellation. Grouping the remaining choices by their distinct index set \(U\), in increasing order, gives \[
1=\sum_{|U|=c}\det(V^U)\det\bigl((H^{\mathsf T})_U\bigr).
\] The signed-permutation formula also gives \(\det(W^{\mathsf T})=\det W\), by replacing a permutation with its inverse. If \(c>n\) every choice repeats, contradicting \(1\ne0\) in the local ring \(\overline B\). Thus \(c\le n\), and at least one \(c\) by \(c\) minor of \(H\) is a unit of \(\overline B\): if all belonged to its maximal ideal the displayed sum could not be one. Choose its variable indices \(U\). The corresponding polynomial \[
\Delta=\det\left(\frac{\partial f_i}{\partial X_\nu}\right)_
{1\le i\le c,\ \nu\in U}\in P
\] is outside \(\mathfrak r\), since its image in \(\kappa(\mathfrak q)\) is nonzero. For \(c=0\) this argument uses the empty determinant \(\Delta=1\) and the empty generator list.

The chosen \(f_i\) generate the whole fibre ideal \(J\) in \(\overline Q\), not merely its conormal quotient. Indeed \(J=(\overline f_1,\ldots,\overline f_c)+J^2\). Its finite quotient by the displayed ideal is equal to \(J\) times that quotient. Since \(J\) is contained in the maximal ideal of the local ring \(\overline Q\), Nakayama makes the quotient zero. Now put \(M=I_{\mathfrak r}/(f_1,\ldots,f_c)Q\), a finite \(Q\)-module. The tensor-injection identification above gives \[
M/\mathfrak pM
=J/(\overline f_1,\ldots,\overline f_c)=0.
\] The ideal \(\mathfrak pQ\) lies in the maximal ideal of \(Q\), so another application of finite Nakayama gives \(M=0\). This is the step where flatness of the local quotient \(B\) is essential.

For each of the finitely many original generators \(a_j\) of \(I\), this equality expresses \(a_j/1\) as a combination of the \(f_i/1\) in \(P_{\mathfrak r}\). Clear the finitely many coefficient denominators and also any multiplier outside \(\mathfrak r\) witnessing each equality in that localization. Their finite product \(v\notin\mathfrak r\) gives \(I_v=(f_1,\ldots,f_c)P_v\). This keeps the killed-cross-difference condition and does not cancel a possible zero divisor. Consequently a neighbourhood of \(\mathfrak q\) is \[
S_{v\Delta}
=R[X_1,\ldots,X_n]_{v\Delta}/(f_1,\ldots,f_c).
\] It is finitely presented over \(R\): add one variable \(Y\) with the relation \((v\Delta)Y-1\), together with the \(c\) equations. These further principal localizations of the earlier \(S_s\) give a principal neighbourhood in the original \(\operatorname{Spec}S\): write their elements with numerators in the original \(S\) and denominators powers of \(s\), and invert \(s\) and those finitely many numerators.

Finally we prove formal smoothness of this neighbourhood by the square-zero calculation. Given an \(R\)-algebra \(D\), an ideal \(K\) with \(K^2=0\), and a map from the displayed algebra to \(D/K\), lift the variable images to \(b_\nu\in D\). The errors \(r_i=f_i(b)\) lie in \(K\). The selected square Jacobian block \(H_U(b)\) has unit determinant, because its determinant is \(\Delta(b)\) and its image modulo \(K\) is a unit. The cofactor calculation in Lemma 3.B gives its inverse. Change only the selected coordinates by \[
(b_\nu')_{\nu\in U}
=(b_\nu)_{\nu\in U}-H_U(b)^{-1}(r_i)_{1\le i\le c},
\] leaving the others fixed. All corrections lie in \(K\). Expanding each polynomial monomial shows \(f_i(b')=f_i(b)+\sum_\nu(\partial f_i/\partial X_\nu)(b)(b_\nu'-b_\nu)\), because every term with two corrections lies in \(K^2\). Thus all \(f_i(b')\) are zero. The image of \(v\Delta(b')\) modulo \(K\) is unchanged and is a unit; it is a unit in \(D\) as well. Explicitly, if a lift of its inverse gives a product \(1+j\) with \(j\in K\), multiply that lift by \(1-j\). Evaluation therefore extends through the quotient and localization to an \(R\)-algebra lift. For \(c=0\) there are no errors to correct; choosing arbitrary variable lifts and using the same unit argument suffices. We have proved formal smoothness for every square-zero test. Together with the finite presentation it is smooth, by the written conormal criterion and its agreement with Definition F.6. The neighbourhood contains \(\mathfrak q\), proving the conclusion. No Noetherian hypothesis has been used. ∎

\phantomsection\label{AG-RG-S05-algebra-lemma-flat-over-regular}

\subsubsection{Lemma F.46: flat over regular}\label{AG-RG-S05-lemma-f.46-flat-over-regular}

Reading source: \href{https://stacks.math.columbia.edu/tag/07DY}{Stacks, Tag 07DY}.

Let \(R\to S\) be a homomorphism of Noetherian local rings. Suppose \(R\) is regular and one regular system of parameters \(x_1,\ldots,x_d\) maps to an \(S\)-regular sequence. Then \(S\) is flat over \(R\). A regular sequence here has proper generated ideal.

\textbf{Proof.} The \(x_i\) generate the maximal ideal of \(R\). Their images belong to the maximal ideal of \(S\), because the ideal they generate is proper in that local ring. Thus the given map is local, even though this was not a separate assumption.

For \(0\leq i\leq d\), put \[
R_i=R/(x_1,\ldots,x_i),\qquad
S_i=S/(x_1,\ldots,x_i)S.
\] These are nonzero Noetherian local rings and their induced maps are local. The written regular-parameter theorem \href{https://kokunoyumeto.github.io/open-math-courses/courses/AG-CA/regular-sequences-depth-and-cohen-macaulay-modules.html}{AG-CA-12, Theorem 6.1} makes the chosen list regular in \(R\). Therefore, for every \(i<d\), multiplication by \(x_{i+1}\) is injective on \(R_i\); the assumed regularity in \(S\) gives the same injection on \(S_i\).

At \(i=d\), the source \(R_d\) is the residue field of \(R\), so \(S_d\) is flat over it. Indeed a vector-space basis identifies tensoring with that space with a direct sum of copies of a module, which preserves injections.

We descend from \(i+1\) to \(i\). Since \(x_{i+1}\) is a nonzerodivisor of \(R_i\), \[
0\longrightarrow R_i\xrightarrow{x_{i+1}}R_i
\longrightarrow R_{i+1}\longrightarrow0
\] is a free resolution. Tensoring it with \(S_i\) computes \[
\operatorname{Tor}_1^{R_i}(R_{i+1},S_i)
=\ker(x_{i+1} : S_i\to S_i)=0.
\] The quotient \(S_i/x_{i+1}S_i=S_{i+1}\) is flat over \(R_{i+1}\) by the descending induction. Apply the complete ideal local criterion \href{https://kokunoyumeto.github.io/open-math-courses/courses/AG-CA/faithful-flatness-and-the-local-criterion-for-flatness.html}{AG-CA-08, Corollary 4.4 and equation (6)} to \(R_i\to S_i\), the proper ideal \((x_{i+1})\subset R_i\), and the module \(S_i\). This module is finite over the target ring \(S_i\), being generated by \(1\); it need not be finite over the source. The Tor vanishing and quotient flatness prove that \(S_i\) is \(R_i\)-flat. The induction reaches \(S_0=S\) over \(R_0=R\). If \(d=0\), the initial field argument already proves the conclusion. ∎

\phantomsection\label{AG-RG-S05-algebra-lemma-flat-over-regular-with-regular-fibre}

\subsubsection{Lemma F.47: flat over regular with regular fibre}\label{AG-RG-S05-lemma-f.47-flat-over-regular-with-regular-fibre}

Reading source: \href{https://stacks.math.columbia.edu/tag/031E}{Stacks, Tag 031E}.

Let \((R,\mathfrak m)\to(S,\mathfrak n)\) be a flat local map of Noetherian local rings. If \(R\) and the closed fibre \(F=S/\mathfrak mS\) are regular, then \(S\) is regular.

\textbf{Proof.} Put \(d=\dim R\) and \(e=\dim F\). Regularity supplies generators \(x_1,\ldots,x_d\) of \(\mathfrak m\) and generators \(\overline y_1,\ldots,\overline y_e\) of the maximal ideal \(\mathfrak n/\mathfrak mS\) of \(F\). Choose lifts \(y_j\in\mathfrak n\). Every \(z\in\mathfrak n\) reduces to a combination of the \(\overline y_j\); lifting the coefficients and subtracting that combination leaves an element of \(\mathfrak mS=(x_1,\ldots,x_d)S\). Consequently \[
\mathfrak n=(x_1,\ldots,x_d,y_1,\ldots,y_e)S.
\] Its cotangent space \(\mathfrak n/\mathfrak n^2\) is therefore spanned by at most \(d+e\) vectors over the residue field of \(S\).

The flat local dimension formula, with precisely the stated Noetherian and flatness hypotheses, is proved in \href{https://kokunoyumeto.github.io/open-math-courses/courses/AG-CA/dimension-theory-of-noetherian-local-rings.html}{AG-CA-11, Theorem 5.1}; it gives \(\dim S=d+e\). Its Proposition 4.2 proves that the dimension of a Noetherian local ring is at most the dimension of its cotangent space: lifts of a cotangent basis generate the maximal ideal by finite Nakayama, and the parameter/height theorem bounds dimension by that generator count. Combining these inequalities gives \[
d+e=\dim S\leq
\dim_{\kappa(\mathfrak n)}\mathfrak n/\mathfrak n^2
\leq d+e.
\] Equality of dimension and embedding dimension is the definition of regularity, so \(S\) is regular. The empty source or fibre parameter lists when \(d=0\) or \(e=0\) obey the same calculation. No finite-type assumption on the ring map is needed. ∎

\phantomsection\label{AG-RG-S05-algebra-lemma-flatness-descends-more-general}

\subsubsection{Lemma F.48: flatness descends more general}\label{AG-RG-S05-lemma-f.48-flatness-descends-more-general}

Reading source: \href{https://stacks.math.columbia.edu/tag/0584}{Stacks, Tag 0584}.

Let \(S\to S'\) be a flat map of \(R\)-algebras, let \(M\) be any \(S\)-module, and put \(M'=S'\otimes_S M\). If \(M\) is flat over \(R\), so is \(M'\). If \(S\to S'\) is faithfully flat, the converse holds as well.

\textbf{Proof.} Test an arbitrary injection \(f:N\hookrightarrow N_1\) of \(R\)-modules. Write \[
Z=\ker\bigl(N\otimes_R M\xrightarrow{f\otimes1}N_1\otimes_R M\bigr).
\] Both tensor products, their map and \(Z\) are \(S\)-modules, using the action on \(M\). Flatness of \(S'\) over \(S\) preserves this kernel. The balancing relations give inverse isomorphisms \[
S'\otimes_S(N\otimes_R M)\ \longleftrightarrow\ N\otimes_R M',
\qquad s'\otimes(n\otimes m)\ \longleftrightarrow\ n\otimes(s'\otimes m),
\] and the same formula for \(N_1\). These maps commute with \(f\), so \[
S'\otimes_S Z\cong
\ker\bigl(N\otimes_R M'\longrightarrow N_1\otimes_R M'\bigr).
\] If \(M\) is \(R\)-flat, \(Z=0\), and the displayed kernel vanishes for every \(f\). If \(M'\) is \(R\)-flat, the right side is zero; faithful flatness then detects \(Z=0\). In either direction the injection characterization of flatness, proved for arbitrary modules in \href{https://kokunoyumeto.github.io/open-math-courses/courses/AG-CA/tor-and-flat-modules.html}{AG-CA-07, Theorem 2.1}, gives the claimed flatness. No module in this calculation is required to be finite. ∎

\phantomsection\label{AG-RG-S05-algebra-lemma-formally-smooth-extensions-easy}

\subsubsection{Lemma F.49: formally smooth extensions easy}\label{AG-RG-S05-lemma-f.49-formally-smooth-extensions-easy}

\emph{Adapted from the Stacks project, algebra.tex, lines 45971--46012.}
\nopagebreak[4]

Let \(K/k\) be an extension of fields.

\begin{enumerate}
\def\labelenumi{\arabic{enumi}.}
\item
  If \(K\) is purely transcendental over \(k\), then \(K\) is formally smooth over \(k\).
\item
  If \(K\) is separable algebraic over \(k\), then \(K\) is formally smooth over \(k\).
\item
  If \(K\) is separable over \(k\), then \(K\) is formally smooth over \(k\).
\end{enumerate}

\textbf{Proof.} For (1) write \(K = k(x_j; j \in J)\). Suppose that \(A\) is a \(k\)-algebra, and \(I \subset A\) is an ideal of square zero. Let \(\varphi : K \to A/I\) be a \(k\)-algebra map. Let \(a_j \in A\) be an element such that \(a_j \mod I = \varphi(x_j)\). Then it is easy to see that there is a unique \(k\)-algebra map \(K \to A\) which maps \(x_j\) to \(a_j\) and which reduces to \(\varphi\) mod \(I\). Hence \(k \subset K\) is formally smooth.

In case (2), each finite subset of \(K\) generates a finite separable extension of \(k\): adjoining its finitely many algebraic generators gives only finitely many finite degrees. These subextensions form a filtered family, by adjoining two finite generating lists together, and their union is \(K\). \href{https://kokunoyumeto.github.io/open-math-courses/courses/AG-CA/formally-smooth-unramified-and-etale-ring-maps.html}{AG-CA-17, Proposition 6.3} proves that every such finite extension is étale; its separable-generator tower proof explicitly gives the unique square-zero lifts, hence formal étaleness. Hence this case follows from Lemma \hyperref[AG-RG-S05-algebra-lemma-colimit-formally-etale]{F.22} and the trivial observation that a formally étale ring map is formally smooth.

In case (3), write \(K = \operatorname*{colim} K_i\) as the filtered colimit of its finitely generated \(k\)-subextensions. By Definition \hyperref[AG-RG-S05-algebra-definition-separable-field-extension]{F.5} each \(K_i\) is separable algebraic over a purely transcendental extension of \(k\). Hence \(K_i/k\) is formally smooth by cases (1) and (2) and Lemma \href{https://kokunoyumeto.github.io/open-math-courses/courses/AG-CA/formally-smooth-unramified-and-etale-ring-maps.html}{AG-CA-17, Theorem 1.2}. Thus \(H_1(L_{K_i/k}) = 0\) by Lemma \hyperref[AG-RG-S05-algebra-lemma-characterize-formally-smooth-field-extension]{F.17}. Hence \(H_1(L_{K/k}) = 0\) by Lemma \hyperref[AG-RG-S05-algebra-lemma-colimits-NL]{F.24}. Hence \(K/k\) is formally smooth by Lemma \hyperref[AG-RG-S05-algebra-lemma-characterize-formally-smooth-field-extension]{F.17} again. ∎

\phantomsection\label{AG-RG-S05-algebra-lemma-geometrically-regular}

\subsubsection{Lemma F.50: geometrically regular}\label{AG-RG-S05-lemma-f.50-geometrically-regular}

\emph{Adapted from the Stacks project, algebra.tex, lines 48753--48798.}
\nopagebreak[4]

For a Noetherian \(k\)-algebra, regularity after every finitely generated field extension is equivalent to regularity after every finite purely inseparable field extension.

\textbf{Proof.} One implication is immediate. For the other, given a finitely generated \(K/k\), choose \(k'/k\) and \(K'/K\) by the explicit separabilization proof \hyperref[AG-RG-S05-algebra-lemma-make-separably-generated]{F.69}. The assumed \(A\otimes_k k'\) is regular and Noetherian.

Since \(K/k\) is finitely generated and \(K'/K\) is finite, \(K'/k'\) is finitely generated too. Choose field generators \(a_1,\ldots,a_n\) with \(K'=k'(a_1,\ldots,a_n)\) and put \(B=k'[a_1,\ldots,a_n]\subset K'\). This is a finite-type domain with fraction field \(K'\). Its generic local ring \(B_{(0)}=K'\) is regular, and its residue field is separably generated over \(k'\) by Lemma F.69. The proved separable-residue criterion \href{https://kokunoyumeto.github.io/open-math-courses/courses/AG-CA/smooth-algebras-over-a-field-and-the-jacobian-criterion.html}{AG-CA-18, Theorem 3.2} therefore makes \(B\) smooth at its generic point. By the definition of smoothness at a point, some \(0\ne g\in B\) has \(B_g\) smooth over \(k'\). This remains a finite-type domain, and \[
K'=(B_g\setminus\{0\})^{-1}B_g.
\] This constructs the smooth field model used here, with the stated finitely generated and separably generated hypotheses.

Base change gives the smooth \((A\otimes_k k')\)-algebra \(C=(A\otimes_k k')\otimes_{k'}B_g\): finite presentation base changes by the polynomial quotient presentation, and formal smoothness base changes by AG-CA-17, Theorem 1.2. The ring \(C\) is Noetherian because it is finite type over the Noetherian ring \(A\otimes_k k'\). Smooth maps are flat by AG-CA-18, Theorem 6.1, and their residue-field fibres are regular by AG-CA-18, Theorem 2.1; regularity ascent proved above makes this base-changed algebra regular. Localization then makes \(A\otimes_k K'\) regular. Choose a \(K\)-basis of the finite extension \(K'/K\) containing \(1\). It identifies \(A\otimes_kK'\) with a positive finite number of copies of \(A\otimes_kK\) as a module. Tensoring with it takes the same positive number of copies of every module, so preserves exact sequences and detects the zero module. Thus \(A\otimes_k K\to A\otimes_kK'\) is faithfully flat. \hyperref[AG-RG-S05-algebra-lemma-descent-regular]{F.38} now gives regularity of \(A\otimes_k K\). These rings are Noetherian because finitely generated field tensor is a localization of a finite-type polynomial extension, the Noetherian field-tensor argument already written in this lesson. ∎

\phantomsection\label{AG-RG-S05-algebra-lemma-geometrically-regular-descent}

\subsubsection{Lemma F.51: geometrically regular descent}\label{AG-RG-S05-lemma-f.51-geometrically-regular-descent}

\emph{Adapted from the Stacks project, algebra.tex, lines 48815--48834.}
\nopagebreak[4]

Geometric regularity descends through faithfully flat maps of Noetherian \(k\)-algebras.

\textbf{Proof.} Let \(A\to B\) be the given faithfully flat map and let \(k'/k\) be finite purely inseparable. Put \(A'=A\otimes_k k'\) and \(B'=B\otimes_k k'\). Both are Noetherian by \hyperref[AG-RG-S05-algebra-lemma-Noetherian-field-extension]{F.12}. Tensor associativity gives \(B'=A'\otimes_A B\) and, for every \(A'\)-module \(M\), a natural identification \[M\otimes_{A'}B'=M\otimes_A B.\] Restriction from \(A'\)-modules to \(A\)-modules preserves exact sequences and the zero module. Since tensoring with \(B\) over \(A\) is exact and detects the zero module, the displayed identification proves that \(A'\to B'\) is faithfully flat. The target \(B'\) is regular by geometric regularity of \(B\). Regularity descent \hyperref[AG-RG-S05-algebra-lemma-descent-regular]{F.38} therefore makes \(A'\) regular. Applying the finite purely inseparable criterion \hyperref[AG-RG-S05-algebra-lemma-geometrically-regular]{F.50} to every such \(k'\) proves that \(A\) is geometrically regular over \(k\). ∎

\phantomsection\label{AG-RG-S05-algebra-lemma-geometrically-regular-over-subfields}

\subsubsection{Lemma F.52: geometrically regular over subfields}\label{AG-RG-S05-lemma-f.52-geometrically-regular-over-subfields}

\emph{Adapted from the Stacks project, algebra.tex, lines 48851--48866.}
\nopagebreak[4]

If \(A\) is Noetherian and geometrically regular over \(k\), then \(A\otimes_kE\) is regular for any field extension \(E/k\) for which this tensor product is Noetherian. If \(A\otimes_kE\) is Noetherian, it is geometrically regular over \(E\). Also if \(k\) is a directed union of subfields \(k_i\) and \(A\) is geometrically regular over every \(k_i\), then it is geometrically regular over \(k\).

\textbf{Proof.} Write \(E\) as the directed union of the intermediate fields \(E_i/k\) generated by finite subsets of \(E\): the union of two finite generating lists gives a common successor. Put \(C_i=A\otimes_k E_i\) and \(C=A\otimes_kE\). The natural inclusion \(C_i\to C\) is injective. Indeed, an \(E_i\)-basis of \(E\) containing \(1\) identifies \(C=C_i\otimes_{E_i}E\) with a direct sum of copies of \(C_i\), with this inclusion as the summand corresponding to \(1\). The same basis shows that \(C\) is a free, hence flat, \(C_i\)-module. Each element of \(C\) is a finite tensor sum, so all its field coefficients lie in one \(E_i\).

Assume \(C\) is Noetherian and fix a prime \(\mathfrak q\subset C\). Its local ring \(T=C_{\mathfrak q}\) is Noetherian; choose a finite list \(x_1,\ldots,x_r\) generating its maximal ideal \(\mathfrak m\). Write \(x_j=a_j/s_j\) with \(a_j\in\mathfrak q\) and \(s_j\notin\mathfrak q\). Including the finite field coefficients of all these numerators and denominators in one \(E_i\) puts every \(a_j,s_j\) in \(C_i\). At the contracted prime \(\mathfrak q_i=\mathfrak q\cap C_i\), put \(B_i=(C_i)_{\mathfrak q_i}\). Every \(s_j\) is invertible in \(B_i\), and each \(a_j\) lies in \(\mathfrak q_i\). The map \(B_i\to T\) is local and flat: localizing the free \(C_i\)-algebra \(C\) first away from \(\mathfrak q_i\) and then at \(\mathfrak q\) preserves flatness. Therefore \(\mathfrak q_iT\) lies in \(\mathfrak m\) and contains each of its chosen generators, giving \[\mathfrak q_iT=\mathfrak m,\qquad T\otimes_{B_i}\kappa(\mathfrak q_i)=T/\mathfrak m.\] The closed fibre is the residue field of \(T\), hence a regular ring. The source \(B_i\) is Noetherian by \hyperref[AG-RG-S05-algebra-lemma-Noetherian-field-extension]{F.12} and localization, and regular by geometric regularity of \(A\) over \(k\), since \(E_i/k\) is finitely generated. Flat local regularity ascent \hyperref[AG-RG-S05-algebra-lemma-flat-over-regular-with-regular-fibre]{F.47} makes \(T\) regular. The empty generating list when \(\mathfrak m=0\) satisfies the same argument. This proves regularity at every prime of \(C\).

For any finite purely inseparable extension \(F/E\), \hyperref[AG-RG-S05-algebra-lemma-Noetherian-field-extension]{F.12} makes \[(A\otimes_kE)\otimes_EF=A\otimes_kF\] Noetherian. Apply the assertion just proved to the field extension \(F/k\); its tensor ring is regular. The criterion \hyperref[AG-RG-S05-algebra-lemma-geometrically-regular]{F.50} now proves geometric regularity of \(A\otimes_kE\) over \(E\).

Finally suppose \(k\) is the directed union of the \(k_i\) and \(A\) is geometrically regular over every \(k_i\). Let \(k'/k\) be finite purely inseparable. If \(k'=k\), take \(k_i'=k_i\) at any index. Otherwise the characteristic is \(p>0\). Choose a \(k\)-basis \(b_1=1,b_2,\ldots,b_d\) of \(k'\). There is a common exponent \(e\) such that every \(b_j^{p^e}\) belongs to \(k\). Write the finitely many multiplication identities \[b_jb_\ell=\sum_{h=1}^d c_{j\ell h}b_h,\qquad c_{j\ell h}\in k.\] Choose one \(k_i\) containing all these coefficients and all the \(b_j^{p^e}\). The subspace \[V=\bigoplus_{j=1}^d k_i b_j\subset k'\] contains \(1\) and is closed under multiplication by the displayed identities. It is a finite-dimensional domain over \(k_i\). For \(0\ne v\in V\), multiplication by \(v\) is an injective endomorphism of this finite-dimensional vector space, hence surjective; a preimage of \(1\) is its inverse. Thus \(V\) is a field, which we call \(k_i'\). For every \(v=\sum_j a_jb_j\in k_i'\) the Frobenius identity gives \[v^{p^e}=\sum_j a_j^{p^e}b_j^{p^e}\in k_i.\] This proves that \(k_i'/k_i\) is finite purely inseparable, including every element of the constructed field. The multiplication map \[k\otimes_{k_i}k_i'\longrightarrow k',\qquad a\otimes v\longmapsto av,\] is an isomorphism: the tensors \(1\otimes b_j\) form a \(k\)-basis of the source and map to the chosen \(k\)-basis of the target. Tensor associativity consequently gives \[A\otimes_kk'\cong A\otimes_{k_i}k_i'.\] The right side is regular by geometric regularity over \(k_i\). Criterion F.50 proves geometric regularity over \(k\). This also covers characteristic zero, in which a finite purely inseparable extension is trivial. ∎

\phantomsection\label{AG-RG-S05-algebra-lemma-grothendieck-regular-sequence}

\subsubsection{Lemma F.53: grothendieck regular sequence}\label{AG-RG-S05-lemma-f.53-grothendieck-regular-sequence}

\emph{Adapted from the Stacks project, algebra.tex, lines 24513--24526.}
\nopagebreak[4]

For a flat local map of Noetherian local rings \(R\to P\), a list whose reduction is regular in the closed fibre is regular in \(P\), and every successive quotient remains flat over \(R\).

\textbf{Proof.} For the first element apply AG-CA-08, Lemma 5.1, to multiplication \(P\xrightarrow{f_1}P\). The target is \(R\)-flat, both modules are finite over the Noetherian local ring \(P\), and the map on the closed fibre is injective. That exact proved lemma gives injectivity and flatness of \(P/(f_1)\). Repeat on the successive local quotient, whose closed fibre is the corresponding regular quotient. Local Nakayama ensures that every quotient remains nonzero, since the elements belong to the maximal ideal. This proves every assertion by induction. ∎

\phantomsection\label{AG-RG-S05-algebra-lemma-henselian-cat-finite-etale}

\subsubsection{Lemma F.54: henselian cat finite etale}\label{AG-RG-S05-lemma-f.54-henselian-cat-finite-etale}

Reading source: \href{https://stacks.math.columbia.edu/tag/04GK}{Stacks, Tag 04GK}.

For a henselian local ring \((R,\mathfrak m,k)\), reduction \(S\mapsto S/\mathfrak mS\) is an equivalence from finite étale \(R\)-algebras to finite products of finite separable extensions of \(k\). Ring maps are unital, and the zero algebra is the empty product.

\textbf{Proof.} The written idempotent and finite-algebra theorem \href{https://kokunoyumeto.github.io/open-math-courses/courses/AG-CA/henselian-local-rings-and-henselization.html}{AG-CA-20, Theorem 2.2} decomposes a finite \(R\)-algebra into finitely many local factors. A factor of an étale algebra is étale, since it is localization at its component idempotent. For each such factor \(S_i\), reduction is an étale \(k\)-algebra and is local: every maximal ideal of a finite algebra contracts to \(\mathfrak m\), and \(S_i\) has only one. The field classification \href{https://kokunoyumeto.github.io/open-math-courses/courses/AG-CA/formally-smooth-unramified-and-etale-ring-maps.html}{AG-CA-17, Theorem 7.2} makes this reduction a finite separable field. In particular \(\mathfrak mS_i\) is its maximal ideal. Every \(S_i\) is henselian by \href{https://kokunoyumeto.github.io/open-math-courses/courses/AG-CA/henselian-local-rings-and-henselization.html}{AG-CA-20, Proposition 3.1}.

Construct a lift of a finite separable field extension \(L/k\) as follows. Choose finitely many field generators and use their successive fields \[
k=L_0\subset L_1\subset\cdots\subset L_a=L,
\qquad L_i=L_{i-1}(\alpha_i).
\] Each minimal polynomial is monic separable. Starting with \(T_0=R\), lift its coefficients to a monic \(h_i\in T_{i-1}[X]\) and define \(T_i=T_{i-1}[X]/(h_i)\). Its bounded powers of \(X\) are a free \(T_{i-1}\)-basis, by monic division. Reduction modulo \(\mathfrak m\) is \(L_i\). Every maximal ideal contracts to \(\mathfrak m\) by integrality, so the field reduction makes \(T_i\) local with maximal ideal \(\mathfrak mT_i\). The derivative \(h_i'\) has nonzero residue in that field and is consequently a unit. The one-equation Jacobian proof \href{https://kokunoyumeto.github.io/open-math-courses/courses/AG-CA/formally-smooth-unramified-and-etale-ring-maps.html}{AG-CA-17, Theorem 4.1 and Section 6} makes \(T_i/T_{i-1}\) étale. Composition makes \(T_a/R\) finite étale, with free rank \([L:k]\) and reduction \(L\). Finite products of these lifts provide every desired residue algebra.

We next show explicitly that every local finite étale \(R\)-algebra \(S\), with residue \(L\), is isomorphic to this constructed \(T_a\). Given the residue identification, lift a map \(T_{i-1}\to S\) through the next step by choosing the root of \(h_i\) reducing to \(\alpha_i\). The simple-root Hensel statement, \href{https://kokunoyumeto.github.io/open-math-courses/courses/AG-CA/henselian-local-rings-and-henselization.html}{AG-CA-20, Theorem 1.2}, gives exactly one such root in the henselian local ring \(S\), and hence exactly one map \(T_i\to S\) at each step. This constructs \(\phi:T_a\to S\) inducing the chosen isomorphism on residues.

By the complete arbitrary-base proof in \href{https://kokunoyumeto.github.io/open-math-courses/courses/AG-RG/AG-RG-S04.html}{S04, Lemma E4.2}, a finite étale algebra is finite locally free. Its proof first obtains a finite module presentation from monic algebra equations, then applies Lemma A1.3 to a finitely presented flat module. Since \(R\) is local, \(S\) is finite free, of rank \(\dim_k(S/\mathfrak mS)=[L:k]\). Both modules in \(\phi\) therefore have this same finite rank. In chosen bases its matrix reduces to an invertible matrix over \(k\), so its determinant lies outside \(\mathfrak m\) and is a unit of \(R\). The adjugate matrix divided by this unit is its inverse. Thus \(\phi\) is an isomorphism of modules, and its bijective algebra map is an algebra isomorphism.

Now let \(S_1,S_2\) be local finite étale algebras and let \(L_1\to L_2\) be a \(k\)-field map between their residues. Replace \(S_1\) by its isomorphic tower \(T_a\). Starting with \(R\to S_2\), lift each successive generator to the unique simple root whose residue is its prescribed image in \(L_2\). The same Hensel calculation supplies a unique \(R\)-algebra map \(S_1\to S_2\) inducing the field map. Conversely, any algebra map has those generator values, so it is this unique lift. The map on residue fields is defined because \(\mathfrak mS_i\) is the maximal ideal.

For products, a map to a local target sends each source component idempotent to \(0\) or \(1\); orthogonality and their sum \(1\) force exactly one component to be chosen. The residue map chooses precisely that component, and the local result lifts its field map uniquely. Doing this independently for each target factor proves a bijection on all Hom sets. The empty-product cases follow from unital maps directly: there is one map to the zero ring and none from it to a nonzero ring. We have proved full faithfulness and essential surjectivity, hence the asserted equivalence, over every henselian local source. ∎

\phantomsection\label{AG-RG-S05-algebra-lemma-huber}

\subsubsection{Lemma F.55: huber}\label{AG-RG-S05-lemma-f.55-huber}

Reading source: \href{https://stacks.math.columbia.edu/tag/07CF}{Stacks, Tag 07CF}.

Suppose a finitely presented \(R\)-algebra \(S\) has a finite polynomial presentation \(S=P/I\), where \(P=R[X_1,\ldots,X_n]\), and \(I/I^2\) is free over \(S\). Then \(S\) has another finite polynomial presentation whose defining equations themselves form a basis of its conormal module.

\textbf{Proof.} \hyperref[AG-RG-S05-algebra-lemma-finite-presentation-independent]{F.44} makes \(I\) finitely generated. Thus \(I/I^2\) has a finite basis: when \(S\ne0\), the supports of a finite generating list in any free basis have finite union, which must contain the entire basis; when \(S=0\), choose its empty conormal basis directly. Lift this basis to \(f_1,\ldots,f_c\in I\) and put \(K=(f_1,\ldots,f_c)\). The basis property gives \(I=K+I^2\). The finite \(P\)-module \(U=I/K\) consequently satisfies \(U=IU\).

Choose generators \(u_1,\ldots,u_t\) of \(U\) and coefficients \(b_{ij}\in I\) with \(u_i=\sum_j b_{ij}u_j\). The adjugate identity proved in \hyperref[AG-RG-S05-algebra-lemma-matrix-left-inverse]{F.70} applied to \(1_t-(b_{ij})\) shows that \[g=\det(1_t-(b_{ij})),\qquad gU=0,\qquad g\in1+I.\] For \(U=0\) take \(g=1\). Hence \(gI\subset K\) and \(I_g=K_g\). The image of \(g\) in \(S\) is exactly \(1\), so \(P_g/I_g=S\). It follows, with both inverse maps given by adjoining or evaluating \(g^{-1}\), that \[S=P[T]/(f_1,\ldots,f_c,gT-1).\]

Let \(L\) be the displayed equation ideal. After localization at \(g\), the classes of \(f_i\) form the given basis of \(I_g/I_g^2\). The explicit inverse-variable conormal calculation \hyperref[AG-RG-S05-algebra-lemma-principal-localization-NL]{F.72}, applied to \(P_g/I_g\) and this presentation, identifies \[L/L^2\cong S^c\oplus S.\] The first basis vectors are \([f_i]\) and the last is \([gT-1]\): the last equation has derivative \(g\), a unit, in its new variable, and the remaining classes are read by evaluation at \(T=g^{-1}\). These are exactly the claimed equation classes. No Noetherian hypothesis or cancellation of a free summand is used. Empty bases use the same determinant convention; if \(S=0\), all the displayed modules and bases are interpreted over the zero ring. ∎

\phantomsection\label{AG-RG-S05-algebra-lemma-isolated-point-fibre}

\subsubsection{Lemma F.56: isolated point fibre}\label{AG-RG-S05-lemma-f.56-isolated-point-fibre}

Reading source: \href{https://stacks.math.columbia.edu/tag/00PK}{Stacks, Tag 00PK}.

Let \(R\to S\) be of finite type and let \(\mathfrak q\in\operatorname{Spec}S\) lie over \(\mathfrak p\). Put \(k=\kappa(\mathfrak p)\), \(T=S\otimes_R k\), and let \(x\in\operatorname{Spec}T\) be the corresponding point. The following conditions are equivalent: \(x\) is isolated; \(T_x=S_{\mathfrak q}/\mathfrak pS_{\mathfrak q}\) is finite-dimensional over \(k\); some \(g\in S\setminus\mathfrak q\) has \(x\) as the only point of its fibre neighbourhood; the fibre point-dimension at \(x\) is zero; \(x\) is closed and \(\dim T_x=0\); and \(\kappa(x)/k\) is finite and \(\dim T_x=0\).

\textbf{Proof.} All the geometry can be checked in the finite-type field algebra \(T\). Such an algebra and its principal localizations are Noetherian by \href{https://kokunoyumeto.github.io/open-math-courses/courses/AG-CA/noetherian-and-artinian-rings.html}{AG-CA-03, Theorem 2.1}. We will use the following consequence of its Theorem 4.2. A zero-dimensional finite-type \(k\)-algebra \(B\) is Artinian, with finitely many local factors and nilpotent radical \(J\). Every residue field is finite over \(k\) by \href{https://kokunoyumeto.github.io/open-math-courses/courses/AG-CA/the-nullstellensatz-and-jacobson-rings.html}{AG-CA-06, Theorem 2.1}. Each \(J^i/J^{i+1}\) is finite over the product \(B/J\) of those fields: \(B\) is Noetherian, so the ideals involved are finite. Its coordinate summands are consequently finite-dimensional over the residue fields, hence over \(k\). The finite filtration by powers of \(J\) proves that \(B\) is finite-dimensional over \(k\). The factor idempotents give a principal singleton neighbourhood of each of its points.

If \(x\) is isolated, choose \(h\in T\) with \(D(h)=\{x\}\). The algebra \(T_h\) has just one prime, so has dimension zero and is finite-dimensional by the preceding calculation. Every element outside its unique maximal ideal is a unit. Its localization at that ideal is therefore itself, and \(T_x=T_h\) is finite-dimensional.

Conversely, a finite-dimensional \(T_x\) is Artinian: chains of its ideals are chains of subspaces of a finite-dimensional \(k\)-space. Each of its prime quotients is a finite-dimensional domain, hence a field, because multiplication by a nonzero element is an injective, and thus surjective, linear map. Thus \(\dim T_x=0\), and its residue field is finite over \(k\). For a prime of a finite-type \(k\)-algebra, finiteness of the residue extension implies closedness: \(T/x\) embeds into the finite-dimensional field \(\kappa(x)\), so is a finite-dimensional domain and the same multiplication argument makes it a field. The converse follows from the cited weak Nullstellensatz.

Suppose now that \(x\) is closed and \(\dim T_x=0\). The prime corresponding to \(x\) is both maximal and minimal. There are finitely many other minimal primes, by AG-CA-03, Proposition 2.2. From each choose an element outside \(x\) and take their product \(h\); use \(h=1\) if there are none. A prime in \(D(h)\) contains none of those other minimal primes. Since every prime contains a minimal prime, it must contain \(x\) and therefore equal \(x\). Hence \(D(h)=\{x\}\) and \(x\) is isolated. This proves the equivalence of isolation, finite local fibre algebra, and the two residue-field/local-dimension conditions.

An isolated point has point-dimension zero, using its singleton open neighbourhood. Conversely, zero point-dimension supplies an open neighbourhood of dimension zero: the infimum defining that dimension is an attained nonnegative integer. Take a principal neighbourhood inside it. Its algebra has dimension zero, so the Artinian factor calculation above supplies a smaller principal singleton neighbourhood of \(x\). Thus zero point-dimension implies isolation as well.

Finally, every \(h\in T=S_{R\setminus\mathfrak p}/\mathfrak pS_{R\setminus\mathfrak p}\) is represented by \(g/s\) with \(g\in S\) and \(s\in R\setminus\mathfrak p\). The image of \(s\) is a fibre unit, so \(D(h)=D(\bar g)\), and \(h\notin x\) means \(g\notin\mathfrak q\). Conversely the fibre of \(D(g)\) is exactly \(D(\bar g)\). This proves the neighbourhood condition over the original ring, with no Noetherian assumption on \(R\). ∎

\phantomsection\label{AG-RG-S05-algebra-lemma-lci}

\subsubsection{Lemma F.57: lci}\label{AG-RG-S05-lemma-f.57-lci}

\emph{Adapted from the Stacks project, algebra.tex, lines 37341--37441.}
\nopagebreak[4]

Let \(k\) be a field. Let \(S\) be a finite type \(k\)-algebra. Let \(\mathfrak q\) be a prime of \(S\). Choose any presentation \(S = k[x_1, \ldots, x_n]/I\). Let \(\mathfrak q'\) be the prime of \(k[x_1, \ldots, x_n]\) corresponding to \(\mathfrak q\). Set \(c = \text{height}(\mathfrak q') - \text{height}(\mathfrak q)\), in other words \(\dim_{\mathfrak q}(S) = n - c\) (see Lemma \hyperref[AG-RG-S05-algebra-lemma-codimension]{F.20}). The following are equivalent

\begin{enumerate}
\def\labelenumi{\arabic{enumi}.}
\item
  There exists a \(g \in S\), \(g \not \in \mathfrak q\) such that \(S_g\) is a global complete intersection over \(k\).
\item
  The ideal \(I_{\mathfrak q'} \subset k[x_1, \ldots, x_n]_{\mathfrak q'}\) can be generated by \(c\) elements.
\item
  The conormal module \((I/I^2)_{\mathfrak q}\) can be generated by \(c\) elements over \(S_{\mathfrak q}\).
\item
  The conormal module \((I/I^2)_{\mathfrak q}\) is a free \(S_{\mathfrak q}\)-module of rank \(c\).
\item
  The ideal \(I_{\mathfrak q'}\) can be generated by a regular sequence in the regular local ring \(k[x_1, \ldots, x_n]_{\mathfrak q'}\).
\end{enumerate}

In this case any \(c\) elements of \(I_{\mathfrak q'}\) which generate \(I_{\mathfrak q'}/\mathfrak q'I_{\mathfrak q'}\) form a regular sequence in the local ring \(k[x_1, \ldots, x_n]_{\mathfrak q'}\).

\textbf{Proof.} Set \(R=k[x_1,\ldots,x_n]_{\mathfrak q'}\). Iterating \href{https://kokunoyumeto.github.io/open-math-courses/courses/AG-CA/regular-local-rings.html}{AG-CA-14, Proposition 3.3}, starting with \(k\), makes this a regular local ring. The full proof of AG-CA-12, Theorem 6.1, makes it Cohen--Macaulay. Its dimension is \(\operatorname{height}(\mathfrak q')\) by prime-localization correspondence. Moreover, \(\overline{R} = R/IR = R/I_{\mathfrak q'} = S_{\mathfrak q}\) is a quotient of dimension \(\text{height}(\mathfrak q)\). Let \(f_1, \ldots, f_c \in I_{\mathfrak q'}\) be elements which generate \((I/I^2)_{\mathfrak q}\). By Lemma \hyperref[AG-RG-S05-algebra-lemma-NAK]{F.10} we see that \(f_1, \ldots, f_c\) generate \(I_{\mathfrak q'}\). Since the dimensions work out, we conclude by Proposition \href{https://kokunoyumeto.github.io/open-math-courses/courses/AG-CA/regular-sequences-depth-and-cohen-macaulay-modules.html}{AG-CA-12, Theorem 4.2 and Corollary 4.3} that \(f_1, \ldots, f_c\) is a regular sequence in \(R\). By Lemma \hyperref[AG-RG-S05-algebra-lemma-regular-quasi-regular]{F.74} we see that \((I/I^2)_{\mathfrak q}\) is free. These arguments show that (2), (3), (4) are equivalent and that they imply the last statement of the lemma, and therefore they imply (5).

If (5) holds, say \(I_{\mathfrak q'}\) is generated by a regular sequence of length \(e\), then \(\text{height}(\mathfrak q) = \dim(S_{\mathfrak q}) =
\dim(k[x_1, \ldots, x_n]_{\mathfrak q'}) - e =
\text{height}(\mathfrak q') - e\) by dimension theory, see Section \hyperref[AG-RG-S05-dimension-context]{the dimension notation}. We conclude that \(e = c\). Thus (5) implies (2).

We continue with the notation introduced in the first paragraph. For each \(f_i\) we may find \(d_i \in k[x_1, \ldots, x_n]\), \(d_i \not \in \mathfrak q'\) such that \(f_i' = d_i f_i \in k[x_1, \ldots, x_n]\). Then it is still true that \(I_{\mathfrak q'} = (f_1', \ldots, f_c')R\). Hence there exists a \(g' \in k[x_1, \ldots, x_n]\), \(g' \not \in \mathfrak q'\) such that \(I_{g'} = (f_1', \ldots, f_c')\). Moreover, pick \(g'' \in k[x_1, \ldots, x_n]\), \(g'' \not \in \mathfrak q'\) such that \(\dim(S_{g''}) = \dim_{\mathfrak q} \operatorname{Spec}(S)\). By Lemma \hyperref[AG-RG-S05-algebra-lemma-codimension]{F.20} this dimension is equal to \(n - c\). Finally, set \(g\) equal to the image of \(g'g''\) in \(S\). Then we see that

\[
S_g \cong k[x_1, \ldots, x_n, x_{n + 1}]
/
(f_1', \ldots, f_c', x_{n + 1}g'g'' - 1)
\]

and by our choice of \(g''\) this ring has dimension \(n - c\). Therefore it is a global complete intersection. Thus each of (2), (3), and (4) implies (1).

Assume (1). Let \(S_g\cong k[y_1,\ldots,y_m]/J\), with \(J=(f_1,\ldots,f_t)\) and \(\dim S_g=m-t\), be a global complete-intersection presentation. Let \(\mathfrak q''\) be the polynomial prime corresponding to \(\mathfrak qS_g\). First justify the dimension at this particular point. Every minimal prime \(\mathfrak a\) of \(J\) has height at most \(t\) by \href{https://kokunoyumeto.github.io/open-math-courses/courses/AG-CA/dimension-theory-of-noetherian-local-rings.html}{AG-CA-11, Theorem 3.1}. Its component has dimension \(m-\operatorname{height}(\mathfrak a)\) by AG-CA-09, Theorem 4.3. The maximum of these finitely many component dimensions is \(\dim S_g=m-t\). Thus each such height is at least \(t\) as well, and every component has dimension exactly \(m-t\).

The complete point-dimension calculation in AG-CA-09, Theorem 6.1, now gives \(\dim_{\mathfrak qS_g}\operatorname{Spec}S_g=m-t\): every component through this point has that dimension. The local rings \((S_g)_{\mathfrak qS_g}\) and \(S_{\mathfrak q}\) are identical, so their heights agree. Applying \hyperref[AG-RG-S05-algebra-lemma-codimension]{F.20} to the polynomial presentation therefore gives \[
t=m-\dim_{\mathfrak qS_g}\operatorname{Spec}S_g
=\operatorname{height}(\mathfrak q'')-\operatorname{height}(\mathfrak q).
\] The ambient polynomial local ring is regular by iterating \href{https://kokunoyumeto.github.io/open-math-courses/courses/AG-CA/regular-local-rings.html}{AG-CA-14, Proposition 3.3}, starting with the field \(k\). It is Cohen--Macaulay by AG-CA-12, Theorem 6.1. Its quotient has dimension \(\operatorname{height}(\mathfrak q)\), so the displayed equality and the expected-dimension criterion AG-CA-12, Theorem 4.2, prove that the actual chosen \(f_1,\ldots,f_t\) form a regular sequence there. This includes the empty-list case. By Lemma \hyperref[AG-RG-S05-algebra-lemma-regular-quasi-regular]{F.74} we see that \((J/J^2)_{\mathfrak q}\) is free of rank \(t\). By Lemma \hyperref[AG-RG-S05-algebra-lemma-conormal-module-localize]{F.33} we have

\[
J/J^2 \oplus S_g^n \cong (I/I^2)_g \oplus S_g^m
\]

Localize this isomorphism at \(\mathfrak q\). Since its left side is finite free, \((I/I^2)_{\mathfrak q}\) is a finite direct summand of a finite free module. It is flat by restricting tensor injections and finitely presented because its complementary summand is finite. The complete local calculation in \href{https://kokunoyumeto.github.io/open-math-courses/courses/AG-RG/AG-RG-S01.html}{AG-RG-S01, Lemma 1.2} makes it finite free. Reduction to the residue field in this same isomorphism computes its rank. Thus \((I/I^2)_{\mathfrak q}\) is free of rank \(t + n - m = m - \dim(S_g) + n - m = n - \dim(S_g) =
\text{height}(\mathfrak q') - \text{height}(\mathfrak q) = c\). Thus we obtain (4). ∎

\phantomsection\label{AG-RG-S05-algebra-lemma-lci-at-prime}

\subsubsection{Lemma F.58: lci at prime}\label{AG-RG-S05-lemma-f.58-lci-at-prime}

Reading source: \href{https://stacks.math.columbia.edu/tag/00SG}{Stacks, Tag 00SG}.

For a finite-type \(k\)-algebra \(S\) and \(\mathfrak q\in\operatorname{Spec}S\), the following are equivalent: \(S_{\mathfrak q}\) is a complete intersection in the sense of \hyperref[AG-RG-S05-algebra-definition-lci-local-ring]{F.3}; it has a global complete-intersection principal neighbourhood; and the kernel of any finite polynomial presentation of \(S\) is generated by a regular sequence in the polynomial local ring above \(\mathfrak q\). Equivalently, there is a principal neighbourhood that is locally a complete intersection over \(k\).

\textbf{Proof.} Write \(C=S_{\mathfrak q}\) and \(P=k[X_1,\ldots,X_n]_{\mathfrak q'}\twoheadrightarrow C\). The ring \(P\) is regular local by \href{https://kokunoyumeto.github.io/open-math-courses/courses/AG-CA/regular-local-rings.html}{AG-CA-14, Proposition 3.3 and Theorem 3.2}. A regular generating list for its kernel therefore gives the ambient ring required by F.3. Conversely suppose F.3 gives \(C=Q/(a_1,\ldots,a_e)\), where \(Q\) is regular local and essentially of finite type over the same field \(k\), and this list is regular in its maximal ideal. We compare this local ambient with \(P\) explicitly.

Present \(Q\) as a quotient of a polynomial local ring \(V=k[Y_1,\ldots,Y_m]_{\mathfrak n}\). Both \(V\) and \(Q\) are regular local. AG-CA-14, Proposition 1.4, makes the kernel of \(V\to Q\) generated by part of a regular parameter system. Append lifts of the \(a_i\). The resulting list generates \(\ker(V\to C)\) and is regular: its first successive quotients are those by the parameter list, and its remaining multiplication injections are precisely those of the given regular list in \(Q\).

Let \(U=k[X,Y]_{\mathfrak r}\), where \(\mathfrak r\) is the preimage of the maximal ideal of \(C\) under the map using both sets of coordinates, and put \(L=\ker(U\to C)\). Lift the images of \(X_i\) to \(V\) and denote them by \(u_i(Y)\); these are fractions with denominators outside \(\mathfrak n\). Those denominators are units in \(U\). The quotient by \(X_i-u_i(Y)\) is \(V\), including its indicated localization: evaluating an element outside \(\mathfrak r\) gives an element outside \(\mathfrak n\), hence a unit, and every fraction defining \(V\) already occurs in \(U\). Translation of polynomial variables shows that the graph equations form a regular sequence. They remain regular after localization by exactness, and their quotient is nonzero. Appending the preceding regular list in \(V\) consequently gives a regular generating list for \(L\).

In the other direction lift the images of \(Y_j\) to fractions \(v_j(X)\) in \(P\). The identical graph calculation gives \[U/(Y_j-v_j(X))\cong P,\qquad L=(\ker(P\to C),Y_j-v_j(X))U.\] The constant and linear coefficient calculation of \hyperref[AG-RG-S05-algebra-lemma-conormal-module-localize]{F.33} works in this localized polynomial ring as well: first invert the finitely many denominators of the \(v_j\), translate \(Y_j\) by \(v_j\), and then localize at \(\mathfrak r\). Modulo \(L^2\) its constant terms are read modulo \(\ker(P\to C)^2\), its graph-linear coefficients modulo \(\ker(P\to C)\), and terms of graph-degree at least two vanish. The resulting mutually inverse maps give \[L/L^2\cong \ker(P\to C)/\ker(P\to C)^2\oplus C^m.\]

The regular list for \(L\) makes its conormal finite free by \hyperref[AG-RG-S05-algebra-lemma-regular-quasi-regular]{F.74}. Its rank is \(\dim U-\dim C\) by \href{https://kokunoyumeto.github.io/open-math-courses/courses/AG-CA/regular-sequences-depth-and-cohen-macaulay-modules.html}{AG-CA-12, Corollary 6.2}. Hence the first summand in the display is finite projective, and is finite free over the local ring \(C\) by \href{https://kokunoyumeto.github.io/open-math-courses/courses/AG-RG/AG-RG-S01.html}{S01, Lemma 1.2}. Reduction to its residue field computes its rank. The polynomial height formula \href{https://kokunoyumeto.github.io/open-math-courses/courses/AG-CA/krull-dimension-and-noether-normalization.html}{AG-CA-09, Theorem 4.3} gives \(\dim U=\dim P+m\), since the two polynomial primes have the same residue field \(\kappa(\mathfrak q)\). The rank is thus \(\dim P-\dim C\). F.57 now supplies a regular generating list for \(\ker(P\to C)\) and a global complete-intersection principal neighbourhood of \(\mathfrak q\). This proves the assertion for every chosen polynomial presentation, including the local ambients in F.3.

Conversely a global complete-intersection neighbourhood has the regular local equation lists proved in F.57, so gives F.3 at \(\mathfrak q\). It is also locally a complete intersection: apply F.57 at each of its points. A locally complete-intersection principal neighbourhood contains a further global complete-intersection principal neighbourhood of the specified point by the defining local condition. If the further inverted element is \(b/g^r\) in \(S_g\), that neighbourhood is \(S_{gb}\), with \(gb\notin\mathfrak q\). This proves the last equivalence. All graph and rank calculations allow empty coordinate or equation lists. ∎

\phantomsection\label{AG-RG-S05-algebra-lemma-limit-essentially-finite-presentation}

\subsubsection{Lemma F.59: limit essentially finite presentation}\label{AG-RG-S05-lemma-f.59-limit-essentially-finite-presentation}

\emph{Adapted from the Stacks project, algebra.tex, lines 33722--33770.}
\nopagebreak[4]

A local essentially finitely presented map \(R\to S\) is a directed colimit of local maps \(R_\lambda\to S_\lambda\) essentially of finite type over \(\mathbb Z\), with transition maps obtained by tensoring the base and localizing at the contracted prime.

\textbf{Proof.} Write \(S=(R[X_1,\ldots,X_n]/(f_1,\ldots,f_c))_{\mathfrak q}\). For each finite subset of \(R\) containing every equation coefficient, let \(A_\lambda\subset R\) be the subring it generates over the image of \(\mathbb Z\), and set \(R_\lambda=(A_\lambda)_{\mathfrak m_R\cap A_\lambda}\). These are local subrings of \(R\) with local inclusions. Their directed union is \(R\), since each element is included in some finite subset. Form \(C_\lambda=R_\lambda[X]/(f_i)\) and localize it at the inverse image \(\mathfrak q_\lambda\) of the chosen target prime, giving \(S_\lambda\). The maps are local and essentially finite type, and every stage is Noetherian by the Hilbert basis theorem.

Tensoring \(C_\lambda\) with \(R_\mu\) gives \(C_\mu\). After initially inverting the complement of \(\mathfrak q_\lambda\), localizing at the contracted prime inverts exactly the additional complement of \(\mathfrak q_\mu\); the fraction universal property identifies this with \(S_\mu\). Every fraction in \(S\) has a numerator and denominator involving finitely many coefficients and hence occurs at one stage. If two fractions become equal, cross multiplication gives an equality in the polynomial quotient after one denominator outside \(\mathfrak q\) is multiplied in. Its finite ideal-membership expression uses only finitely many coefficients, so the equality holds at a later stage. This proves both surjectivity and injectivity of \(\operatorname{colim}S_\lambda\to S\), completing every claimed colimit identification. ∎

\phantomsection\label{AG-RG-S05-algebra-lemma-limit-module-essentially-finite-presentation}

\subsubsection{Lemma F.60: limit module essentially finite presentation}\label{AG-RG-S05-lemma-f.60-limit-module-essentially-finite-presentation}

\emph{Adapted from the Stacks project, algebra.tex, lines 33804--33859.}
\nopagebreak[4]

A finitely presented module \(M\) over the local algebra in the preceding approximation has finitely presented models \(M_\lambda\) at sufficiently large stages, compatible by tensoring with \(S_\mu\), with colimit \(M\).

\textbf{Proof.} Choose a finite matrix presentation \(S^a\xrightarrow{H}S^b\to M\to0\). Its finitely many entries are fractions and occur in one \(S_{\lambda_0}\) by the preceding representative proof. Use their images to define \(M_\lambda=\operatorname{coker}(S_\lambda^a\xrightarrow{H_\lambda}S_\lambda^b)\) at all later stages. Tensor is right exact, so \(M_\lambda\otimes_{S_\lambda}S_\mu=M_\mu\). Taking the colimit of the finite free presentations gives the original cokernel: a vector represents zero precisely when it equals the image of a finite vector, and this finite equality holds at a common later stage. Thus the colimit is \(M\), each model is finitely presented, and all ring-transition assertions remain those proved above. ∎

\phantomsection\label{AG-RG-S05-algebra-lemma-local-dimension-zero-henselian}

\subsubsection{Lemma F.61: local dimension zero henselian}\label{AG-RG-S05-lemma-f.61-local-dimension-zero-henselian}

\emph{Adapted from the Stacks project, algebra.tex, lines 44097--44121.}
\nopagebreak[4]

Let \((R, \mathfrak m)\) be a local ring of dimension \(0\). Then \(R\) is henselian.

\textbf{Proof.} Let \(R \to S\) be a finite ring map. By Lemma \href{https://kokunoyumeto.github.io/open-math-courses/courses/AG-CA/henselian-local-rings-and-henselization.html}{AG-CA-20, Theorem 1.2 and Theorem 2.2} it suffices to show that \(S\) is a product of local rings. By Lemma \href{https://kokunoyumeto.github.io/open-math-courses/courses/AG-CA/integral-extensions-lying-over-going-up-and-going-down.html}{AG-CA-05, Theorem 3.4} \(S\) has finitely many primes \(\mathfrak m_1, \ldots, \mathfrak m_r\) which all lie over \(\mathfrak m\). There are no inclusions among these primes, see Lemma \href{https://kokunoyumeto.github.io/open-math-courses/courses/AG-CA/integral-extensions-lying-over-going-up-and-going-down.html}{AG-CA-05, Theorem 3.3}, hence they are all maximal. Every element of \(\mathfrak m_1 \cap \ldots \cap \mathfrak m_r\) is nilpotent by Lemma \href{https://kokunoyumeto.github.io/open-math-courses/courses/AG-CA/spectra-of-rings.html}{AG-CA-01, Theorem 1.2 and Proposition 2.1}. It follows \(S\) is the product of the localizations of \(S\) at the primes \(\mathfrak m_i\) by Lemma \hyperref[AG-RG-S05-algebra-lemma-product-local]{F.73}. ∎

\phantomsection\label{AG-RG-S05-algebra-lemma-local-syntomic}

\subsubsection{Lemma F.62: local syntomic}\label{AG-RG-S05-lemma-f.62-local-syntomic}

\emph{Adapted from the Stacks project, algebra.tex, lines 37927--37941.}
\nopagebreak[4]

If finitely many principal opens cover \(\operatorname{Spec}S\) and \(R\to S_{g_i}\) is syntomic on each, then \(R\to S\) is syntomic.

\textbf{Proof.} Finite presentation follows from \hyperref[AG-RG-S05-algebra-lemma-cover-upstairs]{F.35}, including its explicit finite substitution calculation. For flatness, tensor any injection of \(R\)-modules with \(S\) and localize its kernel at each \(g_i\). Each kernel is zero by the assumed chart flatness; local detection on this principal cover makes the original kernel zero. In every residue-field fibre these same opens cover, and their local complete-intersection covers combine to give a cover by global complete-intersection neighbourhoods. This is precisely the fibre definition of syntomicness. No Noetherian hypothesis is required for this gluing step. ∎

\phantomsection\label{AG-RG-S05-algebra-lemma-localization-colimit}

\subsubsection{Lemma F.63: localization colimit}\label{AG-RG-S05-lemma-f.63-localization-colimit}

\emph{Adapted from the Stacks project, algebra.tex, lines 1260--1278.}
\nopagebreak[4]

Localization of a module is the filtered colimit of copies of that module with transition maps given by multiplication by denominators. In particular a localization of a ring is the filtered colimit of its principal localizations.

\textbf{Proof.} For the first assertion, index the copies by finite products of elements of the multiplicative set, ordered with multiplication maps (or use the filtered divisibility category). Send an element \(m\) in the copy indexed by \(s\) to \(m/s\). Every fraction has such a representative. Equality \(m/s=n/t\) means \(u(tm-sn)=0\) for some denominator \(u\), precisely the equality at the common successor \(ust\). Thus the map is bijective and respects addition and scalars. For rings use the principal subalgebras \(R_s\): two such stages map to \(R_{st}\); every element and every equality in the full localization is represented or becomes equal at a principal successor by the same fraction criterion. This proves the ring colimit assertion. ∎

\phantomsection\label{AG-RG-S05-algebra-lemma-localization-smooth-separable}

\subsubsection{Lemma F.64: localization smooth separable}\label{AG-RG-S05-lemma-f.64-localization-smooth-separable}

\emph{Adapted from the Stacks project, algebra.tex, lines 46075--46088.}
\nopagebreak[4]

A finitely generated separably generated extension \(K/k\) is the fraction field of a smooth finite-type \(k\)-domain. Conversely a fraction field of a smooth finite-type \(k\)-domain is separably generated.

\textbf{Proof.} For the forward implication choose a finite separating transcendence basis \(T\) and a finite tower of separable algebraic generators \(y_1,\ldots,y_r\). At each step, express the coefficients of its monic minimal polynomial in the standard monomial basis of the preceding field over \(k(T)\). Clear the finitely many rational-function denominators by one \(0\ne h(T)\in k[T]\). This gives monic equations \(F_i\in k[T,1/h,Y_1,\ldots,Y_i]\) and \[C=k[T,1/h,Y_1,\ldots,Y_r]/(F_1,\ldots,F_r).\] Successive monic division gives the standard monomials as a free basis over \(k[T,1/h]\). Over \(k(T)\) the quotient is exactly the chosen field tower \(K\). Since the coefficient ring is a domain and \(C\) is free, \(C\to K\) is injective. Its fraction field is \(K\), because it contains \(T\) and all the \(y_i\).

The evaluated derivatives \(\Delta_i=(\partial F_i/\partial Y_i)(y)\) are nonzero, since each tower step is separable. Put \(\Delta=\prod_i\Delta_i\). The algebra \(C_\Delta\) is a succession of standard étale extensions over \(k[T,1/h]\): at step \(i\) invert \(\Delta_i\) and retain the preceding inverses. AG-CA-17, Proposition 6.1, proves the unique lifts for each such equation. Polynomial algebras and localization are formally smooth by its Theorem 1.2, as are their composites with these maps. All variables, equations and inverted elements are finite in number, so \(C_\Delta\) is finitely presented over \(k\). It is therefore smooth. It remains a domain inside \(K\) with fraction field \(K\). Empty towers use \(\Delta=1\), and the same construction applies.

For the converse take a nonempty principal standard smooth chart \(D\) in the given domain, using AG-CA-17, Theorem 5.1. Its fraction field is still \(K\). Reorder its variables so that the selected square Jacobian uses the pivot variables, and let \(U_1,\ldots,U_t\) be the nonpivot variables. Viewed over \(R_0=k[U_1,\ldots,U_t]\), this presentation has only the pivot variables and the same square invertible Jacobian. The explicit lifting calculation AG-CA-17, Theorem 4.1, gives formal smoothness and zero relative differentials; Theorem 2.2 gives formal étaleness. Its finite presentation makes the map étale, hence flat by AG-CA-18, Corollary 6.2.

The map \(R_0\to D\) is injective. Indeed, multiplication by a nonzero element of the polynomial domain is injective, and flatness preserves this injection after tensoring with \(D\). An element of the structural kernel would act as zero on the nonzero ring \(D\), contradicting that injection. Thus the \(U_i\) are algebraically independent in \(K\). AG-CA-18, Theorem 2.1, at the generic point gives \(t=\operatorname{trdeg}_kK\). They are consequently a transcendence basis. The extension \(K/k(U_1,\ldots,U_t)\) is finite: it is algebraic and has finitely many field generators, whose finite degrees multiply along a finite tower.

Formal étaleness survives base change and localization by AG-CA-17, Theorem 1.2. First invert the nonzero elements of \(R_0\), and then those of \(D\); this gives the formally étale field map \(k(U_1,\ldots,U_t)\to K\). Its differentials vanish by Theorem 2.2. The finite-field converse AG-CA-17, Lemma 7.1, therefore makes it separable. This constructs the required separating transcendence basis directly, without using Proposition F.91. ∎

\phantomsection\label{AG-RG-S05-algebra-lemma-localize-NL}

\subsubsection{Lemma F.65: localize NL}\label{AG-RG-S05-lemma-f.65-localize-nl}

\emph{Adapted from the Stacks project, algebra.tex, lines 37118--37133.}
\nopagebreak[4]

Let \(A \to B\) be a ring map. Let \(S \subset B\) be a multiplicative subset. The canonical map \(\NL_{B/A} \otimes_B S^{-1}B \to \NL_{S^{-1}B/A}\) is a quasi-isomorphism.

\textbf{Proof.} We have \(S^{-1}B = \colim_{g \in S} B_g\) where we think of \(S\) as a directed set (ordering by divisibility), see Lemma \hyperref[AG-RG-S05-algebra-lemma-localization-colimit]{F.63}. By Lemma \hyperref[AG-RG-S05-algebra-lemma-principal-localization-NL]{F.72} each of the maps \(\NL_{B/A} \otimes_B B_g \to \NL_{B_g/A}\) is a quasi-isomorphism. The lemma follows from Lemma \hyperref[AG-RG-S05-algebra-lemma-colimits-NL]{F.24}. ∎

\phantomsection\label{AG-RG-S05-algebra-lemma-localize-relative-complete-intersection}

\subsubsection{Lemma F.66: localize relative complete intersection}\label{AG-RG-S05-lemma-f.66-localize-relative-complete-intersection}

Reading source: \href{https://stacks.math.columbia.edu/tag/00ST}{Stacks, Tag 00ST}.

Let \(R\) be Noetherian and put \(S=R[X_1,\ldots,X_n]/(f_1,\ldots,f_c)\). In each situation below there are \(g\in S\) and a polynomial lift \(h\) of \(g\) for which \[S_g=R[X_1,\ldots,X_n,T]/(f_1,\ldots,f_c,hT-1)\] is a relative global complete intersection as defined in \hyperref[AG-RG-S05-algebra-definition-relative-global-complete-intersection]{F.4}.

\begin{enumerate}
\def\labelenumi{\arabic{enumi}.}
\tightlist
\item
  If every nonempty fibre over \(V(I)\), for an ideal \(I\subset R\), has dimension \(n-c\), one may require \(g\equiv1\pmod{IS}\).
\item
  If the fibre over \(\mathfrak p\in\operatorname{Spec}R\) has dimension \(n-c\), one may require the image of \(g\) in that fibre to be a unit.
\item
  If a point \(\mathfrak q\) of the fibre over \(\mathfrak p\) has fibre point-dimension \(n-c\), one may require \(g\notin\mathfrak q\).
\end{enumerate}

\textbf{Proof.} Set \(d=n-c\). The neighbourhood bound proved in \hyperref[AG-RG-S05-algebra-lemma-dimension-fibres-bounded-open-upstairs]{F.39} gives, around every point whose fibre point-dimension is at most \(d\), a principal neighbourhood on which every fibre point-dimension is at most \(d\). The union of these neighbourhoods is an open \(W\) with that same bound. Each of its nonempty fibres has dimension at most \(d\): a finite-type field fibre is Noetherian, and its component dimensions are its point-dimensions at the corresponding generic points, as in \hyperref[AG-RG-S05-algebra-lemma-codimension]{F.20}. Open subsets retain this bound.

In case (1) choose \(W\) containing \(V(IS)\); each point of those fibres has point-dimension at most their assumed dimension. Write the closed complement of \(W\) as \(V(J)\) in \(\operatorname{Spec}S\). The equality \(V(J+IS)=\varnothing\) implies \(J+IS=S\): a proper ideal would be contained in a maximal, hence prime, ideal. Write \(1=g+b\) with \(g\in J\) and \(b\in IS\). Then \(g\equiv1\pmod{IS}\), and \(D(g)\subset W\).

In case (2) take \(W\) containing the entire indicated fibre, and again write its complement as \(V(J)\). If \(k=\kappa(\mathfrak p)\), the image ideal \(J(S\otimes_R k)\) is the unit ideal. Choose a finite expression \(1=\sum_j\beta_j\bar g_j\), where \(g_j\in J\). Each \(\beta_j\) is represented in \(S_{R\setminus\mathfrak p}/\mathfrak pS_{R\setminus\mathfrak p}\). Clear their finitely many denominators to obtain \(b_j\in S\) and \(s\in R\setminus\mathfrak p\) with \(\beta_j=\bar b_j/\bar s\). Thus \(g=\sum_j b_jg_j\in J\) has fibre image \(\bar s\), a unit. Again \(D(g)\subset W\).

In case (3), one of the principal neighbourhoods supplied directly by F.39 is \(D(g)\) with \(g\notin\mathfrak q\) and the required upper bound. In all three cases select any polynomial lift \(h\) of \(g\). Adjoining the inverse of \(h\) gives the presentation in the statement, by the explicit substitution \(T\mapsto g^{-1}\) and its inverse. Every nonempty fibre of this presentation has dimension at least \[ (n+1)-(c+1)=d.\] Indeed, in its polynomial ring over the residue field a minimal prime over \(c+1\) equations has height at most \(c+1\), by \href{https://kokunoyumeto.github.io/open-math-courses/courses/AG-CA/dimension-theory-of-noetherian-local-rings.html}{AG-CA-11, Theorem 3.1}; \href{https://kokunoyumeto.github.io/open-math-courses/courses/AG-CA/krull-dimension-and-noether-normalization.html}{AG-CA-09, Theorem 4.3} gives that component dimension as \(n+1\) minus its height. The upper bound on \(D(g)\) makes the dimension exactly \(d\). This proves the relative global complete-intersection assertion, without presupposing flatness; its flatness is subsequently proved in F.77. Empty fibres impose no dimension test, and a zero localization is included. ∎

\phantomsection\label{AG-RG-S05-algebra-lemma-locally-smooth}

\subsubsection{Lemma F.67: locally smooth}\label{AG-RG-S05-lemma-f.67-locally-smooth}

Reading source: \href{https://stacks.math.columbia.edu/tag/00TC}{Stacks, Tag 00TC}.

For any ring map \(R\to S\), smoothness is equivalent to the existence, around every prime of \(S\), of a principal localization that is smooth over \(R\).

\textbf{Proof.} Use the finite-presentation and cotangent-complex definition \hyperref[AG-RG-S05-algebra-definition-smooth]{F.6}. If \(S/R\) is smooth, then \(S_g\) is finitely presented by adjoining an inverse variable. The exact calculation \hyperref[AG-RG-S05-algebra-lemma-principal-localization-NL]{F.72}, together with the presentation comparison \hyperref[AG-RG-S05-algebra-lemma-NL-homotopy]{F.11}, identifies its conormal complex with the localized complex plus a contractible summand. Its first homology is zero and its differential module is the localization of a finite projective module. Localization keeps the splitting of a finite free module, so that module is still finite projective. Thus \(S_g/R\) is smooth.

For the converse, the elements defining the given neighbourhoods generate the unit ideal: otherwise a maximal ideal containing them would be an uncovered prime. A finite expression for \(1\) gives a cover by \(D(g_1),\ldots,D(g_r)\). The complete gluing proof \hyperref[AG-RG-S05-algebra-lemma-cover-upstairs]{F.35} makes \(S/R\) finitely presented. Fix \(S=R[X_1,\ldots,X_n]/I\) with \(I\) finite. For this presentation set \[K=\ker(I/I^2\xrightarrow{d}S^n),\qquad M=\operatorname{coker}(I/I^2\xrightarrow{d}S^n)=\Omega_{S/R}.\] The derivatives of a finite list generating \(I\) generate the image of \(d\), so \(M\) is finitely presented. Exact localization and F.72/F.11 identify \(K_{g_i}\) with the zero first homology of each smooth chart. The power-annihilator module calculation in F.35 therefore gives \(K=0\).

It remains to produce a global projective splitting for \(M\). Write \(\pi:S^n\twoheadrightarrow M\). On each chart \(M_{g_i}\) is finite projective. A section \(s_i\) of \(\pi_{g_i}\) exists: write that projective module as a summand of a finite free module, lift the free basis through \(\pi_{g_i}\), and restrict the resulting lift to the summand.

For a finitely presented module, localization of \(\operatorname{Hom}_S(M,S^n)\) is \(\operatorname{Hom}_{S_{g_i}}(M_{g_i},S_{g_i}^n)\). Indeed a finite presentation of \(M\) identifies both modules with the kernel of the same matrix between finite powers of \(S^n\), and localization is exact. Clear the denominator of \(s_i\) to obtain \(t_i:M\to S^n\) whose localized composite is \(g_i^{a_i}1_M\). The error \(\pi t_i-g_i^{a_i}1_M\) vanishes on a finite generating list after localization. A further common power of \(g_i\) kills those finitely many errors. After multiplying \(t_i\) by that power, we have an actual equality \[\pi t_i=g_i^{e_i}1_M.\] Increase \(e_i\) to a positive integer if needed. These powers generate \(1\): expand a finite unit expression in the \(g_i\) to exponent \(1+\sum_i(e_i-1)\); every monomial contains one of the required powers. Choose \(c_i\) with \(\sum_i c_ig_i^{e_i}=1\). Then \(t=\sum_i c_it_i\) satisfies \(\pi t=1_M\). Thus \(S^n=t(M)\oplus\ker\pi\), with both summands generated by projections of its finite basis. This proves finite projectivity of \(M\) directly.

We have proved finite presentation, zero first cotangent homology and finite projective differentials, exactly F.6. If \(S=0\), its finite presentation is \(R/(1)\) and both cotangent modules are zero, so the conclusion also holds for the empty spectrum. No Noetherian hypothesis is needed. ∎

\phantomsection\label{AG-RG-S05-algebra-lemma-make-etale-map-prescribed-residue-field}

\subsubsection{Lemma F.68: make etale map prescribed residue field}\label{AG-RG-S05-lemma-f.68-make-etale-map-prescribed-residue-field}

\emph{Adapted from the Stacks project, algebra.tex, lines 41287--41319.}
\nopagebreak[4]

For a prime \(\mathfrak p\) of \(R\) and a finite separable extension \(L/\kappa(\mathfrak p)\), some finitely presented etale \(R\)-algebra has a point over \(\mathfrak p\) with this prescribed residue field.

\textbf{Proof.} Choose finitely many elements generating \(L\) as a field, and adjoin them one at a time. At a given step their monic minimal polynomial over the preceding residue field has nonzero derivative at the chosen element, since it is separable. Lift its finitely many coefficients to the local ring of the already constructed algebra at the chosen prime. Clear finitely many denominators to obtain that polynomial over one principal neighbourhood. Adjoin a root of this monic polynomial and invert its derivative. The square one-equation Jacobian calculation of AG-CA-17, Proposition 6.1, makes this algebra etale; its chosen residue-field evaluation maps the root to the required element. The kernel of this evaluation at the chosen point has exactly the enlarged residue field, since the localized base residue field and the root generate it. The polynomial coefficients and inverses involved at each of the finitely many steps descend from the local ring by finitely many denominator clearings. Composing the resulting finitely presented etale maps gives the asserted \(R\)-algebra. This construction uses a finite separable tower and does not assume a primitive element theorem. ∎

\phantomsection\label{AG-RG-S05-algebra-lemma-make-separably-generated}

\subsubsection{Lemma F.69: make separably generated}\label{AG-RG-S05-lemma-f.69-make-separably-generated}

\emph{Adapted from the Stacks project, algebra.tex, lines 10182--10253.}
\nopagebreak[4]

For a finitely generated field extension \(K/k\) there are finite purely inseparable extensions \(k'/k\) and \(K'/K\) with a commutative field square and \(K'/k'\) separably generated.

\textbf{Proof.} Choose a finite field-generating list for \(K/k\) and a maximal algebraically independent sublist \(T=(T_1,\ldots,T_d)\). Every remaining generator is algebraic over \(F=k(T)\): a polynomial dependence involving it has a nonzero coefficient after evaluating the variables in \(T\), because those variables are independent. Adjoining the remaining algebraic generators gives a finite tower. At a simple step, division by the monic minimal polynomial gives a unique remainder of smaller degree. The resulting finite-dimensional domain is a field: multiplication by a nonzero element is injective and hence surjective, giving its inverse. It therefore equals the simple field extension, with the powers below the minimal-polynomial degree as a basis. Products of these bases form a basis of the whole tower: they span by successive expansion, and grouping a putative relation by the last generator's powers proves independence inductively. This proves finiteness and multiplication of all degrees used below. In characteristic zero every irreducible polynomial has nonzero derivative, hence is separable; the finite-tower separability proof in AG-CA-17, Lemma 7.1, gives the assertion with \(k'=k\) and \(K'=K\).

Now assume characteristic \(p>0\). When a finite list of \(p\)th roots is needed, adjoin them successively: choose an irreducible factor of each \(X^p-a\) and take its field quotient. Such a factor exists by descent on the positive degree, and the Euclidean algorithm makes the quotient a field. Each polynomial has at most one root in a field extension, because \(u^p=v^p\) gives \((u-v)^p=0\), hence \(u=v\). Retain the existing fields as subfields at every adjunction; all roots and composita below use these compatible inclusions. We first construct the separable subfield of any finite extension \(E/F\). If \(x\in E\) has monic irreducible polynomial \(f\), write \[f(X)=g(X^{p^e})\] with \(e\) as large as possible. The polynomial \(g\) is irreducible, since a factorization would substitute to one of \(f\), and \(g'\ne0\). Euclidean division gives \(\gcd(g,g')=1\): irreducibility and the smaller degree of \(g'\) exclude a nonconstant common factor. Thus no root is repeated, and \(x^{p^e}\) is separable over \(F\).

If \(x\) and \(y\) are separable over \(F\), the minimal polynomial of \(y\) over \(F(x)\) divides its separable polynomial over \(F\), so the two-step tower \(F\subset F(x)\subset F(x,y)\) has separable steps. The complete finite-tower argument AG-CA-17, Lemma 7.1, makes every element of \(F(x,y)\) separable over \(F\). Consequently the separable elements in \(E\) are closed under sums, products and nonzero inverses. They form a field \(E_s\), finite-dimensional over \(F\) as a subspace of \(E\), and every one of its elements is separable over \(F\).

For each of a finite generating list of \(E/F\), the preceding irreducible-polynomial calculation puts some \(p\)-power in \(E_s\). Choose a common exponent \(e\). Frobenius applied to a rational expression in these generators then puts every \(p^e\)th power from \(E\) in \(E_s\). Hence \(E/E_s\) is purely inseparable. Its degree is a power of \(p\): for each generator, adjoin successively its powers in reverse order, beginning with the power already in \(E_s\). Every nontrivial step has the form \(u^p\in H\), \(u\notin H\), and has degree \(p\) by the proper-factor proof in \hyperref[AG-RG-S05-more-algebra-lemma-p-basis]{F.119}. Multiplication of the finite tower degrees proves the assertion.

Apply this construction to \(K/F\), writing its separable subfield as \(K_s\). We use strong induction on \(m=[K:K_s]\) for all finitely generated extensions together with the selected transcendence basis. If \(m=1\), that basis already separates \(K/k\). If \(m>1\), choose \(x\in K\setminus K_s\) and let \(e\ge1\) be the least exponent with \(x^{p^e}\in K_s\). Put \(\beta=x^{p^{e-1}}\) and \(\alpha=\beta^p\). Then \(\beta\notin K_s\) and \(\alpha\in K_s\). The same degree-\(p\) proof gives \[[K_s(\beta):K_s]=p,\qquad [K:K_s(\beta)]=m/p.\]

Write the monic separable minimal polynomial of \(\alpha\) over \(F\) as \[P(X)=X^n+\sum_{i<n}a_iX^i.\] Each \(a_i\) is a rational function in \(T\) involving finitely many coefficients of \(k\). Using the finite-root construction over \(K\), adjoin the \(p\)th roots of all those coefficients. Let \(k'\) be the subfield they generate over \(k\). It is finite and purely inseparable over \(k\), with \((k')^p\subset k\). Next adjoin \(U_j=T_j^{1/p}\) by the same construction. A polynomial relation among the \(U_j\) over \(k'\) would, after taking its \(p\)th power, be a nonzero polynomial relation among the \(T_j\) over \(k\): every coefficient of \(k'\) has its \(p\)th power in \(k\), and Frobenius is injective on this field. Thus the \(U_j\) are algebraically independent over \(k'\).

Set \(F'=k'(U_1,\ldots,U_d)\) and \(L=KF'\). This extension of \(K\) is finite and purely inseparable: it adjoins only the finitely many coefficient roots and the \(U_j\), all with \(p\)th powers in \(K\); the \(p\)th power of any rational expression in them lies in \(K\). For each \(a_i\), take the coefficientwise \(p\)th roots of its numerator and denominator and substitute the \(U_j\). The denominator is nonzero because its \(p\)th power is the original nonzero denominator. We obtain \(b_i\in F'\) with \(b_i^p=a_i\). Put \[Q(X)=X^n+\sum_{i<n}b_iX^i,\qquad Q(X)^p=P(X^p).\] Evaluation at \(\beta\) gives \(Q(\beta)^p=P(\alpha)=0\), hence \(Q(\beta)=0\). Also \[Q'(\beta)^p=P'(\alpha)\ne0.\] If \(h\) is the irreducible minimal polynomial of \(\beta\) over \(F'\), write \(Q=hq\). Evaluating \(Q'=h'q+hq'\) at \(\beta\) makes \(h'(\beta)q(\beta)\ne0\). Thus \(h'\ne0\), so \(\beta\) is separable over \(F'\).

For each separable irreducible minimal polynomial, the gcd calculation above gives a Euclidean identity with its derivative equal to one. That identity remains equal to one after mapping the coefficients to a field extension, so the polynomial still has no repeated root. Each factor, in particular each new minimal polynomial, consequently has distinct roots. Choose finite generators of \(K_s/F\); their separable polynomials therefore make the successive generator tower for \(K_sF'/F'\) separable. Adjoining the now separable \(\beta\) gives a further separable step. The finite-tower proof AG-CA-17, Lemma 7.1, gives \[H=K_s(\beta)F'\subset L_s,\] where \(L_s\) is the separable subfield of \(L/F'\) constructed above.

Here is the needed degree bound after base change. Choose a basis of \(K\) over \(K_s(\beta)\) of size \(m/p\). Its \(H\)-linear span inside \(L\) is a subalgebra: products of basis elements expand with coefficients in \(K_s(\beta)\subset H\). It is a finite-dimensional domain over \(H\), hence a field, since multiplication by a nonzero element is an injective linear map and therefore surjective. It contains both \(K\) and \(F'\), so it is exactly their compositum \(L\). Consequently \[[L:L_s]\le[L:H]\le m/p<m.\] The extension \(L/k'\) is finitely generated, the \(U_j\) are its transcendence basis, and \(L/F'\) is finite. Strong induction applied to this data gives finite purely inseparable \(k''/k'\) and \(L'/L\) with \(L'/k''\) separably generated.

Finally these purely inseparable extensions compose. If \(z\in L'\), some \(p\)-power lies in \(L\) and a further \(p\)-power lies in \(K\), so \(L'/K\) is purely inseparable; the same argument applies to \(k''/k\). Their degrees are finite by the tower basis calculation. Taking the final \(k'=k''\) and \(K'=L'\) gives the required commutative field square. This also covers the empty transcendence basis and trivial root extensions. ∎

\phantomsection\label{AG-RG-S05-algebra-lemma-matrix-left-inverse}

\subsubsection{Lemma F.70: matrix left inverse}\label{AG-RG-S05-lemma-f.70-matrix-left-inverse}

Reading source: \href{https://stacks.math.columbia.edu/tag/07DQ}{Stacks, Tag 07DQ}.

Let \(A\) be an \(n\times m\) matrix over a commutative ring \(R\), where \(0\leq m\leq n\), and let \(J\) be the ideal of its maximal minors. Every \(f\in J\) has an \(m\times n\) matrix \(B\) with \(BA=f1_m\). Conversely, existence of such a \(B\) implies \(f^m\in J\).

\textbf{Proof.} For an increasing \(m\)-element subset \(I\) of \(\{1,\ldots,n\}\), let \(E_I\) select those coordinates from \(R^n\), and put \(A_I=E_IA\). Its determinant is the corresponding minor. By the complete cofactor calculation in \hyperref[AG-RG-S05-algebra-lemma-charpoly]{F.19}, \[
\operatorname{adj}(A_I)E_IA=(\det A_I)1_m.
\] Thus, from an expression \(f=\sum_I c_I\det A_I\), the finite sum \[
B=\sum_I c_I\operatorname{adj}(A_I)E_I
\] constructs the required matrix.

For the converse, here is the determinant product calculation. Expand each column of \(BA\) as the linear combination of columns of \(B\) whose coefficients are entries of that column of \(A\). Multilinearity follows directly from the determinant permutation formula. A term choosing the same column twice is zero by the equal-column cancellation. For any fixed set \(I\) of distinct chosen columns, order those indices increasingly. Reordering them in all possible ways contributes their permutation signs; the signed sum of their coefficient products is exactly \(\det A_I\). The remaining column determinant is \(\det B^I\), where \(B^I\) uses the columns indexed by \(I\). Consequently \[
\det(BA)=\sum_{|I|=m}(\det B^I)(\det A_I)\in J.
\] When \(BA=f1_m\), its determinant is \(f^m\). This proves the converse over arbitrary rings, including rings with nilpotents. If \(m=0\), the sole minor is the empty determinant \(1\), so \(J=R\). The unique map to and from the zero module supplies the first assertion, and \(f^0=1\in J\) supplies the second. ∎

\phantomsection\label{AG-RG-S05-algebra-lemma-presentation-sym-exterior}

\subsubsection{Lemma F.71: presentation sym exterior}\label{AG-RG-S05-lemma-f.71-presentation-sym-exterior}

Reading source: \href{https://stacks.math.columbia.edu/tag/00DO}{Stacks, Tag 00DO}. The quotient maps and both kernels are proved below.

For an exact sequence \(M_2\xrightarrow{u}M_1\xrightarrow{v}M\to0\) of modules over any commutative ring and every integer \(n\geq1\), there are exact sequences \[
M_2\otimes_R\operatorname{Sym}^{n-1}M_1
\longrightarrow\operatorname{Sym}^nM_1
\longrightarrow\operatorname{Sym}^nM\longrightarrow0,
\] \[
M_2\otimes_R\bigwedge^{n-1}M_1
\longrightarrow\bigwedge^nM_1
\longrightarrow\bigwedge^nM\longrightarrow0.
\] The first maps multiply, respectively wedge, with \(u(x)\); the second maps apply \(v\) to every factor. In degree zero both second maps are the identity of \(R\), so both kernels are zero.

\textbf{Proof.} Put \(N=\operatorname{im}u\), so exactness identifies \(M=M_1/N\). Construct the tensor algebra \(T(M_1)=\bigoplus_{d\geq0}M_1^{\otimes d}\) with multiplication by concatenation. Balancing in each tensor factor makes that product well defined. Its symmetric quotient imposes \(xy=yx\) for degree-one generators; its exterior quotient imposes \(x^2=0\) for every degree-one generator. In the latter quotient, expanding \((x+y)^2=0\) gives \(xy=-yx\) over every ring, including characteristic two.

In either algebra let \(L\) be the two-sided homogeneous ideal generated by the degree-one image of \(N\). The factorwise map to the corresponding algebra of \(M_1/N\) is onto: every pure tensor in the target lifts, and pure tensors generate each degree. It kills \(L\). Its inverse on the quotient is explicit. Send a class \(\bar x\in M_1/N\) to the class of any lift \(x\) in the source algebra modulo \(L\). Another lift differs by an element of \(N\), hence gives the same class. This map is linear, so its products extend to the tensor algebra by the tensor balancing relations. Those products obey commutation, respectively square-zero, relations in the indicated quotient; it therefore descends to the symmetric, respectively exterior, algebra. The two composites fix every degree-one generator and the scalar ring, hence every product and sum of generators. They are inverse graded algebra maps.

It remains to identify the degree-\(n\) part of \(L\). It is spanned by products having one selected factor in \(N\), since those generate the ideal and all factors have nonnegative degree. In the symmetric case move that factor to the front by commutation. In the exterior case move it to the front using \(xy=-yx\), retaining the resulting sign. Every remaining product has degree \(n-1\). Thus that part is exactly the image of \(N\otimes_R\operatorname{Sym}^{n-1}M_1\), respectively \(N\otimes_R\bigwedge^{n-1}M_1\), under multiplication. Each selected element of \(N\) lifts through \(u\), so these are also the images of the stated maps from \(M_2\). The inverse quotient maps just constructed prove that these images are exactly the kernels. The ideal has no degree-zero part, proving the separate degree-zero assertion. No module is assumed finite or free. ∎

\phantomsection\label{AG-RG-S05-algebra-lemma-principal-localization-NL}

\subsubsection{Lemma F.72: principal localization NL}\label{AG-RG-S05-lemma-f.72-principal-localization-nl}

\emph{Adapted from the Stacks project, algebra.tex, lines 37031--37116.}
\nopagebreak[4]

Let \(A \to B\) be a ring map. Let \(g \in B\). Suppose \(\alpha : P \to B\) is a presentation with kernel \(I\). Then a presentation of \(B_g\) over \(A\) is the map

\[
\beta : P[x] \longrightarrow B_g
\]

extending \(\alpha\) and sending \(x\) to \(1/g\). The kernel \(J\) of \(\beta\) is generated by \(I\) and the element \(f x - 1\) where \(f \in P\) is an element mapped to \(g \in B\) by \(\alpha\). In this situation we have

\begin{enumerate}
\def\labelenumi{\arabic{enumi}.}
\item
  \(J/J^2 = (I/I^2)_g \oplus B_g (f x - 1)\),
\item
  \(\Omega_{P[x]/A} \otimes_{P[x]} B_g =
  \Omega_{P/A} \otimes_P B_g \oplus B_g \text{d}x\),
\item
  \(\operatorname{NL}(\beta) \cong
  \operatorname{NL}(\alpha) \otimes_B B_g \oplus (B_g \xrightarrow{g} B_g).\)
\end{enumerate}

Hence the canonical map \(\NL_{B/A} \otimes_B B_g \to \NL_{B_g/A}\) is a homotopy equivalence.

\textbf{Proof.} Since \(P[x]/(I, fx - 1) = B[x]/(gx - 1) = B_g\) we get the statement about \(I\) and \(fx - 1\) generating \(J\). Consider the commutative diagram

\[
\begin{gathered}
0 \xrightarrow{} \Omega_{P/A} \otimes_P B_g \\
\Omega_{P/A} \otimes_P B_g \xrightarrow{} \Omega_{P[x]/A} \otimes_{P[x]} B_g \\
\Omega_{P[x]/A} \otimes_{P[x]} B_g \xrightarrow{} \Omega_{B[x]/B} \otimes_{B[x]} B_g \\
\Omega_{B[x]/B} \otimes_{B[x]} B_g \xrightarrow{} 0 \\
(I/I^2)_g \xrightarrow{} J/J^2 \\
(I/I^2)_g \xrightarrow{} \Omega_{P/A} \otimes_P B_g \\
J/J^2 \xrightarrow{} (gx - 1)/(gx - 1)^2 \\
J/J^2 \xrightarrow{} \Omega_{P[x]/A} \otimes_{P[x]} B_g \\
(gx - 1)/(gx - 1)^2 \xrightarrow{} 0 \\
(gx - 1)/(gx - 1)^2 \xrightarrow{} \Omega_{B[x]/B} \otimes_{B[x]} B_g
\end{gathered}
\]

with exact rows of Lemma \hyperref[AG-RG-S05-algebra-lemma-exact-sequence-NL]{F.42}. The \(B_g\)-module \(\Omega_{B[x]/B} \otimes_{B[x]} B_g\) is free of rank \(1\) on \(\text{d}x\). The element \(\text{d}x\) in the \(B_g\)-module \(\Omega_{P[x]/A} \otimes_{P[x]} B_g\) provides a splitting for the top row. The element \(gx - 1 \in (gx - 1)/(gx - 1)^2\) is mapped to \(g\text{d}x\) in \(\Omega_{B[x]/B} \otimes_{B[x]} B_g\) and hence \((gx - 1)/(gx - 1)^2\) is free of rank \(1\) over \(B_g\). Here is a direct conormal calculation that also includes \(B_g=0\). Put \(q=gx-1\). Multiplication by \(q\) on \(B[x]\) is injective: for a nonzero polynomial, the coefficient of its lowest nonzero degree in its product with \(q\) is the negative of that coefficient. The map \[
B[x]/(q)\longrightarrow(q)/(q^2),\qquad [a]\longmapsto[aq],
\] is surjective. If \(aq=bq^2\), injectivity of multiplication by \(q\) gives \(a=bq\), so it is injective as well. Thus \((q)/(q^2)\) is the free rank-one \(B_g\)-module with generator \(q\). When \(g\) is nilpotent, \(B_g=0\) and both modules are zero, so the same calculation applies without a proper regular-sequence hypothesis.

Let us prove \((I/I^2)_g \to J/J^2\) is injective. Consider the \(P\)-algebra map

\[
\pi : P[x] \to (P/I^2)_f = P_f/I_f^2
\]

sending \(x\) to \(1/f\). Since \(J\) is generated by \(I\) and \(fx - 1\) we see that \(\pi(J) \subset (I/I^2)_f = (I/I^2)_g\). Since this is an ideal of square zero we see that \(\pi(J^2) = 0\). If \(a \in I\) maps to an element of \(J^2\) in \(J\), then \(\pi(a) = 0\), which implies that \(a\) maps to zero in \(I_f/I_f^2\). This proves the desired injectivity.

Thus we have a short exact sequence of two term complexes

\[
0 \to \operatorname{NL}(\alpha) \otimes_B B_g \to \operatorname{NL}(\beta)
\to (B_g \xrightarrow{g} B_g) \to 0
\]

For this sequence the following splittings commute with the differentials:

\[
J/J^2 = (I/I^2)_g \oplus B_g (fx - 1)\quad\text{and}\quad
\Omega_{P[x]/A} \otimes B_g = \Omega_{P/A} \otimes B_g \oplus
B_g (g^{-2}\text{d}f + \text{d}x).
\]

This works because \(\text{d}(fx - 1) = x\text{d}f + f \text{d}x =
g(g^{-2}\text{d}f + \text{d}x)\) in \(\Omega_{P[x]/A} \otimes B_g\). The second summand is contractible: multiplication by \(g^{-1}\) from degree zero to degree one is a homotopy, since its two composites with the differential \(g\) are identity. Consequently inclusion and projection on the first summand are homotopy inverses. This proves the claimed homotopy equivalence. ∎

\phantomsection\label{AG-RG-S05-algebra-lemma-product-local}

\subsubsection{Lemma F.73: product local}\label{AG-RG-S05-lemma-f.73-product-local}

\emph{Adapted from the Stacks project, algebra.tex, lines 13035--13062.}
\nopagebreak[4]

A ring with finitely many maximal ideals and locally nilpotent Jacobson radical is the product of its localizations at those ideals; every prime is maximal.

\textbf{Proof.} Write \(I=\bigcap_i\mathfrak m_i\). Every element of \(I\) is nilpotent, so every prime contains \(I\), hence the product of the finitely many \(\mathfrak m_i\), and therefore one maximal ideal. Thus these are all the primes. The Chinese remainder calculation gives \(R/I=\prod_iR/\mathfrak m_i\): for two comaximal ideals choose \(a+b=1\) with \(a\) in the first and \(b\) in the second and use \(ub+va\) to lift a prescribed pair; induction gives the finite version. Idempotents lift uniquely across a nil ideal by AG-CA-01, Lemma 5.1. Lift the coordinate idempotents to \(e_i\). For \(i\ne j\), the product \(e_ie_j\) is an idempotent in the nil ideal, hence zero: an idempotent \(e=e^2\) satisfies \(e=e^N\) for every positive \(N\), so nilpotence makes it zero. The lifted idempotents are therefore pairwise orthogonal, and their sum is idempotent. Hence \(1-\sum_i e_i\) is also an idempotent; its residue is zero, so it lies in the nil ideal and is zero by the same argument. Consequently \(R\to\prod_i e_iR\), \(r\mapsto(e_ir)\), has inverse \((r_i)\mapsto\sum r_i\). The factor \(e_iR\) has just the prime corresponding to \(\mathfrak m_i\), so is local and is the localization \(R_{\mathfrak m_i}\) by the fraction universal property. This proves the claimed product. ∎

\phantomsection\label{AG-RG-S05-algebra-lemma-regular-quasi-regular}

\subsubsection{Lemma F.74: regular quasi regular}\label{AG-RG-S05-lemma-f.74-regular-quasi-regular}

\emph{Adapted from the Stacks project, algebra.tex, lines 17680--17766.}
\nopagebreak[4]

If \(f_1,\ldots,f_c\) is a regular sequence in \(A\), the conormal module of \(I=(f_1,\ldots,f_c)\) is free over \(A/I\) on the classes of the \(f_i\). Only this consequence of quasi-regularity is used here.

\textbf{Proof.} First prove that every relation \(\sum a_if_i=0\) is generated by the pairwise Koszul relations. For \(c=1\), injectivity of \(f_1\) gives \(a_1=0\). For \(c>1\), reduction modulo \((f_1,\ldots,f_{c-1})\) and injectivity of \(f_c\) give \(a_c=\sum_{i<c}b_if_i\). Subtract from the relation the sum of the relations \(b_i(f_ie_c-f_ce_i)\), making its last coefficient zero. Apply induction to the remaining coefficients. In particular all coefficients of every relation lie in \(I\). Now if \(\sum a_if_i\in I^2\), subtract a displayed finite sum of products \(f_if_j\), changing each \(a_i\) by an element of \(I\), to obtain a relation equal to zero. The preceding calculation puts its coefficients in \(I\), hence puts the original \(a_i\) in \(I\). This proves injectivity of \((A/I)^c\to I/I^2\); generation proves surjectivity. ∎

\phantomsection\label{AG-RG-S05-algebra-lemma-regular-sequence-in-polynomial-ring}

\subsubsection{Lemma F.75: regular sequence in polynomial ring}\label{AG-RG-S05-lemma-f.75-regular-sequence-in-polynomial-ring}

\emph{Adapted from the Stacks project, algebra.tex, lines 17576--17624.}
\nopagebreak[4]

Let \(R\) be a ring. Let \(f_1, \ldots, f_r \in R\) which do not generate the unit ideal. The following are equivalent:

\begin{enumerate}
\def\labelenumi{\arabic{enumi}.}
\item
  any permutation of \(f_1, \ldots, f_r\) is a regular sequence,
\item
  any subsequence of \(f_1, \ldots, f_r\) (in the given order) is a regular sequence, and
\item
  \(f_1x_1, \ldots, f_rx_r\) is a regular sequence in the polynomial ring \(R[x_1, \ldots, x_r]\).
\end{enumerate}

\textbf{Proof.} It is clear that (1) implies (2). We prove (2) implies (1) by induction on \(r\). The case \(r = 1\) is trivial. The case \(r = 2\) says that if \(a, b \in R\) are a regular sequence and \(b\) is a nonzerodivisor, then \(b, a\) is a regular sequence. The kernel isomorphism has the following explicit form. If \(ax=by\), send the class of \(x\) modulo \((b)\) to the class of \(y\) modulo \((a)\). The element \(y\) is unique because \(b\) is a nonzerodivisor. Replacing \(x\) by \(x+bz\) replaces \(y\) by \(y+az\), so the map is well-defined; its image is in the kernel of multiplication by \(b\) modulo \((a)\). Conversely, from \(by=ax\) recover the unique \(x\) because \(a\) is a nonzerodivisor. These constructions are inverse and linear. Thus vanishing of the latter kernel gives vanishing of the former. The case \(r > 2\). Assume (2) holds and say we want to prove \(f_{\sigma(1)}, \ldots, f_{\sigma(r)}\) is a regular sequence for some permutation \(\sigma\). We already know that \(f_{\sigma(1)}, \ldots, f_{\sigma(r - 1)}\) is a regular sequence by induction. Hence it suffices to show that \(f_s\) where \(s = \sigma(r)\) is a nonzerodivisor modulo \(f_1, \ldots, \hat f_s, \ldots, f_r\). If \(s = r\) we are done. If \(s < r\), then note that \(f_s\) and \(f_r\) are both nonzerodivisors in the ring \(R/(f_1, \ldots, \hat f_s, \ldots, f_{r - 1})\) (by induction hypothesis again). Since we know \(f_s, f_r\) is a regular sequence in that ring we conclude by the case of sequence of length \(2\) that \(f_r, f_s\) is too.

Note that \(R[x_1, \ldots, x_r]/(f_1x_1, \ldots, f_ix_i)\) as an \(R\)-module is a direct sum of the modules

\[
R/I_E \cdot x_1^{e_1} \ldots x_r^{e_r}
\]

indexed by multi-indices \(E = (e_1, \ldots, e_r)\) where \(I_E\) is the ideal generated by \(f_j\) for \(1 \leq j \leq i\) with \(e_j > 0\). Hence \(f_{i + 1}x_{i + 1}\) is a nonzerodivisor on this if and only if \(f_{i + 1}\) is a nonzerodivisor on \(R/I_E\) for all \(E\). For any subset \(U\subset\{1,\ldots,i\}\), take \(E\) with \(e_j=1\) for \(j\in U\) and \(e_j=0\) otherwise. Then \(I_E=(f_j\mid j\in U)\), and (3) makes \(f_{i+1}\) a nonzerodivisor on this quotient. Apply this successively to the earlier members of any chosen subsequence; it proves every required multiplication injective. Its ideal is proper because it is contained in the original proper ideal. Thus (3) implies (2). Conversely, if (2) holds, then any subsequence of \(f_1, \ldots, f_i, f_{i + 1}\) is a regular sequence in particular \(f_{i + 1}\) is a nonzerodivisor on all \(R/I_E\). In this way we see that (2) implies (3). ∎

\phantomsection\label{AG-RG-S05-algebra-lemma-relative-global-complete-intersection}

\subsubsection{Lemma F.76: relative global complete intersection}\label{AG-RG-S05-lemma-f.76-relative-global-complete-intersection}

\emph{Adapted from the Stacks project, algebra.tex, lines 38280--38292.}
\nopagebreak[4]

A relative global complete intersection over a Noetherian base is syntomic.

\textbf{Proof.} It is finitely presented by its displayed finite equation list. Its fibres are global complete intersections by the dimension definition. Its flatness is the quotient-flatness conclusion of \hyperref[AG-RG-S05-algebra-lemma-relative-global-complete-intersection-conormal]{F.77}, proved immediately below using \hyperref[AG-RG-S05-algebra-lemma-grothendieck-regular-sequence]{F.53}. These are exactly the three syntomic conditions. ∎

\phantomsection\label{AG-RG-S05-algebra-lemma-relative-global-complete-intersection-conormal}

\subsubsection{Lemma F.77: relative global complete intersection conormal}\label{AG-RG-S05-lemma-f.77-relative-global-complete-intersection-conormal}

\emph{Adapted from the Stacks project, algebra.tex, lines 38204--38278.}
\nopagebreak[4]

Let \(R\) be Noetherian and let \(S=R[X_1,\ldots,X_n]/(f_1,\ldots,f_c)\) have all its nonempty fibres of dimension \(n-c\). At every prime above the equation ideal, the \(f_i\) form a regular sequence in the corresponding polynomial local ring, its successive quotients are \(R\)-flat, and \((f)/(f)^2\) is free over \(S\) on their classes. This is the Noetherian form consumed in desingularization.

\textbf{Proof.} Fix a base prime \(\mathfrak p\) and put \(k=\kappa(\mathfrak p)\). In \(k[X]\) every minimal prime of the equation ideal has height at most \(c\), by AG-CA-11, Theorem 3.1. Its component has dimension \(n\) minus that height, by AG-CA-09, Theorem 4.3. The asserted fibre dimension \(n-c\) forces every such height to equal \(c\). At a prime \(\mathfrak r\) containing the equations, \(A=k[X]_{\mathfrak r}\) is regular by iterating \href{https://kokunoyumeto.github.io/open-math-courses/courses/AG-CA/regular-local-rings.html}{AG-CA-14, Proposition 3.3} from the field, hence CM by AG-CA-12, Theorem 6.1. For a minimal equation prime \(\mathfrak a\subset\mathfrak r\), prime-localization correspondence preserves all primes below \(\mathfrak a\), so \(\dim A_{\mathfrak aA}=\operatorname{height}(\mathfrak a)=c\). The CM dimension formula AG-CA-12, Lemma 5.4, applied in \(A\) at \(\mathfrak aA\) therefore gives \(\dim A/\mathfrak aA=\dim A-c\) for this component of the quotient. Taking the maximum gives that same quotient dimension. The expected-dimension criterion AG-CA-12, Theorem 4.2, makes the \(c\) equations a regular sequence in this fibre local ring.

The polynomial local ring over \(R_{\mathfrak p}\) is flat, since polynomial extension is free and localization is flat. Apply \hyperref[AG-RG-S05-algebra-lemma-grothendieck-regular-sequence]{F.53}, whose full Noetherian lifting proof uses AG-CA-08, Lemma 5.1, to obtain regularity before reduction and flatness of every successive quotient. The elementary regular-sequence relation calculation proved at \hyperref[AG-RG-S05-algebra-lemma-regular-quasi-regular]{F.74} makes its conormal module free on the classes of the \(f_i\) at every prime of \(S\). The global map \(S^c\to(f)/(f)^2\) is surjective and has zero kernel after every prime localization, so AG-CA-02, Theorem 5.1, makes it an isomorphism. Flatness of the quotient is likewise detected locally by the ideal injection criterion and local detection. ∎

\phantomsection\label{AG-RG-S05-algebra-lemma-relative-global-complete-intersection-smooth}

\subsubsection{Lemma F.78: relative global complete intersection smooth}\label{AG-RG-S05-lemma-f.78-relative-global-complete-intersection-smooth}

\emph{Adapted from the Stacks project, algebra.tex, lines 39071--39113.}
\nopagebreak[4]

For the preceding Noetherian presentation, smoothness at a prime is equivalent to one \(c\times c\) Jacobian minor being outside that prime.

\textbf{Proof.} The preceding conormal calculation identifies the two-term complex with \(S^c\xrightarrow{J}S^n\). If a chosen maximal minor is invertible, the inverse of that square block gives a left inverse to \(J\); elimination identifies its cokernel with a free module of rank \(n-c\). The split conormal criterion AG-CA-17, Theorem 3.1, gives smoothness on that localization. Conversely smoothness supplies a left inverse after localization by that same criterion. Reducing at the prime preserves the left inverse. Thus the matrix over the residue field has rank \(c\), and some \(c\)-minor is nonzero. Principal localization of the conormal complex is the explicit contractible-summand computation proved at \hyperref[AG-RG-S05-algebra-lemma-principal-localization-NL]{F.72}. ∎

\phantomsection\label{AG-RG-S05-algebra-lemma-section-smooth}

\subsubsection{Lemma F.79: section smooth}\label{AG-RG-S05-lemma-f.79-section-smooth}

\emph{Adapted from the Stacks project, algebra.tex, lines 39859--39914.}
\nopagebreak[4]

For a smooth ring map \(R\to S\) with section \(\sigma:S\to R\), put \(I=\ker\sigma\). Then \(I/I^2\) is finite projective over \(R\). If it is free of rank \(d\), the \(I\)-adic completion is \(R[[T_1,\ldots,T_d]]\).

\textbf{Proof.} The map \(s\mapsto s-\sigma(s)\bmod I^2\) is an \(R\)-derivation to \(I/I^2\) and sends each \(i\in I\) to its class. Universal differentials give the inverse to the natural map \(I/I^2\to\Omega_{S/R}\otimes_SR\): both composites agree on the algebra generators and on the ideal classes. Smoothness makes \(\Omega_{S/R}\) finite projective by the split conormal criterion AG-CA-17, Theorem 3.1, so its base change is finite projective. The ideal \(I\) is finitely generated: for finite algebra generators \(x_j\) of \(S/R\), subtracting their section values shows that \(x_j-\sigma(x_j)\) generate the section kernel.

Choose \(t_i\in I\) representing a free conormal basis. Define a map \(b:R[[T]]\to\widehat S\) by convergent evaluation \(T_i\mapsto t_i\). It is surjective: start with the residue coefficient in \(R\), then successively correct the error in \(I^n/I^{n+1}\) by a homogeneous degree-\(n\) polynomial in the \(t_i\). These monomials generate that quotient because the \(t_i\) generate \(I/I^2\); expanding products proves this in every degree. The corrections converge in the completion.

The chosen conormal identification defines a first-order \(R\)-map \(S\to R[[T]]/(T)^2\) sending \(t_i\) to \(T_i\) and reducing to \(\sigma\). Formal smoothness of \(S/R\), proved equivalent to the split conormal property in AG-CA-17, Theorem 3.1, lifts this map successively through the square-zero kernels \((T)^n/(T)^{n+1}\). Its compatible inverse limit sends \(I\) into \((T)\), hence extends to a continuous map \(a:\widehat S\to R[[T]]\). The composite \(ab\) sends \(T_i\) to \(T_i\) plus terms of degree at least two. Such a substitution is an automorphism: recursively correct each inverse variable one homogeneous degree at a time; the linear term is identity, so the correction at each degree is uniquely its negative error, and the series converge. The same recursion proves that both composites with the resulting inverse are identity in every finite quotient and hence in the separated limit. Replace \(a\) by \((ab)^{-1}a\). Then \(ab=1\). Since \(b\) is already surjective it is also injective and gives the required isomorphism. ∎

\phantomsection\label{AG-RG-S05-algebra-lemma-silly}

\subsubsection{Lemma F.80: silly}\label{AG-RG-S05-lemma-f.80-silly}

Reading source: \href{https://stacks.math.columbia.edu/tag/00DS}{Stacks, Tag 00DS}.

Let \(J,I_1,\ldots,I_r\) be ideals of any commutative ring. Suppose \(J\not\subset I_i\) for every \(i\), and at most two of the \(I_i\) are not prime. There is an element of \(J\) outside all the \(I_i\).

\textbf{Proof.} Remove repetitions and every listed ideal contained in another listed ideal. Avoiding the retained larger ideals also avoids the removed smaller ones. The remaining ideals are pairwise incomparable, and still at most two are not prime. We prove the assertion for this remaining list by induction on its length.

For the empty list choose \(0\in J\); there are no avoidance conditions. For one ideal the hypothesis supplies an element outside it. For two ideals choose \(a\in J\setminus I_1\) and \(b\in J\setminus I_2\). If either avoids both, it suffices. Otherwise \(a\in I_2\), \(b\in I_1\). Their sum belongs to neither: membership in \(I_1\) would force \(a=(a+b)-b\in I_1\), and membership in \(I_2\) would similarly force \(b\in I_2\).

For a list of length at least three, place a prime ideal last, calling it \(\mathfrak p=I_r\). Induction gives \(a\in J\) outside all earlier ideals. If \(a\notin\mathfrak p\), it already suffices. Otherwise choose \(b_0\in J\setminus\mathfrak p\) and, for each \(i<r\), choose \(b_i\in I_i\setminus\mathfrak p\); incomparability supplies the latter choices. The product \(b=b_0b_1\cdots b_{r-1}\) belongs to \(J\) and every earlier \(I_i\). It lies outside \(\mathfrak p\), because iterating the defining implication for a prime ideal shows that a product of elements outside it stays outside. Then \(a+b\in J\) still avoids every earlier ideal by subtraction of \(b\), and avoids \(\mathfrak p\) by subtraction of \(a\). This completes the induction without a Noetherian hypothesis. ∎

\phantomsection\label{AG-RG-S05-algebra-lemma-size-extension-pth-roots}

\subsubsection{Lemma F.81: size extension pth roots}\label{AG-RG-S05-lemma-f.81-size-extension-pth-roots}

\emph{Adapted from the Stacks project, algebra.tex, lines 45787--45812.}
\nopagebreak[4]

If \(a_1,\ldots,a_n\) in a characteristic-\(p\) field have independent absolute differentials, adjoining their \(p\)th roots has degree \(p^n\).

\textbf{Proof.} First prove the absolute differential calculation used here. A maximal \(p\)-independent set is a \(p\)-basis, by the Zorn and degree-\(p\) argument in AG-CA-19S, Lemma 1.1. In its monomial expansion define the partial derivation of one basis element by differentiating the corresponding monomial and killing all \(p\)th-power coefficients. The relations \(T_b^p-b^p\) have derivative zero, so these are well-defined derivations. Their values on the basis elements show that their differentials are linearly independent; differentiating every expansion shows that they span. An expansion has all partial derivatives zero precisely when every exponent is zero, hence precisely when it lies in the subfield of \(p\)th powers. Thus \(da=0\) if and only if \(a\in k^p\).

For \(n=1\) this proves \(a_1\) is not a \(p\)th power. The monic polynomial \(T^p-a_1\) is irreducible: any proper factor over \(k\) is \((T-\alpha)^e\) in an algebraic closure, \(1\le e<p\), and its coefficient \(-e\alpha\) would put \(\alpha\in k\). Suppose by induction the first \(n-1\) roots give degree \(p^{n-1}\), with their standard monomial basis. If \(a_n\) became a \(p\)th power there, write a root in that basis and raise its expansion to \(p\). Then \(a_n=\sum_E\lambda_E^p\prod_{i<n}a_i^{e_i}\). Its differential belongs to the span of the preceding \(da_i\), a contradiction. The same degree-\(p\) irreducibility argument proves the last extension degree \(p\), and multiplication of finite vector-space degrees gives \(p^n\). ∎

\phantomsection\label{AG-RG-S05-algebra-lemma-smooth-over-field}

\subsubsection{Lemma F.82: smooth over field}\label{AG-RG-S05-lemma-f.82-smooth-over-field}

\emph{Adapted from the Stacks project, algebra.tex, lines 38622--38676.}
\nopagebreak[4]

A smooth finite-type algebra over a field is locally a complete intersection.

\textbf{Proof.} At a given prime choose a standard smooth neighbourhood by AG-CA-17, Theorem 5.1. In its polynomial ambient local ring, the classes of its chosen equations in the maximal ideal modulo its square are linearly independent. Indeed, differentiating an element of the square gives zero after reduction to the residue field; a linear relation between those equation classes would therefore give the same relation between their differential rows, contradicting the invertible Jacobian minor. Extend the independent classes to a basis of the maximal ideal modulo its square. The ambient polynomial local ring is regular by iterating \href{https://kokunoyumeto.github.io/open-math-courses/courses/AG-CA/regular-local-rings.html}{AG-CA-14, Proposition 3.3}, starting with the field; therefore AG-CA-12, Theorem 6.1, makes this full minimal generator list a regular sequence; its initial equation list is regular. The complete-intersection neighbourhood criterion just proved gives a global complete-intersection chart around the prime. Such charts cover the spectrum, proving the local statement without any algebraic-closure reduction or omitted diagram calculation. ∎

\phantomsection\label{AG-RG-S05-algebra-lemma-smooth-syntomic}

\subsubsection{Lemma F.83: smooth syntomic}\label{AG-RG-S05-lemma-f.83-smooth-syntomic}

Reading source: \href{https://stacks.math.columbia.edu/tag/00TA}{Stacks, Tag 00TA}.

A smooth map \(R\to S\) admits a cover of \(\operatorname{Spec}S\) by principal opens \(D(g)\) for which \(S_g\) is standard smooth over \(R\). When \(R\) is Noetherian, the map is also syntomic.

\textbf{Proof.} The case \(S=0\) has an empty cover. The zero algebra has the finite presentation \(R/(1)\), is flat as a zero module, and has no nonempty fibres, so it also satisfies the asserted syntomic condition. Hence consider a nonzero \(S\).

Fix a finite presentation \(P=R[x_1,\ldots,x_n]\twoheadrightarrow S\) with kernel \(I=(f_1,\ldots,f_m)\), and put \(M=I/I^2\). The split conormal theorem \href{https://kokunoyumeto.github.io/open-math-courses/courses/AG-CA/formally-smooth-unramified-and-etale-ring-maps.html}{AG-CA-17, Theorem 3.1} puts \(M\) inside \(S^n\) as a direct summand. Its complement is finite, because projections of the coordinate vectors generate it. The splitting therefore presents \(M\) finitely; it also proves flatness by restricting tensor injections for the free module. \href{https://kokunoyumeto.github.io/open-math-courses/courses/AG-RG/AG-RG-S01.html}{S01, Lemma 1.2} makes \(M\) free over each local ring.

For a subset \(E\) of the equation indices, let \(U_E\) be the points where the classes indexed by \(E\) are a basis of \(M\). These loci are open and cover the spectrum. Here are the denominator details. At a prime \(\mathfrak q\), select a basis of the residue vector space from the spanning equation classes, with \(r\) selected members. Their map \(\varphi:S^r\to M\) is surjective over \(S_{\mathfrak q}\) by Nakayama. In a local free basis its matrix is square and invertible modulo the maximal ideal, hence invertible by the determinant and adjugate identities. The finite cokernel vanishes at \(\mathfrak q\); multiplying denominators which kill its generators gives \(a\notin\mathfrak q\) with \(\varphi_a\) surjective. Since \(M_a\) is projective, its kernel is a finite summand of \(S_a^r\). It vanishes at \(\mathfrak q\) too. Multiplying numerators of denominators which kill its generators gives \(b\notin\mathfrak q\) with \(\varphi_{ab}\) an isomorphism. Thus \(D(ab)\) lies in this basis locus. The same reasoning applies at every point of any fixed \(U_E\), proving openness. Selecting a basis from a spanning vector-space list proves covering. For rank zero the selected list is empty and empty denominator products are \(1\).

Choose principal opens inside these loci. On one such open \(D(g)\), adjoining a variable representing \(g^{-1}\) gives a finite polynomial presentation whose conormal module, by \hyperref[AG-RG-S05-algebra-lemma-principal-localization-NL]{F.72}, is \(M_g\oplus S_g\). It is free. \hyperref[AG-RG-S05-algebra-lemma-huber]{F.55} then supplies a polynomial presentation of \(S_g\) whose kernel is generated by a list which is itself a conormal basis. Denote this algebra temporarily by \(T\), and write its presentation as \[
R[y_1,\ldots,y_N]\twoheadrightarrow T,
\qquad K=(a_1,\ldots,a_c),\qquad K/K^2=T^c.
\]

The conormal differential has a left inverse. In these bases let its matrix be \(A:T^c\to T^N\), and let \(B:T^N\to T^c\) satisfy \(BA=1\). At a residue field this forces \(c\le N\). The determinant expansion \hyperref[AG-RG-S05-algebra-lemma-matrix-left-inverse]{F.70(2)} shows that its \(c\)-row minors \(\Delta_E\) generate \(1\). Each selected square map is invertible precisely on \(D(\Delta_E)\), by the adjugate formula, so these principal opens cover \(\operatorname{Spec}T\). The empty minor for \(c=0\) is \(1\).

On a selected minor open, order its \(c\) pivot variables first and let \(h\in R[y]\) represent the inverted element \(u=\Delta_E\in T\). Present \(T_u\) by \(a_1,\ldots,a_c,hz-1\). Its Jacobian on \(y_1,\ldots,y_c,z\) has diagonal blocks the selected Jacobian and \(u\), and upper right block zero. Its determinant is \(u\Delta_E\), a unit. This is a standard smooth presentation. Finally a principal open in \(S_g\) is principal in \(S\): write its defining fraction as \(s/g^k\), so its open corresponds to \(D(gs)\). The charts just constructed are consequently the required principal cover of the original spectrum.

If \(R\) is Noetherian, \hyperref[AG-RG-S05-algebra-lemma-standard-smooth]{F.85} makes each chart a relative global complete intersection. \hyperref[AG-RG-S05-algebra-lemma-relative-global-complete-intersection]{F.76} makes it syntomic, and \hyperref[AG-RG-S05-algebra-lemma-local-syntomic]{F.62} glues that condition on the target. This proves the last assertion. ∎

\phantomsection\label{AG-RG-S05-algebra-lemma-snake}

\subsubsection{Lemma F.84: snake}\label{AG-RG-S05-lemma-f.84-snake}

Reading source: \href{https://stacks.math.columbia.edu/tag/07JW}{Stacks, Tag 07JW}.

Suppose the following diagram of abelian groups commutes and both rows are exact: \[
\begin{array}{ccccccccc}
&&X&\xrightarrow{a}&Y&\xrightarrow{b}&Z&\longrightarrow&0\\
&&\alpha\downarrow&&\beta\downarrow&&\gamma\downarrow\\
0&\longrightarrow&U&\xrightarrow{i}&V&\xrightarrow{j}&W.
\end{array}
\] There is a canonical exact sequence \[
\ker\alpha\xrightarrow{a}\ker\beta\xrightarrow{b}\ker\gamma
\xrightarrow{\partial}\operatorname{coker}\alpha
\xrightarrow{i}\operatorname{coker}\beta
\xrightarrow{j}\operatorname{coker}\gamma.
\] Its first map is injective if \(a\) is injective, and its last map is surjective if \(j\) is surjective.

\textbf{Proof.} For \(z\in\ker\gamma\), choose \(y\) with \(b(y)=z\). Commutativity gives \(j\beta(y)=0\), so exactness supplies \(u\) with \(i(u)=\beta(y)\). The injection \(i\) makes this \(u\) unique for the chosen \(y\). Define \(\partial z=[u]\) modulo \(\alpha(X)\). Replacing \(y\) by another lift changes it by \(a(x)\); commutativity and injectivity of \(i\) then change \(u\) by \(\alpha(x)\). Thus the class is independent of the lift. Adding lifts proves additivity. The other arrows are the restrictions or quotient maps induced by \(a,b,i,j\), which are well defined by the commutative squares.

For exactness at \(\ker\beta\), take \(y\) there with \(b(y)=0\). Write \(y=a(x)\). Since \(i\alpha(x)=\beta a(x)=0\) and \(i\) is injective, \(x\in\ker\alpha\). Conversely the image of such an \(x\) lies in that kernel and has zero \(b\)-image.

At \(\ker\gamma\), the condition \(\partial z=0\) says that the chosen \(u\) is \(\alpha(x)\) for some \(x\). Subtract \(a(x)\) from the chosen \(y\). This still maps to \(z\), and its \(\beta\)-image is zero. Conversely a lift already in \(\ker\beta\) allows \(u=0\), so has zero connecting class.

At \(\operatorname{coker}\alpha\), a class \([u]\) dies in \(\operatorname{coker}\beta\) exactly when \(i(u)=\beta(y)\) for some \(y\). Then \(z=b(y)\) satisfies \(\gamma(z)=j\beta(y)=ji(u)=0\), and the defining lift gives \(\partial z=[u]\). Every connecting class has zero image for this same reason.

At \(\operatorname{coker}\beta\), a class \([v]\) dies in \(\operatorname{coker}\gamma\) exactly when \(j(v)=\gamma(z)\) for some \(z\). Lift \(z\) to \(y\in Y\). The element \(v-\beta(y)\) has zero \(j\)-image, so is \(i(u)\) by bottom-row exactness. Hence \([v]\) comes from \([u]\). Conversely \(ji=0\) kills every such image. These checks prove exactness at every intermediate term. Restriction of an injective \(a\) remains injective. If \(j\) is onto, lift any representative in \(W\) to \(V\) to obtain surjectivity on the last cokernel. ∎

\phantomsection\label{AG-RG-S05-algebra-lemma-standard-smooth}

\subsubsection{Lemma F.85: standard smooth}\label{AG-RG-S05-lemma-f.85-standard-smooth}

Reading source: \href{https://stacks.math.columbia.edu/tag/00T7}{Stacks, Tag 00T7}.

Let \(P=R[x_1,\ldots,x_n]\), \(I=(f_1,\ldots,f_c)\), and \(S=P/I\) be standard smooth, with the first \(c\) variables chosen as pivots. Then:

\begin{enumerate}
\def\labelenumi{\arabic{enumi}.}
\tightlist
\item
  \(R\to S\) is smooth.
\item
  \(\Omega_{S/R}\) has the free basis \(dx_{c+1},\ldots,dx_n\).
\item
  \(I/I^2\) has the free basis \([f_1],\ldots,[f_c]\).
\item
  Every principal localization \(S_g\) is standard smooth over \(R\).
\item
  For every base map \(R\to R'\), the algebra \(R'\otimes_R S\) is standard smooth over \(R'\).
\item
  If \(r\in R\) becomes a unit in \(S\), then \(S\) is standard smooth over \(R_r\).
\item
  This presentation makes \(S\) a relative global complete intersection over \(R\).
\end{enumerate}

\textbf{Proof.} Let \(C=(\partial f_i/\partial x_j)_{1\le i,j\le c}\); its determinant is a unit in \(S\). The map \(S^c\to I/I^2\) taking the standard vectors to the equation classes is onto. Its composite with the conormal differential and projection onto the first \(c\) differential coordinates is an invertible matrix (the transpose of \(C\) under column-vector conventions). Thus this surjection is injective as well. The equation classes are a basis, and the conormal differential is split injective.

Write its image in \(S^n\) as the columns of \(\binom{C^{\mathsf t}}{D}\). An arbitrary vector \((v,w)\) differs from a unique conormal vector by \[
(v,w)-\binom{C^{\mathsf t}}{D}(C^{\mathsf t})^{-1}v
=(0,w-D(C^{\mathsf t})^{-1}v).
\] The conormal image meets the last coordinate summand only in zero. Its cokernel is therefore exactly that summand, with the asserted differential basis. Finite presentation and this split conormal sequence give smoothness by the complete criterion \href{https://kokunoyumeto.github.io/open-math-courses/courses/AG-CA/formally-smooth-unramified-and-etale-ring-maps.html}{AG-CA-17, Theorem 3.1 and Section 4}. This proves (1)--(3), including \(c=0\).

To invert \(g\), choose a polynomial \(h\) representing it and use the equations \(f_1,\ldots,f_c,hz-1\). On pivot variables \(x_1,\ldots,x_c,z\), the Jacobian has blocks \[
\begin{pmatrix}C&0\\z(\partial h/\partial x_1,\ldots,\partial h/\partial x_c)&h\end{pmatrix}.
\] Its determinant \(h\det C\) is a unit in \(S_g\). This proves (4), even when the localization is zero. Under a base map, applying the map to every coefficient presents the tensor product; the same operation takes the determinant to its image, which remains invertible. This proves (5). For (6), the original equations over \(R_r\) present \(S_r=S\), since the image of \(r\) is already invertible; their determinant is unchanged.

For (7), it remains to compute every nonempty fibre's dimension. By (5) such a fibre is a nonzero standard smooth algebra \(B=k[x_1,\ldots,x_n]/(\bar f_1,\ldots,\bar f_c)\) over its residue field \(k\). It is locally a complete intersection by \hyperref[AG-RG-S05-algebra-lemma-smooth-over-field]{F.82}. At any prime \(\mathfrak q\) of \(B\), let \(\mathfrak q'\) be its polynomial preimage. The conormal module is free of rank \(c\) by (3); reducing to the residue field determines that rank uniquely. \hyperref[AG-RG-S05-algebra-lemma-lci]{F.57} identifies it with \(\operatorname{height}(\mathfrak q')-\operatorname{height}(\mathfrak q)\). For a minimal prime \(\mathfrak q\) this gives \(\operatorname{height}(\mathfrak q')=c\). The polynomial height formula \href{https://kokunoyumeto.github.io/open-math-courses/courses/AG-CA/krull-dimension-and-noether-normalization.html}{AG-CA-09, Theorem 4.3} therefore gives dimension \(n-c\) for every component.

The algebra \(B\) is Noetherian by \href{https://kokunoyumeto.github.io/open-math-courses/courses/AG-RG/AG-RG-S01.html}{S01, Lemma 1.3}, so \hyperref[AG-RG-S05-topology-lemma-Noetherian]{F.136} and \hyperref[AG-RG-S05-topology-lemma-irreducible]{F.137} give a finite component cover. Dimension equals the largest component dimension: a finite chain of irreducible closed subsets lies in a component containing its largest member, and every chain inside a component is also a chain in the whole space. Thus \(\dim B=n-c\). This argument works over imperfect fields and for an empty equation list. Empty fibres require no dimension assertion, completing (7). ∎

\phantomsection\label{AG-RG-S05-algebra-lemma-sum-two-terms}

\subsubsection{Lemma F.86: sum two terms}\label{AG-RG-S05-lemma-f.86-sum-two-terms}

Reading source: \href{https://stacks.math.columbia.edu/tag/00S3}{Stacks, Tag 00S3}.

Let \(A_1\xrightarrow{d_A}A_0\) and \(B_1\xrightarrow{d_B}B_0\) be two-term complexes of modules over a ring. If they are homotopy equivalent, then \[
A_1\oplus B_0\cong B_1\oplus A_0.
\]

\textbf{Proof.} Choose chain maps \(\varphi:A_\bullet\to B_\bullet\) and \(\psi:B_\bullet\to A_\bullet\) and homotopies \(h:A_0\to A_1\), \(h':B_0\to B_1\) with \[
1-\psi_1\varphi_1=hd_A,\qquad 1-\psi_0\varphi_0=d_Ah,
\] \[
1-\varphi_1\psi_1=h'd_B,\qquad 1-\varphi_0\psi_0=d_Bh'.
\] Define \(U:B_1\oplus A_0\to A_1\oplus B_0\) and \(V\) in the reverse direction by \[
U=\begin{pmatrix}\psi_1&h\\-d_B&\varphi_0\end{pmatrix},\qquad
V=\begin{pmatrix}\varphi_1&-h'\\d_A&\psi_0\end{pmatrix}.
\] Multiplying the blocks and using the four homotopy identities and the two chain-map identities gives \[
UV=\begin{pmatrix}1&h\psi_0-\psi_1h'\\0&1\end{pmatrix},\qquad
VU=\begin{pmatrix}1&\varphi_1h-h'\varphi_0\\0&1\end{pmatrix}.
\] Each off-diagonal matrix squares to zero, so negating that off-diagonal block gives the inverse product. Consequently \(L=(VU)^{-1}V\) is a left inverse of \(U\), and \(T=V(UV)^{-1}\) is a right inverse. The equality \(L=L(UT)=(LU)T=T\) makes them the same inverse. Hence \(U\) is an isomorphism, proving the claim without any finite-generation or freeness hypothesis. ∎

\phantomsection\label{AG-RG-S05-algebra-lemma-syntomic}

\subsubsection{Lemma F.87: syntomic}\label{AG-RG-S05-lemma-f.87-syntomic}

Reading source: \href{https://stacks.math.columbia.edu/tag/00SY}{Stacks, Tag 00SY}.

Let \(R\) be Noetherian, and let \(\mathfrak q\in\operatorname{Spec}S\) lie over \(\mathfrak p\in\operatorname{Spec}R\). The following neighbourhood conditions are equivalent:

\begin{enumerate}
\def\labelenumi{\arabic{enumi}.}
\tightlist
\item
  Some \(g\notin\mathfrak q\) makes \(R\to S_g\) syntomic.
\item
  Some \(g\notin\mathfrak q\) makes \(S_g\) a relative global complete intersection over \(R\).
\item
  Some \(g\notin\mathfrak q\) makes \(S_g\) finitely presented over \(R\), while \(R_{\mathfrak p}\to S_{\mathfrak q}\) is flat and \(S_{\mathfrak q}/\mathfrak pS_{\mathfrak q}\) is a complete-intersection local ring over \(k=\kappa(\mathfrak p)\) in the sense of \hyperref[AG-RG-S05-algebra-definition-lci-local-ring]{F.3}.
\end{enumerate}

\textbf{Proof.} A syntomic neighbourhood is finitely presented and flat and has locally complete-intersection fibres; \hyperref[AG-RG-S05-algebra-lemma-lci-at-prime]{F.58} gives (3) from (1). \hyperref[AG-RG-S05-algebra-lemma-relative-global-complete-intersection]{F.76} gives (1) from (2). We prove the remaining implication with all rings localized explicitly.

Under (3), first replace \(S\) by the indicated finitely presented neighbourhood, retaining the prime. Write \(P=R[x_1,\ldots,x_n]\twoheadrightarrow S\) with kernel \(I\) and polynomial preimage \(\mathfrak q'\). In the fibre \(\bar P=k[x_1,\ldots,x_n]\), write \(\bar I\) for the image ideal and \(\bar{\mathfrak q}'\) for the induced prime. By \hyperref[AG-RG-S05-algebra-lemma-lci-at-prime]{F.58} and \hyperref[AG-RG-S05-algebra-lemma-lci]{F.57}, the ideal \(\bar I_{\bar{\mathfrak q}'}\) has a regular generating list of length \(c\), its minimal number of generators. Choose \(f_1,\ldots,f_c\in I\) whose images give a basis of \(\bar I_{\bar{\mathfrak q}'}/\bar{\mathfrak q}'\bar I_{\bar{\mathfrak q}'}\). This is possible by selecting a basis from the images of a finite generating list of \(I\). Nakayama gives generation in the fibre local ring, and the last assertion of F.57 makes this very list regular there.

Set \(T=P/(f_1,\ldots,f_c)\) and \(J=\ker(T\to S)\). The polynomial ring is Noetherian by \href{https://kokunoyumeto.github.io/open-math-courses/courses/AG-RG/AG-RG-S01.html}{S01, Lemma 1.3}, so \(J\) is finite. Localize the short exact sequence at \(\mathfrak q'\) as \(P\)-modules: \[
0\longrightarrow J_{\mathfrak q'}\longrightarrow T_{\mathfrak q'}\longrightarrow S_{\mathfrak q}\longrightarrow0.
\] These are \(R_{\mathfrak p}\)-modules. Flatness of the last module makes tensoring with \(k\) injective on the first arrow, by the tensor-kernel calculation \href{https://kokunoyumeto.github.io/open-math-courses/courses/AG-RG/AG-RG-S01.html}{S01, Lemma 1.2 and Lemma 1.A}. The last arrow becomes an isomorphism after tensoring, because the chosen fibre equations generate the localized fibre ideal. Hence \[
J_{\mathfrak q'}/\mathfrak pJ_{\mathfrak q'}=0.
\] Over the local polynomial ring \(P_{\mathfrak q'}\), the ideal \(\mathfrak pP_{\mathfrak q'}\) is contained in its maximal ideal. Applying \hyperref[AG-RG-S05-algebra-lemma-NAK]{F.10} to the finite module \(J_{\mathfrak q'}\) gives zero. Clear annihilating denominators for a finite generating list of \(J\): some \(u\in P\setminus\mathfrak q'\) gives \(J_u=0\), and thus \(T_u=S_u\).

The fibre of \(T\) at the induced prime has the same local quotient by the chosen regular list. F.57 gives its point dimension \(n-c\). Applying \hyperref[AG-RG-S05-algebra-lemma-localize-relative-complete-intersection]{F.66(3)} to this \(c\)-equation presentation produces a relative global complete-intersection principal neighbourhood of that point. Apply this to \(T_u\), using its presentation with the additional equation \(uz-1\); the fibre point dimension remains \(n-c\), and the variable and equation counts both increase by one. The resulting neighbourhood is therefore a principal localization of \(S_u\) with a relative global complete-intersection presentation. Writing its inverted element as a fraction in \(S_u\) makes the composite neighbourhood \(S_g\) for some \(g\notin\mathfrak q\). This is (2). Empty regular lists use the same argument with \(c=0\). ∎

\phantomsection\label{AG-RG-S05-algebra-lemma-syntomic-presentation-ideal-mod-squares}

\subsubsection{Lemma F.88: syntomic presentation ideal mod squares}\label{AG-RG-S05-lemma-f.88-syntomic-presentation-ideal-mod-squares}

Reading source: \href{https://stacks.math.columbia.edu/tag/07BT}{Stacks, Tag 07BT}.

Let \(R\) be Noetherian, let \(S=R[x_1,\ldots,x_n]/I\) with \(I\) finitely generated, and suppose \(S_g\) is syntomic over \(R\). Then \((I/I^2)_g\) is finite projective over \(S_g\).

\textbf{Proof.} Put \(A=S_g\) and \(M=(I/I^2)_g\). If \(A=0\), then \(M=0\) and the assertion follows. Otherwise \(A\) is Noetherian by \href{https://kokunoyumeto.github.io/open-math-courses/courses/AG-RG/AG-RG-S01.html}{S01, Lemma 1.3} and localization, so the finite module \(M\) is finitely presented.

At every prime of \(A\), \hyperref[AG-RG-S05-algebra-lemma-syntomic]{F.87} supplies a principal relative global complete-intersection neighbourhood. \hyperref[AG-RG-S05-algebra-lemma-relative-global-complete-intersection-conormal]{F.77} makes the conormal of its displayed equations finite free. The presentation comparison \hyperref[AG-RG-S05-algebra-lemma-conormal-module-localize]{F.33} identifies \(M\) there, up to finite free summands, with that conormal. Therefore \(M\) is finite projective there: the comparison realizes it as a summand of a finite free module. Choose finitely many of these neighbourhoods \(D(a_i)\) covering \(\operatorname{Spec}A\). The finite choice follows from the principal-open cover criterion: the ideal of all covering elements is not contained in any prime, hence contains \(1\), which is a finite sum.

Here is the global splitting argument. Choose a finite free surjection \(\pi:A^N\to M\). Its localization on \(D(a_i)\) has a right inverse \(s_i\), because \(M_{a_i}\) is projective. A finite presentation \(A^r\to A^s\to M\to0\) realizes \(\operatorname{Hom}_A(M,A^N)\) as the kernel of a map between finite matrix modules. Localization is exact and commutes with these finite modules, so \[
\operatorname{Hom}_A(M,A^N)_{a_i}
=\operatorname{Hom}_{A_{a_i}}(M_{a_i},A_{a_i}^N).
\] Consequently some global \(t_i:M\to A^N\) represents \(a_i^{k_i}s_i\). The map \(\pi t_i-a_i^{k_i}1_M\) becomes zero after inverting \(a_i\). Multiplying by another power of \(a_i\) kills its values on all members of a finite generating list of \(M\). Replacing \(t_i\) by that multiple gives \[
\pi t_i=a_i^{e_i}1_M.
\] The powers \(a_i^{e_i}\) still generate \(1\): if a proper ideal contained them, a containing prime would contain all the corresponding \(a_i\), contradicting the cover. Choose \(b_i\) with \(\sum_i b_i a_i^{e_i}=1\). Then \(\sum_i b_i t_i\) is a global right inverse of \(\pi\). Thus \(M\) is a direct summand of \(A^N\), which proves finite projectivity. All denominator and zero-module cases are included. ∎

\phantomsection\label{AG-RG-S05-algebra-lemma-total-ring-fractions-no-embedded-points}

\subsubsection{Lemma F.89: total ring fractions no embedded points}\label{AG-RG-S05-lemma-f.89-total-ring-fractions-no-embedded-points}

Reading source: \href{https://stacks.math.columbia.edu/tag/02LX}{Stacks, Tag 02LX}.

Suppose a ring \(R\) has finitely many minimal primes \(\mathfrak q_1,\ldots,\mathfrak q_t\), and its zerodivisors are exactly their union. Its total ring of fractions has the canonical decomposition \[
Q(R)\cong\prod_{i=1}^t R_{\mathfrak q_i}.
\]

\textbf{Proof.} If \(t=0\), the assumed zerodivisor set is empty. A nonzero ring would make \(0\) a zerodivisor, since it annihilates \(1\ne0\), so \(R=0\). Both its localization and the empty ring product are zero. Now assume \(t>0\). Let \(D\) be the multiplicative set of nonzerodivisors and put \(A=D^{-1}R\). Every element of \(D\) avoids every \(\mathfrak q_i\), so each map \(R\to R_{\mathfrak q_i}\) extends to \(A\). The prime correspondence for localization identifies the primes of \(A\) with primes of \(R\) disjoint from \(D\). Such a prime is contained in \(\bigcup_i\mathfrak q_i\). By \hyperref[AG-RG-S05-algebra-lemma-silly]{F.80} it is contained in one \(\mathfrak q_i\), and minimality forces equality. Thus the primes of \(A\) are precisely the extensions \(\mathfrak m_i=D^{-1}\mathfrak q_i\). They are distinct and incomparable, hence all maximal.

We construct the product decomposition, including possible nilpotents. For each \(i\) and \(j\ne i\), choose an element in \(\mathfrak m_j\setminus\mathfrak m_i\), and multiply these choices to obtain \(a_i\). Then \(a_i\) avoids \(\mathfrak m_i\) and belongs to every other prime. For \(t=1\) use the empty product \(a_1=1\). Each product \(a_i a_j\) for \(i\ne j\) belongs to every prime, and hence is nilpotent by \href{https://kokunoyumeto.github.io/open-math-courses/courses/AG-CA/spectra-of-rings.html}{AG-CA-01, Theorem 1.2}. Choose a common positive integer \(N\) with \(a_i^N a_j^N=0\) for every pair. The elements \(a_i^N\) generate the unit ideal, since no prime contains all of them. Write \(1=\sum_i b_i a_i^N\) and set \(e_i=b_i a_i^N\). The \(e_i\) are orthogonal and sum to \(1\); multiplying this sum by \(e_i\) gives \(e_i^2=e_i\). Modulo \(\mathfrak m_i\), only the \(i\)th summand survives, so \(e_i\) is \(1\) there and is zero at every other prime.

The maps \(a\mapsto(e_i a)_i\) and \((u_i)_i\mapsto\sum_i u_i\) are inverse ring maps between \(A\) and \(\prod_i e_iA\). Each factor has only the prime induced by \(\mathfrak m_i\), and is therefore local: a nonunit lies in a maximal ideal, which must be that prime. The projection \(A\to e_iA\) consequently inverts every element outside \(\mathfrak m_i\). Conversely, in \(A_{\mathfrak m_i}\) the idempotent \(e_i\) is a unit and hence equals \(1\), so the complementary factor vanishes. These universal properties give \(e_iA\cong A_{\mathfrak m_i}\cong R_{\mathfrak q_i}\). The isomorphism is the canonical map stated above. ∎

\phantomsection\label{AG-RG-S05-algebra-lemma-when-colimit}

\subsubsection{Lemma F.90: when colimit}\label{AG-RG-S05-lemma-f.90-when-colimit}

Reading source: \href{https://stacks.math.columbia.edu/tag/07C3}{Stacks, Tag 07C3}.

For an arbitrary ring map \(R\to\Lambda\), the following are equivalent: \(\Lambda\) is a filtered colimit of smooth \(R\)-algebras; every \(R\)-algebra map from a finitely presented algebra \(A\) to \(\Lambda\) factors through a smooth \(R\)-algebra.

\textbf{Proof.} First suppose such a filtered presentation of \(\Lambda\) is given. Write \(A\) with finitely many polynomial generators and finitely many defining equations. The images of its generators have representatives at a common stage. Each defining equation becomes zero at a later stage, by the equality criterion in a filtered colimit. A common successor of those finitely many stages therefore satisfies every equation. The generator images define the required map to that smooth stage.

For the converse, form the category whose objects are smooth \(R\)-algebras \(B\) with a map \(B\to\Lambda\), and whose arrows commute with those maps. Smooth algebras are finitely presented. This category can be taken small: finite polynomial presentations form a set, and a map from each one to \(\Lambda\) is specified by a finite tuple of elements of \(\Lambda\); use these presentations as representatives. The object \(R\to\Lambda\) makes it nonempty.

Given two objects \(B_1,B_2\), the tensor product \(B_1\otimes_R B_2\) is finitely presented: combine their finite generator lists and equations. Factor its induced map to \(\Lambda\) through a smooth algebra. The two structure maps into that algebra give a common successor in the category. For parallel arrows \(u,v\colon B\rightrightarrows C\), choose finite algebra generators \(b_1,\ldots,b_r\) of \(B\). The quotient \[
C/(u(b_1)-v(b_1),\ldots,u(b_r)-v(b_r))
\] is finitely presented by adding these finitely many equations to a presentation of \(C\). Its map to \(\Lambda\) factors through a smooth algebra by the hypothesis. The resulting arrow from \(C\) equalizes \(u\) and \(v\), because equality on the generators gives equality on every polynomial in them. These two constructions prove filteredness.

Let \(L\) be this category's colimit and \(L\to\Lambda\) its induced map. It is surjective: for any \(\lambda\in\Lambda\), apply the factoring hypothesis to \(R[T]\to\Lambda\) with \(T\mapsto\lambda\); the intermediate object supplies a representative of \(\lambda\) in \(L\). For injectivity, a representative \(b\in B\) with zero image in \(\Lambda\) gives a map \(B/(b)\to\Lambda\). This quotient is finitely presented, so its factorization through a smooth object yields a successor at which \(b=0\). Hence its colimit class is zero. Two representatives can be moved to a common stage and this argument applied to their difference, proving equality detection as well. Thus \(L\to\Lambda\) is an isomorphism of \(R\)-algebras, establishing the filtered presentation without any Noetherian hypothesis. ∎

\phantomsection\label{AG-RG-S05-algebra-proposition-characterize-separable-field-extensions}

\subsubsection{Proposition F.91: characterize separable field extensions}\label{AG-RG-S05-proposition-f.91-characterize-separable-field-extensions}

\emph{Adapted from the Stacks project, algebra.tex, lines 46037--46069.}
\nopagebreak[4]

For any field extension \(K/k\), separability, geometric reducedness, algebraic formal smoothness, vanishing of \(H_1(L_{K/k})\), and injectivity of \(K\otimes_k\Omega_{k/P}\to\Omega_{K/P}\) are equivalent; here \(P\) is the prime field. In characteristic \(p\), separability means that a \(p\)-basis of \(k\) remains \(p\)-independent over \(K^p\). In characteristic zero all these conditions hold. For finitely generated extensions this agrees with having a separating transcendence basis.

\textbf{Proof.} Every field is formally smooth over its prime field by the full arbitrary-field proof AG-CA-19S, Lemma 1.1. Suppose the displayed differential map is injective. In a square-zero lifting test for \(K/k\), first obtain a prime-field lift \(u:K\to E\) by that lemma. Its restriction differs from the given map of \(k\) by a derivation \(D:k\to J\), where \(J\) is a \(K\)-vector space through the map modulo \(J\). The corresponding \(K\)-linear map \(K\otimes_k\Omega_{k/P}\to J\) extends to \(\Omega_{K/P}\) by choosing a basis complement. Subtract the resulting derivation from \(u\). The square-zero product identity shows that this is still a ring map, and its restriction to \(k\) is the prescribed one. Thus \(K/k\) is formally smooth.

Conversely for any \(P\)-derivation \(D:k\to V\), with \(V\) a \(K\)-vector space, use the test algebra \(K\oplus V\) with structural map \(a\mapsto(a,D(a))\) and target quotient \(K\). A lift of the identity gives a \(P\)-derivation of \(K\) extending \(D\). Hence every linear functional on \(K\otimes_k\Omega_{k/P}\) extends to \(\Omega_{K/P}\). If the displayed map had a nonzero kernel vector, a linear functional nonzero on that vector could not extend. It is therefore injective. The conormal/vector-space argument in the preceding formal-field lemma equates formal smoothness with the \(H_1\) condition.

In characteristic zero choose a transcendence basis of \(k/P\) and extend it to one of \(K/P\). Algebraic extensions in characteristic zero are separable. The polynomial and separable algebraic differential formulas AG-CA-16, Theorem 6.2, identify their absolute differentials with the free spaces on those bases, proving the indicated injection. The same formula for arbitrary algebraic unions follows because each finite polynomial expression and relation occurs in a finite subextension. The lifting proof in the first paragraph consequently applies.

In characteristic \(p\), the absolute instance of \hyperref[AG-RG-S05-more-algebra-lemma-p-basis]{F.119} says that an absolute \(p\)-basis \(B\) of \(k\) has differentials forming a basis of \(\Omega_{k/P}\). The same proved equivalence for \(K/P\) says that their images in \(\Omega_{K/P}\) are independent exactly when the monomials in \(B\), with exponents below \(p\), are independent over \(K^p\). This proves equivalence of the differential injection with the \(p\)-basis condition. Its equivalence with the definition in F.5 will be proved below, rather than assumed here.

That condition is equivalent to \(K\otimes_k k^{1/p}\) being reduced. For a finite subset of \(B\) the tensor product is \(K[T_b]/(T_b^p-b)\), free on the standard monomials. It is a field exactly when those roots have degree \(p^{\#B_0}\), which is exactly the indicated monomial independence. If the degree is smaller, its map to the field generated by the roots has a nonzero kernel; every element of that kernel has zero \(p\)th power, so it is not reduced. The whole tensor product is the directed union of these rings; these inclusions are injective by their monomial bases. Reducedness is therefore equivalent to independence of all finite subsets.

The independence persists at every root height. Indeed, a relation \(\sum_a c_a^{p^r}B^a=0\) with \(0\le a_b<p^r\) can be grouped by \(a=pu+v\). Independence of the monomials \(B^v\), \(0\le v_b<p\), over \(K^p\) makes each grouped coefficient, a \(p\)th power, zero. Taking its unique \(p\)th root and induction on \(r\) makes every coefficient zero. Iterating the \(p\)-basis expansion proves that these are the bases for adjoining all \(p^r\)th roots of \(k\). Consequently tensoring \(K\) with any finite purely inseparable extension of \(k\) gives a field: it injects into the corresponding full-root tensor field and is finite-dimensional over \(K\).

For any finitely generated extension \(E/k\), \hyperref[AG-RG-S05-algebra-lemma-make-separably-generated]{F.69} gives finite purely inseparable \(k'/k\) and \(E'/E\) with \(E'/k'\) separably generated. Put \(F=K\otimes_k k'\), a field by the preceding root-height argument. Choose a finite separating basis \(T\) for \(E'/k'\) and its finite separable algebraic generator tower over \(k'(T)\). The ring \(F\otimes_{k'}k'(T)\) is the localization of the polynomial ring \(F[T]\) at the nonzero polynomials from \(k'[T]\), so it is Noetherian and regular by AG-CA-14, Proposition 3.3 and Theorem 3.2. Tensor the separable tower with \(F\). At every step its monic polynomial and invertible derivative give a finite étale algebra, by AG-CA-17, Proposition 6.3 and base change. The residue-field fibre algebras are finite products of finite separable fields by its Theorem 7.2. Thus flatness and regularity ascent \hyperref[AG-RG-S05-algebra-lemma-flat-over-regular-with-regular-fibre]{F.47} make every resulting ring regular; the rings are Noetherian because each step is finite over the preceding one. Regular local rings are domains, hence these rings are reduced by local detection. Consequently \(K\otimes_kE'=F\otimes_{k'}E'\) is reduced.

Tensoring the inclusion \(E\subset E'\) with the \(k\)-vector space \(K\) preserves injectivity: extend a basis of \(E\) to a basis of \(E'\) to see the split vector-space inclusion. Hence \(K\otimes_kE\) is reduced. Every field extension is the directed union of its finitely generated subfields; their tensor maps are injective by the same argument, and every nilpotent element occurs at one stage. This proves geometric reducedness for every field extension. The converse follows by testing \(k^{1/p}\).

Here is the finitely generated bridge to a separating transcendence basis. Suppose first that \(L/k\) is finitely generated and formally smooth. Choose field generators and a domain \(B=k[a_1,\ldots,a_n]\subset L\) with fraction field \(L\). Write \(B=P/I\), where \(P=k[X_1,\ldots,X_n]\). The ideal \(I\) is finitely generated by the Hilbert basis theorem AG-CA-03, Theorem 2.1. Formal smoothness of \(L/k\) splits the conormal sequence at its generic point by AG-CA-17, Theorem 3.1. Choose \(f_1,\ldots,f_c\) among the generators of \(I\) whose conormal classes are a basis over \(L\). Their differentials are independent in \(L^n\), so some \(c\)-column Jacobian determinant \(\Delta\) lies outside \(I\).

The polynomial local ring \(P_I\) has maximal ideal \(IP_I\) and residue field \(L\). Nakayama's lemma \hyperref[AG-RG-S05-algebra-lemma-NAK]{F.10} makes the \(f_i\) generate \(IP_I\). The finite module \(I/(f_1,\ldots,f_c)\) therefore vanishes there. For each of its finite generators choose an annihilating denominator outside \(I\), and multiply those denominators to obtain \(h_0\notin I\) with \(I_{h_0}=(f_1,\ldots,f_c)_{h_0}\). Put \(h=h_0\Delta\) and let \(g\) be its nonzero image in \(B\). Then \[B_g=k[X_1,\ldots,X_n,Z]/(f_1,\ldots,f_c,hZ-1).\] The minor using the selected variable columns and \(Z\) has block determinant \(h\Delta\), a unit. This is a standard smooth domain by AG-CA-17, Theorem 4.1, with fraction field \(L\). The independent converse construction in \hyperref[AG-RG-S05-algebra-lemma-localization-smooth-separable]{F.64} gives a separating transcendence basis. Conversely a finitely generated separably generated field is formally smooth by the transcendental and separable algebraic cases of \hyperref[AG-RG-S05-algebra-lemma-formally-smooth-extensions-easy]{F.49} and formal composition. This proves the finitely generated comparison without identifying separability with the \(p\)-basis condition in advance.

Now use the exact definition \hyperref[AG-RG-S05-algebra-definition-separable-field-extension]{F.5}. If \(K/k\) is separable by that definition, F.49(3) proves formal smoothness, hence all the conditions already compared above. Conversely, in characteristic \(p\), suppose the \(p\)-basis condition holds for \(K/k\). For every finitely generated intermediate field \(L\), it also holds for \(L/k\): a monomial relation with coefficients in \(L^p\subset K^p\) would already be a relation in \(K^p\). Thus the differential map for \(L/k\) is injective, its square-zero lifting proof above gives formal smoothness, and the finitely generated construction makes \(L/k\) separably generated. Definition F.5 therefore makes the original, possibly infinitely generated, \(K/k\) separable. In characteristic zero, every finitely generated intermediate field has a finite algebraic tower over a transcendence basis, and all its irreducible minimal polynomials are separable; the same definition applies. This proves the definition bridge and finishes every comparison in the full stated field scope. ∎

\phantomsection\label{AG-RG-S05-more-algebra-definition-G-ring}

\subsubsection{Definition F.92: G ring}\label{AG-RG-S05-definition-f.92-g-ring}

Reading source: \href{https://stacks.math.columbia.edu/tag/07GH}{Stacks, Tag 07GH}.

The term \emph{G-ring} imposes two requirements on \(R\): it is Noetherian, and at every prime \(\mathfrak p\) its localization-to-completion map is regular. The map in this second requirement is \[
R_{\mathfrak p}\longrightarrow\widehat{R_{\mathfrak p}},
\qquad
\widehat{R_{\mathfrak p}}
=\varprojlim_{n\ge1}R_{\mathfrak p}/(\mathfrak pR_{\mathfrak p})^n.
\] Thus the completion used in the definition is taken at the maximal ideal of the local ring, separately for each prime.

\phantomsection\label{AG-RG-S05-more-algebra-definition-excellent}

\subsubsection{Definition F.93: excellent}\label{AG-RG-S05-definition-f.93-excellent}

Reading source: \href{https://stacks.math.columbia.edu/tag/07QT}{Stacks, Tag 07QT}.

For a Noetherian ring \(R\), \emph{quasi-excellence} means that \(R\) satisfies the G-ring condition of \hyperref[AG-RG-S05-more-algebra-definition-G-ring]{F.92} and the J-2 condition. Here J-2 requires the regular locus \[
\{\mathfrak q\in\operatorname{Spec}B:B_{\mathfrak q}\text{ is regular}\}
\] to be open for each finite-type \(R\)-algebra \(B\); this is the terminology of \href{https://stacks.math.columbia.edu/tag/07P7}{Stacks, Tag 07P7}.

An \emph{excellent} ring has these properties and is universally catenary. To specify that last condition, a ring is catenary when, between each fixed pair of comparable primes, the lengths of finite strict chains have a common finite bound and the chains into which no further prime can be inserted all have equal length. A Noetherian ring is universally catenary when every finite-type algebra over it has this property. These two definitions are \href{https://stacks.math.columbia.edu/tag/00NI}{Stacks, Tag 00NI} and \href{https://stacks.math.columbia.edu/tag/00NL}{Tag 00NL}.

\phantomsection\label{AG-RG-S05-more-algebra-definition-local-complete-intersection}

\subsubsection{Definition F.94: local complete intersection}\label{AG-RG-S05-definition-f.94-local-complete-intersection}

Reading source: \href{https://stacks.math.columbia.edu/tag/07D0}{Stacks, Tag 07D0}.

For a map \(A\to B\) of finite type, choose a surjection \[
P=A[T_1,\ldots,T_n]\twoheadrightarrow B,
\qquad I=\ker(P\to B).
\] The map is a \emph{local complete intersection} when \(I\) is a Koszul-regular ideal, with that ideal condition understood exactly as in \hyperref[AG-RG-S05-more-algebra-definition-regular-ideal]{F.95}. Checking one such presentation suffices: the full arbitrary-ring proof \hyperref[AG-RG-S05-more-algebra-lemma-lci-presentation-independent]{F.94.A} below shows that every other finite polynomial presentation then has a Koszul-regular kernel. Neither this definition nor that comparison assumes \(A\) Noetherian.

\phantomsection\label{AG-RG-S05-more-algebra-lemma-lci-presentation-independent}

\subsubsection{Lemma F.94.A: arbitrary-ring presentation independence}\label{AG-RG-S05-lemma-f.94.a-arbitrary-ring-presentation-independence}

\emph{Adapted from the freely accessible Stacks project, Tag \href{https://stacks.math.columbia.edu/tag/07CZ}{07CZ}, with the finite-generation, change-of-basis and complex arguments written out below. This argument retains the lesson's GNU Free Documentation License terms.}

Let \(A\) be any ring and let \(C\) have two finite polynomial presentations \[
P=A[T_1,\ldots,T_n]\twoheadrightarrow C=P/I,\qquad
Q=A[X_1,\ldots,X_m]\twoheadrightarrow C=Q/J.
\] If \(I\) is a Koszul-regular ideal in the sense of Definition F.95, then so is \(J\). In particular the parenthetical assertion in Definition F.94 holds without a Noetherian hypothesis, and such a finite-type algebra is finitely presented.

\textbf{Proof.} First prove the required finite generation. At every prime containing \(I\), a principal neighbourhood makes \(I\) generated by a finite Koszul-regular list. The corresponding principal opens cover \(V(I)\). If their elements are \(a_\lambda\), no prime contains \(I+(a_\lambda)\), so that ideal is the unit ideal: a proper ideal lies in a prime by the full prime-separation proof in \href{https://kokunoyumeto.github.io/open-math-courses/courses/AG-RG/AG-RG-S04.html\#p-elementary-integral-algebra-before-descent-and-completion}{AG-RG-S04, Lemma P0.1}. Thus finitely many \(a_1,\ldots,a_s\) suffice, and \[
1=\sum_{i=1}^s b_i a_i+d,\qquad d\in I.
\] Choose in \(I\) numerators of the finite generators of each \(I_{a_i}\) and let \(N\) be the ideal generated by those numerators and \(d\). Then \(N_{a_i}=I_{a_i}\). At a prime outside all \(D(a_i)\), the displayed relation makes \(d\) a unit, so \(N\) and \(I\) are both the unit ideal there. Every prime is in one of these cases. Local detection applied to \(I/N\) gives \(I=N\). Hence \(I\) is finite and \(C\) is finitely presented. Lemma F.44, proved for arbitrary rings, now makes \(J\) finite too. If \(C=0\), both kernels are unit ideals and their Koszul-regularity condition is vacuous; this also gives a finite presentation. We may henceforth fix a prime \(\mathfrak q\) of the nonzero \(C\).

Choose \(p_j(T)\in P\) representing the images of \(X_j\) and \(q_i(X)\in Q\) representing the images of \(T_i\). In \(W=A[T,X]\) the common kernel is \[
L=(I,X_j-p_j(T))=(J,T_i-q_i(X)).
\] Quotienting either displayed ideal eliminates its added variables and gives precisely the specified map to \(C\), proving both equalities. Let \(\mathfrak r\) be the inverse image of \(\mathfrak q\) in \(W\) and put \(s_i=T_i-q_i(X)\). Translation identifies \(W\) with \(Q[s_1,\ldots,s_n]\) and \(L\) with \((J,s_1,\ldots,s_n)\). Expanding polynomials by their degree in the \(s_i\) gives \[
L/L^2\simeq (J/J^2)\oplus C^{\oplus n}.
\] Indeed its constant part is \(J/J^2\), its linear coefficients are read modulo \(J\), and all terms with two \(s\) factors lie in \(L^2\). These descriptions define mutually inverse maps, so require no finiteness or flatness assumption.

The Koszul complex on \(f_1,\ldots,f_c\) has free degree-\(k\) basis indexed by the \(k\)-element subsets; its differential deletes one index with the alternating sign and multiplies by the corresponding \(f_i\). The terms of its square cancel in pairs. If its first homology is zero, the classes of the \(f_i\) form a free basis of \((f)/(f)^2\). To prove independence, correct a relation \(\sum b_i f_i\in(f)^2\) by coefficients in \((f)\) so that its scalar product is exactly zero. First-homology vanishing expresses that relation as a sum of the skew Koszul boundaries, whose coefficients all belong to \((f)\). The original \(b_i\) therefore belong to \((f)\) as well. Surjectivity follows because the \(f_i\) generate the ideal.

By the hypothesis on \(I\), choose \(a\in P\) outside the contracted prime and a Koszul-regular list \(f\) generating \(I_a\). Polynomial extension and localization preserve its positive-homology vanishing: they are flat, and the exact flat-tensor homology calculation is proved in AG-RG-S01, Lemma 1.0. Append \(u_j=X_j-p_j(T)\). Modulo \(I\) these are successive polynomial variables translated by constants, and eliminating them leaves \(P/I\). Multiplication by each such variable is injective, as comparison of polynomial coefficients shows, also after localization. Appending one element to a Koszul complex is the cone of its multiplication map; the connecting-homomorphism construction in \href{https://kokunoyumeto.github.io/open-math-courses/courses/AG-CA/resolutions-tor-and-ext.html}{AG-CA-06S, Section 2} then identifies its positive homology with the kernel and higher homology of that map. The higher groups are zero and the degree-zero multiplication is injective. Induction therefore makes the whole list \((f,u)\) Koszul-regular in \(W_a\). Its conormal classes give a free basis at \(\mathfrak r\) by the preceding relation calculation.

Consequently \((J/J^2)_{\mathfrak q}\) is a finite direct summand of that free conormal module. Such a summand is flat by restricting the free tensor injection, and is finitely presented because the complementary summand is generated by the projections of a finite basis. AG-RG-S01, Lemma 1.2 makes it finite free over the local ring \(C_{\mathfrak q}\). Choose from a finite generating list of \(J\) elements \(g_1,\ldots,g_r\) whose conormal residue classes form a basis. Nakayama makes those classes generate the local conormal module; in a free basis their square coordinate matrix has unit determinant, so they form a basis there. They also generate \(J\) in the local ring \(Q_{\mathfrak q_Q}\), where \(\mathfrak q_Q\) is the inverse image of \(\mathfrak q\). In fact \(J=(g)+J^2\) there; the finite module \(J/(g)\) equals \(J\) times itself, and \(J\) is contained in that local maximal ideal. The same proved Nakayama calculation makes this quotient zero.

Thus \((g,s)\) and \((f,u)\) are two generator lists for \(L_{\mathfrak r}\) whose conormal classes are bases. They have the same length, since reduction gives two bases of \(L_{\mathfrak r}/\mathfrak rL_{\mathfrak r}\); here \(L^2\subset\mathfrak rL\). Expressing one list in the other gives a square matrix invertible modulo the maximal ideal, hence invertible by the cofactor identity. Its exterior powers give an isomorphism of their Koszul complexes: the alternating deletion formula above checks commutation with the differential on every basis wedge, and the exterior powers of the inverse give the inverse maps. Permuting the list is the same argument with a permutation matrix. In particular the Koszul complex on \((s,g)\) has zero positive homology at this stage.

We need a principal neighbourhood, not only this stalk. Clear the finitely many denominators in both generator comparisons and in the invertible matrix, and invert its determinant. For equalities in a localization, also include an element outside \(\mathfrak r\) killing each cross-difference; no zero divisor is cancelled. Together with \(a\) and the finite generator comparison for \(J\), their product \(h\in W\setminus\mathfrak r\) makes the list \((s,g)\) generate \(L_h\) and its Koszul complex isomorphic to the Koszul-regular complex on \((f,u)\) over \(W_h\). Thus its positive homology vanishes over the actual neighbourhood \(W_h\).

Finally set \(v=h(q_i(X),X)\in Q\). It is outside \(\mathfrak q_Q\), because evaluation modulo \((s)\) takes \(\mathfrak r\) to \(\mathfrak q_Q\). We have \[
W_h/(s_1,\ldots,s_n)=Q_v.
\] The Koszul complex on \(s\) is a free resolution of this quotient: before localization it is the successive polynomial-variable complex, whose acyclicity follows from the same cone calculation; flat localization preserves it. Tensor this resolution with the finite free terms of the Koszul complex on \(g\). The augmentation of the resulting total complex to the Koszul complex on \(g\) over \(Q_v\) is a homology isomorphism. Here is the finite elimination argument: include the quotient as augmented degree \(-1\) in each column. Those columns are exact. In a total cycle choose its component with largest \(g\) degree, lift its vertical cycle in that exact column, and subtract the total boundary. The largest degree disappears and only smaller degrees are changed. Repetition terminates and makes every cycle a boundary. This proves the augmentation assertion without a spectral-sequence or Noetherian assumption. The total complex is precisely the Koszul complex on \((s,g)\), so \(g\) has zero positive Koszul homology over \(Q_v\). The generator equality also descends to \(J_v=(g)\). We have therefore produced the neighbourhood required by Definition F.95 at every prime. Empty lists obey the same formulas. This proves every assertion. ∎

\phantomsection\label{AG-RG-S05-more-algebra-definition-regular-ideal}

\subsubsection{Definition F.95: regular ideal}\label{AG-RG-S05-definition-f.95-regular-ideal}

\emph{Adapted from the Stacks project, more-algebra.tex, lines 8710--8731.}
\nopagebreak[4]

Let \(R\) be a ring and let \(I \subset R\) be an ideal.

The sequence terms used here have the following meanings. The Koszul complex on a finite list \(f_1,\ldots,f_r\) is the alternating exterior-power complex described in Lemma F.94A. The list is \emph{Koszul-regular} if its positive homology vanishes, and \emph{\(H_1\)-regular} if its first homology vanishes. With \(J=(f_1,\ldots,f_r)\), it is \emph{quasi-regular} if the graded map \[
(R/J)[T_1,\ldots,T_r]\longrightarrow\bigoplus_{n\ge0}J^n/J^{n+1},\qquad
T_1^{e_1}\cdots T_r^{e_r}\longmapsto f_1^{e_1}\cdots f_r^{e_r}\bmod J^{1+\sum e_i}
\] is an isomorphism. In degree one this identifies \((R/J)^r\) with \(J/J^2\) on the displayed generator classes. This is the freely accessible \href{https://stacks.math.columbia.edu/tag/061P}{Stacks project definition, Tag 061P}, written here to specify the exact condition being used.

\begin{enumerate}
\def\labelenumi{\arabic{enumi}.}
\item
  We say \(I\) is a \emph{regular ideal} if for every \(\mathfrak p \in V(I)\) there exists a \(g \in R\), \(g \not \in \mathfrak p\) and a regular sequence \(f_1, \ldots, f_r \in R_g\) such that \(I_g\) is generated by \(f_1, \ldots, f_r\).
\item
  We say \(I\) is a \emph{Koszul-regular ideal} if for every \(\mathfrak p \in V(I)\) there exists a \(g \in R\), \(g \not \in \mathfrak p\) and a Koszul-regular sequence \(f_1, \ldots, f_r \in R_g\) such that \(I_g\) is generated by \(f_1, \ldots, f_r\).
\item
  We say \(I\) is a \emph{\(H_1\)-regular ideal} if for every \(\mathfrak p \in V(I)\) there exists a \(g \in R\), \(g \not \in \mathfrak p\) and an \(H_1\)-regular sequence \(f_1, \ldots, f_r \in R_g\) such that \(I_g\) is generated by \(f_1, \ldots, f_r\).
\item
  We say \(I\) is a \emph{quasi-regular ideal} if for every \(\mathfrak p \in V(I)\) there exists a \(g \in R\), \(g \not \in \mathfrak p\) and a quasi-regular sequence \(f_1, \ldots, f_r \in R_g\) such that \(I_g\) is generated by \(f_1, \ldots, f_r\).
\end{enumerate}

\phantomsection\label{AG-RG-S05-more-algebra-lemma-G-ring-goes-up-quasi-finite}

\subsubsection{Lemma F.96: G ring goes up quasi finite}\label{AG-RG-S05-lemma-f.96-g-ring-goes-up-quasi-finite}

\emph{Adapted from the Stacks project, more-algebra.tex, lines 12823--12870.}
\nopagebreak[4]

Let \(R\to R'\) be a finite type map of Noetherian rings. Let \(\mathfrak q'\subset\mathfrak p'\) be primes of \(R'\) with contractions \(\mathfrak q\subset\mathfrak p\) in \(R\), and assume the map is quasi-finite at \(\mathfrak p'\).

\begin{enumerate}
\def\labelenumi{\arabic{enumi}.}
\tightlist
\item
  If \(\widehat{R_{\mathfrak p}}\otimes_R\kappa(\mathfrak q)\) is geometrically regular over \(\kappa(\mathfrak q)\), then \(\widehat{R'_{\mathfrak p'}}\otimes_{R'}\kappa(\mathfrak q')\) is geometrically regular over \(\kappa(\mathfrak q')\).
\item
  If every formal fibre of \(R_{\mathfrak p}\) is geometrically regular, then every formal fibre of \(R'_{\mathfrak p'}\) is geometrically regular.
\item
  If \(R\to R'\) is quasi-finite and \(R\) is a G-ring, then \(R'\) is a G-ring.
\end{enumerate}

\textbf{Proof.} Put \(\widehat R=\widehat{R_{\mathfrak p}}\) and \(\widehat S=\widehat{R'_{\mathfrak p'}}\). Both completions are Noetherian by AG-CA-19, Theorem 3.3. The finite-type algebra \(T=\widehat R\otimes_R R'\) is Noetherian by AG-CA-03, Theorem 2.1. The explicit quasi-finite completion-factor proof \hyperref[AG-RG-S05-algebra-lemma-completion-at-quasi-finite-prime]{F.26} gives \[T\cong\widehat S\times D,\] with first projection the natural completed-local map.

Write \(k=\kappa(\mathfrak q)\) and \(E=\kappa(\mathfrak q')\). The injective map \(R/\mathfrak q\to R'/\mathfrak q'\) extends to \(k\subset E\): every nonzero element of the source domain has nonzero image and becomes invertible in the target fraction field. Put \(F=\widehat R\otimes_R k\), geometrically regular by the hypothesis of (1). Tensor associativity identifies \[H=T\otimes_{R'}E\cong F\otimes_kE.\] Here \(H\) is Noetherian: if \(S=R'\setminus\mathfrak q'\), then \(E=S^{-1}(R'/\mathfrak q')\), so \(H=S^{-1}(T/\mathfrak q'T)\). Quotients and localizations preserve Noetherianness by the same earlier Hilbert-basis and permanence proof. The full field-extension result \hyperref[AG-RG-S05-algebra-lemma-geometrically-regular-over-subfields]{F.52} therefore makes \(H\) geometrically regular over \(E\).

Tensoring the displayed product decomposition gives \[H\cong(\widehat S\otimes_{R'}E)\times(D\otimes_{R'}E).\] To check geometric regularity of its first factor \(H_1\), let \(E'/E\) be finite purely inseparable. The ring \(H\otimes_EE'\) is regular, and it is the product of the two corresponding factor tensors. The first factor is its localization at the idempotent \((1,0)\), hence regular: its local rings are local rings of that regular product. It is Noetherian as a quotient of the product. The first factor \(H_1\) is itself Noetherian as a quotient of \(H\). Criterion \hyperref[AG-RG-S05-algebra-lemma-geometrically-regular]{F.50} now proves that \(H_1=\widehat S\otimes_{R'}E\) is geometrically regular over \(E\). Zero factors cause no exception, since they have no prime local rings. This proves (1).

For (2), let \(\mathfrak q'\) run over all primes contained in \(\mathfrak p'\); their contractions lie in \(\mathfrak p\), so (1) applies to every formal fibre. Completion is flat by AG-CA-19, Theorem 3.2. Thus these geometrically regular fibres give a regular completion map. Finally, if \(R\) is a G-ring and the map is quasi-finite at every prime, apply this argument at every \(\mathfrak p'\). Definition \hyperref[AG-RG-S05-more-algebra-definition-G-ring]{F.92} says exactly that the resulting regular completion maps make \(R'\) a G-ring. ∎

\phantomsection\label{AG-RG-S05-more-algebra-lemma-another-helper-G-ring}

\subsubsection{Lemma F.97: another helper G ring}\label{AG-RG-S05-lemma-f.97-another-helper-g-ring}

\emph{Adapted from the Stacks project, more-algebra.tex, lines 13103--13159.}
\nopagebreak[4]

Let \(R\) be any Noetherian complete local domain with fraction field \(K\). First, for a prime \(q\) of \(R[X]\), \(\widehat{R[X]_q}\otimes_RK\) is regular. If \(q\) lies over the maximal ideal and \(0\ne r\subset q\) lies over \((0)\subset R\), then \(\widehat{R[X]_q}\otimes_{R[X]}\kappa(r)\) is geometrically regular over \(\kappa(r)\).

\textbf{Proof.} Choose the already proved finite regular complete local subring \(A_0\subset R\), and put \(F=\operatorname{Frac}(A_0)\). Let \(q_0\) contract \(q\) to \(A_0[X]\). The finite-extension completion decomposition gives \[\widehat{A_0[X]_{q_0}}\otimes_{A_0[X]}R[X]=\prod_i\widehat{R[X]_{q_i}},\] with the finitely many primes over \(q_0\), one of which is \(q\). Tensor with \(K\) over \(R\). The left side becomes \(\widehat{A_0[X]_{q_0}}\otimes_{A_0}K\): the finite domain \(R\) over \(A_0\) has \(R\otimes_{A_0}F=K\), since a finite domain over a field is a field. The strong first assertion of the preceding helper makes \(\widehat{A_0[X]_{q_0}}\otimes_{A_0}F\) geometrically regular over \(F\). Its base change to the finite extension \(K/F\) is regular by the field-extension proof. Each factor, including \(\widehat{R[X]_q}\otimes_RK\), is regular. This proves the first assertion without any use of a general formal-smoothness/regularity equivalence.

For the second, \(\kappa(r)/K\) is finite, because the nonzero prime \(rK[X]\) is generated by an irreducible polynomial. Given a finite purely inseparable \(L/\kappa(r)\), choose a finite integral \(R\)-subalgebra \(B\subset L\) with fraction field \(L\), by scaling finitely many algebraic generators into monic equations. As a finite \(R\)-module, \(B\) is Noetherian. Apply Lemma F.27 to the finite map \(R\to B\) at the maximal ideal of \(R\). Since \(R\) is complete, its left side is \(B\), and it gives the canonical isomorphism \(B\cong\prod_i\widehat{B_{\mathfrak n_i}}\), with the finitely many maximal ideals \(\mathfrak n_i\) of \(B\). Each factor has its nonzero residue-field quotient. A domain has only the idempotents zero and one, so this product has exactly one factor. Thus \(B\) has one maximal ideal \(\mathfrak n\), \(B=B_{\mathfrak n}\), and the displayed map identifies \(B\) with its maximal-ideal completion. This proves that \(B\) is a complete local domain directly from the already written finite-completion decomposition. Let \(r'\) be the kernel of \(B[X]\to L\) sending \(X\) to its specified image from \(\kappa(r)\); then \(\kappa(r')=L=\operatorname{Frac}(B)\).

Finite-extension completion and tensor associativity identify the tested fibre with a product of \[\widehat{B[X]_{q_i}}\otimes_{B[X]}\kappa(r'),\] over primes \(q_i\) above \(q\); factors not meeting this fibre are zero. In each remaining factor, \(r'L[X]=(X-a)\) for some \(a\in L\), so this is \[\big(\widehat{B[X]_{q_i}}\otimes_B L\big)/(X-a).\] The parenthesized ring is regular by the first assertion, applied to the complete local domain \(B\). Lemma \hyperref[AG-RG-S05-more-algebra-lemma-derivation-extends]{F.104} extends \(\partial/\partial X\) from \(B[X]\) through its prime localization and completion, and then through localization at \(B\setminus\{0\}\). It kills \(B\), so the fraction formula in that proof makes it zero on \(L=\operatorname{Frac}(B)\) as well. Thus it sends \(X-a\) to one. Lemma \hyperref[AG-RG-S05-more-algebra-lemma-quotient-regular]{F.122} applies to this regular parenthesized ring and makes each quotient factor regular. Thus every finite purely inseparable test extension \(L\) gives a regular fibre, and the geometric-regularity criterion gives the assertion. This is the exact completion argument consumed by polynomial permanence of G-rings. ∎

\phantomsection\label{AG-RG-S05-more-algebra-lemma-cartier-equality}

\subsubsection{Lemma F.98: cartier equality}\label{AG-RG-S05-lemma-f.98-cartier-equality}

Reading source: \href{https://stacks.math.columbia.edu/tag/07E1}{Stacks, Tag 07E1}.

For a finitely generated field extension \(K/k\), the spaces of differentials and first conormal homology are finite-dimensional, and \[
\dim_K\Omega_{K/k}-\dim_K H_1(L_{K/k})
=\operatorname{trdeg}_k K.
\]

\textbf{Proof.} Put \(d=\operatorname{trdeg}_k K\). Apply the explicit monic-tower construction in \hyperref[AG-RG-S05-algebra-lemma-colimit-syntomic]{F.23} to a finite field-generating list for \(K/k\). It gives a domain \(C\subset K\) containing that list and having fraction field \(K\). In that construction there are \(d\) transcendence variables, \(r\) algebraic tower variables and one inverse variable. Consequently its polynomial presentation has \[
N=d+r+1\quad\text{variables},\qquad c=r+1\quad\text{equations}.
\] The first equation is \(hZ-1\) for the chosen nonzero denominator polynomial \(h\). The remaining equations, after imposing this first equation, are the successive monic minimal polynomials in the tower. Multiplication by \(hZ-1\) is injective: the lowest nonzero coefficient of a polynomial in its product with \(hZ-1\) is the negative of that coefficient. Each subsequent monic equation is injective on its polynomial ring by its highest-coefficient calculation. Thus this ordered list is regular, exactly as in the construction of F.23.

Write \(C=k[X_1,\ldots,X_N]/I\). \hyperref[AG-RG-S05-algebra-lemma-regular-quasi-regular]{F.74} identifies \(I/I^2\) with \(C^c\) on the equation classes. Polynomial differentials identify the other term with \(C^N\) on \(dX_i\). Presentation comparison \hyperref[AG-RG-S05-algebra-lemma-NL-homotopy]{F.11} and localization \hyperref[AG-RG-S05-algebra-lemma-localize-NL]{F.65}, applied at all nonzero elements of the domain \(C\), therefore identify the first conormal homology and the differential module of \(K/k\) with the kernel and cokernel of a linear map \[
D:K^c\longrightarrow K^N.
\] Localization is exact by \href{https://kokunoyumeto.github.io/open-math-courses/courses/AG-RG/AG-RG-S01.html}{S01, Lemma 1.1}, so it preserves the kernel and cokernel used in this comparison.

If \(\rho\) is the rank of \(D\), choosing bases for its image and extending them to bases gives \[
\dim_K\ker D=c-\rho,\qquad
\dim_K\operatorname{coker}D=N-\rho.
\] Their difference is \(N-c=d\), and both numbers are finite. This proves the asserted equality without a separability hypothesis. Equivalently the parameter count is \(\dim C=d=\operatorname{trdeg}_k K\), as proved for finite-type domains in \href{https://kokunoyumeto.github.io/open-math-courses/courses/AG-CA/krull-dimension-and-noether-normalization.html}{AG-CA-09, Theorem 4.2}. Empty transcendence or tower lists obey the same count, including \(K=k\). ∎

\phantomsection\label{AG-RG-S05-more-algebra-lemma-check-G-ring-easy}

\subsubsection{Lemma F.99: check G ring easy}\label{AG-RG-S05-lemma-f.99-check-g-ring-easy}

\emph{Adapted from the Stacks project, more-algebra.tex, lines 12803--12821.}
\nopagebreak[4]

A Noetherian ring \(R\) is a G-ring if and only if for every \(\mathfrak q\subset\mathfrak p\) the ring \((R/\mathfrak q)_{\mathfrak p}^{\wedge}\otimes_{R/\mathfrak q}\kappa(\mathfrak q)\) is geometrically regular over \(\kappa(\mathfrak q)\).

\textbf{Proof.} The completion map of each Noetherian local ring is flat by AG-CA-19, Theorem 3.2. Its fibre at a contracted prime \(\mathfrak q\) is \(\widehat{R_{\mathfrak p}}\otimes_R\kappa(\mathfrak q)\). Exactness of finite-module completion, AG-CA-19, Theorem 3.1, gives \(\widehat{R_{\mathfrak p}}/\mathfrak q\widehat{R_{\mathfrak p}}=\widehat{(R/\mathfrak q)_{\mathfrak p}}\). Tensoring this equality with the fraction field of \(R/\mathfrak q\) identifies the displayed fibre with the ring in the assertion. A flat map is regular precisely when all its residue-field fibres are geometrically regular, by the stated definition. Applying this to every local completion gives both directions. ∎

\phantomsection\label{AG-RG-S05-completion-flat-local-map}

\subsubsection{Lemma F.99A: completion of a flat local map}\label{AG-RG-S05-lemma-f.99a-completion-of-a-flat-local-map}

Let \((A,\mathfrak m)\to(B,\mathfrak n)\) be a flat local map of Noetherian local rings. Its induced map on maximal-ideal completions \(\widehat A\to\widehat B\) is faithfully flat. The residue fields need not be equal.

\textbf{Proof.} A local map sends \(\mathfrak m^r\) into \(\mathfrak n^r\); passing to compatible finite quotients therefore defines the map of completions. It remains local because their maximal ideals are \(\mathfrak m\widehat A\) and \(\mathfrak n\widehat B\), and their residue fields are \(A/\mathfrak m\) and \(B/\mathfrak n\), respectively. These quotient and Noetherian assertions are AG-CA-19, Proposition 2.2 and Theorem 3.3.

Put \(\kappa=A/\mathfrak m\). Choose a free \(A\)-resolution \(F_\bullet\to\kappa\), as constructed in \href{https://kokunoyumeto.github.io/open-math-courses/courses/AG-CA/resolutions-tor-and-ext.html}{AG-CA-06S, Section 1}. Completion is flat by AG-CA-19, Theorem 3.2, so \(F_\bullet\otimes_A\widehat A\) is a free \(\widehat A\)-resolution of \(\kappa\otimes_A\widehat A=\kappa\). The ring \(\widehat B\) is flat over \(B\) by the same theorem and hence flat over \(A\) by composition. Consequently

\[
\operatorname{Tor}^{\widehat A}_1(\kappa,\widehat B)
=H_1(F_\bullet\otimes_A\widehat B)
=\operatorname{Tor}^{A}_1(\kappa,\widehat B)=0.
\]

Apply the proved local criterion for flatness, AG-CA-08, Theorem 4.2, to the Noetherian local map \(\widehat A\to\widehat B\) and the finite \(\widehat B\)-module \(\widehat B\). It proves flatness. Locality gives \(\mathfrak m\widehat B\subseteq\mathfrak n\widehat B\), so the closed fibre \(\widehat B/\mathfrak m\widehat B\) surjects onto \(\widehat B/\mathfrak n\widehat B=B/\mathfrak n\). This residue field is nonzero, hence the closed fibre is nonzero. The local faithful-flatness criterion AG-CA-08, Theorem 2.1 now makes the map faithfully flat. This proves the assertion. ∎

\phantomsection\label{AG-RG-S05-more-algebra-lemma-check-G-ring-maximal-ideals}

\subsubsection{Lemma F.100: check G ring maximal ideals}\label{AG-RG-S05-lemma-f.100-check-g-ring-maximal-ideals}

\emph{Adapted from the Stacks project, more-algebra.tex, lines 13001--13037.}
\nopagebreak[4]

Let \(R\) be a Noetherian ring. Then \(R\) is a G-ring if and only if \(R_\mathfrak m\) has geometrically regular formal fibres for every maximal ideal \(\mathfrak m\) of \(R\).

\textbf{Proof.} Assume \(R_\mathfrak m \to R_\mathfrak m^\wedge\) is regular for every maximal ideal \(\mathfrak m\) of \(R\). Let \(\mathfrak p\) be a prime of \(R\) and choose a maximal ideal \(\mathfrak p \subset \mathfrak m\). Since \(R_\mathfrak m \to R_\mathfrak m^\wedge\) is faithfully flat we can choose a prime \(\mathfrak p'\) in \(R_\mathfrak m^\wedge\) lying over \(\mathfrak pR_\mathfrak m\). Consider the commutative diagram

\[
\begin{gathered}
R_\mathfrak m^\wedge \xrightarrow{} (R_\mathfrak m^\wedge)_{\mathfrak p'} \\
(R_\mathfrak m^\wedge)_{\mathfrak p'} \xrightarrow{} (R_\mathfrak m^\wedge)_{\mathfrak p'}^\wedge \\
R_\mathfrak m \xrightarrow{} R_\mathfrak m^\wedge \\
R_\mathfrak m \xrightarrow{} R_\mathfrak p \\
R_\mathfrak p \xrightarrow{} (R_\mathfrak m^\wedge)_{\mathfrak p'} \\
R_\mathfrak p \xrightarrow{} R_\mathfrak p^\wedge \\
R_\mathfrak p^\wedge \xrightarrow{} (R_\mathfrak m^\wedge)_{\mathfrak p'}^\wedge
\end{gathered}
\]

By assumption the ring map \(R_\mathfrak m \to R_\mathfrak m^\wedge\) is regular. By Proposition \hyperref[AG-RG-S05-more-algebra-proposition-Noetherian-complete-G-ring]{F.132} \((R_\mathfrak m^\wedge)_{\mathfrak p'} \to
(R_\mathfrak m^\wedge)_{\mathfrak p'}^\wedge\) is regular. The localization \(R_\mathfrak m^\wedge \to (R_\mathfrak m^\wedge)_{\mathfrak p'}\) is regular. Hence \(R_\mathfrak m \to (R_\mathfrak m^\wedge)_{\mathfrak p'}^\wedge\) is regular by Lemma \hyperref[AG-RG-S05-more-algebra-lemma-regular-composition]{F.124}. Since it factors through the localization \(R_\mathfrak p\), also the ring map \(R_\mathfrak p \to (R_\mathfrak m^\wedge)_{\mathfrak p'}^\wedge\) is regular. The local map \(R_\mathfrak p\to(R_\mathfrak m^\wedge)_{\mathfrak p'}\) is flat, being a localization of the flat completion map. Lemma \hyperref[AG-RG-S05-completion-flat-local-map]{F.99A} makes the induced map \(R_\mathfrak p^\wedge\to
(R_\mathfrak m^\wedge)_{\mathfrak p'}^\wedge\) faithfully flat. Thus Lemma \hyperref[AG-RG-S05-more-algebra-lemma-regular-permanence]{F.125} applies to this actual factorization and proves that \(R_\mathfrak p\to R_\mathfrak p^\wedge\) is regular. ∎

\phantomsection\label{AG-RG-S05-more-algebra-lemma-conormal-sequence-H1-regular-ideal}

\subsubsection{Lemma F.101: conormal sequence H1 regular ideal}\label{AG-RG-S05-lemma-f.101-conormal-sequence-h1-regular-ideal}

Reading source: \href{https://stacks.math.columbia.edu/tag/07CX}{Stacks, Tag 07CX}.

For ideals \(I\subset J\) of an arbitrary ring \(A\), suppose the ideal \(J/I\) of \(A/I\) is \(H_1\)-regular. Then \[
I\cap J^2=IJ.
\]

\textbf{Proof.} The inclusion \(IJ\subset I\cap J^2\) follows from \(I\subset J\). We prove the other inclusion at each prime of \(A\). Exact localization commutes with the intersection, the product and the quotient in this assertion; the fraction proof and the detection of a zero module are \href{https://kokunoyumeto.github.io/open-math-courses/courses/AG-RG/AG-RG-S01.html}{S01, Lemma 1.1}. These facts will turn the local equalities into the global one.

At a prime outside \(V(J)\), the ideal \(J\) becomes the unit ideal, and both sides become \(I\). At a prime containing \(J\), localize \(A\) there. The defining neighbourhood in \hyperref[AG-RG-S05-more-algebra-definition-regular-ideal]{F.95} supplies a finite list \(g_1,\ldots,g_s\in J\) whose images generate \(J/I\) and have zero first Koszul homology over \(A/I\). Indeed we can lift the list from its quotient neighbourhood and localize further at the prime; exact localization preserves cycles, boundaries and their homology. In this local ring \(J=I+(g_1,\ldots,g_s)\).

Take \(x\in I\cap J^2\). Expanding a finite expression for \(x\) as a sum of products of members of \(J\) gives \[
x=u+\sum_{i=1}^s g_i b_i,
\qquad u\in IJ,\quad b_i\in J.
\] Terms having a factor in \(I\) were absorbed in \(u\); all other terms have two factors from the displayed list. Reducing modulo \(I\) makes \((\bar b_i)\) a degree-one Koszul cycle because its scalar product with \((\bar g_i)\) is zero. The zero first homology says this cycle is a boundary. With the convention \[
\partial(e_i\wedge e_j)=\bar g_i e_j-\bar g_j e_i,
\] it is a finite sum of these displayed vectors with coefficients in \(A/I\).

Lift those coefficients to \(A\) and subtract the corresponding vectors \(g_i e_j-g_j e_i\) from \((b_i)\). Their scalar products with \((g_i)\) are exactly zero in \(A\), so this subtraction leaves \(\sum_i g_i b_i\) unchanged. Its new coefficients reduce to zero modulo \(I\), hence belong to \(I\). The sum is therefore in \(IJ\), as is \(u\), proving \(x\in IJ\). An empty list gives the same calculation with sum zero. No finite generation of \(I\) and no Noetherian condition were used. Applying local detection to \((I\cap J^2)/IJ\) completes the proof. ∎

\phantomsection\label{AG-RG-S05-more-algebra-lemma-cotangent-complex-symmetric-algebra}

\subsubsection{Lemma F.102: cotangent complex symmetric algebra}\label{AG-RG-S05-lemma-f.102-cotangent-complex-symmetric-algebra}

Reading source: \href{https://stacks.math.columbia.edu/tag/07EV}{Stacks, Tag 07EV}.

Let \(A\) be any ring and let \[
0\longrightarrow K\longrightarrow A^m\xrightarrow{q}M\longrightarrow0
\] be exact, where \(M\) is finite projective. Set \(C=\operatorname{Sym}_A(M)\) and let \(\alpha:A[y_1,\ldots,y_m]\twoheadrightarrow C\) be the presentation induced by \(q\), with kernel \(J\). The natural map \[
K\otimes_A C\longrightarrow J/J^2,
\qquad (k_1,\ldots,k_m)\otimes c
\longmapsto c\sum_i k_i y_i\bmod J^2
\] is an isomorphism. With this identification the conormal complex and its cokernel are \[
\operatorname{NL}(\alpha)=[K\otimes_A C\longrightarrow C^m],
\qquad \Omega_{C/A}=M\otimes_A C,
\] where the differential is obtained from \(K\hookrightarrow A^m\) by tensoring with \(C\).

\textbf{Proof.} Choose a section \(s:M\to A^m\) using projectivity. The map \((k,m)\mapsto k+s(m)\) identifies \(K\oplus M\) with \(A^m\): its inverse sends \(a\) to \((a-sq(a),q(a))\). We compute the presentation using this decomposition, keeping track of its specified natural map.

For a module \(V\), the symmetric algebra is the tensor algebra modulo the relations that interchange two adjacent factors. A linear map \(V\to D\) into a commutative \(A\)-algebra extends on a pure tensor by multiplying its images; the interchange relations vanish, and these products determine the extension uniquely. This proves the symmetric-algebra universal property from the tensor construction \href{https://kokunoyumeto.github.io/open-math-courses/courses/AG-RG/AG-RG-S01.html}{S01, Lemma 1.0}. It also gives its grading by tensor degree, with degree zero \(A\) and degree one \(V\).

That property and the algebra tensor universal property give inverse maps \[
\operatorname{Sym}_A(K\oplus M)
\ \cong\ \operatorname{Sym}_A(K)\otimes_A C.
\] In one direction send \(k\) to \(k\otimes1\) and \(m\) to \(1\otimes m\). In the other direction use the inclusions of \(K\) and \(M\) into \(K\oplus M\) on the two tensor factors. Both composites fix the algebra generators, hence are identities. Likewise the symmetric algebra of \(A^m\) is the polynomial ring on its coordinate basis. Under these identifications \(\alpha\) is the augmentation on the \(K\) factor and the identity on \(C\).

Grade the tensor product by \(K\) degree. The direct-sum tensor formula of S01, Lemma 1.0 gives \[
J=\bigoplus_{d\ge1}\operatorname{Sym}_A^d(K)\otimes_A C,
\qquad
J^2=\bigoplus_{d\ge2}\operatorname{Sym}_A^d(K)\otimes_A C.
\] For the second equality, a product of positive degrees has degree at least two. Conversely each degree-\(d\) term is generated by products \(k_1\cdots k_d\otimes c\); when \(d\ge2\) separate its first factor from the rest to express it as a product of two elements of \(J\). Quotienting thus leaves precisely \(K\otimes_A C\). Its generator \(k\otimes c\) corresponds in the original coordinates to \(c\sum_i k_i y_i\), proving exactly the natural isomorphism in the statement, independently of the auxiliary section.

Relative differentiation sends this linear polynomial to \(\sum_i k_i\,dy_i\). The polynomial differentials are free on these \(m\) symbols, so the displayed complex has exactly the claimed differential. Tensoring the split sequence with \(C\) remains split exact, and its cokernel is \(M\otimes_A C\). This proves the assertion about \(\Omega_{C/A}\) as well, including zero modules and the empty coordinate basis.

The projectivity hypothesis prevents additional conormal relations. For example let \(A=k[\varepsilon]/(\varepsilon^2)\), let \(M=k^2\) via \(\varepsilon M=0\), and let \(q:A^2\to M\) be reduction. Then \(K=\varepsilon A^2\), \(C=A[y_1,y_2]/(\varepsilon y_1,\varepsilon y_2)\), and \(J=(\varepsilon y_1,\varepsilon y_2)\) has square zero. The source \(K\otimes_A C\) is \(k[y_1,y_2]^2\), whereas the vector \((y_2,-y_1)\) maps to \(\varepsilon y_1y_2-\varepsilon y_2y_1=0\) in \(J/J^2\) and is nonzero in that source. Thus the same conormal identity need not hold after dropping the hypothesis. ∎

\phantomsection\label{AG-RG-S05-more-algebra-lemma-degree-p-extension-regular}

\subsubsection{Lemma F.103: degree p extension regular}\label{AG-RG-S05-lemma-f.103-degree-p-extension-regular}

\emph{Adapted from the Stacks project, more-algebra.tex, lines 12453--12467.}
\nopagebreak[4]

If \(R\) is Noetherian regular and \(D(f)\) is a unit for some derivation, then \(R[Z]/(P(Z)-f)\) is regular for every \(P\in\mathbb Z[Z]\).

\textbf{Proof.} Polynomial extension of a Noetherian regular ring is regular by AG-CA-14, Proposition 3.3. Extend the derivation coefficientwise and set \(D(Z)=0\). Every integer and hence \(P(Z)\) is annihilated. The derivative of \(P(Z)-f\) is \(-D(f)\), still a unit in the quotient, so the preceding hypersurface proof applies at every prime. The special case \(P(Z)=Z^p\) is the one used in the inseparable induction. ∎

\phantomsection\label{AG-RG-S05-more-algebra-lemma-derivation-extends}

\subsubsection{Lemma F.104: derivation extends}\label{AG-RG-S05-lemma-f.104-derivation-extends}

\emph{Adapted from the Stacks project, more-algebra.tex, lines 12363--12405.}
\nopagebreak[4]

A derivation \(D:R\to R\) extends to every localization and canonically to every ideal-adic completion. If \(R\subset R'\) is finite type and \(R_g=R'_g\) with \(g\) a nonzerodivisor in \(R'\), some \(g^ND\) preserves \(R'\).

\textbf{Proof.} In a localization the forced formula is \(D(r/s)=D(r)/s-rD(s)/s^2\). If \(r/s=r'/s'\), a multiplier kills \(rs'-r's\); differentiating this equality and using the original equality after localization proves equality of the two displayed values. Direct expansion verifies addition and the product rule; differentiating \(s(s^{-1})=1\) proves uniqueness.

For completion, Leibniz gives \(D(I^{n+1})\subset I^n\). For a compatible family \(a_n\in R/I^n\), choose a representative of \(a_{n+1}\) and define the \(n\)th coordinate of \(\widehat D(a)\) by its derivative modulo \(I^n\). It is independent of that representative and compatible as \(n\) varies. Represent two families modulo \(I^{n+1}\); the product rule there reduces modulo \(I^n\) to the product rule for the new coordinate. Addition is checked in the same way. Every coordinate therefore satisfies the derivation identities, so their inverse limit does too. This fills the completion verification.

For finite type choose generators \(x_i\) of \(R'/R\) and \(n\) with \(g^nx_i\in R\). Inside \(R'_g\), \[g^{n+1}D(x_i)=-ng^nx_iD(g)+gD(g^nx_i)\in R'.\] Because \(R'\) injects into \(R'_g\), the extended derivation \(g^{n+1}D\) can be restricted to \(R'\): it preserves the coefficient subring and the generators, and the product rule shows that it preserves every polynomial in them. All relations continue to hold in this subring. This proves the assertion. ∎

\phantomsection\label{AG-RG-S05-more-algebra-lemma-find-D}

\subsubsection{Lemma F.105: find D}\label{AG-RG-S05-lemma-f.105-find-d}

\emph{Adapted from the Stacks project, more-algebra.tex, lines 12469--12520.}
\nopagebreak[4]

If a characteristic-\(p\) domain \(B\) is finite type over a Noetherian complete local ring and \(f\in B\) is not a \(p\)th power in its fraction field, there exists a derivation \(D:B\to B\) with \(D(f)\ne0\).

\textbf{Proof.} Replace the complete base by its image in \(B\). It is its quotient by a prime, hence again a complete Noetherian local domain by exact completion AG-CA-19, Theorem 3.1. The finite-over-regular-subring construction already proved above reduces it to a finite extension of \(R=k[[X_1,\ldots,X_n]]\). Thus \(B\) is finite type over \(R\).

Apply Noether normalization AG-CA-09, Corollary 3.2, to \(B\otimes_R\operatorname{Frac}(R)\), a finite-type domain over that field. Clearing finitely many denominators gives algebraically independent \(y_j\in B\) and a nonzero \(g\in R\) such that \(B_g\) is finite over \(R_g[y]\). Here scaling normalized parameters clears their denominators without losing algebraic independence. For each finite generator \(b_i\) of \(B/R[y]\), take its monic integral equation over \(R_g[y]\). For one sufficiently large \(N\), the generator \(g^Nb_i\) satisfies a monic equation over \(R[y]\), because multiplying the coefficient of its degree-\(j\) term by \(g^{N(d-j)}\) clears that coefficient's denominator. The subalgebra \(B'\) generated by these scaled generators is finite over \(R[y]\), by reducing monomials with those monic equations, and \(B'_g=B_g\).

For a sufficiently large \(r\), \(f'=g^{pr}f\) belongs to \(B'\). It is still not a \(p\)th power in its fraction field, since the multiplier is a \(p\)th power. Set \(L=\operatorname{Frac}(B')\) and \(K=\operatorname{Frac}(R[y])\). By the power-series-subfield lemma the finite coefficient subrings \(A_J\) have fraction fields \(K_J\) intersecting in \(K^p\). By the finite-extension intersection lemma, \(\bigcap_JL^pK_J=L^p\). Therefore some \(J\) has \(f'\notin L^pK_J\). The relative \(p\)-basis construction extends \(f'\) to a \(p\)-basis of \(L/K_J\), and its partial derivation gives a \(K_J\)-derivation \(\delta:L\to L\) with \(\delta(f')=1\).

Since \(B'\) is finite over \(R[y]\) and that ring is finite over \(A_J\), it has finitely many \(A_J\)-module generators \(b_1,\ldots,b_t\). Write \(\delta(b_i)\) as fractions of \(B'\) and multiply \(\delta\) by their nonzero denominator product \(h\). Then \(h\delta(B')\subset B'\): differentiate each finite module expression and use that \(\delta\) kills \(A_J\). Its value on \(f'\) is \(h\ne0\). The finite-type derivation-extension proof above now supplies a power \(g^s\) for which \(D=g^sh\delta\) preserves \(B\). The fraction formula gives \(D(f)=g^{s-pr}h\delta(f')=g^{s-pr}h\ne0\) in its fraction field, hence in \(B\). This proves the required derivation without assuming \(k\) perfect or using characteristic zero differentiation. ∎

\phantomsection\label{AG-RG-S05-more-algebra-lemma-formally-smooth}

\subsubsection{Lemma F.106: formally smooth}\label{AG-RG-S05-lemma-f.106-formally-smooth}

Reading source: \href{https://stacks.math.columbia.edu/tag/07EC}{Stacks, Tag 07EC}.

Let \(\varphi:R\to S\) be continuous. If the underlying algebra map is formally smooth, it is formally smooth when \(R\) has any linear topology and \(S\) has a pre-adic topology. Moreover, if the topologies on \(R\) and \(S\) are given by the powers of ideals \(\mathfrak m\) and \(\mathfrak n\), formal smoothness is unchanged when the topology on \(R\) is replaced by the discrete topology.

\textbf{Proof.} We specify the topological test. Give a ring \(D\) and its quotient \(D/J\) the discrete topology, where \(J^2=0\). Fix continuous maps \(a:R\to D\) and \(b:S\to D/J\) such that \(b\varphi\) is the reduction of \(a\). A solution is a continuous \(R\)-algebra map \(c:S\to D\) reducing to \(b\). Topological formal smoothness requires a solution for every such test; algebraic formal smoothness uses the same tests without continuity conditions. The latter definition is \href{https://kokunoyumeto.github.io/open-math-courses/courses/AG-CA/formally-smooth-unramified-and-etale-ring-maps.html}{AG-CA-17, Section 1}.

For a pre-adic topology on \(S\), choose \(\mathfrak n\) whose powers form a neighbourhood basis at zero. A ring map from \(S\) to a discrete ring is continuous precisely when its kernel contains some \(\mathfrak n^r\): continuity at zero gives this condition, and translation gives continuity at every point. Algebraic formal smoothness supplies a lift \(c\) in the test. Continuity of \(b\) gives \(b(\mathfrak n^r)=0\) for some \(r\ge1\). Thus \(c(\mathfrak n^r)\subset J\), and multiplication of two such ideal powers gives \[
c(\mathfrak n^{2r})\subset J^2=0.
\] The lift is therefore continuous. The linear topology on \(R\) needs no further restriction beyond continuity of the stipulated maps.

Now assume both topologies are given by the stated ideal powers. If formal smoothness holds with discrete \(R\), any test with continuous \(a\) is among its tests, so the same property holds with the original topology on \(R\). In the other direction, start with a test having discrete \(R\); the map \(a\) is initially allowed to be arbitrary. Continuity of \(\varphi\) gives an integer \(s\ge1\) with \(\varphi(\mathfrak m^s)\subset\mathfrak n\). For the integer \(r\) supplied by \(b\) as above, compatibility implies \[
a(\mathfrak m^{sr})\subset J,
\qquad a(\mathfrak m^{2sr})=0.
\] Hence \(a\) was continuous for the \(\mathfrak m\)-power topology as well. Formal smoothness for that topology now gives the desired continuous lift. This proves both directions, with no finite-generation or Noetherian assumption. Neither completeness nor separatedness of \(S\) was required. ∎

\phantomsection\label{AG-RG-S05-more-algebra-lemma-formally-smooth-completion}

\subsubsection{Lemma F.107: formally smooth completion}\label{AG-RG-S05-lemma-f.107-formally-smooth-completion}

Reading source: \href{https://stacks.math.columbia.edu/tag/07ED}{Stacks, Tag 07ED}.

Let \(R\to S\) be a continuous map of Noetherian rings with topologies defined by finitely generated ideals \(\mathfrak m\subset R\) and \(\mathfrak n\subset S\). Write \(\widehat R\) and \(\widehat S\) for the corresponding completions. The following three properties are equivalent: \(R\to S\) is formally smooth for the \(\mathfrak n\)-power topology; \(R\to\widehat S\) is formally smooth for the \(\mathfrak n\widehat S\)-power topology; and \(\widehat R\to\widehat S\) is formally smooth for that target topology.

\textbf{Proof.} The continuous map induces the last completion map by applying it to compatible finite-order coordinates. \href{https://kokunoyumeto.github.io/open-math-courses/courses/AG-CA/completion.html}{AG-CA-19, Proposition 2.2} proves the exact quotient identities \[
\widehat R/\mathfrak m^q\widehat R=R/\mathfrak m^q,
\qquad
\widehat S/\mathfrak n^q\widehat S=S/\mathfrak n^q
\quad(q\ge1).
\] These identities use the canonical maps and commute with all transition maps. That proposition also proves completeness for the extended-ideal topologies.

A continuous homomorphism \(u:S\to D\) to a discrete ring kills some \(\mathfrak n^q\). Its factorization through \(S/\mathfrak n^q\) and the second displayed identity define a homomorphism \(\widehat u:\widehat S\to D\) killing \(\mathfrak n^q\widehat S\). It extends \(u\). It is the only continuous extension: any other extension kills a possibly different ideal power; pass to an exponent at least as large as both, where its map is determined by the same surjective finite quotient. Restriction and extension are inverse, and they commute with maps between discrete target rings. The identical argument applies to \(R\) and \(\widehat R\).

Use the tests defined in \hyperref[AG-RG-S05-more-algebra-lemma-formally-smooth]{F.106}, with \(D\to D/J\) square-zero. First compare \(R\to S\) with \(R\to\widehat S\). Keep the structural map from \(R\) fixed. The just-proved correspondence identifies every continuous map to \(D/J\), every continuous lift to \(D\), and their compatibility with \(R\). Thus one test has a solution exactly when its counterpart does.

Next compare \(R\to\widehat S\) with \(\widehat R\to\widehat S\). With the original topology on the base, extension and restriction identify the structural maps to both discrete test rings. Compatibility after restriction implies compatibility after extension: the two continuous maps from \(\widehat R\) agree by uniqueness of extension. The test diagrams and their dotted lifts therefore correspond in both directions. Finally F.106 identifies each of these topological-base tests with the tests using discrete base. In particular a structural map allowed in a discrete-base test is automatically continuous when it is compatible with the continuous target map modulo \(J\); that is precisely the ideal-power calculation in F.106. This proves the three equivalences in the stipulated Noetherian scope. ∎

\phantomsection\label{AG-RG-S05-more-algebra-lemma-gamma-commutative-diagram}

\subsubsection{Lemma F.108: gamma commutative diagram}\label{AG-RG-S05-lemma-f.108-gamma-commutative-diagram}

Reading source: \href{https://stacks.math.columbia.edu/tag/07E3}{Stacks, Tag 07E3}.

Consider a commutative square of fields with arrows \(k\to k'\), \(k\to K\), \(k'\to K'\) and \(K\to K'\). Assume \(k'/k\) and \(K'/K\) are finitely generated. For the natural maps \[
\alpha:\Omega_{K/k}\otimes_KK'\longrightarrow\Omega_{K'/k'},
\qquad
\beta:H_1(L_{K/k})\otimes_KK'\longrightarrow H_1(L_{K'/k'}),
\] all four kernels and cokernels are finite-dimensional, and \[
\dim\ker\alpha-\dim\operatorname{coker}\alpha
-\dim\ker\beta+\dim\operatorname{coker}\beta
=\operatorname{trdeg}_k k'-\operatorname{trdeg}_K K'.
\]

\textbf{Proof.} All vector spaces below are over \(K'\), and all omitted scalar extensions are to \(K'\). Apply the full field-transitivity proof \hyperref[AG-RG-S05-more-algebra-lemma-transitivity-gamma]{F.131} to \(k\subset k'\subset K'\) and \(k\subset K\subset K'\). Let \[
V=H_1(L_{K'/k}),\qquad U=\Omega_{K'/k}.
\] For the first row put \(P_1=H_1(L_{k'/k})\otimes_{k'}K'\), \(W_1=H_1(L_{K'/k'})\), \(Q_1=\Omega_{k'/k}\otimes_{k'}K'\) and \(Z_1=\Omega_{K'/k'}\). For the second put \(P_2=H_1(L_{K/k})\otimes_KK'\), \(W_2=H_1(L_{K'/K})\), \(Q_2=\Omega_{K/k}\otimes_KK'\) and \(Z_2=\Omega_{K'/K}\). The two exact rows are \[
0\longrightarrow P_i\longrightarrow V\longrightarrow W_i
\longrightarrow Q_i\longrightarrow U\longrightarrow Z_i\longrightarrow0
\qquad(i=1,2).
\] Their first injections identify \(P_1,P_2\) with subspaces of \(V\). Naturality of the representative and derivative maps in \hyperref[AG-RG-S05-algebra-lemma-exact-sequence-NL]{F.42} identifies \(\beta\) with \(P_2\to V\to W_1\) and \(\alpha\) with \(Q_2\to U\to Z_1\).

We calculate only finite defects, so no subtraction of infinite dimensions is involved. Set \[
T=P_1\cap P_2,\quad Y=V/(P_1+P_2),\quad
A_i=\ker(Q_i\to U),\quad B_i=\operatorname{im}(Q_i\to U).
\] Write \(i_0=\dim T\), \(r=\dim(P_1/T)\), \(y=\dim Y\), \(a_i=\dim A_i\), \(j=\dim(B_1\cap B_2)\), \(s=\dim(B_1/(B_1\cap B_2))\) and \(z=\dim(U/(B_1+B_2))\). These are finite. Indeed \hyperref[AG-RG-S05-more-algebra-lemma-cartier-equality]{F.98} makes \(P_1,Q_1,W_2,Z_2\) finite. The space \(Y\) is a quotient of \(V/P_2\), which embeds in \(W_2\); \(A_2\) is a quotient of \(W_2\); and the space measured by \(z\) is a quotient of \(Z_2\). All other spaces on the list are subspaces or quotients of \(P_1\) or \(Q_1\).

The kernel of \(\beta\) is \(T\). Exactness of the first row gives an injection \(V/P_1\to W_1\) with quotient \(A_1\). Dividing its image by the image of \(P_2\) therefore yields \[
0\longrightarrow Y\longrightarrow\operatorname{coker}\beta
\longrightarrow A_1\longrightarrow0.
\] For \(\alpha\), its kernel consists of the elements of \(Q_2\) whose image in \(U\) lies in \(B_1\). Mapping them into \(U\) gives a surjection onto \(B_1\cap B_2\) with kernel \(A_2\). Its cokernel is \(U/(B_1+B_2)\). Hence \[
\dim\ker\beta=i_0,\quad \dim\operatorname{coker}\beta=y+a_1,
\qquad \dim\ker\alpha=a_2+j,\quad \dim\operatorname{coker}\alpha=z.
\] The second row identifies \(V/P_2\) with the kernel of \(W_2\to Q_2\). Inside \(V/P_2\), the image of \(P_1\) is \(P_1/T\) and has quotient \(Y\). Similarly \(U/B_2\) contains the image \(B_1/(B_1\cap B_2)\) with quotient \(U/(B_1+B_2)\). The resulting finite exact sequences give \[
\dim P_1=i_0+r,\quad \dim W_2=r+y+a_2,
\qquad \dim Q_1=a_1+j+s,\quad \dim Z_2=s+z.
\] Subtracting these finite numbers shows that the expression in the statement equals \[
(a_2+j)-z-i_0+(y+a_1)
=\dim Q_1-\dim P_1-\dim Z_2+\dim W_2.
\] F.98 applied to \(k'/k\) and \(K'/K\) gives the required transcendence-degree difference. This also proves finiteness of every defect without requiring \(K/k\) or \(K'/k'\) to be finitely generated. ∎

\phantomsection\label{AG-RG-S05-more-algebra-lemma-geometrically-regular-over-field}

\subsubsection{Lemma F.109: geometrically regular over field}\label{AG-RG-S05-lemma-f.109-geometrically-regular-over-field}

\emph{Adapted from the Stacks project, more-algebra.tex, lines 9460--9501.}
\nopagebreak[4]

Let \(k\) be a field of characteristic \(p > 0\). Let \((A, \mathfrak m, K)\) be a Noetherian local \(k\)-algebra. Assume \(A\) is geometrically regular over \(k\). Let \(K/F/k\) be a finitely generated subextension. Let \(\varphi : k[y_1, \ldots, y_m] \to A\) be a \(k\)-algebra map such that \(y_i\) maps to an element of \(F\) in \(K\) and such that \(\text{d}y_1, \ldots, \text{d}y_m\) map to a basis of \(\Omega_{F/k}\). Set \(\mathfrak p = \varphi^{-1}(\mathfrak m)\). Then

\[
k[y_1, \ldots, y_m]_\mathfrak p \to A
\]

is flat and \(A/\mathfrak pA\) is regular.

\textbf{Proof.} Set \(A_0=k[y_1,\ldots,y_m]_{\mathfrak p}\), with maximal ideal \(\mathfrak m_0\). Its residue field \(K_0\) is the fraction field of the image \(k[\overline y_1,\ldots,\overline y_m]\subset F\), so \(K_0=k(\overline y_1,\ldots,\overline y_m)\). Polynomial differentials and localization make \(\Omega_{A_0/k}\) free on the \(dy_i\). These differentials span \(\Omega_{K_0/k}\) by the quotient and fraction rules. They are independent there: any \(K_0\)-linear relation maps to a relation in \(\Omega_{F/k}\), where their images are a basis by hypothesis. Thus \(\Omega_{A_0/k}\otimes K_0\to\Omega_{K_0/k}\) is an isomorphism.

The ring \(A_0\) is regular by the polynomial and localization theorems AG-CA-14, Proposition 3.3 and Theorem 3.2. After every finite purely inseparable \(k'/k\), its tensor product is again a localization of the polynomial ring \(k'[y_1,\ldots,y_m]\), hence regular by those same theorems. The proved field test F.50 therefore makes \(A_0\) geometrically regular. Proposition \hyperref[AG-RG-S05-more-algebra-proposition-characterization-geometrically-regular]{F.133} gives injectivity of \(H_1(L_{K_0/k})\to\mathfrak m_0/\mathfrak m_0^2\). The quotient sequence F.42 makes it surjective, since the following differential map is the isomorphism just proved. Hence it is an isomorphism.

Now consider

\[
\begin{gathered}
H_1(L_{K_0/k}) \otimes K \xrightarrow{} H_1(L_{K/k}) \\
H_1(L_{K_0/k}) \otimes K \xrightarrow{} \mathfrak m_0/\mathfrak m_0^2 \otimes K \\
\mathfrak m_0/\mathfrak m_0^2 \otimes K \xrightarrow{} \mathfrak m/\mathfrak m^2 \\
H_1(L_{K/k}) \xrightarrow{} \mathfrak m/\mathfrak m^2
\end{gathered}
\]

Since the left vertical arrow is injective by Lemma \hyperref[AG-RG-S05-more-algebra-lemma-transitivity-gamma]{F.131} and the lower horizontal by Proposition \hyperref[AG-RG-S05-more-algebra-proposition-characterization-geometrically-regular]{F.133} we conclude that the right vertical one is too. Hence a regular system of parameters in \(A_0\) maps to part of a regular system of parameters in \(A\). The images of the parameters can be completed to a regular system of parameters of \(A\), since their classes are independent in \(\mathfrak m/\mathfrak m^2\); lift a completing residue-field basis and use Nakayama's lemma F.10. AG-CA-12, Theorem 6.1, makes this parameter list regular, so \hyperref[AG-RG-S05-algebra-lemma-flat-over-regular]{F.46} proves flatness of \(A_0\to A\). For the asserted quotient regularity, apply the proved parameter-quotient theorem \href{https://kokunoyumeto.github.io/open-math-courses/courses/AG-CA/regular-local-rings.html}{AG-CA-14, Corollary 1.2} successively to the images of the parameters of \(A_0\). They generate \(\mathfrak pA=\mathfrak m_0A\), so the resulting regular quotient is exactly \(A/\mathfrak pA\). ∎

\phantomsection\label{AG-RG-S05-more-algebra-lemma-helper-G-ring}

\subsubsection{Lemma F.110: helper G ring}\label{AG-RG-S05-lemma-f.110-helper-g-ring}

\emph{Adapted from the Stacks project, more-algebra.tex, lines 12915--12963.}
\nopagebreak[4]

Let \(A_0\) be a regular complete Noetherian local domain and \(P=A_0[Y_1,\ldots,Y_m]\). For every prime \(q\) of \(P\), both \(\widehat{P_q}\otimes_{A_0}\operatorname{Frac}(A_0)\) and \(\widehat{P_q}\otimes_P\operatorname{Frac}(P)\) are geometrically regular over their indicated fields. The second assertion includes the original power-series helper; the first is the extra consumed form that removes a general formal-smoothness theorem from the completion proof.

\textbf{Proof.} The polynomial ring \(P\) is regular by AG-CA-14, Proposition 3.3, and its localization is regular. Completion preserves the cotangent dimension and Krull dimension by AG-CA-19, Theorem 4.1, so \(\widehat{P_q}\) is regular. Both displayed rings are localizations of it and therefore regular. In characteristic zero the finite purely inseparable test is trivial, so they are geometrically regular by that criterion.

In positive characteristic Cohen structure AG-CA-19S, Theorem 6.1, identifies \(A_0\) with a power-series ring over its residue field: its minimal parameter list gives a surjection from that power-series ring, equality of dimensions and the nonzero-prime dimension drop AG-CA-11, Theorem 3.2, make its kernel zero. This is also the regular complete specialization already proved above. Treat either indicated fraction field \(F\), and induct on a finite purely inseparable degree \([L:F]\). At a step \(F\subset M\subset L\), \(L=M(z)\) with \(z^p=f\notin M^p\). Let \(Q\) denote \(A_0\) for the first assertion and \(P\) for the second. Choose a finite \(Q\)-subalgebra \(B\subset M\) with fraction field \(M\) and \(f\in B\). To justify this choice, scale each finite algebraic field generator so it satisfies a monic equation over \(Q\), and scale \(z\) by an element of \(Q\) to put its \(p\)th power \(f\) in that algebra; this does not alter its generated field. Here is the denominator step: write \(f=b/d\) with nonzero \(d\in B\). A monic integral equation for \(d\) may be shortened by cancelling powers of \(d\) until its constant coefficient \(a\in Q\) is nonzero, since \(B\) is a domain. That equation gives \(a/d\in B\). Replacing \(z\) by \(az\) changes its \(p\)th power to \(a^pf=b\,a^{p-1}(a/d)\in B\), as required. The finitely generated integral algebra is finite. Put \(C=B[Z]/(Z^p-f)\subset L\); it is a finite domain with fraction field \(L\).

For the first assertion use the finite extension \(P\subset B[Y]\); for the second use \(P\subset B\). The base \(Q\) is normal: it is regular Noetherian, and AG-CA-15, Corollary 4.5, proves its normality. If \(b\in B\), some \(b^{p^r}\) lies in \(\operatorname{Frac}(Q)\) by pure inseparability and is integral over \(Q\) because \(b\) is. Normality puts that power in \(Q\). The same applies to \(C\). These extensions therefore have every element's sufficiently high \(p\)th power in the smaller ring. Each prime has exactly one prime above it. Indeed membership of an element in a prime is equivalent to membership of that power in its contracted prime; lying over supplies existence. The finite-extension completion proof already written identifies the tensor of \(\widehat{P_q}\) with the finite overring with the completion of its unique local factor. The same holds for \(C\) (with the polynomial \(Y\) variables in the first assertion). Hence the ring at the \(M\) stage is a localization of this completed \(B\)-ring, and the ring at the \(L\) stage is exactly its quotient extension \(E[Z]/(Z^p-f)\). Finite-module completion commutes with this monic finite algebra by AG-CA-19, Theorem 3.1; this is the needed equality rather than an assumed equality of unrelated completions.

The induction hypothesis says \(E\) is regular. The derivation lemma above applies to \(B\), which is finite type over the complete local ring \(A_0\) in either case, and gives \(D:B\to B\) with \(D(f)\ne0\). Extend it coefficientwise over the added \(Y\) variables, then through the prime localization, completion and the localization to \(M\), using the fully checked derivation extension formula. Its value on \(f\) remains the nonzero element \(D(f)\) of the field \(M\), hence a unit of \(E\). The hypersurface derivation proof makes \(E[Z]/(Z^p-f)\) regular. This completes the induction for every finite purely inseparable extension, and the criterion gives both geometric-regularity assertions. ∎

\phantomsection\label{AG-RG-S05-more-algebra-lemma-helper-integral}

\subsubsection{Lemma F.111: helper integral}\label{AG-RG-S05-lemma-f.111-helper-integral}

Reading source: \href{https://stacks.math.columbia.edu/tag/09XG}{Stacks, Tag 09XG}.

Let \(A\to B\) be an integral ring map, \(I\subset A\) an ideal, and \(y\in B\) have idempotent image in \(B/IB\). There are an integer \(d\geq1\) and a monic polynomial \(f(T)\in A[T]\) such that \[
f(y)=0,\qquad f(T)\bmod I=T^d(T-1)^d.
\]

\textbf{Proof.} Write \(z=y^2-y=\sum_{i=1}^h a_ib_i\) with \(a_i\in I\) and \(b_i\in B\). This is a finite sum by the definition of the extended ideal. The subalgebra \(C=A[y,b_1,\ldots,b_h]\subset B\) is a finite \(A\)-module: each of its finitely many generators satisfies a monic equation over \(A\), and successive substitution reduces every monomial to one whose exponent in each generator is below that equation's degree. Those finitely many bounded monomials generate \(C\). This argument uses the image of \(A\); injectivity of \(A\to B\) is unnecessary.

Choose a finite module generating list \(c_1,\ldots,c_d\) for \(C\), including \(1\), so \(d\geq1\). Since \(zC\subset IC\), express each \(zc_i\) as \(\sum_j m_{ij}c_j\) with all \(m_{ij}\in I\). For the column \(c=(c_i)\) these equations say \((z1_d-M)c=0\), where \(M=(m_{ij})\). Multiplying by the adjugate, whose identity was proved completely in \hyperref[AG-RG-S05-algebra-lemma-charpoly]{F.19}, gives \[
g(z)c_i=0\quad\hbox{for all }i,
\qquad g(U)=\det(U1_d-M)\in A[U].
\] The included generator \(1\) makes \(g(z)=0\) in \(C\), hence in \(B\). The determinant expansion makes \(g\) monic of degree \(d\); reducing the entries of \(M\) modulo \(I\) gives \(g(U)\bmod I=U^d\). Set \(f(T)=g(T^2-T)\). Its leading coefficient is \(1\), it annihilates \(y\), and its reduction is \((T^2-T)^d=T^d(T-1)^d\). The same construction permits an empty sum for \(z\), the zero target algebra, and \(I=A\). ∎

\phantomsection\label{AG-RG-S05-more-algebra-lemma-intersection-subfields}

\subsubsection{Lemma F.112: intersection subfields}\label{AG-RG-S05-lemma-f.112-intersection-subfields}

\emph{Adapted from the Stacks project, more-algebra.tex, lines 12067--12128.}
\nopagebreak[4]

Let \(K\) have characteristic \(p\) and downward directed subfields \(F_\lambda\) containing \(K^p\) with intersection \(K^p\). Then the maps \(\Omega_{K/\mathbb F_p}\to\Omega_{K/F_\lambda}\) have zero common kernel. For every finite \(L/K\), the fields \(L^pF_\lambda\) have intersection \(L^p\).

\textbf{Proof.} A nonzero absolute differential is a finite linear combination of differentials of distinct absolute \(p\)-basis elements. If it vanishes over every \(F_\lambda\), those finitely many elements have dependent relative differentials, hence their monomials have a nonzero relation over every \(F_\lambda\) by the preceding \(p\)-basis proof. Apply the finite-subspace intersection lemma to the space of all monomial coefficient relations inside \(K^{p^n}\). It gives a nonzero relation over \(K^p\), contradicting the absolute \(p\)-basis property. This proves the first assertion.

For the second, finite extensions decompose into towers of finite separable simple extensions and degree-\(p\) purely inseparable simple extensions by the repeated-minimal-polynomial and separable-subfield construction given above. It suffices to treat one step and iterate: if the assertion holds for an intermediate \(M\), then the fields \(M^pF_\lambda\) satisfy the same three assumptions inside \(M\).

For a separable simple extension \(L=K(\theta)\) choose its generator in \(L^p\). This is possible because \(K(\theta^p)=K(\theta)\): the intervening extension is purely inseparable and also separable, so trivial. Frobenius identifies \([L^p:K^p]=[L:K]=d\). Thus both \(L^p\) and each \(L^pF_\lambda\) have the common basis \(1,\theta,\ldots,\theta^{d-1}\) over \(K^p\) and \(F_\lambda\), respectively. Independence over \(F_\lambda\) follows from independence over the larger field \(K\). Uniqueness of coefficient expansions makes the intersection precisely the vectors with all coefficients in \(K^p\), namely \(L^p\).

For a degree-\(p\) step \(L=K(\theta)\), \(\theta^p=t\notin K^p\), one has \(L^p=K^p(t)\). Choose an index with \(t\notin F_{\lambda_0}\); it exists by the intersection assumption. On the cofinal family of subfields refined inside \(F_{\lambda_0}\), the elements \(1,t,\ldots,t^{p-1}\) are independent by the same degree-\(p\) argument. Thus \(L^pF_\lambda=F_\lambda(t)\) has this basis. The unique expansion coefficients of an element in the intersection lie in the intersection of that cofinal family, which is \(K^p\), proving it lies in \(K^p(t)\). The opposite inclusion is automatic. This cofinal restriction is necessary: at other indices \(t\) may already belong to \(F_\lambda\), so a degree-\(p\) basis cannot be assumed there. ∎

\phantomsection\label{AG-RG-S05-more-algebra-lemma-intersection-subfields-subspace}

\subsubsection{Lemma F.113: intersection subfields subspace}\label{AG-RG-S05-lemma-f.113-intersection-subfields-subspace}

\emph{Adapted from the Stacks project, more-algebra.tex, lines 12027--12065.}
\nopagebreak[4]

If subfields \(F_\lambda\subset K\) are downward directed and have intersection \(F\), then a \(K\)-linear subspace \(V\subset K^n\) meets every \(F_\lambda^n\) nontrivially if and only if it meets \(F^n\) nontrivially.

\textbf{Proof.} The reverse implication is immediate. Induct on \(n\) for the other. If \(V\cap(0\oplus F^{n-1})\ne0\), the conclusion already holds. Otherwise induction supplies some \(\lambda_0\) for which \(V\cap(0\oplus F_{\lambda_0}^{n-1})=0\). The nonzero intersection with \(F_{\lambda_0}^n\) contains a vector with nonzero first coordinate, which can be normalized to \(v=(1,v_2,\ldots,v_n)\). Subtracting the first coordinate times \(v\) shows that this intersection is exactly the line \(F_{\lambda_0}v\). For every refined subfield inside \(F_{\lambda_0}\), the nonzero intersection similarly has a normalized vector; its uniqueness in the larger line makes it the same \(v\). Thus all its coordinates belong to every refined subfield. Their intersection is \(F\), because downward directedness refines \(\lambda_0\) with each arbitrary index. Hence \(v\in F^n\). The case \(n=1\) is immediate, completing the induction. ∎

\phantomsection\label{AG-RG-S05-more-algebra-lemma-lci-local}

\subsubsection{Lemma F.114: lci local}\label{AG-RG-S05-lemma-f.114-lci-local}

Reading source: \href{https://stacks.math.columbia.edu/tag/07D1}{Stacks, Tag 07D1}.

Let \(R\) be Noetherian and let \(S\) be a finitely presented \(R\)-algebra. A kernel of a finite polynomial presentation of \(S\) is a regular ideal if and only if this property holds for polynomial presentations on a finite principal cover of \(\operatorname{Spec}S\). The condition does not depend on the chosen presentations.

\textbf{Proof.} The arbitrary-ring comparison \hyperref[AG-RG-S05-more-algebra-lemma-lci-presentation-independent]{F.94.A} proves presentation independence for Koszul-regular kernels. The finite polynomial rings and their localizations are Noetherian by \href{https://kokunoyumeto.github.io/open-math-courses/courses/AG-CA/noetherian-and-artinian-rings.html}{AG-CA-03, Theorem 2.1}; its Theorem 1.1 and Proposition 1.2 make their finite-module submodules finite. In these rings, \hyperref[AG-RG-S05-more-algebra-lemma-noetherian-finite-all-equivalent]{F.118} converts a Koszul-regular generating list in a local maximal ideal into a regular list. To obtain the neighbourhood condition in \hyperref[AG-RG-S05-more-algebra-definition-regular-ideal]{F.95}, we will explicitly spread the successive injections, rather than infer global regularity merely from Koszul acyclicity.

Fix \(P=R[X_1,\ldots,X_n]\twoheadrightarrow S\) with kernel \(I\), and a prime \(\mathfrak q\supset I\) of \(P\). Suppose a chart \(S_g\) containing its image has a regular presentation. Choose \(u\in P\) lifting \(g\), so \(u\notin\mathfrak q\). A presentation of that chart is \[
Q=P[Z]\twoheadrightarrow S_g,\qquad H=(I,t),\quad t=uZ-1.
\] The prime \(\mathfrak q'\) of \(Q\) given by \(\mathfrak q\) and \(Z=1/u\) satisfies \[
Q_{\mathfrak q'}/(t)=P_{\mathfrak q}.
\] Indeed quotienting first gives \(P_u\), and then localization at the prime over \(\mathfrak q\) gives \(P_{\mathfrak q}\). Presentation comparison and F.118 give a regular generating list for \(H_{\mathfrak q'}\). Its conormal classes are a basis by \hyperref[AG-RG-S05-algebra-lemma-regular-quasi-regular]{F.74}. Meanwhile \hyperref[AG-RG-S05-algebra-lemma-principal-localization-NL]{F.72} gives, at this point, \[
(H/H^2)_{\mathfrak q'}
\cong (I/I^2)_{\mathfrak q}\oplus S_{\mathfrak q/I}\,t.
\] The first summand is a finite direct summand of a finite free module. It is flat and finitely presented, hence free over the local quotient by \href{https://kokunoyumeto.github.io/open-math-courses/courses/AG-RG/AG-RG-S01.html}{S01, Lemma 1.2}. Choose \(f_1,\ldots,f_c\in I\) whose classes are a basis there; denominators may be cleared since their values outside \(\mathfrak q\) are units. The list \(t,f_1,\ldots,f_c\) is a conormal basis. Nakayama applied to \(H_{\mathfrak q'}/(t,f_1,\ldots,f_c)\) makes it a generator list for \(H_{\mathfrak q'}\).

Compare it with the regular generator list just obtained. Both have the same length, equal to the conormal rank. Express the new list in the old one. Modulo the maximal ideal the coefficient matrix takes one basis of the conormal fibre to the other, so its determinant is a unit in \(Q_{\mathfrak q'}\). Its exterior powers identify the Koszul complexes: the differential on a wedge deletes one entry with its alternating sign, and this formula commutes with the change of basis. Thus \(t,f_1,\ldots,f_c\) has zero first homology. F.118 makes this very order regular, with \(t\) first. Quotienting by \(t\) proves that \(f_1,\ldots,f_c\) is a regular list generating \(I_{\mathfrak q}\).

Here is the required principal neighbourhood. The module \(I/(f_1,\ldots,f_c)\) is finite and vanishes at \(\mathfrak q\). For each \(i\), the kernel of multiplication by \(f_i\) on \(P/(f_1,\ldots,f_{i-1})\) is also finite, because \(P\) is Noetherian; it vanishes at \(\mathfrak q\) by the just-proved regularity. Exact localization and the finite-generator fraction-zero argument \href{https://kokunoyumeto.github.io/open-math-courses/courses/AG-RG/AG-RG-S01.html}{S01, Lemma 1.1} give one element \(h\notin\mathfrak q\) annihilating all these finitely many modules. In \(P_h\) the list therefore generates \(I_h\) and every successive multiplication is injective. Its final quotient is nonzero because the prime \(\mathfrak qP_h\) still contains \(I_h\). The list is regular on this actual neighbourhood, as required by F.95.

The argument applies at every point lying in one of the stipulated charts, hence at every prime above \(I\). Conversely, localization preserves the successive injections of a regular list. For a chart where a given neighbourhood also inverts \(u\), present the localization by adjoining \(Z\) and put \(uZ-1\) first. Multiplication by this polynomial is injective by its lowest-coefficient calculation, and its quotient is \(P_u\); the remaining localized list is regular wherever the quotient is nonzero. Charts with zero quotient have no primes to check. Thus the kernel condition restricts to every principal chart. The local neighbourhoods cover the target, and a finite subcover exists by the finite unit-ideal calculation \href{https://kokunoyumeto.github.io/open-math-courses/courses/AG-CA/spectra-of-rings.html}{AG-CA-01, Proposition 2.3}. Applying the same comparison to a second initial presentation proves the claimed independence. If \(S=0\), all these prime conditions are vacuous, including the empty-cover case. ∎

\phantomsection\label{AG-RG-S05-more-algebra-lemma-lift-factorization-monic}

\subsubsection{Lemma F.115: lift factorization monic}\label{AG-RG-S05-lemma-f.115-lift-factorization-monic}

Reading source: \href{https://stacks.math.columbia.edu/tag/0ALH}{Stacks, Tag 0ALH}.

Let \(I\subset A\) be any ideal. Suppose the reduction of a monic \(f\in A[T]\) has a factorization \(\bar f=\bar G\bar H\) into coprime monic polynomials over \(A/I\). There is a finitely presented étale map \(A\to A'\) such that its canonical map \(A/I\to A'/IA'\) is an isomorphism and \(f=G'H'\) in \(A'[T]\), with monic factors reducing to the prescribed ones.

\textbf{Proof.} If \(I=A\), take \(A'=0\), the principal localization at zero; its closed fibre and the prescribed factors are zero, so all requirements hold. Assume \(A/I\ne0\). Put \(d=\deg\bar G\) and \(e=\deg\bar H\). Introduce \(d+e\) variables for the lower coefficients of monic polynomials \(G\) and \(H\) of those degrees. The \(d+e\) lower coefficients of \(GH-f\) are equations in these variables. Write \(C_0\) for their quotient algebra and \(\Delta\) for the determinant of their square Jacobian matrix.

Over \(D=A/I\), the derivative of multiplication at the given factors is \[
D[T]_{<d}\oplus D[T]_{<e}\longrightarrow D[T]_{<d+e},
\qquad (v,w)\longmapsto\bar H v+\bar G w.
\] Here \(D[T]_{<r}\) is the free module on \(1,T,\ldots,T^{r-1}\). Choose a Bezout identity \(a\bar G+b\bar H=1\). For a polynomial \(q\) in the right module, let \(v\) be the remainder of \(bq\) on monic division by \(\bar G\). Then \(q-\bar H v\) is divisible by \(\bar G\); its quotient \(w\) has degree below \(e\). This constructs a preimage of \(q\). If the image of \((v,w)\) is zero, reduction modulo \(\bar G\) and invertibility of \(\bar H\) there give \(v=0\), by its degree bound, and monic multiplication then gives \(w=0\). Thus this is an isomorphism even when \(D\) has zero divisors. Degree-zero factors are covered by empty polynomial modules.

Consequently \(\Delta\) has a unit value at the prescribed coefficient tuple. Set \(C=(C_0)_\Delta\). The square Jacobian calculation \href{https://kokunoyumeto.github.io/open-math-courses/courses/AG-CA/formally-smooth-unramified-and-etale-ring-maps.html}{AG-CA-17, Theorem 4.1} gives formal smoothness, and there are no remaining differential coordinates. Theorem 2.2 of that lesson therefore gives formal unramifiedness. The finite equations and principal localization give finite presentation, so \(C/A\) is étale. The coefficient tuple defines a \(D\)-algebra retraction \(\sigma:E=C/IC\to D\) of its structural map.

We select exactly this section in the closed fibre. Let \(L=\ker\sigma\). Finite algebra generators \(x_i\) of \(E/D\) give finite ideal generators \(x_i-\sigma(x_i)\) for \(L\): replace each variable by its constant value in a polynomial and telescope the difference one variable at a time. The map \[
\delta:E\longrightarrow L/L^2,\qquad
x\longmapsto x-\sigma(x)\pmod {L^2}
\] is a \(D\)-derivation, with \(E\) acting on \(L/L^2\) through \(\sigma\). To check products, expand \(x=\sigma(x)+\ell\) and \(y=\sigma(y)+m\) and discard \(\ell m\in L^2\). Since \(E/D\) is étale by base change, its differentials vanish by the same Theorem 2.2. The universal differential property makes \(\delta=0\). Its restriction to \(L\) is the quotient map, so \(L=L^2\).

Choose finite ideal generators \(\ell_i\) and write \(\ell_i=\sum_j m_{ij}\ell_j\) with \(m_{ij}\in L\). The adjugate identity of \href{https://kokunoyumeto.github.io/open-math-courses/courses/AG-RG/AG-RG-S01.html}{S01, Lemma 1.2} shows that \(a=\det(1-M)\) annihilates \(L\), while \(a\equiv1\pmod L\). Put \(p=1-a\in L\). Then \(p\ell=\ell\) for every \(\ell\in L\), so \(p^2=p\) and \(L=pE\). Inverting \(1-p\) kills precisely \(L\) and induces \(E_{1-p}\cong E/L=D\), compatibly with its structural map. Lift \(1-p\) to \(h\in C\) and put \(A'=C_h\). Its closed fibre is this very copy of \(D\), and principal localization preserves finite presentation and étaleness by AG-CA-17, Theorem 1.2 and Section 4. The universal polynomials in \(A'[T]\) supply \(G',H'\). If both degrees are zero, one can directly take \(A'=A\) and both factors equal to one. ∎

\phantomsection\label{AG-RG-S05-more-algebra-lemma-lift-idempotent-upstairs}

\subsubsection{Lemma F.116: lift idempotent upstairs}\label{AG-RG-S05-lemma-f.116-lift-idempotent-upstairs}

Reading source: \href{https://stacks.math.columbia.edu/tag/07M4}{Stacks, Tag 07M4}.

Suppose \(A\to B\) is integral and \(I\subset A\) is an ideal. Every idempotent \(\bar e\in B/IB\) lifts to an idempotent of \(B\otimes_A A'\) after a finitely presented étale map \(A\to A'\) for which \(A/I\to A'/IA'\) is an isomorphism.

\textbf{Proof.} Choose \(y\in B\) reducing to \(\bar e\). The determinant proof \hyperref[AG-RG-S05-more-algebra-lemma-helper-integral]{F.111} gives an integer \(d\ge1\) and a monic \(f\in A[T]\) with \[
f(y)=0,\qquad \bar f=T^d(T-1)^d.
\] The two factors generate the unit ideal: in the expansion of \((T+(1-T))^{2d-1}=1\), every term contains either \(T^d\) or \((1-T)^d\). Apply \hyperref[AG-RG-S05-more-algebra-lemma-lift-factorization-monic]{F.115}. Over its étale algebra \(A_1\) we have \(f=GH\) with reductions \(T^d\) and \((T-1)^d\). Put \(B_1=B\otimes_AA_1\), write \(y_1\) for the image of \(y\), and set \(b_1=G(y_1)\), \(b_2=H(y_1)\). Then \[
b_1b_2=0,\qquad
\bar b_1=\bar e,\quad \bar b_2=(-1)^d(1-\bar e).
\] Thus \(b_1,b_2\) generate the unit ideal modulo \(IB_1\).

The map \(A_1\to B_1\) remains integral. Indeed each tensor element uses finitely many elements of \(B\); their monic equations generate a finite module algebra after base change, and its elements are integral by the determinant test \href{https://kokunoyumeto.github.io/open-math-courses/courses/AG-CA/integral-extensions-lying-over-going-up-and-going-down.html}{AG-CA-05, Theorems 1.1--1.2}. Let \(K\) be the inverse image in \(A_1\) of the ideal \((b_1,b_2)B_1\). The closed-image and lying-over proof in Theorem 3.4 of that lesson gives \[
\operatorname{image}V_{B_1}(b_1,b_2)=V_{A_1}(K).
\] No prime in this image can contain \(IA_1\), since its prime upstairs would contain \(IB_1\) as well as \(b_1,b_2\). Hence \(V(K+IA_1)\) is empty and \(K+IA_1=A_1\), by the prime-separation proof \href{https://kokunoyumeto.github.io/open-math-courses/courses/AG-CA/spectra-of-rings.html}{AG-CA-01, Lemma 1.1}. Choose \(a\in K\) with \(a\equiv1\pmod {IA_1}\) and set \(A'=(A_1)_a\).

In \(B'=B_1\otimes_{A_1}A'\), the image of \(a\) is a unit belonging to \((b_1,b_2)\). Choose \(u,v\in B'\) with \(ub_1+vb_2=1\) and put \(e'=ub_1\). The equations give \[
e'(1-e')=uvb_1b_2=0,
\] so \(e'\) is idempotent. Reducing \(ub_1+vb_2=1\) and multiplying by \(\bar e\) gives \(\bar u\bar e=\bar e\), proving that \(e'\) reduces to the specified idempotent. Finally \(a\) reduces to one, so the last localization leaves the quotient \(A_1/IA_1=A/I\) unchanged. Its composition with \(A\to A_1\) is finitely presented étale by the proved stability in AG-CA-17, Theorem 1.2 and Section 4. This supplies every stated property, including the zero closed fibre. ∎

\phantomsection\label{AG-RG-S05-more-algebra-lemma-lift-projective-module}

\subsubsection{Lemma F.117: lift projective module}\label{AG-RG-S05-lemma-f.117-lift-projective-module}

Reading source: \href{https://stacks.math.columbia.edu/tag/07M5}{Stacks, Tag 07M5}.

For an ideal \(I\) of any ring \(A\) and a finite projective \(A/I\)-module \(\bar P\), there is a finitely presented étale map \(A\to A'\) inducing \(A/I\cong A'/IA'\), and a finite projective \(A'\)-module \(P'\) with \(P'/IP'\cong\bar P\) under that identification.

\textbf{Proof.} Realize \(\bar P\) as a direct summand of \((A/I)^n\) and let \(\bar p\) be its projector. Such a realization follows from projectivity by splitting a finite free surjection, as in \href{https://kokunoyumeto.github.io/open-math-courses/courses/AG-RG/AG-RG-S01.html}{S01, Lemma 1.0}. Lift the finitely many entries of \(\bar p\) to a matrix \(\varphi\) over \(A\). Its characteristic polynomial \(f(T)=\det(T1-\varphi)\) is monic. The full coefficient proof of Cayley--Hamilton \hyperref[AG-RG-S05-algebra-lemma-charpoly]{F.19} gives an algebra map \[
B=A[T]/(f)\longrightarrow\operatorname{End}_A(A^n),
\qquad x\longmapsto\varphi,
\] where \(x\) is the class of \(T\). Monic division makes \(B\) a finite free \(A\)-module, so \(A\to B\) is integral.

In each residue field of \(A/I\), an idempotent linear operator decomposes the vector space into its image and kernel: \(v=\bar p v+(1-\bar p)v\), and applying \(\bar p\) proves the intersection zero. Choose bases of the two summands. Its characteristic polynomial there is \(T^{n-r}(T-1)^r\) for the dimension \(r\) of its image. Consequently at every prime of \(B/IB\) the image of \(x\) is zero or one: pass to the field of that prime and use this factorization. The nilradical calculation \href{https://kokunoyumeto.github.io/open-math-courses/courses/AG-CA/spectra-of-rings.html}{AG-CA-01, Theorem 1.2} then gives an integer \(N\ge1\) such that \[
x^N(1-x)^N=0\quad\text{in }E=B/IB.
\] This uses a single nilpotent element, with no Noetherian condition.

Set \(a=x^N\), \(b=(1-x)^N\) in \(E\). They generate the unit ideal: expand \((x+(1-x))^{2N-1}\), whose every term is divisible by \(a\) or \(b\). Also \(ab=0\). Choose \(u,v\) with \(ua+vb=1\). Then \((a+b)(u^2a+v^2b)=u^2a^2+v^2b^2=(ua+vb)^2=1\). Thus \(a+b\) is a unit, and \[
\epsilon=\frac{a}{a+b},\qquad
\epsilon(1-\epsilon)=\frac{ab}{(a+b)^2}=0
\] is an idempotent. Under the reduction of the displayed endomorphism map, \(a\) goes to \(\bar p^N=\bar p\), \(b\) to \((1-\bar p)^N=1-\bar p\), and their sum to the identity. Hence \(\epsilon\) maps to exactly \(\bar p\).

Apply \hyperref[AG-RG-S05-more-algebra-lemma-lift-idempotent-upstairs]{F.116} to the integral algebra \(B/A\) and this idempotent of \(B/IB\). It supplies the required étale \(A'\) and an idempotent of \(B\otimes_AA'\) lifting \(\epsilon\). Its image in \(\operatorname{End}_{A'}((A')^n)\) is a projector \(p'\) reducing to \(\bar p\). Put \(P'=\operatorname{im}p'\). The decomposition \[
(A')^n=\operatorname{im}p'\oplus\operatorname{im}(1-p')
\] follows by writing each vector as the sum of its two projections; applying \(p'\) makes their intersection zero. The projections of the finite coordinate basis generate each summand, so \(P'\) is finite projective. Reduction preserves this split decomposition and identifies its first summand with \(\operatorname{im}\bar p=\bar P\). If \(n=0\), take \(A'=A\) and \(P'=0\) directly. This also handles \(A/I=0\) by choosing its zero module with \(n=0\). ∎

\phantomsection\label{AG-RG-S05-more-algebra-lemma-noetherian-finite-all-equivalent}

\subsubsection{Lemma F.118: noetherian finite all equivalent}\label{AG-RG-S05-lemma-f.118-noetherian-finite-all-equivalent}

Reading source: \href{https://stacks.math.columbia.edu/tag/09CC}{Stacks, Tag 09CC}, together with the regular-sequence and Koszul arguments cited there. The proof below is independently written for finite modules.

For a nonzero finite module over a Noetherian local ring and a list in its maximal ideal, regularity, vanishing of all positive Koszul homology, vanishing of first Koszul homology, and the polynomial associated-graded condition are equivalent.

\textbf{Proof.} Write \((A,\mathfrak m)\) for the ring, \(M\) for the module, and \(J=(f_1,\ldots,f_r)\). The fourth assertion means that the natural graded map \[\theta:(M/JM)[T_1,\ldots,T_r]\longrightarrow\bigoplus_{n\geq0}J^nM/J^{n+1}M,\qquad \bar mT^a\longmapsto[mf^a],\] is an isomorphism. It is always surjective. If \(r=0\), the Koszul complex has only \(M\) in degree zero, \(J=0\), and \(\theta\) is the identity in degree zero. The empty list is regular because \(M\ne0\). Thus all four assertions hold in that case; assume \(r>0\).

Use the exterior complex on the symbols \(e_1,\ldots,e_r\), with coefficients in \(M\), differential \(d(e_i)=f_i\), and the signed product rule. Separating wedges which contain \(e_r\) identifies it with the cone of multiplication by \(f_r\) on the complex for the first \(r-1\) entries. Its homology sequence contains \[H_1(K_{r-1})\xrightarrow{f_r}H_1(K_{r-1})\longrightarrow H_1(K_r)\longrightarrow H_0(K_{r-1})\xrightarrow{f_r}H_0(K_{r-1}).\] For the two consequences needed here, exactness can also be checked directly. A degree-one cycle of \(K_{r-1}\) becomes a boundary in \(K_r\) when \(H_1(K_r)=0\). Writing that boundary as the sum of a wedge without \(e_r\) and a wedge containing \(e_r\) shows that its class in \(H_1(K_{r-1})\) is \(f_r\) times a class represented by a cycle. This gives surjectivity of the first arrow. If \(f_rm\in(f_1,\ldots,f_{r-1})M\), choose coefficients for this equality to obtain a degree-one cycle in \(K_r\) with last coefficient \(m\). A degree-two boundary has last coefficient in \((f_1,\ldots,f_{r-1})M\), so \([m]=0\) in \(H_0(K_{r-1})\). This gives injectivity of the last arrow.

If \(H_1(K_r)=0\), the preceding \(H_1\) is a finite \(A\)-module, since its complex consists of finite modules over a Noetherian ring. Nakayama, with \(f_r\in\mathfrak m\), makes it zero by the surjectivity just proved. Induction therefore gives all successive multiplication injections, including the final one displayed above. Their last quotient is nonzero by Nakayama applied to \(M/JM\). This proves regularity. Conversely, assume the list regular and the preceding complex acyclic in positive degrees. A cycle in \(K_r\) of degree \(k\) has a component \(v\) multiplying \(e_r\), and \(v\) is a cycle of degree \(k-1\) in \(K_{r-1}\). For \(k\geq2\), preceding acyclicity makes \(v\) a boundary. For \(k=1\), the cycle identity says \(f_r[v]=0\) in \(H_0(K_{r-1})\); regularity again makes \(v\) a boundary. Subtract the differential of a suitable wedge containing \(e_r\) to remove \(v\). What remains is a degree-\(k\) cycle in \(K_{r-1}\) and is a boundary by preceding acyclicity. This proves acyclicity in every positive degree by induction. In particular it implies vanishing of \(H_1\).

We next prove that a regular list makes \(\theta\) injective. In degree zero this is the defining quotient \(M/JM\). In degree one, suppose \(\sum_i f_im_i\in J^2M\). Express this element as \(\sum_i f_iu_i\) with \(u_i\in JM\), and subtract. The vector \((m_i-u_i)\) is a degree-one Koszul cycle. By the preceding acyclicity it is a boundary, whose coefficients all belong to \(JM\), since the degree-two differential has coefficients \(f_i\). Consequently every \(m_i\in JM\). This proves directly the module version of the conormal calculation in \hyperref[AG-RG-S05-algebra-lemma-regular-quasi-regular]{F.74}.

Here is a coefficient calculation for any degree \(n\geq1\). In a relation \[\sum_{|a|=n}m_af^a\in J^{n+1}M,\] collect every term involving one of \(f_1,\ldots,f_{r-1}\) into their generated submodule. The only remaining term on the left is \(m_{(0,\ldots,0,n)}f_r^n\); the only remaining type on the right is \(u f_r^{n+1}\). Thus \[\bigl(m_{(0,\ldots,0,n)}-uf_r\bigr)f_r^n+\sum_{i<r}f_iv_i=0\] for suitable \(u,v_i\in M\). The list \(f_1,\ldots,f_{r-1},f_r^n\) is regular by AG-CA-12, Proposition 1.3, whose power argument applies to modules. Its first Koszul homology is zero by the cone argument. Hence the last coefficient of this relation belongs to \((f_1,\ldots,f_{r-1})M\), so \(m_{(0,\ldots,0,n)}\in JM\).

To recover the coefficients of all monomials, put \(B=A[U_1,\ldots,U_r,U_r^{-1}]\), \(N=M\otimes_AB\), and \[g_i=f_i-(U_i/U_r)f_r\quad(i<r),\qquad g_r=f_r/U_r.\] The ring \(B\) is free and faithfully flat over \(A\), with its Laurent monomials as a basis. Tensoring preserves Koszul acyclicity. The displayed change of generators is invertible, so its exterior powers identify the two Koszul complexes and give \(H_1(K(g;N))=0\).

Fix a prime \(\mathfrak p\) of \(B\). If \(N_{\mathfrak p}=0\), every term of every Koszul complex with these coefficients is zero. If \(N_{\mathfrak p}\ne0\) and \(JB\subset\mathfrak p\), the local criterion already proved makes \(g_1,\ldots,g_r\) regular on the nonzero finite \(B_{\mathfrak p}\)-module \(N_{\mathfrak p}\). Its power list \(g_1,\ldots,g_{r-1},g_r^n\) is regular and has zero \(H_1\). Finally, if \(JB\not\subset\mathfrak p\), at least one \(g_i\) is a unit in \(B_{\mathfrak p}\). The power list still contains a unit, so choose \(a_i\) with \(\sum_{i<r}a_ig_i+a_rg_r^n=1\). Wedge multiplication by \(\sum_i a_ie_i\) satisfies \(dh+hd=1\) by the signed product rule, and contracts its Koszul complex. These three cases show that the power list has zero first homology at every prime. Localization commutes with the complex and its homology; a module whose localizations at all primes vanish is zero. The global \(H_1\) therefore vanishes as well.

The preceding pure-last-coefficient calculation now applies to \(g\) on \(N\). Substituting \(f_i=g_i+U_ig_r\) for \(i<r\) and \(f_r=U_rg_r\) makes the coefficient of \(g_r^n\) equal to \(\sum_{|a|=n}m_aU^a\). Thus this element belongs to \(JN\). Since \(N/JN=(M/JM)\otimes_AB\) is the direct sum of its Laurent monomial copies of \(M/JM\), each coefficient \(m_a\) belongs to \(JM\). This proves injectivity of \(\theta\) in every degree.

For the converse, suppose \(\theta\) is an isomorphism. Krull intersection, AG-CA-03, Theorem 6.1, gives \(\bigcap_{n\geq0}J^nM=0\). Every nonzero \(m\in M\) therefore has a finite order \(v\), characterized by \(m\in J^vM\setminus J^{v+1}M\). Its initial class is nonzero; multiplication by \(T_1\) is injective on \((M/JM)[T_1,\ldots,T_r]\), so \(f_1m\) is nonzero and has order \(v+1\). In particular \(f_1\) acts injectively on \(M\). The same calculation gives \[f_1M\cap J^nM=f_1J^{n-1}M\qquad(n\geq1):\] if \(f_1m\in J^nM\) and \(m\ne0\), its order forces \(m\in J^{n-1}M\); the reverse inclusion follows from \(f_1\in J\).

Set \(M'=M/f_1M\) and let \(J'\) be the ideal induced by \(J\). In degree \(n\geq1\), the last equality identifies the kernel of \(J^nM/J^{n+1}M\to J'^nM'/J'^{n+1}M'\) with multiplication by \(T_1\) from degree \(n-1\). Indeed an element in that kernel is \(y+f_1m\), with \(y\in J^{n+1}M\); then \(f_1m\in J^nM\), so \(m\in J^{n-1}M\). Degree zero is unchanged. It follows that \[\operatorname{gr}_{J'}M'\simeq(M/JM)[T_2,\ldots,T_r].\] Repeat the order and quotient argument for the remaining entries. Nakayama ensures that each finite intermediate module and the final quotient are nonzero. This proves regularity and hence all four equivalences. ∎

\phantomsection\label{AG-RG-S05-more-algebra-lemma-p-basis}

\subsubsection{Lemma F.119: p basis}\label{AG-RG-S05-lemma-f.119-p-basis}

\emph{Adapted from the Stacks project, more-algebra.tex, lines 11972--12025.}
\nopagebreak[4]

For a characteristic-\(p\) extension \(L/F\), a subset is \(p\)-independent over \(F\) if its monomials of exponents below \(p\) are independent over \(FL^p\). This is equivalent to independence of its relative differentials. Every such subset extends to a \(p\)-basis, whose differentials are a basis of \(\Omega_{L/F}\).

\textbf{Proof.} Put \(E=FL^p\). A chain of \(p\)-independent subsets has independent union, because every monomial relation uses only finitely many elements and coefficients. Zorn's lemma therefore extends any such subset to a maximal one \(B\). The subalgebra \(E[B]\) is a field: every element belongs to a subalgebra on finitely many basis elements, which is a finite-dimensional domain over \(E\) and hence a field. Here the last assertion follows directly: multiplication by a nonzero element is an injective linear map of a finite-dimensional space, hence surjective, so that element has an inverse.

If \(x\in L\) lies outside \(E[B]\), then \(x^p\in E\). The polynomial \(T^p-x^p\) is irreducible over \(E[B]\). Inside \(L[T]\) it is \((T-x)^p\). In a factorization of this polynomial, evaluation at \(T=x\) makes at least one factor vanish; division of that factor by \(T-x\) and induction on the exponent show that every monic factor is a power of \(T-x\). A monic proper factor over \(E[B]\) would therefore be \((T-x)^r\) with \(1\le r<p\); its coefficient of \(T^{r-1}\) is \(-rx\), so \(x\) would belong to \(E[B]\), since \(r\) is invertible in characteristic \(p\). Thus adjoining \(x\) has degree \(p\) and preserves monomial independence, contradicting maximality. Consequently \(E[B]=L\). Reducing exponents by the monic relations, and then using monomial independence, proves the presentation \[
L=E[T_b:b\in B]/(T_b^p-b^p:b\in B).
\] All expressions use finitely many variables, so this argument also covers infinite \(B\).

The product and inverse rules make every \(F\)-derivation, including the universal one, vanish on \(E\). In the displayed presentation, monomial differentiation defines a partial \(F\)-derivation for each \(b\in B\): it kills the coefficient field \(E\) and all the defining relations. These derivations send the basis elements to their respective coordinate values, proving that the \(db\) are independent. Differentiating the finite monomial expansion of each element proves that they span \(\Omega_{L/F}\). In particular, every \(p\)-independent subset has independent differentials after extending it to \(B\).

For the converse, suppose a finite list \(db_1,\ldots,db_r\) is independent. Here is the required extension of its coordinate functionals, including when \(\Omega_{L/F}\) is infinite-dimensional. Prescribe a coordinate functional on their span. Order pairs consisting of a larger subspace and a linear functional extending this prescription by extension. A chain has an upper bound obtained by taking the union of its subspaces and compatible functionals. Zorn's lemma gives a maximal pair. If its subspace omits a vector \(v\), extend to its direct sum with \(Lv\) by assigning value zero to \(v\); independence of \(v\) makes this extension well-defined, contradicting maximality. The maximal subspace is therefore all of \(\Omega_{L/F}\). By the universal differential property, these extended coordinate functionals give \(F\)-derivations \(D_i:L\to L\) with \(D_i(b_j)=\delta_{ij}\), and they kill \(E\).

If the \(b_j\) have a nonzero monomial relation over \(E\) with all exponents below \(p\), choose one of least total degree. Its degree is positive, since a nonzero constant cannot vanish in a field. Some partial derivative of its polynomial is nonzero: choose a variable having positive exponent in a nonconstant term; that exponent is between one and \(p-1\), and different monomials give different derivative monomials. This derivative has smaller total degree. Applying the corresponding \(D_i\) to the original relation makes its evaluation zero, contradicting minimality. Thus the finite list is \(p\)-independent. For an arbitrary subset, every relation would involve a finite list, so the same conclusion holds. This proves the converse, the equivalence, and the basis assertion in the full stated field scope. ∎

\phantomsection\label{AG-RG-S05-more-algebra-lemma-power-series-ring-subfields}

\subsubsection{Lemma F.120: power series ring subfields}\label{AG-RG-S05-lemma-f.120-power-series-ring-subfields}

\emph{Adapted from the Stacks project, more-algebra.tex, lines 12130--12171.}
\nopagebreak[4]

For \(A=k[[X_1,\ldots,X_n]][Y_1,\ldots,Y_m]\) in characteristic \(p\), choose a \(p\)-basis \(B\) of \(k\). For each finite subset \(J\subset B\), let \(k_J=k^p(B\setminus J)\) and \(A_J=k_J[[X_1^p,\ldots,X_n^p]][Y_1^p,\ldots,Y_m^p]\). Then \(A\) is finite free over \(A_J\) and the fraction fields \(K_J\) are downward directed, contain \(K^p\), and intersect in \(K^p\).

\textbf{Proof.} The monomials in the finite set \(J\), in the \(X_i\), and in the \(Y_j\), all with exponents below \(p\), form a basis of \(A\) over \(A_J\). Group power-series terms by the \(X\) exponents modulo \(p\), polynomial terms by the \(Y\) exponents modulo \(p\), and each coefficient by the finite basis of \(k/k_J\); uniqueness of each expansion proves both spanning and independence. This proves finite freeness directly, including imperfect \(k\).

Every \(p\)th power lies in each \(A_J\), hence \(K^p\subset K_J\), and adjoining more excluded basis elements shrinks \(K_J\), proving directedness. To compute the intersection, write \(u=f/g^p\) with \(f,g\in A\), \(g\ne0\); every fraction admits this form by multiplying its numerator by \(g^{p-1}\). If \(u\in K_J\), write it as \(a/b^p\) with \(a\in A_J\) and nonzero \(b\in A\): the inverse of any denominator \(h\in A_J\) is \(h^{p-1}/h^p\), so this form is available. Then \(b^pf=ag^p\in A_J\). Any \(A_J\)-derivation of \(A\) therefore kills \(f\), since differentiating gives \(b^pD(f)=0\) in the domain. The power-series and polynomial partial derivatives imply every nonzero exponent of \(f\) in the \(X,Y\) variables is divisible by \(p\). The partial derivations on the excluded coefficient \(p\)-basis elements, extended coefficientwise to the power series, imply its coefficients lie in \(k_J\): their monomial basis expansion shows that their common kernel is \(k_J\). For every finite \(J\) this holds, so each coefficient lies in \(\bigcap_Jk_J=k^p\), again by its finite \(p\)-basis monomial expansion. Thus \(f\in A^p\) (take coefficient roots and divide all exponents by \(p\)), and \(u\in K^p\). This proves the intersection assertion. ∎

\phantomsection\label{AG-RG-S05-more-algebra-lemma-quasi-regular-ideal-finite-projective}

\subsubsection{Lemma F.121: quasi regular ideal finite projective}\label{AG-RG-S05-lemma-f.121-quasi-regular-ideal-finite-projective}

\emph{Adapted from the Stacks project, more-algebra.tex, lines 8769--8778.}
\nopagebreak[4]

Let \(I \subset R\) be a quasi-regular ideal of a ring. Then \(I/I^2\) is a finite projective \(R/I\)-module.

\textbf{Proof.} Put \(S=R/I\) and \(M=I/I^2\). If \(S=0\), then \(I=R\) and \(M=0\), which is finite projective. Otherwise Definition F.95 gives, around each prime of \(S\), a principal open \(D(\bar g)\) where \(I_g\) is generated by a quasi-regular list. The degree-one isomorphism in that definition and exact localization identify \(M_{\bar g}\) with a finite free \(S_{\bar g}\)-module. The elements \(\bar g\) generate the unit ideal: a proper ideal would lie in a prime missing every one of these opens, contrary to their cover. Hence finitely many \(g_1,\ldots,g_s\in S\) suffice, with \((g_1,\ldots,g_s)=S\).

We use the following finite-cover detection calculation for arbitrary modules. If \(N_{g_i}=0\) for every \(i\) and \(x\in N\), then \(g_i^{e_i}x=0\) for some positive \(e_i\). The powers \(g_i^{e_i}\) still generate the unit ideal, since a prime containing all of them would contain every \(g_i\). Writing \(1=\sum b_i g_i^{e_i}\) gives \(x=0\), so \(N=0\).

For each local basis of \(M_{g_i}\), clear its finitely many denominators to choose elements of \(M\) whose localized images are a basis; multiplying basis elements by powers of the unit \(g_i\) preserves that property. Let \(M_0\subset M\) be generated by all these finitely many elements. The quotient \(M/M_0\) vanishes on every open, so the detection calculation gives \(M=M_0\). Choose a finite free surjection \(F=S^n\to M\) and let \(K\) be its kernel. On each open the quotient \(M_{g_i}\) is free, so the localized surjection splits. Thus \(K_{g_i}\) is a direct summand of \(S_{g_i}^n\), generated by the projections of its finite basis. Clear denominators in those kernel generators to obtain finitely many elements of \(K\); the submodule they generate has the same localization as \(K\) on every open. Apply the same detection calculation to their quotient to conclude that \(K\) is finitely generated. Therefore \(M\) is finitely presented.

At every prime one of the \(g_i\) is invertible, and \(M\) localizes there to a free module. The proved local-flatness criterion \href{https://kokunoyumeto.github.io/open-math-courses/courses/AG-CA/tor-and-flat-modules.html}{AG-CA-07, Theorem 3.3} makes \(M\) flat. Its finite presentation now supplies the full hypothesis of \href{https://kokunoyumeto.github.io/open-math-courses/courses/AG-CA/tor-and-flat-modules.html}{AG-CA-07, Theorem 5.3}, giving finite projectivity. No Noetherian assumption, or finite presentation of the initial relation module, was used. ∎

\phantomsection\label{AG-RG-S05-more-algebra-lemma-quotient-regular}

\subsubsection{Lemma F.122: quotient regular}\label{AG-RG-S05-lemma-f.122-quotient-regular}

\emph{Adapted from the Stacks project, more-algebra.tex, lines 12407--12424.}
\nopagebreak[4]

If \(R\) is Noetherian regular and a derivation \(D:R\to R\) sends \(f\) to a unit modulo \((f)\), then \(R/(f)\) is regular.

\textbf{Proof.} At a prime containing \(f\), extend \(D\) to the local ring by the fraction calculation above. If \(f\in\mathfrak m^2\), Leibniz on a finite expression \(f=\sum a_ib_i\), \(a_i,b_i\in\mathfrak m\), gives \(D(f)\in\mathfrak m\), contradicting that its class is a unit. Hence \(f\notin\mathfrak m^2\). In the regular local ring its class extends to a regular parameter system, and AG-CA-14, Corollary 1.2, gives regularity of the parameter quotient. These quotients are exactly all local rings of \(R/(f)\). ∎

\phantomsection\label{AG-RG-S05-more-algebra-lemma-regular-base-change}

\subsubsection{Lemma F.123: regular base change}\label{AG-RG-S05-lemma-f.123-regular-base-change}

\emph{Adapted from the Stacks project, more-algebra.tex, lines 11022--11047.}
\nopagebreak[4]

Let \(R \to \Lambda\) be a regular ring map. For any finite type ring map \(R \to R'\) the base change \(R' \to \Lambda \otimes_R R'\) is regular too.

\textbf{Proof.} Flatness survives base change by AG-CA-07, Proposition 3.2. Fix \(\mathfrak p'\subset R'\) with contraction \(\mathfrak p\subset R\). Set \(k=\kappa(\mathfrak p)\) and \(E=\kappa(\mathfrak p')\). The field extension \(E/k\) is finitely generated, as follows directly from the finite-type hypothesis. Choose algebra generators \(u_1,\ldots,u_n\) of \(R'/R\). The injection \(R/\mathfrak p\to R'/\mathfrak p'\) makes \[(R\setminus\mathfrak p)^{-1}(R'/\mathfrak p')=k[\bar u_1,\ldots,\bar u_n]\] a domain with fraction field \(E\). Thus \(E=k(\bar u_1,\ldots,\bar u_n)\).

The original fibre \(F=\Lambda\otimes_Rk\) is Noetherian and geometrically regular. Tensor associativity identifies the new fibre with \[(\Lambda\otimes_RR')\otimes_{R'}E\cong F\otimes_kE.\] It is Noetherian by \hyperref[AG-RG-S05-algebra-lemma-Noetherian-field-extension]{F.12}, using the field generation just proved. Its geometric regularity over \(E\) is the full base-field conclusion of \hyperref[AG-RG-S05-algebra-lemma-geometrically-regular-over-subfields]{F.52}. Flatness and these fibre assertions are precisely the regularity conditions. No assumption that \(E/k\) is algebraic or separable has been added. ∎

\phantomsection\label{AG-RG-S05-more-algebra-lemma-regular-composition}

\subsubsection{Lemma F.124: regular composition}\label{AG-RG-S05-lemma-f.124-regular-composition}

\emph{Adapted from the Stacks project, more-algebra.tex, lines 11049--11065.}
\nopagebreak[4]

A composition of regular maps of Noetherian rings is regular.

\textbf{Proof.} Let \(A\to B\to C\) be the maps. Their composite is flat by AG-CA-07, Proposition 3.2. For \(\mathfrak a\subset A\), put \(k_0=\kappa(\mathfrak a)\) and \(S=A\setminus\mathfrak a\). Both original fibres are Noetherian: for example \[C\otimes_Ak_0=S^{-1}(C/\mathfrak aC).\] The right side is a localization of a quotient of the Noetherian ring \(C\). Thus the Noetherian composite-fibre condition is automatic here.

Let \(K/k_0\) be finite purely inseparable, and set \(D=B\otimes_AK\) and \(E=C\otimes_AK\). Lemma \hyperref[AG-RG-S05-algebra-lemma-Noetherian-field-extension]{F.12} makes them Noetherian, since each is the finite field tensor of its original Noetherian fibre. Regularity of \(A\to B\) makes \(D\) regular. The identification \(E=C\otimes_BD\) and the same flat base-change proposition make \(D\to E\) flat.

Fix \(\mathfrak r\subset E\) and let \(\mathfrak s\subset D\) and \(\mathfrak b\subset B\) be its successive contractions. The contraction of \(\mathfrak b\) to \(A\) is \(\mathfrak a\): elements of \(\mathfrak a\) vanish in \(D\), whereas every element outside \(\mathfrak a\) becomes invertible through \(K\). There is a field inclusion \(k_0\subset\kappa(\mathfrak b)\). Tensoring \(D\) with this residue field gives \[Q=D\otimes_B\kappa(\mathfrak b)\cong\kappa(\mathfrak b)\otimes_{k_0}K.\] A finite \(k_0\)-basis of \(K\) makes \(Q\) a finite-dimensional \(\kappa(\mathfrak b)\)-algebra. The prime \(\mathfrak s\) descends to a prime \(\overline{\mathfrak s}\) of \(Q\): first quotient by \(\mathfrak bD\) and then invert \(B\setminus\mathfrak b\). Its quotient \(Q/\overline{\mathfrak s}\) is a finite-dimensional domain over \(\kappa(\mathfrak b)\), hence a field. Indeed, multiplication by a nonzero element is injective and therefore surjective, giving its inverse. Localizing this quotient at its nonzero elements changes nothing. The quotient-localization construction therefore identifies it with \(\kappa(\mathfrak s)\), proving \[[\kappa(\mathfrak s):\kappa(\mathfrak b)]<\infty.\]

The fibre of \(D\to E\) at \(\mathfrak s\) is \[E\otimes_D\kappa(\mathfrak s)\cong(C\otimes_B\kappa(\mathfrak b))\otimes_{\kappa(\mathfrak b)}\kappa(\mathfrak s).\] The original \(B\)-fibre is Noetherian and geometrically regular. Its displayed finite field tensor is Noetherian by F.12 and regular by \hyperref[AG-RG-S05-algebra-lemma-geometrically-regular-over-subfields]{F.52}. The closed fibre of the flat local map \(D_{\mathfrak s}\to E_{\mathfrak r}\) is a localization of this regular ring, hence regular. The source is a regular local ring and both local rings are Noetherian. \hyperref[AG-RG-S05-algebra-lemma-flat-over-regular-with-regular-fibre]{F.47} makes \(E_{\mathfrak r}\) regular. This holds for every target prime, so \(E\) is regular. Criterion \hyperref[AG-RG-S05-algebra-lemma-geometrically-regular]{F.50}, applied to every finite purely inseparable \(K/k_0\), proves geometric regularity of the composite fibre. Doing this at every \(\mathfrak a\) proves the composition assertion. Empty fibres have no target primes and satisfy the same criterion. ∎

\phantomsection\label{AG-RG-S05-more-algebra-lemma-regular-permanence}

\subsubsection{Lemma F.125: regular permanence}\label{AG-RG-S05-lemma-f.125-regular-permanence}

\emph{Adapted from the Stacks project, more-algebra.tex, lines 11116--11129.}
\nopagebreak[4]

If \(A\to B\to C\) are maps of Noetherian rings, \(A\to C\) is regular and \(B\to C\) faithfully flat, then \(A\to B\) is regular.

\textbf{Proof.} For an injection \(N_1\to N_2\) of \(A\)-modules, put \(K=\ker(N_1\otimes_AB\to N_2\otimes_AB)\). Flatness of \(C\) over \(B\) preserves the exact sequence \[0\longrightarrow K\longrightarrow N_1\otimes_AB\longrightarrow N_2\otimes_AB.\] Tensor associativity identifies the last map after tensoring with \(C\) with \(N_1\otimes_AC\to N_2\otimes_AC\). It is injective by \(A\)-flatness of \(C\). Consequently \(K\otimes_BC=0\), and faithful flatness makes \(K=0\). This proves that \(B\) is flat over \(A\).

For \(\mathfrak a\subset A\), set \(k=\kappa(\mathfrak a)\), \(F=B\otimes_Ak\) and \(G=C\otimes_Ak\). These rings are Noetherian: they are respectively localizations of \(B/\mathfrak aB\) and \(C/\mathfrak aC\) at the image of \(A\setminus\mathfrak a\). The identification \(G=F\otimes_BC\) proves faithful flatness of \(F\to G\) directly. For every \(F\)-module \(M\), \[M\otimes_FG\cong M\otimes_BC.\] Restriction of scalars is exact, while the right-hand tensor functor is exact and detects the zero module. The left-hand functor therefore has those properties too. The target \(G\) is geometrically regular over \(k\), since \(A\to C\) is regular. Geometric regularity descent \hyperref[AG-RG-S05-algebra-lemma-geometrically-regular-descent]{F.51} applies to these Noetherian \(k\)-algebras and makes \(F\) geometrically regular. These are all the fibre conditions for \(A\to B\), completing the proof. ∎

\phantomsection\label{AG-RG-S05-more-algebra-lemma-relative-global-complete-intersection-koszul}

\subsubsection{Lemma F.126: relative global complete intersection koszul}\label{AG-RG-S05-lemma-f.126-relative-global-complete-intersection-koszul}

Reading source: \href{https://stacks.math.columbia.edu/tag/07D2}{Stacks, Tag 07D2}. The following proof gives the Noetherian form used in this lesson.

Let \(R\) be Noetherian, \(P=R[X_1,\ldots,X_n]\), and \(S=P/(f_1,\ldots,f_c)\). Suppose each nonempty fibre of \(R\to S\) has dimension \(n-c\), as in Definition F.4. The Koszul complex \(K_P(f_1,\ldots,f_c)\) has zero homology in every positive degree.

\textbf{Proof.} Let \(I=(f_1,\ldots,f_c)\) and fix a prime \(\mathfrak p\) of \(P\). If \(I\subset\mathfrak p\), the complete relative conormal argument \hyperref[AG-RG-S05-algebra-lemma-relative-global-complete-intersection-conormal]{F.77} makes this ordered list regular in \(P_{\mathfrak p}\). The cone induction in \hyperref[AG-RG-S05-more-algebra-lemma-noetherian-finite-all-equivalent]{F.118} then makes its localized Koszul complex acyclic in positive degrees. That implication only needs the successive multiplication injections; it imposes no assertion of global ordered regularity outside \(V(I)\).

If \(I\not\subset\mathfrak p\), write \(\sum_i a_i f_i=1\) in \(P_{\mathfrak p}\). On the exterior complex with basis \(e_i\) and differential \(d(e_i)=f_i\), set \(h(v)=(\sum_i a_i e_i)\wedge v\). The signed product rule gives \[d h(v)+h d(v)=\Big(\sum_i a_i f_i\Big)v=v.\] Thus every localized cycle is a boundary in this case too. Localization is exact and therefore commutes with kernels, images and homology, by \href{https://kokunoyumeto.github.io/open-math-courses/courses/AG-RG/AG-RG-S01.html\#1-the-algebra-of-localization}{AG-RG-S01, Lemma 1.1}. Each positive homology module now has zero localization at every prime. The local-detection proof in that same lemma makes it zero. When \(c=0\), there are no positive complex terms. When \(S=0\), every prime is in the contracting-homotopy case. Both cases satisfy the assertion. ∎

\phantomsection\label{AG-RG-S05-more-algebra-lemma-relative-regular-immersion-algebra}

\subsubsection{Lemma F.127: relative regular immersion algebra}\label{AG-RG-S05-lemma-f.127-relative-regular-immersion-algebra}

Reading source: \href{https://stacks.math.columbia.edu/tag/0CEP}{Stacks, Tag 0CEP}. All tensor steps are specified below.

Take ring maps \(A\to B\) and \(A\to A'\), a finite list \(f_1,\ldots,f_r\in B\), and \(J=(f_1,\ldots,f_r)\). Assume \(B/J\) is flat over \(A\). Put \(B'=B\otimes_A A'\), let \(f_i'\) be the images, and put \(J'=(f_1',\ldots,f_r')\). If the original list is quasi-regular, so is the new list. If its first Koszul homology vanishes, the new first Koszul homology vanishes. No flatness of \(B\) is assumed.

\textbf{Proof.} We use the exact quotient-flatness statement \href{https://kokunoyumeto.github.io/open-math-courses/courses/AG-CA/tor-and-flat-modules.html}{AG-CA-07, Theorem 6.1}: a short exact sequence with flat quotient remains short exact under any tensor functor, and an extension of two flat modules is flat.

In the quasi-regular case Definition \hyperref[AG-RG-S05-more-algebra-definition-regular-ideal]{F.95} identifies \[J^n/J^{n+1}\simeq (B/J)^{\oplus\{a\in\mathbf N^r:|a|=n\}}.\] These modules are \(A\)-flat. Starting with \(B/J\), the sequences \[0\to J^n/J^{n+1}\to B/J^{n+1}\to B/J^n\to0\] therefore prove by induction that \(B/J^n\) is flat for every \(n\geq1\). Apply quotient-flatness to \(0\to J^n\to B\to B/J^n\to0\). Its tensor identifies \(J^n\otimes_A A'\) with its image in \(B'\). That image is exactly \((J')^n\): the finite products of the \(f_i\) generate \(J^n\), and their images generate \((J')^n\). Right exactness applied to \(J^{n+1}\to J^n\to J^n/J^{n+1}\to0\) now gives \[ (J')^n/(J')^{n+1}\simeq (J^n/J^{n+1})\otimes_A A'. \] In degree zero it gives \(B'/J'=(B/J)\otimes_A A'\). The identifications commute with multiplication, since every one is induced by the images of the same products of generators. They carry each \(T_i\) to \(f_i'\) and prove the required polynomial graded-ring identity.

For the first-homology assertion, put \(Z=\ker(B^r\xrightarrow{(f_i)}J)\). Quotient-flatness first gives an injection \(J\otimes_A A'\hookrightarrow B'\), whose image is \(J'\). Right exactness of \[Z\longrightarrow B^r\longrightarrow J\longrightarrow0\] then makes \(Z\otimes_A A'\) surject onto \(Z'=\ker((B')^r\xrightarrow{(f_i')}B')\). The original Koszul boundaries are the image of \(\bigwedge^2 B^r\to Z\). If \(H_1=0\), this map is surjective. Tensoring preserves its surjectivity, and the natural identification \((\bigwedge^2 B^r)\otimes_A A'=\bigwedge^2(B')^r\) carries its alternating boundary formula to the new formula. Consequently every element of \(Z'\) is a new boundary, as required. For an empty list \(J=0\) and all first complex terms vanish; both conclusions still hold. Unit ideals and zero rings cause no change to these exact-sequence arguments. ∎

\phantomsection\label{AG-RG-S05-more-algebra-lemma-symmetric-algebra-smooth}

\subsubsection{Lemma F.128: symmetric algebra smooth}\label{AG-RG-S05-lemma-f.128-symmetric-algebra-smooth}

Reading source: \href{https://stacks.math.columbia.edu/tag/07M6}{Stacks, Tag 07M6}. The universal-property and lifting argument is written here for arbitrary modules and rings.

For an \(A\)-module \(M\), its symmetric algebra \(C=\operatorname{Sym}_A(M)\) is smooth over \(A\) precisely when \(M\) is finite projective. The base-change, direct-sum and free-module identities used in D.10 hold without finiteness hypotheses.

\textbf{Proof.} Start with \(\bigoplus_{d\geq0}M^{\otimes_A d}\), where degree zero is \(A\). Concatenation of tensors defines its product: the balancing relations defining a tensor product are preserved by concatenation, and associativity holds on pure tensors, which generate. Quotient by the homogeneous relations \(m\otimes n-n\otimes m\). Interchanging adjacent factors makes the quotient commutative. This constructs \(C\), with \(C_0=A\), \(C_1=M\), and all its pieces generated by products from degree one. Projection \(\epsilon:C\to A\) is an algebra map. If \(J=\ker\epsilon\), the homogeneous degree calculation gives \(J/J^2=M\).

For a commutative \(A\)-algebra \(D\), an \(A\)-linear map \(u:M\to D\) extends by the formula \[m_1\otimes\cdots\otimes m_d\longmapsto u(m_1)\cdots u(m_d).\] It respects balancing and the displayed commutation relations, hence defines an algebra map \(C\to D\). Every algebra map has this restriction and is determined by it. These constructions are inverse, proving the universal property.

That property gives actual mutually inverse maps \[\begin{aligned}
B\otimes_A\operatorname{Sym}_A(M)&\simeq\operatorname{Sym}_B(B\otimes_A M),\\
\operatorname{Sym}_A(M\oplus N)&\simeq\operatorname{Sym}_A(M)\otimes_A\operatorname{Sym}_A(N),\\
\operatorname{Sym}_A(A^n)&\simeq A[X_1,\ldots,X_n].
\end{aligned}\] In the first pair of maps send \(b\otimes m\) to \(b\otimes m\), and send \(m\) back to \(1\otimes m\); the tensor-algebra coproduct calculation in \href{https://kokunoyumeto.github.io/open-math-courses/courses/AG-RG/AG-RG-S01.html\#1-the-algebra-of-localization}{AG-RG-S01, Lemma 1.0} combines the latter map with the base map from \(B\). For the second, send \((m,n)\) to \(m\otimes1+1\otimes n\); the inverse comes from the two inclusions into \(M\oplus N\) and the same coproduct formula. For the last, send basis vectors to variables and variables back to basis vectors. Each composite fixes the base ring and every degree-one generator, so is identity. The maps preserve degree. These verifications prove the identities, including reduction modulo a base ideal, used later in the lesson.

Suppose first \(C\) is smooth. It is finitely presented, hence has algebra generators \(z_1,\ldots,z_s\). The differences \(y_i=z_i-\epsilon(z_i)\) generate \(J\): for a polynomial \(P\), successive subtraction in each variable expresses \(P(z)-P(\epsilon(z))\) as a sum of multiples of those differences. For an element of \(J\) the second term is zero. Thus \(M=J/J^2\) is finite. Choose a finite free surjection \(\pi:F\to M\).

For any module \(V\), form the square-zero algebra \(A\oplus V\), with multiplication \((a,v)(b,w)=(ab,aw+bv)\). The map \(A\oplus F\to A\oplus M\) induced by \(\pi\) is surjective with square-zero kernel. The linear map \(m\mapsto(0,m)\) extends to an algebra map \(C\to A\oplus M\). The formal-smoothness comparison proved in \href{https://kokunoyumeto.github.io/open-math-courses/courses/AG-CA/formally-smooth-unramified-and-etale-ring-maps.html}{AG-CA-17, Theorem 3.1 and Section 4} gives a lift to \(A\oplus F\). Its first coordinate on \(m\) is zero, because reduction leaves that coordinate unchanged. Write its second coordinate as \(s(m)\). Restriction of the algebra map makes \(s\) linear and \(\pi s=1_M\). Hence \[F=s(M)\oplus\ker\pi,\] as the decomposition \(f=s(\pi f)+(f-s(\pi f))\) and application of \(\pi\) verify. This is exactly finite projectivity of \(M\).

Conversely write \(F=A^n=M\oplus K\) with maps \(i:M\to F\) and \(\pi:F\to M\) satisfying \(\pi i=1_M\). The projections of a basis of \(F\) generate \(K\). By the universal property, \(C\) is the quotient of \(A[X_1,\ldots,X_n]\) by the linear forms corresponding to \(K\): maps from that quotient are precisely linear maps \(F\to D\) killing \(K\), and hence precisely linear maps \(M\to D\). The generator maps give the inverse algebra identifications. Since \(K\) has the finite generating list just exhibited, this is a finite presentation of \(C\).

Now take a square-zero surjection \(D\to D/H\) of \(A\)-algebras and a map \(C\to D/H\). Its restriction is a linear \(u:M\to D/H\). Lift the finitely many values of \(u\pi\) on the free basis of \(F\) to elements of \(D\). The resulting linear map \(v:F\to D\) reduces to \(u\pi\), so \(vi\) lifts \(u\). Its algebra extension lifts the original map from \(C\). Thus \(C\) is formally smooth; its finite presentation and the same proved formal/conormal comparison make it smooth. The constructions include \(M=0\), \(n=0\), and the zero ring. ∎

\phantomsection\label{AG-RG-S05-more-algebra-lemma-syntomic-lci}

\subsubsection{Lemma F.129: syntomic lci}\label{AG-RG-S05-lemma-f.129-syntomic-lci}

Reading source: \href{https://stacks.math.columbia.edu/tag/07D3}{Stacks, Tag 07D3}. We use the Noetherian-base version and specify the syntomic conditions directly.

Let \(R\) be Noetherian. A map \(R\to S\) is syntomic if and only if it is flat and a local complete intersection homomorphism in Definition \hyperref[AG-RG-S05-more-algebra-definition-local-complete-intersection]{F.94}. Here syntomic means flat, finitely presented, with local complete intersection fibres as in Definition F.2.

\textbf{Proof.} If \(R\to S\) is syntomic, the complete principal-neighbourhood construction \hyperref[AG-RG-S05-algebra-lemma-syntomic]{F.87} covers \(\operatorname{Spec}S\) by charts that have relative global complete intersection presentations. On each chart \hyperref[AG-RG-S05-more-algebra-lemma-relative-global-complete-intersection-koszul]{F.126} makes the presenting equations Koszul-regular. Thus each chart is a local complete intersection homomorphism. A finite collection of the principal charts suffices, and \hyperref[AG-RG-S05-more-algebra-lemma-lci-local]{F.114} proves their local complete intersection property descends to \(R\to S\). Flatness was one of the syntomic conditions.

Conversely suppose \(R\to S\) is flat and a local complete intersection. The finite-kernel and arbitrary-presentation proof \hyperref[AG-RG-S05-more-algebra-lemma-lci-presentation-independent]{F.94.A} makes it finitely presented. Choose \(S=P/I\) with \(P=R[X_1,\ldots,X_n]\). For a field-valued base map \(R\to k\), put \(P'=k[X_1,\ldots,X_n]\) and \(S'=S\otimes_R k=P'/I'\), with \(I'\) the image of \(I\).

Fix a prime \(\mathfrak q'\) of \(S'\) and its inverse image \(\mathfrak p'\) in \(P'\). Contract to a prime \(\mathfrak p\) of \(P\). Definition \hyperref[AG-RG-S05-more-algebra-definition-regular-ideal]{F.95} and the presentation comparison give \(g\notin\mathfrak p\) such that \(I_g\) has a finite Koszul-regular generating list \(f_1,\ldots,f_r\). Its first Koszul homology is zero. The quotient \(P_g/I_g=S_g\) is \(R\)-flat, since localization is flat and flat maps compose. Apply \hyperref[AG-RG-S05-more-algebra-lemma-relative-regular-immersion-algebra]{F.127} to \(R\to P_g\) and \(R\to k\). It proves first-homology vanishing for the image list in \(P'_g\); this list generates \(I'_g\) by tensor right exactness. The image of \(g\) is outside \(\mathfrak p'\) by contraction. Localize at \(\mathfrak p'\). This is a Noetherian polynomial local ring, the list lies in its maximal ideal, and \hyperref[AG-RG-S05-more-algebra-lemma-noetherian-finite-all-equivalent]{F.118} makes it regular there. The full field complete-intersection neighbourhood argument \hyperref[AG-RG-S05-algebra-lemma-lci]{F.57} therefore gives a global complete intersection chart of \(S'\) around \(\mathfrak q'\). Since every prime of every fibre has such a chart, the fibres meet Definition F.2. Together with flatness and finite presentation this proves syntomicity. If a fibre or \(S\) is zero there are no primes requiring charts; the same conditions and argument remain valid. ∎

\phantomsection\label{AG-RG-S05-more-algebra-lemma-transitive-lci-at-end}

\subsubsection{Lemma F.130: transitive lci at end}\label{AG-RG-S05-lemma-f.130-transitive-lci-at-end}

Reading source: \href{https://stacks.math.columbia.edu/tag/07D4}{Stacks, Tag 07D4}. The two rows and their maps are written explicitly.

Let \(A\to B\to C\) be maps, with \(B\to C\) a local complete intersection homomorphism. Choose \(P=A[X_s:s\in\Sigma]\twoheadrightarrow B\) with kernel \(I\), allowing an arbitrary set \(\Sigma\). Choose \(B[Y_1,\ldots,Y_m]\twoheadrightarrow C\) with kernel \(J\). Put \(Q=P[Y_1,\ldots,Y_m]\) and let \(K\) be the kernel of \(Q\to C\). There are natural short exact rows of \(C\)-modules \[\begin{gathered}
0\to(I/I^2)\otimes_B C\to K/K^2\to J/J^2\to0,\\
0\to\Omega_{P/A}\otimes_P C\to\Omega_{Q/A}\otimes_Q C
\to\Omega_{B[Y]/B}\otimes_{B[Y]} C\to0.
\end{gathered}\] The conormal differential maps from the first row to the second commute. In particular the cotangent transitivity sequence is exact with zero at the left: \[\begin{gathered}
0\to H_1(\NL_{B/A}\otimes_B C)\to H_1(L_{C/A})\to H_1(L_{C/B})\\
\to\Omega_{B/A}\otimes_B C\to\Omega_{C/A}\to\Omega_{C/B}\to0.
\end{gathered}\]

\textbf{Proof.} The kernel of \(Q\to B[Y]\) is \(IQ\), by comparing polynomial coefficients, and \(J=K/IQ\). The arbitrary-ring presentation comparison \hyperref[AG-RG-S05-more-algebra-lemma-lci-presentation-independent]{F.94.A} makes \(J\) a Koszul-regular ideal in this chosen finite polynomial presentation. In particular it is \(H_1\)-regular. Apply the intersection identity \hyperref[AG-RG-S05-more-algebra-lemma-conormal-sequence-H1-regular-ideal]{F.101} to \(Q\to Q/IQ\) and the ideal \(K\). It gives \[IQ\cap K^2=(IQ)K.\] Thus the kernel of \(K/K^2\to J/J^2\) is \(IQ/(IQ)K\), embedded in \(K/K^2\). On the other hand \[IQ/(IQ)K\simeq (I/I^2)\otimes_B C.\] To check this identity directly, polynomial coefficients identify \(IQ/I^2Q=(I/I^2)\otimes_B B[Y]\). Quotienting by the products with \(J=K/IQ\) gives the right-hand tensor with \(C=B[Y]/J\), and gives \(IQ/(IQ)K\) on the left; \(I^2Q\subset(IQ)K\) already. The maps send a generator class of \(I\) to its class in \(K/K^2\). This proves exactness of the first row, including its left injection and right surjection.

The second row is the split sequence of free differential modules on the \(X\) symbols, on both the \(X\) and \(Y\) symbols, and on the \(Y\) symbols. Inclusion and projection give the maps, and inclusion of the \(Y\) symbols is a splitting. These statements remain true for an infinite set \(\Sigma\), since each polynomial and differential involves finitely many symbols. The vertical maps send a relation \(f\) to \(df\) after reduction in \(C\). Including a relation from \(I\), or reducing one from \(K\) modulo \(IQ\), commutes with this differentiation; hence the square on each pair of maps commutes.

Apply the complete snake argument \hyperref[AG-RG-S05-algebra-lemma-snake]{F.84} to these actual short exact rows. Their kernel and cokernel sequence begins with zero and ends with zero. The polynomial-presentation homotopy comparison \hyperref[AG-RG-S05-algebra-lemma-NL-homotopy]{F.11} identifies the middle and right kernels with the indicated first cotangent homologies and their cokernels with the corresponding differential modules. For the left vertical map, right exactness identifies its source and target with the terms of \(\NL_{B/A}\otimes_B C\). Its kernel is therefore \(H_1(\NL_{B/A}\otimes_B C)\), and its cokernel is \(\Omega_{B/A}\otimes_B C\). No flatness of \(B\to C\) is needed or inferred in this identification. The representative description of the snake connecting map agrees with the cotangent map in \hyperref[AG-RG-S05-algebra-lemma-exact-sequence-NL]{F.42}: lift a relation, differentiate and take the class in the left cokernel. Thus the displayed sequence has the specified canonical maps. Empty variable lists and the zero target obey the same quotient and complex formulas. ∎

\phantomsection\label{AG-RG-S05-more-algebra-lemma-transitivity-gamma}

\subsubsection{Lemma F.131: transitivity gamma}\label{AG-RG-S05-lemma-f.131-transitivity-gamma}

Reading source: \href{https://stacks.math.columbia.edu/tag/07E2}{Stacks, Tag 07E2}.

For any tower of fields \(K\subset L\subset M\), the natural sequence \[
\begin{gathered}
0\longrightarrow H_1(L_{L/K})\otimes_LM
\longrightarrow H_1(L_{M/K})\longrightarrow H_1(L_{M/L})\\
\longrightarrow\Omega_{L/K}\otimes_LM
\longrightarrow\Omega_{M/K}\longrightarrow\Omega_{M/L}\longrightarrow0.
\end{gathered}
\] is exact, including its first injection. No finite-generation assumption is imposed on the fields.

\textbf{Proof.} First replace \(M\) by a global complete-intersection \(L\)-algebra \(C=L[Y_1,\ldots,Y_r]/J\), with \(J\) generated by a regular list. Choose a polynomial presentation \(P=K[X_s]\twoheadrightarrow L\) with kernel \(I\), allowing an arbitrary set of variables. Let \(Q=P[Y_1,\ldots,Y_r]\) and let \(H=\ker(Q\to C)\); the image of \(H\) in \(Q/IQ=L[Y]\) is \(J\).

The regular-list Koszul argument in \hyperref[AG-RG-S05-more-algebra-lemma-noetherian-finite-all-equivalent]{F.118} proves vanishing of its first homology; that direction uses only the successive multiplication injections. Thus \hyperref[AG-RG-S05-more-algebra-lemma-conormal-sequence-H1-regular-ideal]{F.101}, applied to \(IQ\subset H\) in \(Q\), gives \(IQ\cap H^2=IH\). Polynomial coefficient expansion identifies \[
IQ/IH=(I/I^2)\otimes_L C.
\] Indeed first reduce \(I\) modulo \(I^2\), extend polynomial coefficients to \(L[Y]\), and then impose the equations of \(J\); these are precisely the relations obtained by dividing \(IQ\) by \(IH\). Lifting representatives of \(J\) to \(H\) gives the short exact conormal row \[
0\longrightarrow (I/I^2)\otimes_L C
\longrightarrow H/H^2\longrightarrow J/J^2\longrightarrow0.
\] The first injection is exactly the intersection equality just proved. Differentiation, with coefficients evaluated in \(C\), maps this row to the split exact row \[
0\longrightarrow\Omega_{P/K}\otimes_P C
\longrightarrow\Omega_{Q/K}\otimes_Q C
\longrightarrow\Omega_{L[Y]/L}\otimes_{L[Y]}C\longrightarrow0.
\] Its maps include the free \(dX_s\) coordinates and project to the free \(dY_i\) coordinates. Differentiating any lifted polynomial relation verifies commutativity. The complete snake proof \hyperref[AG-RG-S05-algebra-lemma-snake]{F.84} now gives the asserted six-term kernel/cokernel sequence for \(C\), with zero at both ends. Presentation comparison \hyperref[AG-RG-S05-algebra-lemma-NL-homotopy]{F.11} identifies its kernels and cokernels with the named first conormal homologies and differential modules. For the left terms, tensoring over the field \(L\) is exact, so they are \(H_1(L_{L/K})\otimes_L C\) and \(\Omega_{L/K}\otimes_L C\).

Finally \hyperref[AG-RG-S05-algebra-lemma-colimit-syntomic]{F.23} writes \(M\) as a directed union of global complete-intersection \(L\)-subalgebras of this kind. The relation and differential comparison \hyperref[AG-RG-S05-algebra-lemma-colimits-NL]{F.24} identifies their cotangent complexes over both \(K\) and \(L\) with those of the union. The sequence maps agree at each inclusion because they were defined by the same polynomial representatives and their derivatives. Filtered colimits are exact: a representative killed in a colimit is killed at some common later stage, and finite sums and equalities can all be placed at one such stage, as proved in F.24. The tensor generators and their balanced relations likewise give \(\varinjlim(C\otimes_L V)=M\otimes_L V\) for each \(L\)-vector space \(V\). Taking the colimit of the exact sequences therefore proves the full sequence for \(M\), including its first zero. Empty equation lists and infinitely generated intermediate fields obey the same argument. ∎

\phantomsection\label{AG-RG-S05-more-algebra-proposition-Noetherian-complete-G-ring}

\subsubsection{Proposition F.132: Noetherian complete G ring}\label{AG-RG-S05-proposition-f.132-noetherian-complete-g-ring}

\emph{Adapted from the Stacks project, more-algebra.tex, lines 12965--12999.}
\nopagebreak[4]

A Noetherian complete local ring is a G-ring.

\textbf{Proof.} Let \(A\) be a Noetherian complete local ring. By Lemma \hyperref[AG-RG-S05-more-algebra-lemma-check-G-ring-easy]{F.99} it suffices to check that \(B = A/\mathfrak q\) has geometrically regular formal fibres over the minimal prime \((0)\) of \(B\). Thus we may assume that \(A\) is a domain and it suffices to check the condition for the formal fibres over the minimal prime \((0)\) of \(A\). Let \(K\) be the fraction field of \(A\).

We can choose a subring \(A_0 \subset A\) which is a regular complete local ring such that \(A\) is finite over \(A_0\), see Algebra, Lemma \hyperref[AG-RG-S05-algebra-lemma-complete-local-Noetherian-domain-finite-over-regular]{F.25}. Moreover, we may assume that \(A_0\) is a power series ring over a field or a Cohen ring. By Lemma \hyperref[AG-RG-S05-more-algebra-lemma-G-ring-goes-up-quasi-finite]{F.96} we see that it suffices to prove the result for \(A_0\).

Assume that \(A\) is a power series ring over a field or a Cohen ring. Since \(A\) is regular local, each \(A_\mathfrak p\) is regular local by the written localization proof \href{https://kokunoyumeto.github.io/open-math-courses/courses/AG-CA/regular-local-rings.html\#3-regularity-at-more-general-points}{AG-CA-14, Theorem 3.2}. Hence the completions \(A_\mathfrak p^\wedge\) are regular, see Lemma \href{https://kokunoyumeto.github.io/open-math-courses/courses/AG-CA/completion.html}{AG-CA-19, Theorem 4.1}. Hence the fibre \(A_{\mathfrak p}^\wedge \otimes_A K\) is, as a localization of \(A_\mathfrak p^\wedge\), also regular. Thus we are done if the characteristic of \(K\) is \(0\). The positive characteristic case is the case \(A = k[[x_1, \ldots, x_d]]\) which is a special case of Lemma \hyperref[AG-RG-S05-more-algebra-lemma-helper-G-ring]{F.110}. ∎

\phantomsection\label{AG-RG-S05-more-algebra-proposition-characterization-geometrically-regular}

\subsubsection{Proposition F.133: characterization geometrically regular}\label{AG-RG-S05-proposition-f.133-characterization-geometrically-regular}

Reading source: \href{https://stacks.math.columbia.edu/tag/07E5}{Stacks, Tag 07E5}. The proof below includes both finite-defect dimension arguments.

Let \(k\) have characteristic \(p>0\), and let \((A,\mathfrak m,K)\) be a Noetherian local \(k\)-algebra. The following conditions are equivalent:

\begin{enumerate}
\def\labelenumi{\arabic{enumi}.}
\tightlist
\item
  \(A\) is geometrically regular over \(k\).
\item
  \(A\otimes_k k'\) is regular whenever \(k\subset k'\subset k^{1/p}\) and \([k':k]<\infty\).
\item
  \(A\) is regular and \(H_1(L_{K/k})\to\mathfrak m/\mathfrak m^2\) is injective.
\item
  \(A\) is regular and \(\Omega_{k/\mathbf F_p}\otimes_k K\to\Omega_{A/\mathbf F_p}\otimes_A K\) is injective.
\end{enumerate}

\textbf{Proof.} Write \(E=\mathfrak m/\mathfrak m^2\), \(H=H_1(L_{K/k})\), \(U=\Omega_{k/\mathbf F_p}\otimes_k K\) and \(V=\Omega_{A/\mathbf F_p}\otimes_A K\). The quotient presentation of \(K\) over \(A\) has no relative polynomial variables, so its naive complex is \(E\to0\). Every field is formally smooth over its prime field by \href{https://kokunoyumeto.github.io/open-math-courses/courses/AG-CA/coefficient-rings-and-cohen-structure.html}{AG-CA-19S, Lemma 1.1}, and \hyperref[AG-RG-S05-algebra-lemma-characterize-formally-smooth-field-extension]{F.17} consequently gives \(H_1(L_{K/\mathbf F_p})=0\). The complete quotient transitivity proof \hyperref[AG-RG-S05-algebra-lemma-exact-sequence-NL]{F.42} and the field version \hyperref[AG-RG-S05-more-algebra-lemma-transitivity-gamma]{F.131} give \[\begin{gathered}
0\to E\to V\to\Omega_{K/\mathbf F_p}\to0,\\
0\to H\to U\to\Omega_{K/\mathbf F_p}\to\Omega_{K/k}\to0.
\end{gathered}\] Naturality maps the second row into the first, with identity on \(\Omega_{K/\mathbf F_p}\). If \(u\in U\) maps to zero in \(V\), it maps to zero in \(\Omega_{K/\mathbf F_p}\) and hence comes from a unique \(h\in H\). The first row embeds \(E\) in \(V\), so this \(h\) maps to zero in \(E\). Conversely an \(h\) killed in \(E\) maps to a vector of \(U\) killed in \(V\). Thus the two kernels are canonically isomorphic, proving (3) equivalent to (4).

Assume (3), and take any finite purely inseparable \(k'/k\). Set \(A'=A\otimes_k k'\). A basis of \(k'\) containing \(1\) makes \(A'\) finite free over \(A\) and the structural map injective. It is integral by the determinant proof \href{https://kokunoyumeto.github.io/open-math-courses/courses/AG-CA/integral-extensions-lying-over-going-up-and-going-down.html}{AG-CA-05, Theorem 1.2}. For some \(e\), \((k')^{p^e}\subset k\); applying Frobenius to a finite tensor sum shows \(x^{p^e}\in A\) for every \(x\in A'\). For a prime \(\mathfrak q\) over \(\mathfrak p\), the equivalence \(x\in\mathfrak q\Longleftrightarrow x^{p^e}\in\mathfrak p\) determines \(\mathfrak q\) uniquely. Lying over supplies it. Integral maximal-ideal correspondence, AG-CA-05, Theorem 3.3, makes \(A'\) local with maximal ideal \(\mathfrak m'\) over \(\mathfrak m\). It is Noetherian, since it is finite over \(A\). Write \(K'=A'/\mathfrak m'\) and \(E'=\mathfrak m'/\mathfrak m'^2\). The surjection \(K\otimes_k k'\to K'\) proves \([K':K]<\infty\). The integral inclusion preserves dimension by AG-CA-05, Theorem 5.1.

Put \(H'=H_1(L_{K'/k'})\). Relative quotient transitivity gives rows \[\begin{gathered}
0\to H\otimes_K K'\to E\otimes_K K'\to V_0\to W\to0,\\
H'\to E'\to V_0\to W'\to0,
\end{gathered}\] where \(W=\Omega_{K/k}\otimes_K K'\), \(W'=\Omega_{K'/k'}\), and differential base change \href{https://kokunoyumeto.github.io/open-math-courses/courses/AG-CA/kahler-differentials.html}{AG-CA-16, Theorem 2.1} identifies the two middle differential terms \(V_0\). Let \(\alpha:W\to W'\) and \(\beta:H\otimes_K K'\to H'\) be the natural maps. Surjectivity from the common \(V_0\) makes \(\alpha\) surjective. Both extensions \(k'/k\) and \(K'/K\) are finite, so \hyperref[AG-RG-S05-more-algebra-lemma-gamma-commutative-diagram]{F.108} applies with transcendence degrees zero and proves \[\dim\ker\alpha-\dim\ker\beta+\dim\operatorname{coker}\beta=0.\] All these defects are finite. Also \(H\) embeds in the finite space \(E\), and finite cokernel of \(\beta\) therefore makes \(H'\) finite.

Let \(L\) and \(L'\) be the images of the two conormal terms in \(V_0\). Then \(L\subset L'\) and \(L'/L=\ker\alpha\). The first row gives \(\dim E=\dim H+\dim L\); the second gives \(\dim E'\leq\dim H'+\dim L'\). Consequently \[\begin{aligned}
\dim E'-\dim E
&\leq\dim H'-\dim H+\dim\ker\alpha\\
&=\dim\operatorname{coker}\beta-\dim\ker\beta+\dim\ker\alpha=0.
\end{aligned}\] Here only finite spaces are subtracted. Since \(A\) is regular and \(\dim A'=\dim A\), this says \(\dim E'\leq\dim A'\). Conversely a basis of \(E'\) lifts to generators of \(\mathfrak m'\) by \hyperref[AG-RG-S05-algebra-lemma-NAK]{F.10}, and \href{https://kokunoyumeto.github.io/open-math-courses/courses/AG-CA/dimension-theory-of-noetherian-local-rings.html}{AG-CA-11, Theorem 3.1 and Proposition 4.2} give \(\dim A'\leq\dim E'\). Equality makes \(A'\) regular. This holds for every finite purely inseparable extension, so \hyperref[AG-RG-S05-algebra-lemma-geometrically-regular]{F.50} proves (1).

It remains to prove (2) implies (4). Taking \(k'=k\) first makes \(A\) regular. Choose \(a_1,\ldots,a_n\in k\) whose absolute differentials are independent. The full root-degree proof \hyperref[AG-RG-S05-algebra-lemma-size-extension-pth-roots]{F.81} identifies \[k'=k(a_1^{1/p},\ldots,a_n^{1/p})=k[T_1,\ldots,T_n]/(T_i^p-a_i),\qquad [k':k]=p^n.\] The preceding prime and dimension arguments apply to \(A'=A\otimes_k k'\); by (2) it is regular. Its monic equations are a regular list, since successive monic quotient division makes each next monic polynomial injective by its leading coefficient. Their conormal classes are a free basis by \hyperref[AG-RG-S05-algebra-lemma-regular-quasi-regular]{F.74}. All relative derivatives are zero. Thus \(H_1(L_{A'/A})\) and \(\Omega_{A'/A}\) are free of rank \(n\). The representative connecting-map construction in F.42 sends the \(i\)th homology basis vector to \(da_i\otimes1\), up to an immaterial common sign, and yields \[A'^n\to\Omega_{A/\mathbf F_p}\otimes_A A'\to\Omega_{A'/\mathbf F_p}\to A'^n\to0.\] This row stays exact after tensoring with \(K'\). Indeed the last surjection splits because its target is free. Its kernel therefore embeds after tensoring, while tensor right exactness preserves the preceding cokernel presentation. Combining these two statements gives a map \[\delta:V\otimes_K K'\longrightarrow V'=\Omega_{A'/\mathbf F_p}\otimes_{A'}K'\] whose kernel is spanned by the \(n\) vectors \(da_i\otimes1\) and whose cokernel is \(K'^n\).

Use the two absolute short exact rows \[\begin{gathered}
0\to E\otimes_K K'\to V\otimes_K K'\to\Omega_{K/\mathbf F_p}\otimes_K K'\to0,\\
0\to E'\to V'\to\Omega_{K'/\mathbf F_p}\to0.
\end{gathered}\] The left vertical map has finite equally dimensional domain and target: both dimensions equal \(\dim A=\dim A'\). Its kernel and cokernel thus have equal dimensions. For the right vertical map, F.131 identifies the kernel with \(H_1(L_{K'/K})\) and the cokernel with \(\Omega_{K'/K}\). Their finite dimensions agree by \hyperref[AG-RG-S05-more-algebra-lemma-cartier-equality]{F.98}, applied to the proved finite extension \(K'/K\). The full snake sequence \hyperref[AG-RG-S05-algebra-lemma-snake]{F.84} now gives \[\dim\ker\delta-\dim\operatorname{coker}\delta=0.\] Every defect in this computation is finite, although the absolute differential spaces need not be finite. Since the cokernel has dimension \(n\), the \(n\) vectors spanning the kernel are independent. The \(p\)-basis proof in AG-CA-19S, Lemma 1.1, and its differential calculation F.81 give a basis of \(\Omega_{k/\mathbf F_p}\) consisting of such \(da_i\). Every prospective relation uses a finite subset. Applying the preceding calculation to that subset proves the full injection in (4). For \(n=0\) the argument is empty. Finally (1) implies (2) directly from geometric regularity. These implications prove the equivalence. ∎

\phantomsection\label{AG-RG-S05-more-algebra-proposition-finite-type-over-G-ring}

\subsubsection{Proposition F.134: finite type over G ring}\label{AG-RG-S05-proposition-f.134-finite-type-over-g-ring}

Reading source: \href{https://stacks.math.columbia.edu/tag/07PV}{Stacks, Tag 07PV}. The completed-local descent and each individual fibre reduction are explicit below.

If \(R\) is a G-ring and \(S\) is an essentially finite-type \(R\)-algebra, then \(S\) is a G-ring.

\textbf{Proof.} G-rings are Noetherian by Definition \hyperref[AG-RG-S05-more-algebra-definition-G-ring]{F.92}; finite-type algebras and their localizations are Noetherian by \href{https://kokunoyumeto.github.io/open-math-courses/courses/AG-CA/noetherian-and-artinian-rings.html}{AG-CA-03, Theorem 2.1 and Proposition 1.2}. A localization has the same prime local rings as the original ring at the surviving primes. Thus it suffices to handle finite-type algebras. Such an algebra is a quotient of a finite-variable polynomial ring. Quotient maps are quasi-finite, and \hyperref[AG-RG-S05-more-algebra-lemma-G-ring-goes-up-quasi-finite]{F.96} proves G-ring ascent for them. Induction on the number of variables reduces the problem to \(R[X]\).

By \hyperref[AG-RG-S05-more-algebra-lemma-check-G-ring-maximal-ideals]{F.100}, fix a maximal ideal \(\mathfrak q\) of \(R[X]\) and prove its local completion map regular. Let \(\mathfrak p=\mathfrak q\cap R\). Replacing \(R\) by \(R_{\mathfrak p}\) leaves \(R[X]_{\mathfrak q}\) unchanged, and this new base is a local G-ring. Write its maximal ideal as \(\mathfrak m\), its completion as \(\widehat R\), and put \[A=R[X]_{\mathfrak q},\qquad B=\widehat R[X]_{\mathfrak q'}.\] There is exactly one prime \(\mathfrak q'\) above \(\mathfrak q\): the quotient identity \(\widehat R/\mathfrak m\widehat R=R/\mathfrak m\) gives \(\widehat R[X]/\mathfrak q\widehat R[X]=R[X]/\mathfrak q\), a field. This uses \href{https://kokunoyumeto.github.io/open-math-courses/courses/AG-CA/completion.html}{AG-CA-19, Proposition 2.2} and the fact \(\mathfrak m\subset\mathfrak q\). The map \(A\to B\) is local and flat. It is also regular, since it is a localization of the base change of the regular G-ring map \(R\to\widehat R\); these operations preserve regularity by \hyperref[AG-RG-S05-more-algebra-lemma-regular-base-change]{F.123}.

If \(B\to\widehat B\) is regular, \hyperref[AG-RG-S05-more-algebra-lemma-regular-composition]{F.124} makes \(A\to\widehat B\) regular. The completed map \(\widehat A\to\widehat B\) is faithfully flat by the full proof \hyperref[AG-RG-S05-completion-flat-local-map]{F.99A} for this actual flat local \(A\to B\). Therefore \hyperref[AG-RG-S05-more-algebra-lemma-regular-permanence]{F.125} makes \(A\to\widehat A\) regular. This reduces us to a complete Noetherian local base, with \(\mathfrak q\) above its maximal ideal; it does not assume unrelated completions are equal.

For this complete base, fix \(\mathfrak r\subset\mathfrak q\) in \(R[X]\). Its completed fibre is \[\widehat{R[X]_{\mathfrak q}}\otimes_{R[X]}\kappa(\mathfrak r).\] Let \(\mathfrak a=\mathfrak r\cap R\). Finite-module completion commutes with quotient by \(\mathfrak a\), by AG-CA-19, Theorem 3.1. Since \(\mathfrak a\) vanishes in the displayed residue field, replacing \(R\) by \(R/\mathfrak a\) gives exactly the same fibre. The quotient is still complete and Noetherian, by that theorem and Theorem 3.3. Thus assume \(R\) is a complete local domain and \(\mathfrak r\cap R=0\).

Choose the finite regular complete local subring \(R_0\subset R\) from \hyperref[AG-RG-S05-algebra-lemma-complete-local-Noetherian-domain-finite-over-regular]{F.25}. The map \(R_0[X]\to R[X]\) is finite, hence quasi-finite at the contracted primes \(\mathfrak r_0\subset\mathfrak q_0\). The prime \(\mathfrak q_0\) is maximal by integral maximal-ideal correspondence; it lies above the maximal ideal of \(R_0\), since \(R\) is finite local over the local subring with the same residue field. Also \(\mathfrak r_0\cap R_0=0\). Assertion (1) of F.96 transfers geometric regularity of this individual completed fibre from \(R_0[X]\) to \(R[X]\). Its tensor and completion-product proof requires only that fibre, without any prior G-ring assertion for \(R_0[X]\). We may therefore assume the base domain is regular and complete.

Let \(K=\operatorname{Frac}(R)\). If \(\mathfrak r=0\), the fibre is \(\widehat{R[X]_{\mathfrak q}}\otimes_{R[X]}\operatorname{Frac}(R[X])\), geometrically regular by the second full assertion of \hyperref[AG-RG-S05-more-algebra-lemma-helper-G-ring]{F.110}. If \(\mathfrak r\ne0\), its extension to the principal ideal domain \(K[X]\) is a nonzero prime, hence generated by an irreducible polynomial. Its residue field is finite over \(K\). The second complete assertion of \hyperref[AG-RG-S05-more-algebra-lemma-another-helper-G-ring]{F.97} proves geometric regularity of this fibre: for each finite purely inseparable test field it constructs a finite complete local domain \(D\) inside that field and identifies the tested fibre with factors of \[\big(\widehat{D[X]_{\mathfrak q_i}}\otimes_D\operatorname{Frac}(D)\big)/(X-a).\] Its first assertion makes the parenthesized algebra regular. The extended \(\partial/\partial X\) sends \(X-a\) to \(1\), and the proved hypersurface argument then makes the quotient regular. This quotient is used as a quotient; no localization assertion about it is required. The two prime cases exhaust the fibres. Completion is flat by AG-CA-19, Theorem 3.2, so the geometrically regular fibres make its map regular. This proves polynomial permanence and all the reductions to \(S\). Zero target rings have no local rings to check. ∎

\phantomsection\label{AG-RG-S05-more-algebra-proposition-ubiquity-G-ring}

\subsubsection{Proposition F.135: ubiquity G ring}\label{AG-RG-S05-proposition-f.135-ubiquity-g-ring}

Reading source: \href{https://stacks.math.columbia.edu/tag/07PX}{Stacks, Tag 07PX}. We prove precisely the finite-type-over-\(\mathbf Z\) assertion used in this programme.

Every ring essentially of finite type over \(\mathbf Z\) is a G-ring.

\textbf{Proof.} The prime ideals of \(\mathbf Z\) are \((0)\) and \((\ell)\) for prime integers \(\ell\). At \((0)\) its local ring is \(\mathbf Q\). Completion at the zero maximal ideal is identity, so this local completion map is regular.

At \((\ell)\) put \(A=\mathbf Z_{(\ell)}\). Factoring an integer into its power of \(\ell\) and its remaining factor shows every nonzero element of \(A\) is \(\ell^n u\) with \(u\) a unit. Every nonzero ideal is generated by the element of smallest such order. Hence \(A\) is Noetherian local with maximal ideal \((\ell)\) and only the primes \((0)\) and \((\ell)\). These are also the DVR computations in \href{https://kokunoyumeto.github.io/open-math-courses/courses/AG-CA/discrete-valuation-rings-normal-rings-and-serres-criterion.html}{AG-CA-15, Section 1}.

Write \(\widehat A\) for its maximal-ideal completion. The exact quotient identities and local-ring proof \href{https://kokunoyumeto.github.io/open-math-courses/courses/AG-CA/completion.html}{AG-CA-19, Proposition 2.2, Proposition 2.3 and Theorems 3.2--3.3} make it Noetherian, complete and separated, with maximal ideal \(\ell\widehat A\) and residue field \(\mathbf F_\ell\). They also make \(A\to\widehat A\) flat, so multiplication by \(\ell\) remains injective. If \(0\ne x\in\widehat A\), separatedness gives a finite largest \(n\) with \(x\in\ell^n\widehat A\). Thus \(x=\ell^n u\), and \(u\) cannot belong to \(\ell\widehat A\), or maximality would fail. The local-ring unit criterion makes \(u\) invertible. For two nonzero elements their product is a unit times a power of \(\ell\), which is nonzero by its multiplication injection. This proves directly that \(\widehat A\) is a domain. The same order argument makes its nonzero ideals powers of \(\ell\) and makes \(\widehat A[1/\ell]\) its fraction field.

There are just two fibres of the flat map \(A\to\widehat A\). The closed fibre is \(\widehat A/\ell\widehat A=\mathbf F_\ell\) over itself; the generic fibre is the field \(\widehat A[1/\ell]\) over \(\mathbf Q\). A field of characteristic zero is geometrically regular over any characteristic-zero subfield by \hyperref[AG-RG-S05-algebra-lemma-geometrically-regular]{F.50}: the only finite purely inseparable test extension is the identity, and the tested ring is a field, hence regular. The closed fibre satisfies the same criterion over itself, since a finite purely inseparable extension tensored with its base field is again that extension field. Thus both fibres are geometrically regular. Definition \hyperref[AG-RG-S05-more-algebra-definition-G-ring]{F.92} now makes \(\mathbf Z\) a G-ring. Apply the fully proved finite-type permanence \hyperref[AG-RG-S05-more-algebra-proposition-finite-type-over-G-ring]{F.134} and its localization step to obtain the assertion. The zero algebra is included by its empty spectrum. ∎

\phantomsection\label{AG-RG-S05-topology-lemma-Noetherian}

\subsubsection{Lemma F.136: Noetherian}\label{AG-RG-S05-lemma-f.136-noetherian}

Reading source: \href{https://stacks.math.columbia.edu/tag/0052}{Stacks, Tag 0052}.

Let \(X\) be a Noetherian topological space. Every subspace is Noetherian, there are finitely many irreducible components, and each component contains a nonempty subset open in \(X\).

\textbf{Proof.} Recall that Noetherianity means stabilization of descending sequences of closed subsets. For a subspace \(T\), take a descending sequence \(C_j\) closed in \(T\) and choose closed \(D_j\subset X\) with \(C_j=T\cap D_j\). Put \(E_j=D_1\cap\cdots\cap D_j\). Because the \(C_j\) descend, \(T\cap E_j=C_j\). The descending sequence \(E_j\) stabilizes in \(X\), so the sequence \(C_j\) stabilizes in \(T\). This proves the assertion for every subspace, including the empty one.

We first show that every closed subset of \(X\) is a finite union of irreducible closed subsets. The empty subset is the union of an empty family. Otherwise, if there were closed subsets without such an expression, the descending-chain condition would supply a minimal one \(C\): lacking a minimal member would allow a strictly descending sequence of counterexamples. This \(C\) is nonempty and cannot itself be irreducible. Thus \(C=C_1\cup C_2\), where both \(C_i\) are proper subsets closed in \(C\), and hence closed in \(X\). Minimality supplies a finite irreducible closed cover of each \(C_i\), and their combined covers supply one of \(C\), a contradiction.

Apply the result to \(X\). From the finite cover remove any member contained in the union of the others; this operation terminates and still leaves a cover. By \hyperref[AG-RG-S05-topology-lemma-pick-irreducible-components]{F.138}, its members are exactly the irreducible components. There are consequently only finitely many of them.

Fix a component \(C\) and list the others as \(C_1,\ldots,C_s\). They are closed by \hyperref[AG-RG-S05-topology-lemma-irreducible]{F.137}. The open set \(U=X\setminus\bigcup_{i=1}^s C_i\) lies in \(C\), since the components cover \(X\). It is nonempty: otherwise \(C\) would be the finite union of the closed subsets \(C\cap C_i\). Binary irreducibility, applied repeatedly to a finite union, would force \(C\subset C_i\) for some \(i\), contradicting the distinctness and maximality of these components. If there are no other components, \(U=X=C\), which is nonempty by the definition of an irreducible component. ∎

\phantomsection\label{AG-RG-S05-topology-lemma-irreducible}

\subsubsection{Lemma F.137: irreducible}\label{AG-RG-S05-lemma-f.137-irreducible}

Reading source: \href{https://stacks.math.columbia.edu/tag/004W}{Stacks, Tag 004W}.

In an arbitrary topological space, the closure of every irreducible subset is irreducible. Every irreducible subset is contained in an irreducible component; components are closed and together cover the space.

\textbf{Proof.} We use the nonempty convention for irreducibility. A nonempty space is irreducible exactly when any two nonempty open subsets intersect: taking complements converts this into the assertion that two proper closed subsets cannot cover it.

Let \(Y\subset X\) be irreducible. An open subset of \(X\) which meets \(\overline Y\) meets \(Y\), by the definition of closure. If two relatively open subsets of \(\overline Y\) are nonempty, represent them by ambient opens. Their intersections with \(Y\) are nonempty and relatively open in \(Y\), so they intersect. The original two subsets therefore intersect as well. This proves irreducibility of \(\overline Y\).

Consider the set of all irreducible subsets of \(X\) containing \(Y\), ordered by inclusion. It is nonempty because it contains \(Y\). For a nonempty chain, its union \(V\) is nonempty. Two nonempty relatively open subsets of \(V\) meet some members of the chain. One of those members contains the other, so that member meets both opens. Its irreducibility supplies a point in their intersection. Thus \(V\) is irreducible and is an upper bound in the set. The empty chain has \(Y\) as an upper bound. Zorn's lemma now gives a maximal member \(C\). Any irreducible subset containing \(C\) also contains \(Y\), so this is a maximal irreducible subset of \(X\), namely a component.

The first part shows that \(\overline C\) is irreducible; maximality forces \(C=\overline C\), so every component is closed. Each singleton is irreducible, since its only nonempty open subset is itself. Applying the existence argument to every singleton proves that the components cover \(X\). For \(X=\varnothing\) both the component family and its union are empty. ∎

\phantomsection\label{AG-RG-S05-topology-lemma-pick-irreducible-components}

\subsubsection{Lemma F.138: pick irreducible components}\label{AG-RG-S05-lemma-f.138-pick-irreducible-components}

Reading source: \href{https://stacks.math.columbia.edu/tag/0G2Y}{Stacks, Tag 0G2Y}.

Suppose \(X=C_1\cup\cdots\cup C_n\), where every \(C_i\) is irreducible and closed in \(X\), and no \(C_i\) lies in the union of the remaining members. Then this family is exactly the family of irreducible components of \(X\).

\textbf{Proof.} Let \(D\) be any irreducible component. The sets \(D\cap C_i\) form a finite closed cover of \(D\). Irreducibility implies that one member of a finite closed cover is the whole space: for two members this is its definition, and grouping the last member against the preceding union proves the assertion by induction. Consequently \(D\subset C_i\) for some \(i\). Because \(C_i\) is itself irreducible, maximality of \(D\) gives \(D=C_i\).

Conversely fix \(C_i\). The general existence argument \hyperref[AG-RG-S05-topology-lemma-irreducible]{F.137} puts it in a component \(D\). The preceding paragraph identifies \(D\) with some \(C_j\). If \(j\ne i\), the containment \(C_i\subset C_j\) would put \(C_i\) in the union of all the other members, contrary to the hypothesis. Thus \(j=i\), and \(C_i\) is a component. These two arguments prove equality of the families. When \(n=0\), the covering assumption says \(X\) is empty, and it has no nonempty irreducible components, giving the same conclusion. ∎

\section{6. Exact earlier proof dependencies}\label{AG-RG-S05-exact-earlier-proof-dependencies}

The statements in Sections 1--5 specify the scope actually used in approximation. The following earlier proofs supply their elementary algebra inputs. \href{https://kokunoyumeto.github.io/open-math-courses/courses/AG-RG/AG-RG-S04.html}{Descent and Zariski Main}, Lemma D2.2, Theorem C4.3 and Corollary B1.2, precedes the henselian section argument. The broader formal-smoothness/regularity equivalence is not consumed: the explicit completion arguments F.99A, F.110 and F.132--F.135 prove the required G-ring inputs. This dependency list points to the proofs of the results it names; results that it does not name are not covered by it. Finite type and finite presentation on every compatible affine pair, used in the target-descent step of AG-GS-05, Lemma 7.6, have their complete morphism argument in \href{https://kokunoyumeto.github.io/open-math-courses/courses/AG-RG/AG-RG-S04.html\#affine-presentation-communication}{AG-RG-S04, Lemma A1.8}, using the preceding finite-generator and finite-relation algebra patching proof. In the application of AG-GS-05, Theorem 7.7, the fixed-character-group splitting scheme is globally a disjoint union on the splitting cover, so this application uses that case of Lemma 7.6.

\subsection{Exact earlier programme proofs used}\label{AG-RG-S05-exact-earlier-programme-proofs-used}

These locators refer to written programme proofs. They do not claim that an unused, broader source statement has been proved.

{\def\LTcaptype{none} % do not increment counter
\begin{longtable}[]{@{}
  >{\raggedright\arraybackslash}p{(\linewidth - 2\tabcolsep) * \real{0.5000}}
  >{\raggedright\arraybackslash}p{(\linewidth - 2\tabcolsep) * \real{0.5000}}@{}}
\toprule\noalign{}
\begin{minipage}[b]{\linewidth}\raggedright
Consumed source label
\end{minipage} & \begin{minipage}[b]{\linewidth}\raggedright
Earlier programme proof
\end{minipage} \\
\midrule\noalign{}
\endhead
\bottomrule\noalign{}
\endlastfoot
definition regular & \href{https://kokunoyumeto.github.io/open-math-courses/courses/AG-CA/regular-local-rings.html}{AG-CA-14, Section 1, definition} \\
definition regular local & \href{https://kokunoyumeto.github.io/open-math-courses/courses/AG-CA/regular-local-rings.html}{AG-CA-14, Section 1, definition} \\
definition standard smooth & \href{https://kokunoyumeto.github.io/open-math-courses/courses/AG-CA/formally-smooth-unramified-and-etale-ring-maps.html}{AG-CA-17, Section 4, equations (6) and (7)} \\
lemma Artin Rees & \href{https://kokunoyumeto.github.io/open-math-courses/courses/AG-CA/noetherian-and-artinian-rings.html}{AG-CA-03, Theorem 5.1} \\
lemma Noetherian irreducible components & \href{https://kokunoyumeto.github.io/open-math-courses/courses/AG-CA/noetherian-and-artinian-rings.html}{AG-CA-03, Proposition 2.2} \\
lemma Noetherian topology & \href{https://kokunoyumeto.github.io/open-math-courses/courses/AG-CA/noetherian-and-artinian-rings.html}{AG-CA-03, Proposition 2.2} \\
lemma Zariski topology & \href{https://kokunoyumeto.github.io/open-math-courses/courses/AG-CA/spectra-of-rings.html}{AG-CA-01, Theorem 1.2 and Proposition 2.1} \\
lemma artinian finite length & \href{https://kokunoyumeto.github.io/open-math-courses/courses/AG-CA/noetherian-and-artinian-rings.html}{AG-CA-03, Theorem 4.2} \\
lemma characterize henselian & \href{https://kokunoyumeto.github.io/open-math-courses/courses/AG-CA/henselian-local-rings-and-henselization.html}{AG-CA-20, Theorem 1.2 and Theorem 2.2} \\
lemma characterize projective & \href{https://kokunoyumeto.github.io/open-math-courses/courses/AG-CA/resolutions-tor-and-ext.html}{AG-CA-06S, Section 1} \\
lemma cokernel flat & \href{https://kokunoyumeto.github.io/open-math-courses/courses/AG-CA/faithful-flatness-and-the-local-criterion-for-flatness.html}{AG-CA-08, Lemma 5.1} \\
lemma compose formally smooth & \href{https://kokunoyumeto.github.io/open-math-courses/courses/AG-CA/formally-smooth-unramified-and-etale-ring-maps.html}{AG-CA-17, Theorem 1.2} \\
lemma differentials base change & \href{https://kokunoyumeto.github.io/open-math-courses/courses/AG-CA/kahler-differentials.html}{AG-CA-16, Theorem 2.1} \\
lemma differentials finitely presented & \href{https://kokunoyumeto.github.io/open-math-courses/courses/AG-CA/kahler-differentials.html}{AG-CA-16, Theorem 5.1, finite polynomial presentation} \\
lemma dimension base fibre equals total & \href{https://kokunoyumeto.github.io/open-math-courses/courses/AG-CA/dimension-theory-of-noetherian-local-rings.html}{AG-CA-11, Theorem 5.1} \\
lemma dimension base fibre total & \href{https://kokunoyumeto.github.io/open-math-courses/courses/AG-CA/dimension-theory-of-noetherian-local-rings.html}{AG-CA-11, Theorem 5.1} \\
lemma dimension prime polynomial ring & \href{https://kokunoyumeto.github.io/open-math-courses/courses/AG-CA/krull-dimension-and-noether-normalization.html}{AG-CA-09, Theorem 4.3} \\
lemma disjoint implies product & \href{https://kokunoyumeto.github.io/open-math-courses/courses/AG-CA/spectra-of-rings.html}{AG-CA-01, Theorem 5.2 and Corollary 5.3} \\
lemma etale at prime & \href{https://kokunoyumeto.github.io/open-math-courses/courses/AG-CA/formally-smooth-unramified-and-etale-ring-maps.html}{AG-CA-17, Theorem 7.2 and Corollary 7.3} \\
lemma exact sequence differentials & \href{https://kokunoyumeto.github.io/open-math-courses/courses/AG-CA/kahler-differentials.html}{AG-CA-16, Theorem 3.1} \\
lemma faithfully flat universally injective & \href{https://kokunoyumeto.github.io/open-math-courses/courses/AG-CA/faithful-flatness-and-the-local-criterion-for-flatness.html}{AG-CA-08, Theorem 2.2} \\
lemma finite finite fibres & \href{https://kokunoyumeto.github.io/open-math-courses/courses/AG-CA/integral-extensions-lying-over-going-up-and-going-down.html}{AG-CA-05, Theorem 3.4} \\
lemma finite projective & \href{https://kokunoyumeto.github.io/open-math-courses/courses/AG-CA/tor-and-flat-modules.html}{AG-CA-07, Theorem 5.3} \\
lemma finite transitive & \href{https://kokunoyumeto.github.io/open-math-courses/courses/AG-CA/integral-extensions-lying-over-going-up-and-going-down.html}{AG-CA-05, Theorem 1.2} \\
lemma flat & \href{https://kokunoyumeto.github.io/open-math-courses/courses/AG-CA/tor-and-flat-modules.html}{AG-CA-07, Theorem 2.1} \\
lemma flat base change & \href{https://kokunoyumeto.github.io/open-math-courses/courses/AG-CA/tor-and-flat-modules.html}{AG-CA-07, Proposition 3.2} \\
lemma flat eq & \href{https://kokunoyumeto.github.io/open-math-courses/courses/AG-CA/tor-and-flat-modules.html}{AG-CA-07, Theorem 5.1} \\
lemma flat going down & \href{https://kokunoyumeto.github.io/open-math-courses/courses/AG-CA/tor-and-flat-modules.html}{AG-CA-07, Theorem 6.3} \\
lemma flat tor zero & \href{https://kokunoyumeto.github.io/open-math-courses/courses/AG-CA/tor-and-flat-modules.html}{AG-CA-07, Theorem 2.1} \\
lemma formally etale etale & \href{https://kokunoyumeto.github.io/open-math-courses/courses/AG-CA/formally-smooth-unramified-and-etale-ring-maps.html}{AG-CA-17, Theorem 2.2 and Proposition 6.1} \\
lemma integral no inclusion & \href{https://kokunoyumeto.github.io/open-math-courses/courses/AG-CA/integral-extensions-lying-over-going-up-and-going-down.html}{AG-CA-05, Theorem 3.3} \\
lemma intersect powers ideal module zero & \href{https://kokunoyumeto.github.io/open-math-courses/courses/AG-CA/noetherian-and-artinian-rings.html}{AG-CA-03, Theorem 6.1} \\
lemma lci CM & \href{https://kokunoyumeto.github.io/open-math-courses/courses/AG-CA/regular-sequences-depth-and-cohen-macaulay-modules.html}{AG-CA-12, Corollary 6.2} \\
lemma maximal chain CM & \href{https://kokunoyumeto.github.io/open-math-courses/courses/AG-CA/regular-sequences-depth-and-cohen-macaulay-modules.html}{AG-CA-12, Lemma 5.4 and Theorem 5.5} \\
lemma quasi compact & \href{https://kokunoyumeto.github.io/open-math-courses/courses/AG-CA/spectra-of-rings.html}{AG-CA-01, Proposition 2.3} \\
lemma rank omega & \href{https://kokunoyumeto.github.io/open-math-courses/courses/AG-CA/smooth-algebras-over-a-field-and-the-jacobian-criterion.html}{AG-CA-18, Proposition 3.1} \\
lemma reformulate CM & \href{https://kokunoyumeto.github.io/open-math-courses/courses/AG-CA/regular-sequences-depth-and-cohen-macaulay-modules.html}{AG-CA-12, Theorem 4.2 and Corollary 4.3} \\
lemma regular domain & \href{https://kokunoyumeto.github.io/open-math-courses/courses/AG-CA/regular-sequences-depth-and-cohen-macaulay-modules.html}{AG-CA-12, Theorem 6.1} \\
lemma regular goes up & \href{https://kokunoyumeto.github.io/open-math-courses/courses/AG-CA/regular-local-rings.html}{AG-CA-14, Proposition 3.3} \\
lemma regular ring CM & \href{https://kokunoyumeto.github.io/open-math-courses/courses/AG-CA/regular-sequences-depth-and-cohen-macaulay-modules.html}{AG-CA-12, Theorem 6.1} \\
lemma regular sequence powers & \href{https://kokunoyumeto.github.io/open-math-courses/courses/AG-CA/regular-sequences-depth-and-cohen-macaulay-modules.html}{AG-CA-12, Proposition 1.3} \\
lemma spec localization & \href{https://kokunoyumeto.github.io/open-math-courses/courses/AG-CA/localization-local-properties-and-support.html}{AG-CA-02, Theorem 3.1} \\
lemma variant local criterion flatness & \href{https://kokunoyumeto.github.io/open-math-courses/courses/AG-CA/faithful-flatness-and-the-local-criterion-for-flatness.html}{AG-CA-08, Corollary 4.4 and equation (6)} \\
proposition CM module & \href{https://kokunoyumeto.github.io/open-math-courses/courses/AG-CA/regular-sequences-depth-and-cohen-macaulay-modules.html}{AG-CA-12, Theorem 4.2 and Corollary 4.3} \\
proposition characterize formally smooth & \href{https://kokunoyumeto.github.io/open-math-courses/courses/AG-CA/formally-smooth-unramified-and-etale-ring-maps.html}{AG-CA-17, Theorem 3.1} \\
proposition smooth formally smooth & \href{https://kokunoyumeto.github.io/open-math-courses/courses/AG-CA/formally-smooth-unramified-and-etale-ring-maps.html}{AG-CA-17, Theorem 3.1 and the split conormal definition of smoothness in this draft} \\
lemma completion regular & \href{https://kokunoyumeto.github.io/open-math-courses/courses/AG-CA/completion.html}{AG-CA-19, Theorem 4.1} \\
complete-local simple-root lifting & \href{https://kokunoyumeto.github.io/open-math-courses/courses/AG-CA/completion.html\#5-hensel-lifting-in-every-complete-local-ring}{AG-CA-19, Theorem 5.1} \\
finite residue fields in finite-type fibres & \href{https://kokunoyumeto.github.io/open-math-courses/courses/AG-CA/the-nullstellensatz-and-jacobson-rings.html\#2-closed-points-detect-radical-equations}{AG-CA-06, Theorem 2.1} \\
minimal resolutions and Tor detection & \href{https://kokunoyumeto.github.io/open-math-courses/courses/AG-CA/projective-dimension-and-the-auslander-buchsbaum-formula.html\#2-minimal-resolutions-over-a-local-ring}{AG-CA-13, Theorem 2.2} \\
regular Noetherian rings are normal & \href{https://kokunoyumeto.github.io/open-math-courses/courses/AG-CA/discrete-valuation-rings-normal-rings-and-serres-criterion.html\#4-serre-conditions-and-the-normality-criterion}{AG-CA-15, Corollary 4.5} \\
arbitrary fields lift over their prime field & \href{https://kokunoyumeto.github.io/open-math-courses/courses/AG-CA/coefficient-rings-and-cohen-structure.html\#1-lifting-maps-from-the-prime-field}{AG-CA-19S, Lemma 1.1} \\
complete local coefficient rings & \href{https://kokunoyumeto.github.io/open-math-courses/courses/AG-CA/coefficient-rings-and-cohen-structure.html\#5-coefficient-subrings-in-every-complete-local-ring}{AG-CA-19S, Theorem 5.1} \\
complete local power-series presentations & \href{https://kokunoyumeto.github.io/open-math-courses/courses/AG-CA/coefficient-rings-and-cohen-structure.html\#6-the-power-series-presentation}{AG-CA-19S, Theorem 6.1} \\
\end{longtable}
}

\section{Sources, attribution, and history}\label{AG-RG-S05-sources-attribution-and-history}

The mathematical exposition reused here is the Stacks Project authors' free \emph{Smoothing Ring Maps}, \emph{Algebra}, and \emph{More on Algebra}, with categorical, field, and topological support from the same pinned edition. The consulted public source is the \href{https://github.com/KokunoYumeto/unofficial-stacks-project-ai-drafts/tree/565b10e987aba5969b21145a0833f42d69f96790}{AI Integrated Stacks fork at the pinned commit}. The \href{https://github.com/KokunoYumeto/unofficial-stacks-project-ai-drafts/blob/565b10e987aba5969b21145a0833f42d69f96790/README.md\#attribution-and-licensing}{source licensing statement} identifies GNU Free Documentation License 1.2.

Original source edition: \emph{Smoothing Ring Maps} and supporting Stacks Project chapters, Stacks Project Authors and maintainers of the consulted fork. Modified version: \emph{Artin approximation for polynomial equations and its desingularization proof}, programme supporting draft, 5 October 2026. Adaptation, finite-system application, and explicit corrections written by GPT-6.1 Sol (OpenAI) in Codex at Ultra. Self-checked by the writing AI. The authors of the source are not responsible for this adaptation.

This document contains adapted GFDL expression and is retained under GFDL-1.2-only, with no Invariant Sections, no Front-Cover Texts and no Back-Cover Texts. It must not be labelled CC0. The complete unchanged \href{assets/GFDL-1.2.txt}{GNU FDL 1.2 licence} accompanies this edition and its PDF. The free editable source and its copying notice remain accessible at the pinned links above.
\appendix
\chapter*{GNU Free Documentation License, Version 1.2}
\addcontentsline{toc}{chapter}{GNU Free Documentation License, Version 1.2}
\begingroup\scriptsize
\begin{verbatim}
		GNU Free Documentation License
		  Version 1.2, November 2002


 Copyright (C) 2000,2001,2002  Free Software Foundation, Inc.
     51 Franklin St, Fifth Floor, Boston, MA  02110-1301  USA
 Everyone is permitted to copy and distribute verbatim copies
 of this license document, but changing it is not allowed.


0. PREAMBLE

The purpose of this License is to make a manual, textbook, or other
functional and useful document "free" in the sense of freedom: to
assure everyone the effective freedom to copy and redistribute it,
with or without modifying it, either commercially or noncommercially.
Secondarily, this License preserves for the author and publisher a way
to get credit for their work, while not being considered responsible
for modifications made by others.

This License is a kind of "copyleft", which means that derivative
works of the document must themselves be free in the same sense.  It
complements the GNU General Public License, which is a copyleft
license designed for free software.

We have designed this License in order to use it for manuals for free
software, because free software needs free documentation: a free
program should come with manuals providing the same freedoms that the
software does.  But this License is not limited to software manuals;
it can be used for any textual work, regardless of subject matter or
whether it is published as a printed book.  We recommend this License
principally for works whose purpose is instruction or reference.


1. APPLICABILITY AND DEFINITIONS

This License applies to any manual or other work, in any medium, that
contains a notice placed by the copyright holder saying it can be
distributed under the terms of this License.  Such a notice grants a
world-wide, royalty-free license, unlimited in duration, to use that
work under the conditions stated herein.  The "Document", below,
refers to any such manual or work.  Any member of the public is a
licensee, and is addressed as "you".  You accept the license if you
copy, modify or distribute the work in a way requiring permission
under copyright law.

A "Modified Version" of the Document means any work containing the
Document or a portion of it, either copied verbatim, or with
modifications and/or translated into another language.

A "Secondary Section" is a named appendix or a front-matter section of
the Document that deals exclusively with the relationship of the
publishers or authors of the Document to the Document's overall subject
(or to related matters) and contains nothing that could fall directly
within that overall subject.  (Thus, if the Document is in part a
textbook of mathematics, a Secondary Section may not explain any
mathematics.)  The relationship could be a matter of historical
connection with the subject or with related matters, or of legal,
commercial, philosophical, ethical or political position regarding
them.

The "Invariant Sections" are certain Secondary Sections whose titles
are designated, as being those of Invariant Sections, in the notice
that says that the Document is released under this License.  If a
section does not fit the above definition of Secondary then it is not
allowed to be designated as Invariant.  The Document may contain zero
Invariant Sections.  If the Document does not identify any Invariant
Sections then there are none.

The "Cover Texts" are certain short passages of text that are listed,
as Front-Cover Texts or Back-Cover Texts, in the notice that says that
the Document is released under this License.  A Front-Cover Text may
be at most 5 words, and a Back-Cover Text may be at most 25 words.

A "Transparent" copy of the Document means a machine-readable copy,
represented in a format whose specification is available to the
general public, that is suitable for revising the document
straightforwardly with generic text editors or (for images composed of
pixels) generic paint programs or (for drawings) some widely available
drawing editor, and that is suitable for input to text formatters or
for automatic translation to a variety of formats suitable for input
to text formatters.  A copy made in an otherwise Transparent file
format whose markup, or absence of markup, has been arranged to thwart
or discourage subsequent modification by readers is not Transparent.
An image format is not Transparent if used for any substantial amount
of text.  A copy that is not "Transparent" is called "Opaque".

Examples of suitable formats for Transparent copies include plain
ASCII without markup, Texinfo input format, LaTeX input format, SGML
or XML using a publicly available DTD, and standard-conforming simple
HTML, PostScript or PDF designed for human modification.  Examples of
transparent image formats include PNG, XCF and JPG.  Opaque formats
include proprietary formats that can be read and edited only by
proprietary word processors, SGML or XML for which the DTD and/or
processing tools are not generally available, and the
machine-generated HTML, PostScript or PDF produced by some word
processors for output purposes only.

The "Title Page" means, for a printed book, the title page itself,
plus such following pages as are needed to hold, legibly, the material
this License requires to appear in the title page.  For works in
formats which do not have any title page as such, "Title Page" means
the text near the most prominent appearance of the work's title,
preceding the beginning of the body of the text.

A section "Entitled XYZ" means a named subunit of the Document whose
title either is precisely XYZ or contains XYZ in parentheses following
text that translates XYZ in another language.  (Here XYZ stands for a
specific section name mentioned below, such as "Acknowledgements",
"Dedications", "Endorsements", or "History".)  To "Preserve the Title"
of such a section when you modify the Document means that it remains a
section "Entitled XYZ" according to this definition.

The Document may include Warranty Disclaimers next to the notice which
states that this License applies to the Document.  These Warranty
Disclaimers are considered to be included by reference in this
License, but only as regards disclaiming warranties: any other
implication that these Warranty Disclaimers may have is void and has
no effect on the meaning of this License.


2. VERBATIM COPYING

You may copy and distribute the Document in any medium, either
commercially or noncommercially, provided that this License, the
copyright notices, and the license notice saying this License applies
to the Document are reproduced in all copies, and that you add no other
conditions whatsoever to those of this License.  You may not use
technical measures to obstruct or control the reading or further
copying of the copies you make or distribute.  However, you may accept
compensation in exchange for copies.  If you distribute a large enough
number of copies you must also follow the conditions in section 3.

You may also lend copies, under the same conditions stated above, and
you may publicly display copies.


3. COPYING IN QUANTITY

If you publish printed copies (or copies in media that commonly have
printed covers) of the Document, numbering more than 100, and the
Document's license notice requires Cover Texts, you must enclose the
copies in covers that carry, clearly and legibly, all these Cover
Texts: Front-Cover Texts on the front cover, and Back-Cover Texts on
the back cover.  Both covers must also clearly and legibly identify
you as the publisher of these copies.  The front cover must present
the full title with all words of the title equally prominent and
visible.  You may add other material on the covers in addition.
Copying with changes limited to the covers, as long as they preserve
the title of the Document and satisfy these conditions, can be treated
as verbatim copying in other respects.

If the required texts for either cover are too voluminous to fit
legibly, you should put the first ones listed (as many as fit
reasonably) on the actual cover, and continue the rest onto adjacent
pages.

If you publish or distribute Opaque copies of the Document numbering
more than 100, you must either include a machine-readable Transparent
copy along with each Opaque copy, or state in or with each Opaque copy
a computer-network location from which the general network-using
public has access to download using public-standard network protocols
a complete Transparent copy of the Document, free of added material.
If you use the latter option, you must take reasonably prudent steps,
when you begin distribution of Opaque copies in quantity, to ensure
that this Transparent copy will remain thus accessible at the stated
location until at least one year after the last time you distribute an
Opaque copy (directly or through your agents or retailers) of that
edition to the public.

It is requested, but not required, that you contact the authors of the
Document well before redistributing any large number of copies, to give
them a chance to provide you with an updated version of the Document.


4. MODIFICATIONS

You may copy and distribute a Modified Version of the Document under
the conditions of sections 2 and 3 above, provided that you release
the Modified Version under precisely this License, with the Modified
Version filling the role of the Document, thus licensing distribution
and modification of the Modified Version to whoever possesses a copy
of it.  In addition, you must do these things in the Modified Version:

A. Use in the Title Page (and on the covers, if any) a title distinct
   from that of the Document, and from those of previous versions
   (which should, if there were any, be listed in the History section
   of the Document).  You may use the same title as a previous version
   if the original publisher of that version gives permission.
B. List on the Title Page, as authors, one or more persons or entities
   responsible for authorship of the modifications in the Modified
   Version, together with at least five of the principal authors of the
   Document (all of its principal authors, if it has fewer than five),
   unless they release you from this requirement.
C. State on the Title page the name of the publisher of the
   Modified Version, as the publisher.
D. Preserve all the copyright notices of the Document.
E. Add an appropriate copyright notice for your modifications
   adjacent to the other copyright notices.
F. Include, immediately after the copyright notices, a license notice
   giving the public permission to use the Modified Version under the
   terms of this License, in the form shown in the Addendum below.
G. Preserve in that license notice the full lists of Invariant Sections
   and required Cover Texts given in the Document's license notice.
H. Include an unaltered copy of this License.
I. Preserve the section Entitled "History", Preserve its Title, and add
   to it an item stating at least the title, year, new authors, and
   publisher of the Modified Version as given on the Title Page.  If
   there is no section Entitled "History" in the Document, create one
   stating the title, year, authors, and publisher of the Document as
   given on its Title Page, then add an item describing the Modified
   Version as stated in the previous sentence.
J. Preserve the network location, if any, given in the Document for
   public access to a Transparent copy of the Document, and likewise
   the network locations given in the Document for previous versions
   it was based on.  These may be placed in the "History" section.
   You may omit a network location for a work that was published at
   least four years before the Document itself, or if the original
   publisher of the version it refers to gives permission.
K. For any section Entitled "Acknowledgements" or "Dedications",
   Preserve the Title of the section, and preserve in the section all
   the substance and tone of each of the contributor acknowledgements
   and/or dedications given therein.
L. Preserve all the Invariant Sections of the Document,
   unaltered in their text and in their titles.  Section numbers
   or the equivalent are not considered part of the section titles.
M. Delete any section Entitled "Endorsements".  Such a section
   may not be included in the Modified Version.
N. Do not retitle any existing section to be Entitled "Endorsements"
   or to conflict in title with any Invariant Section.
O. Preserve any Warranty Disclaimers.

If the Modified Version includes new front-matter sections or
appendices that qualify as Secondary Sections and contain no material
copied from the Document, you may at your option designate some or all
of these sections as invariant.  To do this, add their titles to the
list of Invariant Sections in the Modified Version's license notice.
These titles must be distinct from any other section titles.

You may add a section Entitled "Endorsements", provided it contains
nothing but endorsements of your Modified Version by various
parties--for example, statements of peer review or that the text has
been approved by an organization as the authoritative definition of a
standard.

You may add a passage of up to five words as a Front-Cover Text, and a
passage of up to 25 words as a Back-Cover Text, to the end of the list
of Cover Texts in the Modified Version.  Only one passage of
Front-Cover Text and one of Back-Cover Text may be added by (or
through arrangements made by) any one entity.  If the Document already
includes a cover text for the same cover, previously added by you or
by arrangement made by the same entity you are acting on behalf of,
you may not add another; but you may replace the old one, on explicit
permission from the previous publisher that added the old one.

The author(s) and publisher(s) of the Document do not by this License
give permission to use their names for publicity for or to assert or
imply endorsement of any Modified Version.


5. COMBINING DOCUMENTS

You may combine the Document with other documents released under this
License, under the terms defined in section 4 above for modified
versions, provided that you include in the combination all of the
Invariant Sections of all of the original documents, unmodified, and
list them all as Invariant Sections of your combined work in its
license notice, and that you preserve all their Warranty Disclaimers.

The combined work need only contain one copy of this License, and
multiple identical Invariant Sections may be replaced with a single
copy.  If there are multiple Invariant Sections with the same name but
different contents, make the title of each such section unique by
adding at the end of it, in parentheses, the name of the original
author or publisher of that section if known, or else a unique number.
Make the same adjustment to the section titles in the list of
Invariant Sections in the license notice of the combined work.

In the combination, you must combine any sections Entitled "History"
in the various original documents, forming one section Entitled
"History"; likewise combine any sections Entitled "Acknowledgements",
and any sections Entitled "Dedications".  You must delete all sections
Entitled "Endorsements".


6. COLLECTIONS OF DOCUMENTS

You may make a collection consisting of the Document and other documents
released under this License, and replace the individual copies of this
License in the various documents with a single copy that is included in
the collection, provided that you follow the rules of this License for
verbatim copying of each of the documents in all other respects.

You may extract a single document from such a collection, and distribute
it individually under this License, provided you insert a copy of this
License into the extracted document, and follow this License in all
other respects regarding verbatim copying of that document.


7. AGGREGATION WITH INDEPENDENT WORKS

A compilation of the Document or its derivatives with other separate
and independent documents or works, in or on a volume of a storage or
distribution medium, is called an "aggregate" if the copyright
resulting from the compilation is not used to limit the legal rights
of the compilation's users beyond what the individual works permit.
When the Document is included in an aggregate, this License does not
apply to the other works in the aggregate which are not themselves
derivative works of the Document.

If the Cover Text requirement of section 3 is applicable to these
copies of the Document, then if the Document is less than one half of
the entire aggregate, the Document's Cover Texts may be placed on
covers that bracket the Document within the aggregate, or the
electronic equivalent of covers if the Document is in electronic form.
Otherwise they must appear on printed covers that bracket the whole
aggregate.


8. TRANSLATION

Translation is considered a kind of modification, so you may
distribute translations of the Document under the terms of section 4.
Replacing Invariant Sections with translations requires special
permission from their copyright holders, but you may include
translations of some or all Invariant Sections in addition to the
original versions of these Invariant Sections.  You may include a
translation of this License, and all the license notices in the
Document, and any Warranty Disclaimers, provided that you also include
the original English version of this License and the original versions
of those notices and disclaimers.  In case of a disagreement between
the translation and the original version of this License or a notice
or disclaimer, the original version will prevail.

If a section in the Document is Entitled "Acknowledgements",
"Dedications", or "History", the requirement (section 4) to Preserve
its Title (section 1) will typically require changing the actual
title.


9. TERMINATION

You may not copy, modify, sublicense, or distribute the Document except
as expressly provided for under this License.  Any other attempt to
copy, modify, sublicense or distribute the Document is void, and will
automatically terminate your rights under this License.  However,
parties who have received copies, or rights, from you under this
License will not have their licenses terminated so long as such
parties remain in full compliance.


10. FUTURE REVISIONS OF THIS LICENSE

The Free Software Foundation may publish new, revised versions
of the GNU Free Documentation License from time to time.  Such new
versions will be similar in spirit to the present version, but may
differ in detail to address new problems or concerns.  See
https://www.gnu.org/copyleft/.

Each version of the License is given a distinguishing version number.
If the Document specifies that a particular numbered version of this
License "or any later version" applies to it, you have the option of
following the terms and conditions either of that specified version or
of any later version that has been published (not as a draft) by the
Free Software Foundation.  If the Document does not specify a version
number of this License, you may choose any version ever published (not
as a draft) by the Free Software Foundation.


ADDENDUM: How to use this License for your documents

To use this License in a document you have written, include a copy of
the License in the document and put the following copyright and
license notices just after the title page:

    Copyright (c)  YEAR  YOUR NAME.
    Permission is granted to copy, distribute and/or modify this document
    under the terms of the GNU Free Documentation License, Version 1.2
    or any later version published by the Free Software Foundation;
    with no Invariant Sections, no Front-Cover Texts, and no Back-Cover Texts.
    A copy of the license is included in the section entitled "GNU
    Free Documentation License".

If you have Invariant Sections, Front-Cover Texts and Back-Cover Texts,
replace the "with...Texts." line with this:

    with the Invariant Sections being LIST THEIR TITLES, with the
    Front-Cover Texts being LIST, and with the Back-Cover Texts being LIST.

If you have Invariant Sections without Cover Texts, or some other
combination of the three, merge those two alternatives to suit the
situation.

If your document contains nontrivial examples of program code, we
recommend releasing these examples in parallel under your choice of
free software license, such as the GNU General Public License,
to permit their use in free software.

\end{verbatim}
\endgroup
\end{document}
